%!TeX program = lualatex
%
% Appunti di Algoritmi Avanzati -- trascrizione LaTeX degli appunti manoscritti.
% Compilare con:  make        (oppure: lualatex main, due volte)
%
\documentclass[oneside]{academica}
%%%%%%%%%%%%%%%%%%%%%%%%%%%%%%%%%%%%%%%%%%%%%%%%%%%%%%%%%%%%%%%%%%%%%%%%%%%%%%
%% Preambolo -- Algoritmi Avanzati
%% Classe: academica (Airy).  Compilare con lualatex.
%%%%%%%%%%%%%%%%%%%%%%%%%%%%%%%%%%%%%%%%%%%%%%%%%%%%%%%%%%%%%%%%%%%%%%%%%%%%%%

% academica carica gia': tikz, pgfplots, mathtools, amssymb, tcolorbox, xcolor,
% graphicx, float, caption/subcaption, listings, hyperref, xspace, unicode-math.

\usetikzlibrary{automata,positioning,trees,arrows.meta,calc,fit,shapes.geometric,
                decorations.pathreplacing,decorations.markings,matrix,backgrounds,chains}
\usepackage{booktabs}
\usepackage{multirow}
\usepackage{array}
\usepackage{algorithm}
\usepackage{algpseudocode}
\usepackage{marginnote}

% Localizzazione italiana.
% NB: non si puo' usare babel/polyglossia perche' academica carica biblatex
% (che pretende di essere caricato DOPO il pacchetto di lingua). Rinominiamo
% quindi a mano le poche stringhe che compaiono nel documento.
\renewcommand{\chaptername}{Capitolo}
\renewcommand{\contentsname}{Indice}
\renewcommand{\listfigurename}{Elenco delle figure}
\renewcommand{\listtablename}{Elenco delle tabelle}
\renewcommand{\figurename}{Figura}
\renewcommand{\tablename}{Tabella}
\renewcommand{\bibname}{Bibliografia}
\renewcommand{\indexname}{Indice analitico}
\renewcommand{\appendixname}{Appendice}
\renewcommand{\partname}{Parte}

\providecommand{\qedhere}{}

% academica scrive "List of Figures"/"List of Tables" dentro una tikzpicture,
% quindi \listfigurename non basta: va ridefinito il comando.
\makeatletter
\renewcommand{\listoffigures}{%
  \chapter*{%
    \vspace*{-20\p@}%
    \begin{tikzpicture}[overlay]
      \draw[fill=accent,draw=accent,rounded corners] (10,-0.65) rectangle (20,1);
      \pgftext[right,x=16.7cm,y=0.2cm]{\Huge\sc\bfseries\textcolor{white}{\listfigurename}}
    \end{tikzpicture}%
  }%
  \@starttoc{lof}%
}
\renewcommand{\listoftables}{%
  \chapter*{%
    \vspace*{-20\p@}%
    \begin{tikzpicture}[overlay]
      \draw[fill=accent,draw=accent,rounded corners] (10.5,-0.65) rectangle (20,1);
      \pgftext[right,x=16.7cm,y=0.2cm]{\Huge\sc\bfseries\textcolor{white}{\listtablename}}
    \end{tikzpicture}%
  }%
  \@starttoc{lot}%
}
\makeatother

\setlength{\parindent}{0pt}

%%%%%%%%%%%%%%%%%%%%%%%%%%%%%%%%%%%%%%%%%%%%%%%%%%%%%%%%%%%%%%%%%%%%%%%%%%%%%%
%% AMBIENTI IN ITALIANO
%% academica definisce gli ambienti colorati in inglese (Definition, Theorem,
%% ...). Qui aggiungiamo gli alias italiani, numerati e non, riusando gli
%% stessi contatori e colori della classe.
%%%%%%%%%%%%%%%%%%%%%%%%%%%%%%%%%%%%%%%%%%%%%%%%%%%%%%%%%%%%%%%%%%%%%%%%%%%%%%

\makeatletter
% versione numerata (riusa i contatori della classe: stessa numerazione
% degli ambienti inglesi, così i due stili non divergono)
\newcommand{\aa@num}[4]{% #1 nome env, #2 contatore, #3 etichetta, #4 colore
  \newenvironment{#1}[1][]{%
    \refstepcounter{#2}\def\@currentlabelname{##1}%
    \begin{accentbox}[##1]{#3}{\csname the#2\endcsname}{#4}%
  }{\end{accentbox}}%
  \expandafter\newcommand\csname #1autorefname\endcsname{#3}%
}
\aa@num{definizione}{definition}{Definizione}{lilla}
\aa@num{teorema}{theorem}{Teorema}{blue}
\aa@num{proposizione}{proposition}{Proposizione}{deepsea}
\aa@num{corollario}{corollary}{Corollario}{deepsea}
\aa@num{dimostrazione}{proof}{Dimostrazione}{pink}
\aa@num{esempio}{example}{Esempio}{orange}
\aa@num{esercizio}{exercise}{Esercizio}{orange}
\aa@num{domanda}{question}{Domanda}{green}
\makeatother

% versioni non numerate
\newenvironment{definizione*}[1][]{\begin{accentbox*}[#1]{Definizione}{lilla}}{\end{accentbox*}}
\newenvironment{teorema*}[1][]{\begin{accentbox*}[#1]{Teorema}{blue}}{\end{accentbox*}}
\newenvironment{proposizione*}[1][]{\begin{accentbox*}[#1]{Proposizione}{deepsea}}{\end{accentbox*}}
\newenvironment{corollario*}[1][]{\begin{accentbox*}[#1]{Corollario}{deepsea}}{\end{accentbox*}}
\newenvironment{dimostrazione*}[1][]{\begin{accentbox*}[#1]{Dimostrazione}{pink}}{\end{accentbox*}}
\newenvironment{esempio*}[1][]{\begin{accentbox*}[#1]{Esempio}{orange}}{\end{accentbox*}}
\newenvironment{esercizio*}[1][]{\begin{accentbox*}[#1]{Esercizio}{orange}}{\end{accentbox*}}
\newenvironment{domanda*}[1][]{\begin{accentbox*}[#1]{Domanda}{green}}{\end{accentbox*}}
\newenvironment{lemma*it}[1][]{\begin{accentbox*}[#1]{Lemma}{deepsea}}{\end{accentbox*}}

% "Osservazione" non numerata (equivalente italiano di \begin{note})
\newenvironment{osservazione}[1][]{\begin{accentbox*}[#1]{Osservazione}{gray}}{\end{accentbox*}}

% autoref in italiano per gli ambienti della classe
\renewcommand{\lemmaautorefname}{Lemma}
\renewcommand{\chapterautorefname}{Capitolo}
\renewcommand{\sectionautorefname}{Sezione}

\hyphenation{suf-fis-so suf-fis-si pre-fis-so pre-fis-si strin-ga strin-ghe}
\hyphenation{al-li-nea-men-to ran-do-miz-za-to ran-do-miz-za-ti com-pres-se}

%%%%%%%%%%%%%%%%%%%%%%%%%%%%%%%%%%%%%%%%%%%%%%%%%%%%%%%%%%%%%%%%%%%%%%%%%%%%%%
%% MARCATORI EDITORIALI
%% Usati SOLO quando il dubbio resta irrisolto. Ogni scelta e' comunque
%% registrata in NOTE-EDITORIALI.md.
%%%%%%%%%%%%%%%%%%%%%%%%%%%%%%%%%%%%%%%%%%%%%%%%%%%%%%%%%%%%%%%%%%%%%%%%%%%%%%

% Riferimento alla pagina della scansione originale (appunti.pdf).
% Usa \marginnote e non \marginpar: quest'ultimo e' un float e va perso se
% invocato dentro un tcolorbox (cioe' dentro definizione/teorema/...).
%
% NB: academica non riserva spazio a margine (\marginparwidth vale 0pt), per cui
% senza queste due righe ogni \scan produce un Overfull \hbox di ~20pt. Il
% margine destro utile e' 44.8pt (597.5 - 72.27 - 2.42 - 478.01), quindi si sta
% larghi con 34+4 = 38pt.
\setlength{\marginparwidth}{34pt}
\setlength{\marginparsep}{4pt}

\newcommand{\scan}[1]{%
  \marginnote{%
    \raggedright\scriptsize\color{gray!140}\textsf{scan\\p.\,#1}%
  }%
}

% Blocco: lettura del manoscritto non determinabile con certezza
\newenvironment{dubbio}[1][Lettura incerta]{%
  \begin{accentbox*}[#1]{Dubbio}{red}%
}{%
  \end{accentbox*}%
}

% Blocco: buco negli appunti, ricostruzione non conclusiva
\newenvironment{lacuna}[1][Lacuna negli appunti]{%
  \begin{accentbox*}[#1]{Lacuna}{yellow}%
}{%
  \end{accentbox*}%
}

% Versioni inline
\newcommand{\dubbioinl}[1]{{\color{red}[?\,\textit{#1}]}}
\newcommand{\lacunainl}[1]{{\color{orange}[\dots\,\textit{#1}]}}

%%%%%%%%%%%%%%%%%%%%%%%%%%%%%%%%%%%%%%%%%%%%%%%%%%%%%%%%%%%%%%%%%%%%%%%%%%%%%%
%% STRINGHE E STRING MATCHING
%%%%%%%%%%%%%%%%%%%%%%%%%%%%%%%%%%%%%%%%%%%%%%%%%%%%%%%%%%%%%%%%%%%%%%%%%%%%%%

\newcommand{\alfa}{\S}                    % alfabeto (\S = \Sigma da academica)
\newcommand{\str}[1]{\texttt{#1}}         % stringa letterale
\newcommand{\sub}[3]{#1[#2\,..\,#3]}      % sottostringa S[i..j]
\newcommand{\suf}[2]{#1[#2\,..\,]}        % suffisso S[i..]
\newcommand{\pre}[2]{#1[1\,..\,#2]}       % prefisso S[1..i]
\newcommand{\conc}{\cdot}                 % concatenazione
\newcommand{\vuota}{\varepsilon}          % stringa vuota

\DeclareMathOperator{\lcp}{lcp}
\DeclareMathOperator{\lcs}{lcs}
\DeclareMathOperator{\lca}{lca}
\DeclareMathOperator{\lcey}{lce}

% KMP / Z
\newcommand{\spv}[1]{\mathit{sp}_{#1}}     % sp_i
\newcommand{\spp}[1]{\mathit{sp}'_{#1}}    % sp'_i
\newcommand{\Zv}[1]{Z_{#1}}                % Z_i
\newcommand{\Zbox}{Z\text{-box}\xspace}
\newcommand{\fail}{\mathit{fail}}

% Boyer-Moore
\newcommand{\bc}[1]{R(#1)}                 % bad character rule
\newcommand{\Lp}[1]{L(#1)}                 % good suffix (strong) rule
\newcommand{\lp}[1]{\ell(#1)}

% Aho-Corasick
\newcommand{\goto}{\mathit{goto}}
\newcommand{\out}{\mathit{output}}
\newcommand{\nv}[1]{n_{#1}}                % failure link

% Suffix tree / trie
\newcommand{\STrie}{\textsf{STrie}\xspace}
\newcommand{\STree}{\textsf{STree}\xspace}
\newcommand{\albero}[1]{\mathcal{T}(#1)}
\newcommand{\slink}{\mathit{sl}}           % suffix link
\newcommand{\etich}[1]{\mathit{label}(#1)}
\newcommand{\cammino}[1]{\mathit{path}(#1)}

% Allineamento / distanze
\newcommand{\dedit}{d_{\mathrm{edit}}}
\newcommand{\dham}{d_{\mathrm{H}}}
\newcommand{\Dm}[2]{D(#1,#2)}
\newcommand{\Vm}[2]{V(#1,#2)}
\newcommand{\gap}{\mathit{gap}}
\newcommand{\indel}{\mathit{indel}}

%%%%%%%%%%%%%%%%%%%%%%%%%%%%%%%%%%%%%%%%%%%%%%%%%%%%%%%%%%%%%%%%%%%%%%%%%%%%%%
%% ALGORITMI RANDOMIZZATI
%%%%%%%%%%%%%%%%%%%%%%%%%%%%%%%%%%%%%%%%%%%%%%%%%%%%%%%%%%%%%%%%%%%%%%%%%%%%%%

\newcommand{\Prob}{\mathbb{P}}
\newcommand{\E}{\mathbb{E}}
\renewcommand{\Pr}{\operatorname{Pr}}
\DeclareMathOperator{\Var}{Var}
\DeclareMathOperator*{\argmin}{arg\,min}
\DeclareMathOperator*{\argmax}{arg\,max}
\newcommand{\indic}[1]{\mathbf{1}_{\{#1\}}}
\newcommand{\RandQS}{\textsc{RandQS}\xspace}
\newcommand{\LasVegas}{Las~Vegas\xspace}
\newcommand{\MonteCarlo}{Monte~Carlo\xspace}
\newcommand{\mincut}{\textsc{MinCut}\xspace}
\newcommand{\contract}{\textsc{Contract}\xspace}
% classi di complessita' randomizzate
\newcommand{\RP}{\mathsf{RP}}
\newcommand{\coRP}{\mathsf{coRP}}
\newcommand{\ZPP}{\mathsf{ZPP}}
\newcommand{\BPP}{\mathsf{BPP}}
\newcommand{\PP}{\mathsf{PP}}
\newcommand{\PSPACE}{\mathsf{PSPACE}}
\newcommand{\NP}{\mathsf{NP}}
\renewcommand{\P}{\mathsf{P}}   % nota: academica definiva \P = \mathcal{P}

%%%%%%%%%%%%%%%%%%%%%%%%%%%%%%%%%%%%%%%%%%%%%%%%%%%%%%%%%%%%%%%%%%%%%%%%%%%%%%
%% STRUTTURE DATI COMPRESSE
%%%%%%%%%%%%%%%%%%%%%%%%%%%%%%%%%%%%%%%%%%%%%%%%%%%%%%%%%%%%%%%%%%%%%%%%%%%%%%

\DeclareMathOperator{\rank}{rank}
\DeclareMathOperator{\select}{select}
\DeclareMathOperator{\access}{access}
\newcommand{\BWT}{\textsf{BWT}\xspace}
\newcommand{\FMindex}{\textsf{FM-index}\xspace}
\newcommand{\SA}{\mathit{SA}}
\newcommand{\LF}{\mathit{LF}}
\newcommand{\entropia}[1]{H_{#1}}

%%%%%%%%%%%%%%%%%%%%%%%%%%%%%%%%%%%%%%%%%%%%%%%%%%%%%%%%%%%%%%%%%%%%%%%%%%%%%%
%% VARIE
%%%%%%%%%%%%%%%%%%%%%%%%%%%%%%%%%%%%%%%%%%%%%%%%%%%%%%%%%%%%%%%%%%%%%%%%%%%%%%

\newcommand{\Oh}{\mathcal{O}}
\newcommand{\Th}{\Theta}
\newcommand{\Om}{\Omega}
\newcommand{\defeq}{\vcentcolon=}
\renewcommand{\emptyset}{\varnothing}

% Ambienti algoritmo in italiano
\algrenewcommand\algorithmicrequire{\textbf{Input:}}
\algrenewcommand\algorithmicensure{\textbf{Output:}}
\algrenewcommand\algorithmicwhile{\textbf{while}}
\algrenewcommand\algorithmicfor{\textbf{for}}
\algrenewcommand\algorithmicif{\textbf{if}}
\algrenewcommand\algorithmicreturn{\textbf{return}}
\floatname{algorithm}{Algoritmo}


\overtitle{appunti di}
\title{Algoritmi Avanzati}
\author{rw-r-r-0644}
\date{}

\begin{document}

\frontmatter
\maketitle

\chapter*{Avvertenza}
\markboth{Avvertenza}{Avvertenza}

\begin{note}[Trascrizione automatica, non ancora rivista]
Questi appunti sono stati trascritti da un LLM a partire da un quaderno scritto a
mano. Non sono materiale ufficiale del corso e nessuno li ha ancora riletti.

Il punto importante: dove pensava di esserne certo, il modello ha \textbf{corretto
errori e ricostruito pezzi mancanti senza segnalarlo}. Ci sono quindi passaggi ---
e figure --- che non erano negli appunti originali, e che possono divergere da
quanto è stato detto a lezione. In più la calligrafia era difficile, quindi
qualche simbolo può essere stato letto male.

\textbf{Nessuna garanzia}, insomma: prima di fidarvi di un passaggio, a maggior
ragione se state preparando l'esame, verificatelo sui libri o sul materiale del
corso.

Per controllare: la nota a margine \textsf{scan p.\,N} rimanda alla pagina della
scansione da cui quel testo viene, e \texttt{note-trascrizione/NOTE-EDITORIALI.md}
elenca tutte le scelte fatte. I dubbi rimasti aperti --- una quarantina --- sono
marcati qui dentro con i riquadri \textbf{Dubbio} e \textbf{Lacuna}; occhio però a
non fraintenderli, perché il resto è \emph{non marcato}, non \emph{verificato}.
\end{note}

\tableofcontents

\mainmatter

\part{String matching esatto}
\chapter{Introduzione e pattern matching esatto}
\label{cap:introduzione}

\section{Presentazione del corso}
\label{sec:presentazione}

\scan{1}
Il corso di Algoritmi Avanzati è tenuto dal prof.\ Alberto Policriti.

\subsection{Le tre parti del corso}

Il corso si articola in tre parti, a ciascuna delle quali corrisponde un libro di
riferimento: \emph{string matching}, \emph{algoritmi randomizzati} e
\emph{strutture dati compresse}.

\begin{figure}[H]
  \centering
  % corso-tre-parti -- mappa delle tre parti del corso (quaderno, pag. 1).
\begin{tikzpicture}[
    >={Stealth[length=4.5pt,width=3.5pt]},
    font=\small,
    ramo/.style={->,semithick,shorten >=1.5pt},
    voce/.style={anchor=west,inner sep=1pt},
    glossa/.style={anchor=west,inner sep=1pt,font=\footnotesize},
  ]

  % --- livello 0: la radice ------------------------------------------------
  \node[anchor=west,inner sep=1pt] (radice) at (0,0) {3 parti};
  \coordinate (fork) at ($(radice.east)+(0.20,0)$);

  % --- livello 1: le tre parti del corso -----------------------------------
  \node[voce] (r1) at (2.00, 1.10) {\underline{string matching}};
  \node[voce] (r2) at (2.00, 0.00) {\underline{algoritmi randomizzati}};
  \node[voce] (r3) at (2.00,-1.10) {\underline{strutture dati compresse}};

  \draw[ramo] (fork) -- (r1.west);
  \draw[ramo] (fork) -- (r2.west);
  \draw[ramo] (fork) -- (r3.west);

  % --- diramazione: le due classi di algoritmi randomizzati -----------------
  \coordinate (v) at ($(r2.east)+(0.22,0)$);
  \node[glossa] (s1) at ($(v)+(0.40, 0.50)$) {Monte Carlo --- non sempre corretti};
  \node[glossa] (s2) at ($(v)+(0.40,-0.50)$) {Las Vegas --- sempre corretti};
  \draw[semithick] (v) -- ($(s1.west)+(-0.08,0)$);
  \draw[semithick] (v) -- ($(s2.west)+(-0.08,0)$);

  % --- glossa sull'obiettivo delle strutture dati compresse -----------------
  \node[glossa] (tilde) at ($(r3.east)+(0.14,0)$) {$\sim$};
  \node[glossa] (g3) at ($(tilde.east)+(0.08,0)$)
        {comprimere dati anche molto complessi};

\end{tikzpicture}

  \caption{Le tre parti in cui è suddiviso il corso, con le annotazioni a margine
  sulla dicotomia Monte Carlo / Las Vegas e sull'obiettivo delle strutture dati compresse.}
  \label{fig:corso-tre-parti}
\end{figure}

\begin{osservazione}
La distinzione fra i due tipi di algoritmi randomizzati riguarda dove viene «scaricata»
l'aleatorietà: un algoritmo \MonteCarlo ha tempo di esecuzione deterministico ma può
restituire una risposta errata (con probabilità piccola e controllabile), mentre un
algoritmo \LasVegas restituisce sempre la risposta corretta, e ad essere aleatorio è
invece il suo tempo di esecuzione.
\end{osservazione}

L'obiettivo della terza parte, infine, è comprimere dati anche molto complessi
mantenendoli interrogabili.

\subsection{Testi di riferimento e materiale}

\begin{enumerate}
  \item D.~Gusfield, \emph{Algorithms on Strings, Trees, and Sequences};
  \item R.~Motwani, P.~Raghavan, \emph{Randomized Algorithms};
  \item G.~Navarro, \emph{Compact Data Structures}.
\end{enumerate}

La numerazione dei testi ricalca quella delle tre parti del corso: a ogni modulo
corrisponde un testo canonico distinto. A questi si aggiungono gli appunti
(non ufficiali) del docente e i nostri appunti presi a lezione. Tutto il materiale
del corso è reso disponibile su Teams.

\subsection{Modalità d'esame}

L'esame è orale. Durante il corso vengono proposti spunti per approfondimenti;
l'esame consiste quindi in un (eventuale) approfondimento seguito dal colloquio orale.

\section{Il problema del pattern matching esatto}
\label{sec:pattern-matching-esatto}

\scan{2}
La prima delle tre parti del corso, quella dedicata al \emph{string matching}, comincia
dal problema che ne costituisce il nucleo.

\begin{definizione}[Pattern matching esatto]\label{def:pattern-matching-esatto}
\emph{Input.} Un alfabeto $\S$ di cardinalità finita e due stringhe $P,T\in\S^{*}$, con
$\abs{P}=n$ e $\abs{T}=m$ (di norma $n\le m$).

\emph{Output.} L'insieme delle posizioni
\[
  \set{i \mid \sub{T}{i}{i+n-1}=P}.
\]
\end{definizione}

Si vogliono cioè \emph{tutte} le occorrenze di $P$ (il \emph{pattern}) all'interno di $T$
(il \emph{testo}), non soltanto la prima.

\begin{osservazione}
Le stringhe sono indicizzate a partire da $1$ e $\sub{S}{i}{j}$ denota il fattore
$S[i]S[i+1]\cdots S[j]$. Perché il confronto abbia senso deve essere $1\le i\le m-n+1$:
le posizioni di allineamento del pattern sul testo sono dunque $m-n+1$.
\end{osservazione}

\begin{figure}[H]
\centering
% pm-allineamento -- il pattern P appoggiato sul testo T (quaderno, pag. 2).
\begin{tikzpicture}[
    >={Stealth[length=4pt,width=3pt]},
    stringa/.style={line width=1pt},
    car/.style={anchor=base,inner sep=0pt,font=\small},
  ]

  % --- il testo T ----------------------------------------------------------
  \draw[stringa] (0,0) -- (8.4,0);
  \node[car] at (0.30, 0.16) {$a$};
  \node[car] at (0.90, 0.16) {$b$};
  \node[car] at (1.50, 0.16) {$c$};
  \node[car] at (2.15, 0.16) {$\cdots$};
  \node[anchor=west,font=\small] at (8.6,0) {$T$};

  % --- il pattern P --------------------------------------------------------
  \draw[stringa] (0,-1.3) -- (3.4,-1.3);
  \node[car] at (0.30,-1.14) {$a$};
  \node[car] at (0.90,-1.14) {$b$};
  \node[car] at (1.50,-1.14) {$c$};
  \node[car] at (2.15,-1.14) {$\cdots$};
  \node[anchor=west,font=\small] (etP) at (3.6,-1.3) {$P$};

  % --- posizione corrente di confronto: stessa ascissa in T e in P ----------
  \draw[->] (0.30,-0.52) -- (0.30,-0.10);
  \draw[->] (0.30,-1.82) -- (0.30,-1.40);

\end{tikzpicture}

\caption{Il pattern $P$ allineato sul testo $T$: le due stringhe vengono confrontate
carattere per carattere a partire dalla posizione corrente (le frecce), e poi $P$ scorre
verso destra di una posizione alla volta.}
\label{fig:pm-allineamento}
\end{figure}

\section{L'algoritmo naive}
\label{sec:naive}

L'idea è provare tutti gli allineamenti, uno per uno, con due loop innestati: quello
esterno sceglie la posizione $i$ in cui il pattern viene appoggiato sul testo, quello
interno confronta $P$ con il fattore $\sub{T}{i}{i+n-1}$ carattere per carattere.

\begin{algorithm}[H]
\caption{Algoritmo naive per il pattern matching esatto}
\label{alg:naive}
\begin{algorithmic}[1]
\Require $T$ con $\abs{T}=m$, $P$ con $\abs{P}=n$
\Ensure tutte le posizioni $i$ tali che $\sub{T}{i}{i+n-1}=P$
\For{$i \gets 1$ \textbf{to} $m-n+1$}
    \Comment{allineamenti possibili}
    \State $j \gets 1$
    \While{$j \le n$ \textbf{and} $T[i+j-1] = P[j]$}
        \Comment{\textsc{compare}$(P,\sub{T}{i}{i+n-1})$}
        \State $j \gets j+1$
    \EndWhile
    \If{$j = n+1$}
        \State \textbf{output} $i$
        \Comment{$n$ confronti positivi: occorrenza in $i$}
    \EndIf
\EndFor
\end{algorithmic}
\end{algorithm}

Il ciclo interno --- la procedura \textsc{compare} --- scorre da sinistra a destra: o si
arresta al primo mismatch, e allora in posizione $i$ non c'è un'occorrenza, oppure arriva
in fondo dopo $\abs{P}=n$ confronti tutti positivi, e in tal caso $i$ viene emesso in
output. In entrambi i casi il pattern viene poi fatto scorrere di una sola posizione verso
destra: nessuna informazione raccolta durante i confronti viene riutilizzata.

\subsection{Complessità}
\label{sec:naive-complessita}

Il ciclo esterno esegue $m-n+1$ iterazioni e ogni chiamata a \textsc{compare} costa al più
$n$ confronti, da cui
\[
  \Oh(m\cdot n).
\]

Il limite è stretto (\emph{tight}): basta considerare il caso
\[
  P=a^{\,n-1}b, \qquad T=a^{\,m},
\]
cioè
\[
  T=\overbrace{a\,a\,a\,\cdots\,a\,a}^{m\ \text{volte}},
  \qquad
  P=\underbrace{a\,a\,\cdots\,a}_{n-1}\,b .
\]
Per ogni allineamento i primi $n-1$ confronti hanno successo e l'$n$-esimo fallisce: si
spendono esattamente $n(m-n+1)$ confronti --- cioè $\Th(m^{2})$ quando $n$ e $m-n$ sono
entrambi $\Th(m)$, con il massimo per $n\approx m/2$ --- e ciò nonostante l'output è vuoto.

D'altra parte non si può scendere sotto il costo della lettura dell'input. Per esempio, con
$n=1$ il naive costa $\Oh(m\cdot 1)=\Oh(m)$ ed è già ottimo: qualunque algoritmo deve almeno
leggere il tutto, e dunque
\[
  \Om(n+m).
\]

Per questo motivo la complessità $\Oh(m\cdot n)$ dell'algoritmo naive è detta
\emph{quadratica}. Fra questo upper bound e il lower bound banale $\Om(n+m)$ resta un
divario: la domanda naturale --- esiste un algoritmo sub-quadratico? --- è il punto di
partenza del capitolo successivo.

\chapter{Knuth--Morris--Pratt e la funzione \texorpdfstring{$Z$}{Z}}
\label{cap:kmp-z}

\scan{3}
Fissiamo una volta per tutte le notazioni. $P$ è il \emph{pattern} e $T$ il
\emph{testo}, entrambi stringhe sull'alfabeto $\alfa$; poniamo $n=\abs{P}$ e
$m=\abs{T}$, con $n\le m$, e indicizziamo le stringhe a partire da~$1$, così che
$\sub{S}{i}{j}$ denoti il fattore di $S$ compreso fra le posizioni $i$ e $j$
estremi inclusi. Il problema dell'\emph{exact string matching} chiede di trovare
tutte le occorrenze di $P$ in $T$.

\section{Il problema e i suoi due limiti}

I due termini del confronto sono
\[
  \underbrace{\Oh(n\cdot m)}_{\text{\emph{upper bound}: l'algoritmo naive}}
  \qquad\text{contro}\qquad
  \underbrace{\Om(n+m)}_{\text{\emph{lower bound} del problema}} .
\]
Il primo è quadratico ed è il costo dell'algoritmo naive, che prova tutti gli
allineamenti di $P$ su $T$; il secondo è lineare e non è un limite inferiore del
naive, bensì del \emph{problema} --- l'input va comunque letto per intero. Resta
quindi aperta la domanda:

\begin{domanda}\label{dom:subquadratico}
  Esiste un algoritmo \underline{sub}-quadratico per l'exact string matching?
\end{domanda}

\noindent Il bersaglio è il lower bound stesso, cioè un algoritmo che lavori in
tempo
\[
  \Th(n+m),
\]
ed è esattamente ciò che Knuth--Morris--Pratt riesce a ottenere.

\section{L'idea di Knuth--Morris--Pratt}

L'idea è di non buttare via l'informazione raccolta dai confronti già eseguiti.
Supponiamo di aver allineato $P$ su $T$, di aver riconosciuto correttamente il
prefisso $\pre{P}{i}$ e di trovare un mismatch al carattere successivo: il testo
presenta un carattere $y$, il pattern presenta $x=P[i+1]$, e $x\neq y$.

Il tratto già verificato è noto \emph{a priori} --- è un prefisso di $P$,
indipendente da $T$ --- e quindi tutta l'informazione riutilizzabile può essere
ricavata dal solo pattern.

Sia $\alpha$ il più lungo \emph{prefisso-suffisso} proprio di $\pre{P}{i}$, cioè
la più lunga stringa che ne sia simultaneamente prefisso proprio e suffisso
proprio. Poiché $\alpha$ compare sia in testa a $P$ sia in coda al tratto già
riconosciuto, posso far scivolare il pattern verso destra finché il suo prefisso
$\alpha$ non si sovrappone all'occorrenza di $\alpha$ che nel testo termina
appena prima del mismatch: lo spostamento è di $i-\abs{\alpha}$ posizioni e il
confronto riprende dalla \emph{stessa} posizione del testo, ora contro
$P[\abs{\alpha}+1]$. In particolare l'indice sul testo non torna mai indietro.

\begin{figure}[H]
  \centering
  % kmp-shift -- regola di spostamento di KMP (quaderno, pag. 3).
% Tre barre sovrapposte: il testo T, il pattern P nell'allineamento corrente,
% il pattern P dopo lo scorrimento di i-|alpha| posizioni.
% Proporzioni prese dal quaderno: i = 5 unita', |alpha| = 3, spostamento = 2.
% Le due barre di P hanno la stessa lunghezza (11 unita'): e' la stessa stringa.
\begin{tikzpicture}[
    x=0.52cm, y=1cm,
    font=\footnotesize,
    >={Stealth[length=4pt,width=3pt]},
    barra/.style={line width=0.7pt},
    sep/.style={line width=0.5pt},
    guida/.style={gray!60!black, line width=0.5pt},
  ]

  % ===== colonna del mismatch (attraversa tutte e tre le righe) =========
  \draw[guida] (6,-0.45) -- (6,2.62);
  \draw[guida] (7,-0.45) -- (7,2.62);

  % ===== RIGA 1 : il testo T ===========================================
  \draw[barra] (0,2.05) rectangle (20,2.50);
  % porzione gia' scandita / non specificata
  \draw[guida] plot[domain=0.45:2.55,samples=50,smooth]
        ({\x},{2.275+0.075*sin(\x*360)});
  \draw[sep] (3,2.05) -- (3,2.50);
  \node at (4.5,2.275) {$\alpha$};
  \node at (6.5,2.275)  {$y$};
  \node[right,inner sep=3pt] at (20,2.275) {$T$};
  % punto di confronto corrente sul testo
  \draw[->,line width=0.6pt] (6.5,3.12) -- (6.5,2.68);

  % ===== mismatch ======================================================
  \node[fill=white,inner sep=1.5pt] at (6.5,1.725) {$\neq$};

  % ===== RIGA 2 : il pattern P nell'allineamento corrente ===============
  \draw[barra] (1,0.95) rectangle (12,1.40);
  % la porzione di P che lo scorrimento scarta (lunga i-|alpha|), tratteggiata
  \begin{scope}
    \clip (1,0.95) rectangle (3,1.40);
    \foreach \i in {-2,...,13}{%
      \draw[guida,line width=0.4pt]
        ({1+\i*0.30},0.90) -- ({1+\i*0.30+0.42},1.45);
    }
  \end{scope}
  \draw[sep] (3,0.95) -- (3,1.40);
  \node at (4.5,1.175) {$\alpha$};
  \node at (6.5,1.175)  {$x$};
  \node[right,inner sep=3pt] at (12,1.175) {$P$};

  % ===== RIGA 3 : il pattern P dopo lo spostamento ======================
  \draw[barra] (3,-0.35) rectangle (14,0.10);
  \node at (4.5,-0.125) {$\alpha$};
  \node at (6.5,-0.125)  {$y$};
  \node[right,inner sep=3pt] at (14,-0.125) {$P$};
  % il confronto riprende dalla stessa posizione del testo
  \draw[->,line width=0.6pt] (6.5,-1.00) -- (6.5,-0.53);

  % ===== annotazioni dello spostamento =================================
  % il prefisso alpha di P: lungo |alpha|, oltrepassa la zona tratteggiata ...
  \draw[guida,->] (1.05,0.78) to[bend right=18] (3.95,0.78);
  \draw[guida,->] (2.00,-0.02) -- (2.00,0.55);
  \node[right,inner sep=1pt] at (2.10,-0.02) {$\alpha$};
  % ... e scivola di i-|alpha| fino a sovrapporsi all'occorrenza di alpha in T
  \draw[guida,dashed,->] (1.00,-0.58) to[out=-35,in=205] (2.92,-0.30);
  % dopo lo shift il tratto alpha e' gia' verificato
  \draw[guida,->] (3.10,-0.58) to[bend right=15] (5.90,-0.58);

\end{tikzpicture}

  \caption{Regola di spostamento di KMP. Dopo aver riconosciuto $\pre{P}{i}$ il
  confronto fallisce fra il carattere $y$ del testo e $x=P[i+1]$: il pattern
  scivola di $i-\abs{\alpha}$ posizioni, fino a sovrapporre il proprio prefisso
  $\alpha$ all'occorrenza di $\alpha$ che in $T$ termina appena prima del
  mismatch. Mi sposto cioè in base al più lungo prefisso-suffisso.}
  \label{fig:kmp-shift}
\end{figure}

\paragraph{Schema dell'algoritmo.}
\begin{enumerate}
  \item \emph{Pre-processing del pattern} (che di solito è molto più piccolo del
    testo), che mi consenta, per ogni prefisso $\pre{P}{i}$, di determinare il più
    lungo prefisso-suffisso proprio, cioè
    \[
      \max\bigl\{\, k \;\bigm|\; 0\le k<i,\ \pre{P}{k} = \sub{P}{i-k+1}{i} \,\bigr\}.
    \]
    (\emph{Come farlo?} È la domanda a cui risponde la sezione seguente.)
  \item Uso il pre-processing del punto~1 per eseguire una scansione di $T$ in
    cui, ad ogni passo, posso
    \begin{itemize}
      \item o avanzare il punto di confronto,
      \item oppure avanzare il pattern.
    \end{itemize}
\end{enumerate}

\noindent Poiché ad ogni passo almeno una delle due quantità avanza, e nessuna
delle due può superare la lunghezza del testo, il numero totale di passi è
lineare e il costo complessivo è
\[
  \Th(n+m).
\]

\begin{osservazione}
  È più efficiente, per quanto possa sembrare contro\-intuitivo, preprocessare
  il \emph{testo} piuttosto che il pattern: così il tempo di ricerca di pattern
  diversi sullo stesso testo è proporzionale alla lunghezza del pattern, e non a
  quella del testo. È l'idea che sta dietro ai suffix tree e ai suffix array.
\end{osservazione}

\section{Il pre-processing del pattern: i valori
  \texorpdfstring{$\spv{i}$}{sp\_i}}
\label{sec:kmp-preprocessing}

\scan{4}
Il pre-processing di $P$ consiste nel calcolare, per ogni prefisso di $P$, il più
lungo $\alpha$ che ne sia contemporaneamente prefisso e suffisso \emph{proprio}.
Conviene enunciarlo per una stringa generica $S$ --- che nell'applicazione a KMP
sarà il pattern, cosicché $\abs{S}=n$ --- scrivendo $S_i \defeq \pre{S}{i}$ per
il suo prefisso $i$-esimo.

\begin{definizione}[Prefisso-suffisso massimo]\label{def:sp-i}
  Sia $S$ una stringa e sia $1\le i\le\abs{S}$. Poniamo
  \[
    \spv{i}(S) \defeq \max\set{\,k \;:\; 0\le k<i,\ \pre{S}{k}=\sub{S}{i-k+1}{i}\,} ,
  \]
  cioè $\spv{i}(S)$ è la lunghezza del massimo prefisso \emph{proprio} di $S_i$
  che ne sia anche suffisso; tale prefisso-suffisso si dice anche \emph{bordo}
  di $S_i$. Il vincolo $k<i$ è essenziale: senza di esso ogni
  prefisso sarebbe banalmente prefisso-suffisso di sé stesso e si avrebbe
  $\spv{i}(S)=i$. Il massimo è sempre ben definito perché $k=0$ (il bordo vuoto
  $\vuota$) soddisfa la condizione; in particolare $\spv{1}(S)=0$.
\end{definizione}

\begin{figure}[H]
  \centering
  % kmp-alfa-S -- il massimo prefisso-suffisso alpha del prefisso S_i (pag. 4).
% alpha compare una volta in testa a S e una volta in coda a S_i, e |alpha|=sp_i(S).
\begin{tikzpicture}[
    x=0.85cm, y=1cm,
    font=\footnotesize,
    linea/.style={line width=0.8pt},
    tacca/.style={gray!70!black, line width=0.5pt},
    graffa/.style={decorate, decoration={brace, amplitude=5pt}, gray!70!black},
  ]

  % --- la stringa S -----------------------------------------------------
  \draw[linea] (0,0) -- (10,0);
  \node[right,inner sep=3pt] at (10,0) {$S$};

  % --- confini fra caratteri --------------------------------------------
  \foreach \x in {0,2.4,4.4,6.8}{\draw[tacca] (\x,-0.18) -- (\x,0.18);}

  % --- le due occorrenze di alpha, della stessa lunghezza ---------------
  \draw[graffa] (0,0.26) -- (2.4,0.26)
        node[midway,above=6pt,black] {$\alpha$};
  \draw[graffa] (4.4,0.26) -- (6.8,0.26)
        node[midway,above=6pt,black] {$\alpha$};

  % --- posizioni ---------------------------------------------------------
  \node[below,inner sep=4pt] at (0,-0.18)   {$1$};
  \node[below,inner sep=4pt] at (2.4,-0.18) {$\spv{i}(S)$};
  \node[below,inner sep=4pt] at (4.4,-0.18) {$i-\spv{i}(S)+1$};
  \node[below,inner sep=4pt] at (6.8,-0.18) {$i$};

  % --- il prefisso S_i ---------------------------------------------------
  \draw[graffa,decoration={mirror}] (0,-0.90) -- (6.8,-0.90)
        node[midway,below=6pt,black] {$S_i = \pre{S}{i}$};

\end{tikzpicture}

  \caption{Per ogni prefisso $S_i=\pre{S}{i}$ si cerca il massimo $\alpha$ che ne
  sia al tempo stesso prefisso e suffisso: $\alpha$ compare una volta in testa a
  $S$ e una volta in coda a $S_i$, terminando in posizione $i$, e
  $\abs{\alpha}=\spv{i}(S)$.}
  \label{fig:kmp-alfa-S}
\end{figure}

\begin{osservazione}
  Il pre-processing prende tempo \emph{lineare}. È l'affermazione che ci
  accingiamo a giustificare, e non è affatto ovvia.
\end{osservazione}

\noindent Per calcolare $\spv{i}(S)$ ci sono due casi. Posto $k=\spv{i-1}(S)$ e
$\alpha=\pre{S}{k}$:
\begin{enumerate}
  \item (\emph{caso base}) se $S[i]=S[k+1]$ posso estendere $\alpha$ di un
    carattere, e
    \[
      \spv{i}(S) \;=\; \spv{i-1}(S)+1 ;
    \]
  \item altrimenti $\spv{i}(S)$ si ottiene ricorsivamente: il candidato
    successivo è $\beta$, il massimo bordo di $\alpha$, cioè
    $\abs{\beta}=\spv{k}(S)$, e si ripete il confronto fra $S[i]$ e
    $S[\abs{\beta}+1]$. Il procedimento si itera lungo la catena dei bordi
    $\alpha,\beta,\dots$ finché un carattere coincide con $S[i]$, e allora si
    ricade nel caso~1, oppure finché la catena si esaurisce nella stringa vuota:
    in tal caso $\spv{i}(S)=1$ se $S[i]=S[1]$, e $\spv{i}(S)=0$ altrimenti.
\end{enumerate}

\begin{figure}[H]
  \centering
  % kmp-bordo-del-bordo -- il passo ricorsivo del pre-processing (pag. 4).
% alpha = massimo bordo di S_{i-1} (|alpha|=k=sp_{i-1}(S)),
% beta  = massimo bordo di alpha. Qui S[i]=a diverso da b=S[k+1]:
% alpha non si estende e si ricade su beta.
\begin{tikzpicture}[
    x=0.80cm, y=1cm,
    font=\footnotesize,
    >={Stealth[length=4pt,width=3pt]},
    linea/.style={line width=0.8pt},
    tacca/.style={gray!70!black, line width=0.5pt},
    segm/.style={line width=0.6pt},
  ]

  % --- la stringa S ------------------------------------------------------
  \draw[linea] (0,0) -- (12.6,0);
  \node[right,inner sep=3pt] at (12.6,0) {$S$};
  \draw[tacca] (0,-0.16) -- (0,0.16);

  % --- le due celle di un carattere ---------------------------------------
  % b = S[k+1] segue la copia-prefisso di alpha
  \draw[segm] (4,0) rectangle (5,0.42);
  \node at (4.5,0.21) {$b$};
  % a = S[i] segue la copia-suffisso di alpha
  \draw[segm] (10,0) rectangle (11,0.42);
  \node at (10.5,0.21) {$a$};

  % --- livello 1 : le due occorrenze di alpha ------------------------------
  \foreach \s/\e in {0/4, 6/10}{%
    \draw[segm] (\s,0.62) -- (\e,0.62);
    \draw[segm] (\s,0.52) -- (\s,0.72);
    \draw[segm] (\e,0.52) -- (\e,0.72);
  }
  \node[above,inner sep=2pt] at (3.0,0.66) {$\alpha$};
  \node[above,inner sep=2pt] at (7.0,0.66) {$\alpha$};

  % --- livello 2 : beta, bordo di alpha ------------------------------------
  % beta e' prefisso di alpha (copia sinistra) e suffisso di alpha (copia destra)
  \foreach \s/\e in {0/2, 8/10}{%
    \draw[segm] (\s,1.28) -- (\e,1.28);
    \draw[segm] (\s,1.18) -- (\s,1.38);
    \draw[segm] (\e,1.18) -- (\e,1.38);
  }
  \node[above,inner sep=2pt] at (1.0,1.32) {$\beta$};
  \node[above,inner sep=2pt] at (9.0,1.32) {$\beta$};

  % il confronto del passo ricorsivo riprende dal carattere S[i]
  \draw[->,line width=0.6pt] (10.5,1.10) -- (10.5,0.52);

  % --- posizioni ----------------------------------------------------------
  \foreach \x in {4,5,10,11}{\draw[tacca] (\x,-0.16) -- (\x,0);}
  \node[below,inner sep=4pt] at (3.6,-0.10)  {$k$};
  \node[below,inner sep=4pt] at (4.5,-0.10)  {$k+1$};
  \node[below,inner sep=4pt] at (9.5,-0.10)  {$i-1$};
  \node[below,inner sep=4pt] at (10.5,-0.10) {$i$};

\end{tikzpicture}

  \caption{Il passo ricorsivo: $\alpha$ è il massimo prefisso-suffisso di
  $S_{i-1}$ e $\beta$ è il massimo prefisso-suffisso di $\alpha$. Nel disegno
  $S[i]=a \ne b=S[k+1]$, quindi $\alpha$ non si estende e si ricade su $\beta$.}
  \label{fig:kmp-bordo-del-bordo}
\end{figure}

\noindent Quante chiamate ricorsive innesca il caso~2)?

\begin{osservazione}[Problema]
  Il caso~2) può innescare fino a $i$ chiamate ricorsive: per
  $i\in\set{1,\dots,n}$ l'ordine del caso pessimo farebbe concludere che il
  pre-processing è \emph{quadratico}!
\end{osservazione}

\noindent Con l'analisi ammortizzata si vede invece che la somma delle operazioni
su \emph{tutte} le chiamate ricorsive, per $i\in\set{1,\dots,n}$, è lineare.

\subsection{Analisi ammortizzata}
\label{sec:analisi-ammortizzata}

\scan{5}
Lo schema è il seguente: ogni volta che applico il caso~1) --- quello in cui il
massimo prefisso-suffisso si estende di un carattere --- se il costo effettivo è
$k$ ne \emph{alloco} $2k$. Raddoppiando l'addebito metto da parte un credito con
cui «pago» in anticipo ogni chiamata ricorsiva del caso~2) che eseguirò per
quella posizione.

\begin{osservazione}
  Detto con una funzione potenziale: sia $\Phi$ la lunghezza del bordo corrente,
  cioè il valore $\spv{i-1}(S)$ da cui la ricorsione riparte. Il caso~1) fa
  crescere $\Phi$ di $1$, mentre ogni chiamata ricorsiva del caso~2) esegue
  $\Phi \leftarrow \spv{\Phi}(S)$ e quindi la fa \emph{decrescere strettamente}.
  Poiché $\Phi\ge 0$ e il caso~1) viene applicato al più $n$ volte in tutto, il
  numero complessivo di chiamate ricorsive del caso~2) su tutti gli
  $i\in\set{1,\dots,n}$ è a sua volta al più $n$: il credito accantonato basta a
  coprirle tutte. Il pre-processing costa dunque $\Th(n)$, e non il $\Oh(n^{2})$
  che si otterrebbe maggiorando il caso pessimo di ogni singolo passo.
\end{osservazione}

\begin{esempio}[analisi ammortizzata su uno stack]
  Sullo stack consideriamo due operazioni:
  \begin{itemize}
    \item \textsc{Push}: inserisco un elemento nello stack --- costo $\Oh(1)$;
    \item \textsc{Multi-Pop}: rimuovo tutti gli elementi sovrastanti nello
      stack --- costo proporzionale al numero di elementi rimossi, e quindi
      lineare nell'altezza dello stack.
  \end{itemize}
  Quanto costa una sequenza arbitraria di $t$ operazioni di \textsc{Push} e
  \textsc{Multi-Pop}? $\Oh(t^{2})$?

  No! Prima di poter rimuovere un elemento devo prima averlo inserito: ogni
  elemento entra nello stack una volta sola ed esce al più una volta. Il costo
  complessivo di tutta la sequenza è quindi al più $2t$, dunque lineare.
\end{esempio}

\begin{figure}[H]
  \centering
  % stack-push-multipop -- le due operazioni sullo stack (quaderno, pag. 5).
% A sinistra PUSH: una sola cella in cima al contenuto.
% A destra MULTI-POP: tutto il blocco degli elementi sovrastanti, in un colpo solo.
\begin{tikzpicture}[
    font=\footnotesize,
    >={Stealth[length=4.5pt,width=3.5pt]},
    parete/.style={line width=0.9pt},
    cella/.style={line width=0.5pt},
    guida/.style={gray!60!black, line width=0.6pt},
  ]

  % ============ STACK SINISTRO : PUSH ==================================
  % contenitore aperto in alto
  \draw[parete] (0,3.2) -- (0,0) -- (1.6,0) -- (1.6,3.2);
  % contenuto gia' presente: 4 celle
  \foreach \y in {0.4,0.8,1.2,1.6}{\draw[cella] (0,\y) -- (1.6,\y);}
  % l'elemento appena inserito
  \begin{scope}
    \clip (0,1.6) rectangle (1.6,2.0);
    \foreach \i in {-2,...,9}{%
      \draw[guida,line width=0.4pt] ({\i*0.28},1.55) -- ({\i*0.28+0.5},2.05);
    }
  \end{scope}
  \draw[cella] (0,2.0) -- (1.6,2.0);
  % freccia di inserimento
  \draw[->,line width=0.7pt] (-0.95,1.8) -- (-0.10,1.8);
  \node[left,inner sep=3pt] at (-0.95,1.8) {\textsc{Push}};

  % ============ STACK DESTRO : MULTI-POP ===============================
  \begin{scope}[xshift=4.0cm]
    \draw[parete] (0,3.2) -- (0,0) -- (1.6,0) -- (1.6,3.2);
    % il blocco che viene rimosso tutto insieme
    \begin{scope}
      \clip (0,0.8) rectangle (1.6,2.0);
      \foreach \i in {-4,...,12}{%
        \draw[guida,line width=0.4pt] ({\i*0.28},0.75) -- ({\i*0.28+1.3},2.05);
      }
    \end{scope}
    % celle del contenuto
    \foreach \y in {0.4,0.8,1.2,1.6,2.0}{\draw[cella] (0,\y) -- (1.6,\y);}
    % freccia di rimozione
    \draw[->,line width=0.7pt] (-1.85,-0.20) to[out=25,in=200] (-0.05,0.80);
    \node[below,inner sep=2pt] at (-1.05,-0.45) {\textsc{Multi-Pop}};
  \end{scope}

\end{tikzpicture}

  \caption{Le due operazioni sullo stack. \textsc{Push} deposita un solo elemento
  in cima al contenuto già presente (costo costante); \textsc{Multi-Pop} rimuove
  in un colpo solo tutto il blocco degli elementi sovrastanti (costo
  proporzionale alla sua altezza).}
  \label{fig:stack-push-multipop}
\end{figure}

\noindent Il confronto fra i due modi di contare è riassunto qui sotto.

\begin{center}
\begin{tabular}{@{}p{0.44\linewidth}p{0.46\linewidth}@{}}
\toprule
\textbf{analisi semplice} & \textbf{analisi ammortizzata}\\
\midrule
singolo \textsc{Push}: $k$        & singolo \textsc{Push}: $2k$\\[4pt]
singolo \textsc{Multi-Pop}: $k\cdot i$ & singolo \textsc{Multi-Pop}: $0$\\[6pt]
\small dove $i$ è il numero di elementi sopra l'elemento da eliminare
& \small (ho già pagato durante il \textsc{Push} degli elementi in alto)\\
\bottomrule
\end{tabular}
\end{center}

\noindent L'analisi semplice conclude che il costo è quadratico soltanto perché
maggiora $i$ con il suo caso pessimo $i\in\Oh(t)$, addebitando a \emph{ogni}
\textsc{Multi-Pop} il prezzo pieno; l'analisi ammortizzata fa pagare la rimozione
già all'inserimento e ottiene la maggiorazione lineare $2t$. È esattamente il
rapporto fra il caso~1) e il caso~2) del pre-processing di KMP.

\section{La funzione \texorpdfstring{$Z$}{Z}}
\label{sec:funzione-z}

\scan{5}
Siano $P,T\in\alfa^{*}$: voglio tutte le occorrenze esatte di $P$ in $T$. Invece
di lavorare direttamente sulla coppia pattern/testo, introduciamo un approccio
alternativo, basato su una funzione definita su una stringa generica.

\begin{definizione}[Funzione $Z$]\label{def:funzione-z}
  Sia $S\in\alfa^{*}$ e sia $2\le i\le\abs{S}$. Si pone
  \[
    \Zv{i}(S)\;\defeq\;\max\set{\,j \ge 0 \;:\; \pre{S}{j}=\sub{S}{i}{i+j-1}\,} .
  \]
  L'insieme contiene sempre $j=0$ (la stringa vuota è prefisso di qualunque
  stringa), quindi il massimo è ben definito e $\Zv{i}(S)\ge 0$: $\Zv{i}(S)$ è la
  lunghezza del più lungo prefisso di $S$ che ricorre in $S$ anche a partire
  dalla posizione $i$. Se $\Zv{i}(S)>0$, il fattore
  $\sub{S}{i}{i+\Zv{i}(S)-1}$ si dice \Zbox di inizio $i$.
\end{definizione}

\begin{figure}[H]
  \centering
  % z-definizione -- Z_i(S)=j: il prefisso alpha=S[1..j] ricompare in posizione i.
\begin{tikzpicture}[
    x=0.90cm, y=1cm,
    font=\footnotesize,
    tacca/.style={gray!80!black, line width=0.5pt},
    linea/.style={line width=0.7pt},
  ]

  % --- la stringa S ---------------------------------------------------
  \draw[linea] (0,0) -- (11.2,0);
  \node[right, inner sep=2pt] at (11.35,0) {$S$};

  % --- tacche: 1, j, i, i+j-1  (|1..j| = |i..i+j-1| = j) --------------
  \foreach \x in {0,3,5.5,8.5}{%
    \draw[tacca] (\x,-0.22) -- (\x,0.22);
  }

  % --- etichette di posizione -----------------------------------------
  \node[below, inner sep=3pt] at (3,-0.22)   {$j$};
  \node[below, inner sep=3pt] at (5.5,-0.22) {$i$};
  \node[below, inner sep=3pt] at (8.5,-0.22) {$i+j-1$};

  % --- graffe: le due occorrenze di alpha ------------------------------
  \draw[decorate, decoration={brace, amplitude=5pt, mirror}, gray!70!black]
    (0,-0.90) -- (3,-0.90)
    node[midway, below=6pt, black] {$\alpha$};
  \draw[decorate, decoration={brace, amplitude=5pt, mirror}, gray!70!black]
    (5.5,-0.90) -- (8.5,-0.90)
    node[midway, below=6pt, black] {$\alpha$};

  % --- Z_i(S) misura la lunghezza del prefisso alpha -------------------
  \node (zeta) at (6.35,-1.95) {$\Zv{i}(S)$};
  \draw[-{Stealth[length=4pt]}, gray!70!black]
    (zeta.west) to[out=175, in=-15] (3.10,-1.30);

\end{tikzpicture}

  \caption{$\Zv{i}(S)=j$: il prefisso $\alpha=\pre{S}{j}$ di $S$ ricompare
  identico a partire dalla posizione $i$, e $j$ è la lunghezza massima per cui
  ciò accade.}
  \label{fig:z-definizione}
\end{figure}

\begin{osservazione}
  La definizione si trova spesso scritta a due rami: $\Zv{i}(S)=0$ se
  $S[i]\neq S[1]$, e il massimo di cui sopra altrimenti. È la stessa cosa: la
  prima riga è semplicemente il caso in cui l'unico $j$ ammissibile è $j=0$. La
  posizione $i=1$ va invece esclusa a mano, come si è fatto: la definizione
  darebbe $\Zv{1}(S)=\abs{S}$, un valore privo di informazione.
\end{osservazione}

\subsection{Riduzione dell'exact matching alla funzione
  \texorpdfstring{$Z$}{Z}}
\label{sec:z-riduzione}

\scan{6}
Avendo un algoritmo per $Z$ siamo a posto. Per usare la funzione $Z$ come
algoritmo di ricerca c'è bisogno di avere $\Zv{i}(A)$ per ogni
$i\in\set{2,\dots,\abs{A}}$, e la scelta giusta della stringa $A$ è la
concatenazione di pattern e testo separati da un carattere fresco:
\[
  A \;=\; P\,\$\,T, \qquad \$ \notin \alfa .
\]
Il separatore non compare né in $P$ né in $T$, quindi nessuna copia del prefisso
di $A$ può oltrepassarlo e vale $\Zv{i}(A)\le\abs{P}$ per ogni $i>\abs{P}+1$.
Posto $j=\Zv{i}(A)$ si ha allora
\[
  j = \abs{P}
  \quad\Longleftrightarrow\quad
  P \ \text{occorre in } T \text{ nella posizione corrispondente a } i ,
\]
dove $i$ è contata su $A$: il confronto ha senso solo per $i \ge \abs{P}+2$ e in
tal caso l'occorrenza comincia in $T$ alla posizione $i-\abs{P}-1$.

Resta dunque una sola domanda: \emph{come calcolo $Z$?}

\subsection{Calcolo di \texorpdfstring{$Z$}{Z}}
\label{sec:z-calcolo}

\scan{6}
La risposta è: \emph{per ricorsione}, con un'analisi ammortizzata della
complessità. L'idea è di mantenere --- inizializzandolo e aggiornandolo a ogni
passo --- un intervallo $[\ell,r]$ tale che $\sub{S}{\ell}{r}$ sia un prefisso di
$S$, con $r$ \emph{massimo} fra tutti quelli calcolati fino a quel momento.

\begin{figure}[H]
  \centering
  % z-blocco-lr -- l'invariante: il blocco [l,r] con S[l..r] prefisso di S e i dentro.
\begin{tikzpicture}[
    x=1.00cm, y=1cm,
    font=\footnotesize,
    tacca/.style={gray!80!black, line width=0.5pt},
    linea/.style={line width=0.7pt},
  ]

  % --- la stringa S ---------------------------------------------------
  \draw[linea] (0,0) -- (10,0);
  \node[right, inner sep=2pt] at (10.15,0) {$S$};

  % --- tacche: 1, ell, i, r -------------------------------------------
  \foreach \x in {0,5.0,5.6,6.9}{%
    \draw[tacca] (\x,-0.22) -- (\x,0.22);
  }

  % --- il blocco [ell,r], evidenziato da un tratto sopra la retta ------
  \draw[gray!70!black, line width=0.5pt]
    (5.0,0.26) -- (5.0,0.46) -- (6.9,0.46) -- (6.9,0.26);
  \node[above, inner sep=2pt] at (5.0,0.46) {$\ell$};
  \node[above, inner sep=2pt] at (6.9,0.46) {$r$};

  % --- l'indice corrente, dentro il blocco ----------------------------
  \node[below, inner sep=3pt] at (5.6,-0.22) {$i$};

\end{tikzpicture}

  \caption{L'invariante dell'algoritmo: il blocco $[\ell,r]$, con
  $\sub{S}{\ell}{r}$ prefisso di $S$ e $r$ massimo fra quelli già calcolati;
  l'indice corrente $i$ cade dentro il blocco.}
  \label{fig:z-blocco-lr}
\end{figure}

\scan{7}
Scandiamo $i=2,3,\dots,\abs{S}$ e supponiamo di aver già calcolato
$\Zv{2},\dots,\Zv{i-1}$. Se $r<i$ il blocco non dà alcuna informazione sulla
posizione $i$: si calcola $\Zv{i}$ per confronti diretti, allineando $\suf{S}{i}$
con $S$, e se $\Zv{i}>0$ si pone $[\ell,r]\leftarrow[i,\,i+\Zv{i}-1]$. Il caso
interessante è quello in cui $\ell \le i \le r$. Poniamo allora
\[
  \alpha \;=\; \sub{S}{\ell}{r}\;=\;\pre{S}{r-\ell+1},
  \qquad \abs{\alpha}=r-\ell+1=\Zv{\ell},
\]
e sia
\[
  i' \;\defeq\; i-\ell+1
\]
la posizione del prefisso che corrisponde a $i$ dentro il blocco. Le due
occorrenze di $\alpha$ --- quella iniziale, che è un prefisso di $S$, e quella nel
blocco --- sono copie l'una dell'altra: è questa corrispondenza a \emph{suggerire}
il valore di $\Zv{i}$, perché $\Zv{i'}$ è già stato calcolato. Chiamiamo
\[
  \beta\;=\;\sub{S}{i'}{\,i'+\Zv{i'}-1},
  \qquad \abs{\beta}=\Zv{i'},
\]
la \Zbox che comincia in $i'$. Poiché il blocco $\sub{S}{\ell}{r}$ è una copia del
prefisso, la parte di $\beta$ contenuta in $\alpha$ ricompare a partire dalla
posizione $i$; l'intera $\beta$ soltanto quando $\Zv{i'}\le r-i+1$, cioè quando
$\beta$ non oltrepassa la fine di $\alpha$. Indichiamo infine con
\[
  x \;=\; S[\abs{\alpha}+1], \qquad y \;=\; S[r+1]
\]
i due caratteri che seguono le due occorrenze di $\alpha$: la massimalità di
$\Zv{\ell}$ dice esattamente che $x\neq y$.

\begin{figure}[H]
  \centering
  % z-alfa-beta -- il prefisso alpha e la sua copia S[l..r]; i' corrisponde a i.
\begin{tikzpicture}[
    x=0.78cm, y=1cm,
    font=\footnotesize,
    tacca/.style={gray!80!black, line width=0.5pt},
    linea/.style={line width=0.7pt},
    graffa/.style={decorate, decoration={brace, amplitude=4pt}, gray!70!black},
    carattere/.style={fill=white, inner sep=1pt},
  ]

  % --- la stringa S ---------------------------------------------------
  \draw[linea] (0,0) -- (14,0);
  \node[right, inner sep=2pt] at (14.15,0) {$S$};

  % --- tacche -----------------------------------------------------------
  % gruppo sinistro: 1, i', fine di alpha, fine della cella di x
  % gruppo destro:   ell, i, r,            fine della cella di y
  \foreach \x in {0,1.2,3.2,3.8,7.0,8.2,10.2,10.8}{%
    \draw[tacca] (\x,-0.20) -- (\x,0.20);
  }

  % --- i caratteri che seguono le due copie di alpha --------------------
  \node[carattere] at (3.5,0)  {$x$};
  \node[carattere] at (10.5,0) {$y$};

  % --- posizioni --------------------------------------------------------
  \node[above, inner sep=2pt] at (7.0,0.18)  {$\ell$};
  \node[above, inner sep=2pt] at (10.2,0.18) {$r$};
  \node[below, inner sep=3pt] at (1.2,-0.20) {$i'=i-\ell+1$};
  \node[below, inner sep=3pt] at (8.2,-0.20) {$i$};

  % --- graffe beta (la Z-box che comincia in i', e la sua copia in i) ---
  \draw[graffa] (2.6,0.95) -- (1.2,0.95) node[midway, above=4pt, black] {$\beta$};
  \draw[graffa] (9.6,0.95) -- (8.2,0.95) node[midway, above=4pt, black] {$\beta$};

  % --- graffe alpha (prefisso di S, e la copia dentro il blocco) --------
  \draw[graffa] (3.2,1.85)  -- (0,1.85)   node[midway, above=4pt, black] {$\alpha$};
  \draw[graffa] (10.2,1.85) -- (7.0,1.85) node[midway, above=4pt, black] {$\alpha$};

  % --- i' corrisponde a i: le due copie coincidono ----------------------
  \draw[gray!70!black, line width=0.5pt]
    (1.2,2.40) to[out=65, in=180] (4.35,3.00);
  \node at (4.70,3.00) {$=$};
  \draw[gray!70!black, line width=0.5pt]
    (5.05,3.00) to[out=0, in=115] (8.2,2.40);

  % --- massimalita' di Z_ell: x e y sono diversi ------------------------
  \draw[gray!60, line width=0.5pt, dashed]
    (3.5,-0.30) .. controls (5.6,-1.40) and (8.4,-1.40) .. (10.5,-0.30);
  \node[fill=white, inner sep=1.5pt] at (7.0,-1.07) {$\neq$};

\end{tikzpicture}

  \caption{Il prefisso $\alpha=\pre{S}{r-\ell+1}$ e la sua copia
  $\sub{S}{\ell}{r}$. La posizione $i'=i-\ell+1$ della copia di sinistra
  corrisponde alla posizione $i$ della copia di destra: il valore già noto
  $\Zv{i'}$ suggerisce $\Zv{i}$. I caratteri $x$ e $y$ che seguono le due copie
  sono diversi, per la massimalità di $\Zv{\ell}$.}
  \label{fig:z-alfa-beta}
\end{figure}

\noindent Si presentano allora due casi, a seconda di dove termina $\beta$
rispetto alla fine di $\alpha$.

\begin{description}
  \item[Caso 1 \textnormal{($\Zv{i'}<r-i+1$).}] La \Zbox $\beta$ termina
    strettamente prima della fine di $\alpha$: tutto ciò che serve sta dentro la
    copia già nota e si ha subito il valore cercato,
    \[
      \Zv{i}\;=\;\Zv{i'} .
    \]
    L'intervallo $[\ell,r]$ resta invariato.
  \item[Caso 2 \textnormal{($\Zv{i'}\ge r-i+1$).}] $\beta$ arriva almeno fino
    alla fine di $\alpha$: la copia garantisce soltanto che $\sub{S}{i}{r}$
    coincide con $\pre{S}{r-i+1}$, mentre di ciò che accade oltre $r$ non si sa
    nulla. Si riprendono allora i confronti espliciti fra le posizioni $r+1$ e
    $r-i+2$ (i primi $r-i+1$ caratteri sono già garantiti dal blocco), e
    \[
      \Zv{i}\;=\;(r-i+1)+q,
      \qquad
      q=\max\set{\,j\ge 0 \;:\; \sub{S}{r+1}{r+j}=\sub{S}{r-i+2}{\,r-i+1+j}\,} .
    \]
    Si aggiorna poi $[\ell,r]\leftarrow[i,\,r+q]$.
\end{description}

\begin{figure}[H]
  \centering
  % z-caso1 -- Caso 1: beta finisce strettamente prima della fine di alpha.
\begin{tikzpicture}[
    x=0.78cm, y=1cm,
    font=\footnotesize,
    tacca/.style={gray!80!black, line width=0.5pt},
    linea/.style={line width=0.7pt},
    graffasu/.style={decorate, decoration={brace, amplitude=4pt}, gray!70!black},
    graffagiu/.style={decorate, decoration={brace, amplitude=4pt, mirror}, gray!70!black},
    carattere/.style={fill=white, inner sep=1pt},
  ]

  % --- titolo -----------------------------------------------------------
  \node[right, inner sep=0pt] at (0,1.75) {\small Caso 1};

  % --- la stringa S -----------------------------------------------------
  \draw[linea] (0,0) -- (14,0);
  \node[right, inner sep=2pt] at (14.15,0) {$S$};

  % --- tacche -----------------------------------------------------------
  % gruppo sinistro: 1, i', fine di beta, fine di alpha, fine della cella di x
  % gruppo destro:   ell, i, fine di beta, r,            fine della cella di y
  \foreach \x in {0,1.1,2.3,3.0,3.6,8.0,9.1,10.3,11.0,11.6}{%
    \draw[tacca] (\x,-0.20) -- (\x,0.20);
  }

  % --- caratteri che seguono le due copie di alpha ----------------------
  \node[carattere] at (3.3,0)  {$x$};
  \node[carattere] at (11.3,0) {$y$};

  % --- posizioni --------------------------------------------------------
  \node[below, inner sep=3pt] at (0,-0.20)    {$1$};
  \node[below, inner sep=3pt] at (1.1,-0.20)  {$i'$};
  \node[below, inner sep=3pt] at (8.0,-0.20)  {$\ell$};
  \node[below, inner sep=3pt] at (9.1,-0.20)  {$i$};
  \node[below, inner sep=3pt] at (11.0,-0.20) {$r$};

  % --- graffe alpha, sopra la retta -------------------------------------
  \draw[graffasu] (3.0,0.60) -- (0,0.60)
    node[midway, above=4pt, black] {$\alpha$};
  \draw[graffasu] (11.0,0.60) -- (8.0,0.60)
    node[midway, above=4pt, black] {$\alpha$};

  % --- graffe beta, sotto la retta (strettamente dentro alpha) ----------
  \draw[graffagiu] (1.1,-0.80) -- (2.3,-0.80)
    node[midway, below=4pt, black] {$\beta$};
  \draw[graffagiu] (9.1,-0.80) -- (10.3,-0.80)
    node[midway, below=4pt, black] {$\beta$};

\end{tikzpicture}

  \caption{Caso 1: $\abs{\beta}=\Zv{i'}<r-i+1$. La \Zbox $\beta$ è strettamente
  contenuta in $\alpha$, quindi si copia tale e quale in posizione $i$:
  $\Zv{i}=\Zv{i'}$.}
  \label{fig:z-caso1}
\end{figure}

\begin{figure}[H]
  \centering
  % Caso 2 dell'algoritmo Z: |beta| = Z_{i'} >= r-i+1.
% Copia sinistra: beta compare sia come prefisso di S sia a partire da i',
% dove raggiunge esattamente la fine di alpha.  Copia destra: beta arriva
% fino a r, oltre il quale nulla e' garantito -> confronto a =? y.
\begin{tikzpicture}[x=0.78cm,y=1cm,font=\footnotesize,
                    >={Stealth[length=4pt,width=3pt]}]

  % ---------------------------------------------------------------- la stringa S
  \draw[thick,->] (0,0) -- (13.5,0);
  \node[anchor=west,inner sep=2pt] at (13.5,0) {$S$};
  \foreach \x in {0,1.3,2.0,2.6,3.3,3.9,8.0,9.3,11.3,11.9}
    \draw[thick] (\x,-0.13) -- (\x,0.13);

  % caratteri notevoli, scritti sopra la retta nella loro cella
  \node[anchor=south,inner sep=1pt] at (2.3,0.12)  {$a$};
  \node[anchor=south,inner sep=1pt] at (3.6,0.12)  {$x$};
  \node[anchor=south,inner sep=1pt] at (11.6,0.12) {$y$};

  % etichette di posizione
  \node[anchor=north west,inner sep=2pt] at (0,-0.14)    {\scriptsize $1$};
  \node[anchor=north west,inner sep=2pt] at (1.3,-0.14)  {\scriptsize $i'$};
  \node[anchor=north west,inner sep=2pt] at (8.0,-0.14)  {\scriptsize $\ell$};
  \node[anchor=north west,inner sep=2pt] at (9.3,-0.14)  {\scriptsize $i$};
  \node[anchor=north east,inner sep=2pt] at (11.3,-0.14) {\scriptsize $r$};

  % ------------------------------------------------- gruppo sinistro: prefisso
  % beta come prefisso di S (livello alto)
  \draw[gray!60!black] (0,1.16) -- (2.0,1.16);
  \draw[gray!60!black] (0,1.05) -- (0,1.27);
  \draw[gray!60!black] (2.0,1.05) -- (2.0,1.27);
  \node[anchor=south,inner sep=1pt] at (0.6,1.20) {$\beta$};

  % beta = Z-box che comincia in i' e finisce alla fine di alpha (livello basso)
  \draw[gray!60!black] (1.3,0.62) -- (3.3,0.62);
  \draw[gray!60!black] (1.3,0.51) -- (1.3,0.73);
  \draw[gray!60!black] (3.3,0.51) -- (3.3,0.73);
  \node[anchor=south,inner sep=1pt] at (2.9,0.66) {$\beta$};

  % graffa: alpha = prefisso S[1..r-l+1]
  \draw[decorate,decoration={brace,amplitude=4pt}] (0,1.70) -- (3.3,1.70);
  \node[anchor=south,inner sep=2pt] at (1.65,1.86) {$\alpha$};

  % --------------------------------------------------- gruppo destro: il blocco
  % beta copiata a partire da i; oltre r non si sa nulla (tratteggio)
  \draw[gray!60!black] (9.3,0.62) -- (11.3,0.62);
  \draw[gray!60!black] (9.3,0.51) -- (9.3,0.73);
  \draw[gray!60!black] (11.3,0.51) -- (11.3,0.73);
  \draw[gray!60!black,dashed,->] (11.35,0.62) -- (12.15,0.62);
  \node[anchor=south,inner sep=1pt] at (10.3,0.66) {$\beta$};

  % graffa: alpha = S[l..r]
  \draw[decorate,decoration={brace,amplitude=4pt}] (8.0,1.70) -- (11.3,1.70);
  \node[anchor=south,inner sep=2pt] at (9.65,1.86) {$\alpha$};

  % ------------------------------------------------------ confronto da eseguire
  \draw[gray!60!black,thin,<->,rounded corners=3pt]
        (2.3,-0.17) -- (2.3,-0.88) -- (11.6,-0.88) -- (11.6,-0.17);
  \node[anchor=north,inner sep=3pt] at (6.95,-0.90)
        {$a\overset{?}{=}y$};

  % ------------------------------------------------------------------- formula
  \node[anchor=north west,inner sep=2pt] at (0,-1.62)
        {$\abs{\beta}=\Zv{i'}\ge r-i+1,\qquad i'=i-\ell+1$};

\end{tikzpicture}

  \caption{Caso 2: $\abs{\beta}=\Zv{i'}\ge r-i+1$. Nella copia di sinistra
  $\beta$ compare sia come prefisso di $S$ sia a partire da $i'$, dove raggiunge
  la fine di $\alpha$; nella copia di destra $\beta$ arriva fino a $r$, ma se
  prosegua oltre è incognito ($a\overset{?}{=}y$).}
  \label{fig:z-caso2}
\end{figure}

\begin{osservazione}
  Nel caso disegnato vale $\Zv{i'}=r-i+1$ esattamente: $\beta$ è il suffisso di
  $\alpha$ che comincia in $i'$, la posizione da cui riprendere i confronti è
  $\abs{\beta}+1=\Zv{i'}+1$ e il carattere da confrontare con $y=S[r+1]$ è
  $a=S[\abs{\beta}+1]$; è il confronto marcato con «\,?\,» nella figura. Se
  invece $\Zv{i'}>r-i+1$ strettamente, la \Zbox $\beta$ oltrepassa la fine di
  $\alpha$, cioè $S[r-\ell+2]=S[r-i+2]$; il primo confronto mette allora
  $y=S[r+1]$ contro $S[r-i+2]=S[\abs{\alpha}+1]=x$ e fallisce, perché $x\neq y$.
  In tal caso $q=0$ e $\Zv{i}=r-i+1$: il valore corretto \emph{non} è $\Zv{i'}$.
  È per questo che nel caso~2 si riparte da $r-i+1$ e non da $\Zv{i'}$.
\end{osservazione}

\subsection{L'algoritmo e la sua complessità}
\label{sec:z-complessita}

\begin{algorithm}[H]
  \caption{\textsc{Zcalc}$(S)$ --- calcolo di tutti i valori $\Zv{i}(S)$}
  \label{alg:zcalc}
  \begin{algorithmic}[1]
    \Require una stringa $S$ con $\abs{S}\ge 2$
    \Ensure i valori $\Zv{2}(S),\dots,\Zv{\abs{S}}(S)$
    \State $\Zv{2}\gets$ lunghezza del più lungo prefisso comune fra $S$ e
      $\suf{S}{2}$ \Comment{confronti diretti}
    \If{$\Zv{2}>0$}
      \State $\ell\gets 2$, \quad $r\gets \Zv{2}+1$
    \Else
      \State $\ell\gets 0$, \quad $r\gets 0$
    \EndIf
    \For{$i\gets 3$ \textbf{to} $\abs{S}$}
      \If{$i>r$} \Comment{nessun blocco utile: confronti diretti}
        \State $\Zv{i}\gets$ lunghezza del più lungo prefisso comune fra $S$ e
          $\suf{S}{i}$
        \If{$\Zv{i}>0$}
          \State $\ell\gets i$, \quad $r\gets i+\Zv{i}-1$
        \EndIf
      \Else
        \State $i'\gets i-\ell+1$
        \If{$\Zv{i'}<r-i+1$} \Comment{Caso 1}
          \State $\Zv{i}\gets \Zv{i'}$
        \Else \Comment{Caso 2}
          \State $q\gets$ numero di confronti riusciti a partire da $S[r+1]$
            contro $S[r-i+2]$
          \State $\Zv{i}\gets (r-i+1)+q$, \quad $\ell\gets i$, \quad
            $r\gets r+q$
        \EndIf
      \EndIf
    \EndFor
  \end{algorithmic}
\end{algorithm}

\begin{osservazione}[Complessità]
  L'estremo destro $r$ non decresce mai. Ogni confronto esplicito eseguito
  dall'algoritmo o ha successo --- e allora fa crescere $r$ di almeno una
  posizione --- oppure fallisce, e in tal caso chiude l'iterazione corrente. I
  confronti riusciti sono quindi al più $\abs{S}$ in tutto e quelli falliti al
  più uno per iterazione: il calcolo di \emph{tutti} i valori $\Zv{i}(S)$ costa
  $\Th(\abs{S})$. Applicato ad $A=P\,\$\,T$ questo dà $\Th(n+m)$ per l'exact
  string matching, cioè esattamente il lower bound del problema: la
  Domanda~\ref{dom:subquadratico} ha risposta affermativa.
\end{osservazione}

\paragraph{Riferimenti.}
\scan{6}
La trattazione segue la Parte~I di Gusfield; le note del corso sono disponibili
online nel file \texttt{alg2.pdf}.

\section{Dalla funzione \texorpdfstring{$Z$}{Z} ai valori
  \texorpdfstring{$\spv{i}$}{sp\_i}}
\label{sec:z-verso-sp}

\scan{8}
La funzione $Z$ non serve soltanto a risolvere direttamente l'exact string
matching: permette anche di calcolare i valori $\spv{i}(P)$ della
Definizione~\ref{def:sp-i}, cioè i prefissi-suffissi su cui si regge KMP.
Vediamo come.

Ricordiamo il ruolo di $\spv{i}(P)$: serve a riposizionare $P$ durante KMP senza
perdere possibili occorrenze. Dopo aver riconosciuto $\pre{P}{i}$ nel testo, il
pattern viene fatto scorrere di $i-\spv{i}(P)$ posizioni, in modo che il prefisso
$\alpha=\pre{P}{\spv{i}(P)}$ vada a sovrapporsi all'occorrenza di $\alpha$ che
termina in $i$. Massimizzare $\spv{i}(P)$ equivale a minimizzare lo shift, ed è
proprio la massimalità a garantire che nessuna occorrenza venga saltata.

\begin{definizione}[$j$ mappa su $i$]\label{def:mappa-su}
  Dati $P$ e un indice $j$, diciamo che \emph{$j$ mappa su $i$} se e solo se
  \[
    i \;=\; j+\Zv{j}(P)-1 ,
  \]
  cioè se la \Zbox che comincia in $j$ termina esattamente in $i$.
\end{definizione}

Se so calcolare i valori $\Zv{j}$, so allora calcolare anche gli $\spv{i}(P)$.

\begin{osservazione}
  Il legame è il seguente: se $j\ge 2$ mappa su $i$, allora $\Zv{j}(P)=i-j+1$ e
  $\sub{P}{j}{i}=\pre{P}{i-j+1}$, cioè $\Zv{j}(P)$ è la lunghezza di un
  prefisso-suffisso proprio di $\pre{P}{i}$; fra tutti i $j$ che mappano su $i$,
  quello \emph{minimo} è quello che dà il prefisso-suffisso più lungo. Attenzione
  però: il valore così ottenuto è esattamente $\spp{i}(P)$ (definito nella
  sezione seguente) e non $\spv{i}(P)$, perché la massimalità della \Zbox impone
  in più $P[i+1]\ne P[\Zv{j}(P)+1]$. Dai valori $\spp{i}(P)$ si risale poi agli
  $\spv{i}(P)$ con una scansione all'indietro:
  \[
    \spv{\abs{P}}(P)=\spp{\abs{P}}(P),
    \qquad
    \spv{i}(P)=\max\bigl(\spp{i}(P),\ \spv{i+1}(P)-1\bigr)
    \quad\text{per } i=\abs{P}-1,\dots,1 .
  \]
  Un controesempio minimale al criterio ingenuo è $P=\str{aaa}$, per cui
  $\Zv{2}=2$ e $\Zv{3}=1$: il criterio dà $\spp{1}=0$, $\spp{2}=0$, $\spp{3}=2$,
  mentre i valori corretti sono $\spv{1}=0$, $\spv{2}=1$, $\spv{3}=2$.
\end{osservazione}

\section{La variante \texorpdfstring{$\spp{i}$}{sp'\_i}}
\label{sec:sp-primo}

\scan{8}
Passiamo ora al punto successivo. Supponiamo di aver riconosciuto $\pre{P}{i}$
nel testo e di trovare un mismatch alla posizione successiva: $T$ presenta il
carattere $y$, mentre $P$ presenta $x=P[i+1]$, con $x\ne y$. Lo shift dettato da
$\spv{i}(P)$ porta il prefisso $\alpha$ a sovrapporsi all'occorrenza di $\alpha$
che termina in $i$; il carattere di $P$ che va a finire sotto $y$ è quindi
$z=P[\spv{i}(P)+1]$.

\begin{figure}[H]
  \centering
  % KMP: mismatch fra y (in T) e x = P[i+1].  Lo shift di i - sp_i(P) porta il
% prefisso alpha = P[1..sp_i(P)] sotto l'occorrenza di alpha che termina in i:
% il carattere di P che finisce sotto y e' allora z = P[sp_i(P)+1].
\begin{tikzpicture}[x=0.9cm,y=1cm,font=\footnotesize,
                    >={Stealth[length=4pt,width=3pt]}]

  % =============================================================== il testo T
  \draw[thick,->] (-0.6,2) -- (10.9,2);
  \node[anchor=west,inner sep=2pt] at (10.9,2) {$T$};
  \foreach \x in {5.0,7.2,8.0} \draw[thick] (\x,1.87) -- (\x,2.13);
  \node[anchor=south,inner sep=1pt] at (7.6,2.12) {$y$};

  % posizione corrente di confronto nel testo
  \draw[->,thick] (7.6,3.00) -- (7.6,2.56);

  % ============================================================== il pattern P
  \draw[thick] (0,0) -- (9.5,0);
  \foreach \x in {0,2.2,3.0,5.0,7.2,8.0,9.5} \draw[thick] (\x,-0.13) -- (\x,0.13);
  \node[anchor=south,inner sep=1pt] at (2.6,0.12) {$z$};
  \node[anchor=south,inner sep=1pt] at (7.6,0.12) {$x$};

  % le due occorrenze di alpha: prefisso e suffisso di P[1..i]
  \draw[decorate,decoration={brace,amplitude=4pt}] (0,0.45) -- (2.2,0.45);
  \node[anchor=south,inner sep=2pt] at (1.1,0.61) {$\alpha$};
  \draw[decorate,decoration={brace,amplitude=4pt}] (5.0,0.45) -- (7.2,0.45);
  \node[anchor=south,inner sep=2pt] at (6.1,0.61) {$\alpha$};

  % etichette di posizione
  \node[anchor=north west,inner sep=2pt] at (0,-0.14)    {\scriptsize $1$};
  \node[anchor=north east,inner sep=2pt] at (2.2,-0.14)  {\scriptsize $\spv{i}(P)$};
  \node[anchor=north east,inner sep=2pt] at (5.0,-0.14)  {\scriptsize $i-\spv{i}(P)$};
  \node[anchor=north east,inner sep=2pt] at (7.2,-0.14)  {\scriptsize $i$};
  \node[anchor=north,inner sep=2pt]      at (7.6,-0.14)  {\scriptsize $i{+}1$};
  \node[anchor=north east,inner sep=2pt] at (9.5,-0.14)  {\scriptsize $\abs{P}$};

  % ================================================== il confronto che fallisce
  \draw[->,dashed,gray!60!black] (7.6,1.78) -- (7.6,0.52);
  \node[anchor=west,inner sep=3pt] at (7.72,1.15) {$x\neq y$};

  % ============================================== effetto dello shift su z
  \draw[->,gray!60!black]
        (4.30,-1.15) .. controls (3.50,-1.15) and (2.60,-0.95) .. (2.60,-0.40);
  \node[anchor=west,align=left,inner sep=3pt] at (4.40,-1.15)
        {shift di $i-\spv{i}(P)$:\\ $z$ finisce sotto $y$};

\end{tikzpicture}

  \caption{Mismatch fra $y$ (in $T$) e $x=P[i+1]$: dopo lo shift il carattere di
  $P$ allineato a $y$ è $z=P[\spv{i}(P)+1]$. Se $z=x$ il mismatch si ripresenta
  con certezza.}
  \label{fig:kmp-sp-primo}
\end{figure}

Se $z$ non è $y$, dopo lo shift di $P$ avrò di nuovo un mismatch e lo spostamento
sarà stato inutile. Di $y$ non so nulla, ma so che $x\ne y$: quantomeno voglio
allora che $z\ne x$, perché se $z=x$ il mismatch dopo lo shift è garantito. Da
qui la definizione della variante $\spp{i}$.

\begin{definizione}[$\spp{i}(P)$]\label{def:sp-primo}
  Per $1\le i\le\abs{P}$,
  \begin{align*}
    \spp{i}(P)\;\defeq\;\max\bigl\{\,k \;:\;{} & 0\le k<i,\ \ \pre{P}{k}=\sub{P}{i-k+1}{i},\\
    & i=\abs{P}\ \text{oppure}\ P[k+1]\ne P[i+1] \,\bigr\} ,
  \end{align*}
  cioè il massimo prefisso-suffisso proprio di $\pre{P}{i}$ seguito da un
  carattere diverso da $P[i+1]$ (per $i=\abs{P}$ il carattere $P[i+1]$ non esiste
  e la condizione è vuota), con la convenzione $\spp{i}(P)=0$ quando l'insieme è
  vuoto --- cioè quando nessun prefisso-suffisso proprio di $\pre{P}{i}$, nemmeno
  quello vuoto, è seguito da un carattere diverso da $P[i+1]$: è il caso di
  $P=\str{aaa}$ con $i=2$. In particolare $\spp{i}(P)\le\spv{i}(P)$ e
  $\spp{\abs{P}}(P)=\spv{\abs{P}}(P)$.
\end{definizione}

Si può però ottimizzare ancora: con lo stesso costo posso chiedere una funzione
che mi faccia avanzare $P$ più velocemente di $\spp{i}(P)$.

\section{KMP real-time: \texorpdfstring{$\spv{i,c}$}{sp\_\{i,c\}}}
\label{sec:kmp-real-time}

\scan{9}
\begin{definizione}[$\spv{i,c}(P)$]\label{def:sp-ic}
  Siano $i\in\set{0,1,\dots,n}$ e $c\in\alfa$ (lo stato $i=0$ corrisponde al
  prefisso vuoto). Allora $\spv{i,c}(P)$ è la lunghezza del massimo
  prefisso-suffisso proprio di $\pre{P}{i}$ \emph{seguito da $c$}, cioè
  \[
    \spv{i,c}(P) \defeq \max\set{k \mid 0\le k<i,\ \ \pre{P}{k}=\sub{P}{i-k+1}{i},\ \ P[k+1]=c} .
  \]
  Se nessun prefisso-suffisso di $\pre{P}{i}$ è seguito da $c$ l'insieme è vuoto:
  si pone allora per convenzione $\spv{i,c}(P)=-1$.
\end{definizione}

\noindent Si noti l'assenza dell'apice. L'apice di $\spp{i}$ marca la variante in
cui il bordo è seguito da un carattere \emph{diverso} da $P[i+1]$; qui invece la
condizione è di uguaglianza a $c$.

Chiedere una versione \emph{real-time} di KMP significa \textbf{non} assumere che
$T$ sia dato completamente in input: i caratteri del testo arrivano uno per
volta, come in uno \emph{stream}, e per ciascuno di essi si vuole rispondere in
tempo costante \emph{nel caso peggiore}.

Usando $\spv{i,c}(P)$ questo è immediato: se sono già stati riconosciuti i primi
$i$ caratteri di $P$ e arriva il carattere $c$, aggiorno la posizione di $P$ nel
confronto con $T$ in $\Oh(1)$, passando allo stato
\[
  \delta(i,c) \;=\;
  \begin{cases}
    i+1 & \text{se } i<n \ \text{ e } \ P[i+1]=c,\\[3pt]
    \spv{i,c}(P)+1 & \text{altrimenti}
  \end{cases}
\]
(con la convenzione appena fissata, quando nessun prefisso-suffisso è seguito da
$c$ si ricade nello stato $0$, che è anche lo stato iniziale; per $i=n$ il
carattere $P[n+1]$ non esiste e si ricade sempre nel secondo ramo, che dopo
un'occorrenza completa di $P$ fa ripartire il confronto da $\spv{n,c}(P)+1$).

Con il solo $\spv{i}(P)$, invece, l'arrivo di un singolo carattere può innescare
la cascata di fallimenti $i\to\spv{i}\to\spv{\spv{i}}\to\cdots$ e potrei
impiegare $\Oh(n)$: il costo per carattere resta $\Oh(1)$ \emph{ammortizzato},
ma non nel caso peggiore.

\begin{osservazione}
  Il prezzo si paga in spazio: la tabella dei valori $\spv{i,c}(P)$ ha
  $\Th(n\abs{\alfa})$ celle, contro le $\Th(n)$ del solo vettore $\spv{i}(P)$. È,
  di fatto, la tabella di transizione dell'automa deterministico associato a $P$.
\end{osservazione}

\chapter{Boyer--Moore}

\scan{9}
L'algoritmo di Boyer--Moore (1977) risolve lo stesso problema di Knuth--Morris--Pratt
--- trovare tutte le occorrenze di un pattern $P$ dentro un testo $T$ --- ma cambia la
direzione dei confronti all'interno di un singolo allineamento. Da questa sola idea
discendono due euristiche di riposizionamento che, nei casi favorevoli, permettono di
non leggere affatto alcuni caratteri del testo. Per tutto il capitolo poniamo
\[
  m=\abs{T},\qquad n=\abs{P},
\]
con indicizzazione a partire da $1$; $\alfa$ è l'alfabeto e $P^r$ denota il pattern
\emph{rovesciato}.

\section{L'idea: confronti da destra a sinistra}

Il testo $T$ viene percorso da sinistra a destra, cioè gli allineamenti di $P$ su $T$
avanzano come nell'algoritmo naive e in KMP; \emph{dentro} un allineamento, però, il
pattern $P$ viene confrontato con il testo \emph{da destra a sinistra}. È questa
l'idea che si aggiunge rispetto agli algoritmi visti finora.

\begin{figure}[H]
\centering
% Un passo di Boyer--Moore: confronti da destra a sinistra, mismatch x/y, shift.
\begin{tikzpicture}[x=0.38cm,y=0.38cm,
    testo/.style   = {line width=0.9pt},
    patt/.style    = {line width=0.9pt,deepsea},
    taccaT/.style  = {line width=0.8pt},
    taccaP/.style  = {line width=0.8pt,deepsea},
    forte/.style   = {line width=1.8pt,deepsea},
    et/.style      = {font=\footnotesize},
    ets/.style     = {font=\scriptsize,graydark}]

  %% ------------------------------------------------- riga 1: il testo T
  \draw[testo] (0,0) -- (25,0);
  \node[et,anchor=west] at (25.5,0) {$T$};
  \draw[taccaT] (11,-0.7) -- (11,0.7);
  \draw[taccaT] (12,-0.7) -- (12,0.7);
  \draw[taccaT] (15,-0.7) -- (15,0.7);
  \node[et,anchor=south] at (11.5,0.75) {$x$};
  \node[et,anchor=south] at (13.5,0.75) {$\alpha$};

  %% ------------------------- direzione dei confronti (da destra a sinistra)
  \draw[graydark,line width=0.6pt,-{Latex[length=4pt]}] (15,-2.2) -- (11,-2.2);
  \node[ets,anchor=south] at (13,-2.1) {confronti};

  %% ---------------------------- riga 2: il pattern P, allineamento corrente
  \draw[patt] (6,-4.6) -- (15,-4.6);
  \draw[taccaP] (6,-5.25) -- (6,-3.95);
  \draw[taccaP] (11,-5.25) -- (11,-3.95);
  \draw[taccaP] (12,-5.25) -- (12,-3.95);
  \draw[forte]  (15,-5.5) -- (15,-3.7);
  \node[et,anchor=south] at (11.5,-3.9) {$y$};
  \node[et,anchor=south] at (13.5,-3.9) {$\alpha$};
  \node[et,anchor=north] at (15,-5.75) {$P$};

  %% --------------------------------- porzione di P non ancora confrontata
  \draw[graydark,line width=0.6pt]
        (6.8,-5.9) .. controls (8.2,-6.8) and (9.6,-5.2) .. (11,-5.9);
  \node[ets,anchor=north] at (8.9,-6.9) {non ancora confrontata};

  %% ------------------------------- riga 2: il pattern P dopo lo shift
  \draw[patt] (15.8,-4.6) -- (24.8,-4.6);
  \draw[taccaP] (15.8,-5.25) -- (15.8,-3.95);
  \draw[forte]  (24.8,-5.5) -- (24.8,-3.7);
  \node[et,anchor=north] at (24.8,-5.75) {$P$};

  %% ------------------------------------------------- annotazione a destra
  \node[et,anchor=west,align=left] at (26.1,-2.4)
       {ci sono casi\\[2pt] $\Oh(m/n)$};
\end{tikzpicture}

\caption{Un passo di Boyer--Moore: i confronti procedono da destra a sinistra,
$\alpha$ è il suffisso di $P$ già riconosciuto e il mismatch è fra il carattere
$y$ del pattern e il corrispondente carattere $x$ del testo; il pattern viene poi
fatto scorrere verso destra, nel disegno addirittura di tutta la propria lunghezza.}
\end{figure}

Eseguo dunque i confronti da destra e, quando trovo un mismatch, applico la
\emph{più conveniente} fra le due regole della sezione seguente, cioè quella che
produce il \emph{massimo shift}.

\begin{osservazione}
Poiché dopo un mismatch il pattern può scorrere anche di tutta la propria lunghezza,
ci sono casi in cui i confronti eseguiti --- e quindi i caratteri di $T$
effettivamente letti --- sono soltanto $\Th(m/n)$: l'algoritmo può essere
\emph{sublineare} nella lunghezza del testo, cioè terminare senza aver mai letto
alcuni caratteri di $T$. Il preprocessing del pattern, $\Th(n)$, resta ovviamente
da pagare.
\end{osservazione}

\section{Le due regole}\label{sec:bm-regole}

\scan{10}
Fissiamo la notazione di un allineamento. Supponiamo che $P$ sia allineato con la
finestra di testo che termina in $T[k]$ e che i confronti da destra abbiano già
riconosciuto il suffisso
\[
  \alpha=\sub{P}{i}{n}
\]
(con $\alpha=\vuota$ quando $i=n+1$, cioè quando il primissimo confronto fallisce).
Il primo confronto che fallisce è allora quello fra il carattere $y=P[i-1]$ del
pattern e il carattere $x$ del testo che gli sta di fronte, con $x\neq y$: in tutto
il capitolo $x$ è dunque il carattere «cattivo» del \emph{testo} --- i due ruoli
sono scambiati rispetto al capitolo precedente, dove $x$ era il carattere del
pattern. Al mismatch il pattern viene spostato a destra applicando la più
conveniente fra le due regole seguenti.

\begin{definizione}[Bad character rule (BCR)]\label{def:bcr}
Dopo un mismatch con il carattere $x$ di $T$, si riposiziona il pattern cercando in
$P$ l'occorrenza di $x$ \emph{più a destra fra quelle a sinistra} del punto di
mismatch.
\end{definizione}

\begin{definizione}[Good suffix rule (GSR)]\label{def:gsr}
Dopo un mismatch che segue il suffisso $\alpha$ già riconosciuto in $P$, si
riposiziona $P$ cercando l'occorrenza di $\alpha$ più a destra che non sia un
suffisso di $P$.
\end{definizione}

Le due regole hanno costi di preprocessing molto diversi: la BCR si implementa
direttamente, mentre il preprocessing della GSR si riduce all'algoritmo $Z$
(\autoref{sec:bm-gsr}).

\begin{osservazione}[Di quanto si sposta il pattern]
Per ogni posizione $j$ e ogni $x\in\alfa$ poniamo
\[
  R_j(x) \defeq \max\set{\, p \;:\; 1\le p<j \ \text{e}\ P[p]=x \,},
  \qquad R_j(x)=0 \ \text{se un tale } p \text{ non esiste}.
\]
Con la notazione fissata sopra, la BCR sposta $P$ di $(i-1)-R_{i-1}(x)$ posizioni,
mentre la GSR lo sposta di $n-j_\alpha$, dove $j_\alpha$ è la posizione finale
dell'occorrenza di $\alpha$ individuata dalla regola. Nella \autoref{sec:bm-gsr}
rafforzeremo la regola, chiedendo in più che il carattere che precede
l'occorrenza interna sia diverso da quello che ha provocato il mismatch, e vedremo
che allora $j_\alpha$ è esattamente il valore $L'(i)$.
Quando una tale occorrenza \emph{non} esiste --- cioè quando $L'(i)=0$ --- lo
spostamento sarebbe di $n$ posizioni intere, e così com'è non è \emph{safe}: su
$P=\str{aca}$ e $T=\str{bcaca}$ farebbe saltare l'occorrenza in posizione~3. Il
caso va trattato a parte, come si vedrà nella \autoref{sec:bm-ricerca}.
Con questa avvertenza entrambi gli spostamenti sono $\ge 1$, e si applica sempre
il massimo dei due,
\[
  \mathit{shift} \;=\; \max\set{\,\mathit{shift}_{\mathrm{BCR}},\ \mathit{shift}_{\mathrm{GSR}}\,},
\]
il che è lecito perché entrambe le regole sono \emph{safe} --- nessuna occorrenza
viene persa --- e quindi lo è anche il loro massimo.
\end{osservazione}

\begin{esempio*}[Comportamento sublineare]
Siano $T=a^{m}$ e $P=b^{n}$. Ogni allineamento fallisce al primo confronto e il
carattere cattivo $a$ non compare affatto in $P$: la BCR dà $R_n(a)=0$ e fa scorrere
il pattern di tutta la sua lunghezza. Si esaminano così $m/n$ allineamenti con un
solo confronto ciascuno, cioè $\Th(m/n)$ caratteri di $T$ in tutto.
\end{esempio*}

\begin{osservazione}
All'estremo opposto, se $T=a^{m}$ e $P=a^{n}$ ogni allineamento produce
un'occorrenza completa e nessuna delle due regole permette di avanzare di più di una
posizione: l'algoritmo degenera nell'algoritmo naive, con costo quadratico
$\Th(mn)$.
\end{osservazione}

\textbf{N.B.} È possibile, aggiungendo una struttura dati, rendere l'algoritmo
lineare nel caso peggiore (Apostolico--Giancarlo) --- \emph{da approfondire}.

\section{Implementazione della BCR}

Il preprocessing riguarda il solo pattern $P$.

\begin{figure}[H]
\centering
% Preprocessing della BCR: da j si guarda all'indietro dentro P cercando x.
\begin{tikzpicture}[font=\footnotesize,
    tacca/.style = {line width=1pt},
    ets/.style   = {font=\scriptsize,graydark}]

  %% il pattern P
  \draw[deepsea,line width=2.2pt,-{Latex[length=5pt,width=4.5pt]}] (0,0) -- (10,0);
  \node[anchor=west] at (10.3,0) {$P$};
  \node[ets,anchor=north] at (0,-0.45) {$1$};
  \node[ets,anchor=north] at (9.8,-0.45) {$n$};

  %% le due posizioni: p (occorrenza trovata) e j (mismatch corrente)
  \draw[tacca] (6.4,-0.32) -- (6.4,0.32);
  \draw[tacca] (7.6,-0.32) -- (7.6,0.32);
  \node[anchor=north] at (6.4,-0.36) {$p$};
  \node[anchor=south] at (7.6,0.36) {$j$};
  \node[anchor=south] at (6.8,0.3) {$x$};

  %% si guarda all'indietro, da j verso p
  \draw[gray,dashed,line width=0.7pt,-{Latex[length=4.5pt]}]
        (7.6,-0.9) .. controls (7.6,-1.5) and (6.4,-1.5) .. (6.4,-0.9);
  \node[ets,anchor=north] at (7.0,-1.45) {$j-R_j(x)$};
\end{tikzpicture}

\caption{Preprocessing della BCR: dalla posizione $j$ si guarda all'indietro dentro
$P$ per individuare la posizione $p<j$ più a destra con $P[p]=x$.}
\end{figure}

Per ogni $j\in\set{1,\dots,n}$ e per ogni $x\in\alfa$ serve il massimo indice $p$
tale che $p<j$ e $P[p]=x$, cioè esattamente il valore $R_j(x)$ definito sopra. Ci
sono due modi di memorizzarlo.

\begin{itemize}
  \item \textbf{Matrice.} Si tabula direttamente $R_j(x)$ in una matrice
  $n\times\abs{\alfa}$: costa $\Th(n\,\abs{\alfa})$ in spazio e in tempo di
  preprocessing, ma risponde a ogni interrogazione in tempo costante.

  \item \textbf{Liste.} Per ogni carattere $x$ che compare in $P$ si mantiene la
  lista delle posizioni di $P$ in cui $x$ occorre, in ordine decrescente. Lo spazio
  scende a $\Th(n+\abs{\alfa})$, ma per calcolare $R_j(x)$ bisogna scorrere la
  lista.
\end{itemize}

La seconda soluzione sembra quadratica; facendo bene i conti, però, con un'analisi a
complessità ammortizzata si vede che in realtà è lineare.

\begin{osservazione}[Perché il costo ammortizzato è lineare]
Poiché la lista di $x$ è ordinata per posizione decrescente, per calcolare $R_j(x)$
basta scorrerla dalla testa scartando le posizioni $\ge j$. Se il mismatch cade in
$P[j]$, in quel momento sono già stati eseguiti $n-j+1$ confronti (le posizioni
$n,\dots,j$ di $P$), e le posizioni scartate sono tutte comprese fra $j$ e $n$:
sono quindi al più $n-j+1$. Il costo della scansione della lista è perciò
proporzionale al numero di confronti appena eseguiti, e sull'intera esecuzione il
sovraccosto complessivo resta lineare nel numero totale di confronti:
asintoticamente le due implementazioni si equivalgono.
\end{osservazione}

\section{Implementazione della GSR}\label{sec:bm-gsr}

\scan{11}
Resta da capire quale sia il preprocessing sufficiente a implementare la
\emph{good suffix rule}. La risposta è: l'algoritmo $Z$ applicato a $P^r$.

\begin{definizione}[$L'(i)$]\label{def:bm-lprimo}
Siano $P\in\alfa^*$ e $i\in\set{2,\dots,n}$. Allora $L'(i)$ è il \emph{massimo}
indice $j$ tale che
\begin{enumerate}
  \item $\sub{P}{i}{n}$ è un suffisso di $\sub{P}{1}{j}$;
  \item $P[i-1]\neq P[\,j-(n-i)-1\,]$;
\end{enumerate}
e $L'(i)=0$ se nessun indice soddisfa 1) e 2).
\end{definizione}

\begin{osservazione}
La sola condizione~1), con il vincolo $j\le n-1$, definirebbe $L(i)$, cioè la
versione «debole» della regola: l'occorrenza interna più a destra del suffisso già
riconosciuto (il vincolo è indispensabile: $\sub{P}{i}{n}$ è sempre suffisso di
$P$, quindi senza di esso il massimo sarebbe banalmente $j=n$). È la condizione~2)
a giustificare l'apice in $L'(i)$: essa impone che il carattere che precede
l'occorrenza interna sia \emph{diverso} da $P[i-1]$, cioè dal carattere che
nell'allineamento corrente ha provocato il mismatch; senza questo vincolo lo
spostamento riproporrebbe lo stesso mismatch e sarebbe quindi sprecato.

Gli indici della condizione~2) si leggono così: se $\sub{P}{i}{n}$, che ha lunghezza
$n-i+1$, è un suffisso di $\sub{P}{1}{j}$, allora occupa le posizioni
$j-(n-i),\dots,j$, e il carattere che la precede sta in posizione $j-(n-i)-1$.

Restano sottintese due convenzioni: se $j-(n-i)-1=0$ la posizione non esiste e la
condizione~2) si intende soddisfatta; $L'(i)$ non è definito per $i=1$, perché la
condizione~2) chiamerebbe in causa $P[0]$. Si noti infine che la~2) esclude
automaticamente il valore $j=n$: in tal caso diventerebbe $P[i-1]\neq P[i-1]$, che è
falsa. Dunque $L'(i)\le n-1$ senza bisogno di imporlo, ed è per questo che lo
spostamento $n-L'(i)$ della GSR è sempre positivo.
\end{osservazione}

\begin{figure}[H]
\centering
% Lettura grafica di L'(i): la stessa stringa P[i..n] occorre due volte in P
% (come suffisso e come occorrenza interna che termina in L'(i)) e i due
% caratteri che la precedono sono diversi.
\begin{tikzpicture}[x=1.3cm,font=\footnotesize,
    tacca/.style = {line width=1pt},
    barra/.style = {line width=2.4pt,graydark},
    ets/.style   = {font=\scriptsize,graydark}]

  %% il pattern P: posizione 1 in x=0, posizione n in x=10.4
  \draw[deepsea,line width=2.2pt,-{Latex[length=5pt,width=4.5pt]}] (0,0) -- (10.6,0);
  \node[anchor=west] at (10.9,0) {$P$};

  %% tacche
  \foreach \t in {2.4,3.6,6.0,6.8,8.0} \draw[tacca] (\t,-0.3) -- (\t,0.3);

  %% i due caratteri che precedono le due occorrenze
  \node[anchor=south] at (3.0,0.16) {$z$};
  \node[anchor=south] at (7.4,0.16) {$y$};

  %% le due occorrenze della stessa stringa P[i..n], di uguale lunghezza
  \draw[barra] (3.6,0.55) -- (6.0,0.55);
  \draw[tacca,graydark] (3.6,0.38) -- (3.6,0.72);
  \draw[tacca,graydark] (6.0,0.38) -- (6.0,0.72);
  \draw[barra] (8.0,0.55) -- (10.4,0.55);
  \draw[tacca,graydark] (8.0,0.38) -- (8.0,0.72);
  \draw[tacca,graydark] (10.4,0.38) -- (10.4,0.72);
  \node[anchor=south] at (4.8,0.68) {$\sub{P}{i}{n}$};
  \node[anchor=south] at (9.2,0.68) {$\sub{P}{i}{n}$};

  %% condizione 2): i caratteri che precedono le due occorrenze sono diversi
  \draw[gray,dashed,line width=0.7pt]
        (3.0,0.6) .. controls (3.9,2.6) and (6.5,2.6) .. (7.4,0.6);
  \node[fill=white,inner sep=3pt] at (5.2,2.1) {$\neq$};

  %% etichette delle posizioni
  \node[ets,anchor=north] at (3.0,-0.34) {$L'(i)-(n-i)-1$};
  \node[anchor=north] at (6.0,-0.34) {$L'(i)$};
  \node[anchor=north] at (7.4,-0.34) {$i-1$};
  \node[anchor=north west] at (8.0,-0.34) {$i$};

  %% condizione 1): P[i..n] e' un suffisso di P[1..L'(i)]
  \draw[decorate,decoration={brace,mirror,amplitude=5pt},gray,line width=0.7pt]
        (0,-1.05) -- (6.0,-1.05);
  \node[ets,anchor=north] at (3.0,-1.32) {$\sub{P}{1}{L'(i)}$};
\end{tikzpicture}

\caption{Lettura grafica della definizione di $L'(i)$: la stessa stringa
$\sub{P}{i}{n}$ compare due volte in $P$ --- come suffisso e come occorrenza interna
che termina in $L'(i)$ --- e i due caratteri che la precedono sono diversi.}
\end{figure}

\begin{definizione}[{{$N_j[P]$}}]\label{def:bm-nj}
$N_j[P]$ è la lunghezza del massimo suffisso di $\sub{P}{1}{j}$ che è anche un
suffisso di $P$, ossia
\[
  N_j[P]=\max\set{\ell\ge 0 \;:\; \sub{P}{j-\ell+1}{j}=\sub{P}{n-\ell+1}{n}},
  \qquad j=1,\dots,n .
\]
\end{definizione}

I valori $N_j$ sono l'analogo «riflesso» dei valori $Z$: $\Zv{k}$ guarda i prefissi
da sinistra, $N_j$ guarda i suffissi da destra. In particolare $N_n[P]=n$.

\begin{figure}[H]
\centering
% N_j[P] = |alpha|, con alpha suffisso di P[1..j] e insieme suffisso di P.
\begin{tikzpicture}[font=\footnotesize,
    tacca/.style = {line width=1pt},
    ets/.style   = {font=\scriptsize,graydark}]

  %% il pattern P: posizione 1 in x=0, posizione n in x=10.4
  \draw[deepsea,line width=2.2pt,-{Latex[length=5pt,width=4.5pt]}] (0,0) -- (10.8,0);
  \node[anchor=west] at (11.1,0) {$P$};

  %% tacche agli estremi delle due occorrenze di alpha
  \foreach \t in {2.3,4.3,8.4,10.4} \draw[tacca] (\t,-0.3) -- (\t,0.3);

  %% occorrenza che termina in j
  \draw[decorate,decoration={brace,amplitude=5pt,raise=3pt},graydark,line width=0.7pt]
        (2.3,0) -- (4.3,0);
  \node[anchor=south] at (3.3,0.36) {$\alpha$};

  %% occorrenza che termina in n: alpha e' anche un suffisso di P
  \draw[decorate,decoration={brace,amplitude=5pt,raise=3pt},graydark,line width=0.7pt]
        (8.4,0) -- (10.4,0);
  \node[anchor=south] at (9.4,0.36) {$\alpha$};

  %% le due occorrenze sono la stessa stringa
  \draw[gray,dashed,line width=0.7pt]
        (3.3,0.78) .. controls (4.4,1.9) and (8.3,1.9) .. (9.4,0.78);
  \node[fill=white,inner sep=3pt] at (6.35,1.62) {$=$};

  %% etichette delle posizioni
  \node[ets,anchor=north] at (0,-0.3) {$1$};
  \node[anchor=north] at (4.3,-0.3) {$j$};
  \node[ets,anchor=north] at (10.4,-0.3) {$n$};
\end{tikzpicture}

\caption{$N_j[P]=\abs{\alpha}$, dove $\alpha$ è la più lunga stringa che è
simultaneamente suffisso di $\sub{P}{1}{j}$ e suffisso di $P$.}
\end{figure}

\begin{teorema}\label{thm:bm-lprimo-nj}
Per ogni $i\in\set{2,\dots,n}$, $L'(i)$ è il massimo $j$ tale che
\[
  N_j[P]=\abs{\sub{P}{i}{n}}=n-i+1 .
\]
\end{teorema}

\begin{dimostrazione}
Basta mostrare che, per ogni indice $j$, le condizioni 1) e 2) della
Definizione~\ref{def:bm-lprimo} valgono per quel $j$ se e solo se $N_j[P]=n-i+1$.

Se valgono 1) e 2), la 1) dice che $\sub{P}{j-(n-i)}{j}=\sub{P}{i}{n}$, dunque
$\sub{P}{1}{j}$ ha un suffisso di lunghezza $n-i+1$ che è anche suffisso di $P$ e
quindi $N_j[P]\ge n-i+1$; se fosse $N_j[P]\ge n-i+2$ il match si estenderebbe di un
carattere verso sinistra, cioè $P[j-(n-i)-1]=P[i-1]$, contro la 2).

Viceversa, se $N_j[P]=n-i+1$ allora il suffisso comune di quella lunghezza è
esattamente $\sub{P}{j-(n-i)}{j}=\sub{P}{i}{n}$, cioè vale la 1); e la
\emph{massimalità} di $N_j$ impedisce l'estensione a sinistra, cioè
$P[j-(n-i)-1]\neq P[i-1]$, che è la 2) (se $j-(n-i)-1=0$ la condizione è soddisfatta
per convenzione, perché il match non può estendersi oltre l'inizio di $P$).

Le due condizioni sono equivalenti indice per indice, quindi hanno lo stesso massimo.
Si noti che per $i\ge 2$ il valore $j=n$ è automaticamente escluso, poiché
$N_n[P]=n>n-i+1$.
\end{dimostrazione}

Da qui, e dalla definizione stessa di $N_j$, segue la riduzione all'algoritmo $Z$:
per ogni $1\le j\le n-1$,
\[
  \boxed{\;N_j[P]=\Zv{n-j+1}[P^r]\;}
\]
Il solo valore escluso è $j=n$: chiamerebbe in causa $\Zv{1}[P^r]$, che la
definizione della funzione $Z$ non contempla. Non è una perdita, perché
$N_n[P]=n$ è noto a priori e nel seguito $N_j$ serve soltanto per $j\le n-1$.

\begin{osservazione}
Rovesciando $P$, la posizione $j$ di $P$ diventa la posizione $n-j+1$ di $P^r$; un
suffisso di $\sub{P}{1}{j}$ diventa un prefisso di $\sub{P^r}{n-j+1}{n}$ e un
suffisso di $P$ diventa un prefisso di $P^r$. Quindi il più lungo suffisso di
$\sub{P}{1}{j}$ che è anche suffisso di $P$ ha la stessa lunghezza del più lungo
prefisso di $\sub{P^r}{n-j+1}{n}$ che è anche prefisso di $P^r$, cioè
$\Zv{n-j+1}[P^r]$.

Poiché l'algoritmo $Z$ su una stringa di lunghezza $n$ costa $\Th(n)$, tutto il
preprocessing della GSR costa $\Th(n)$ in tempo e in spazio.
\end{osservazione}

\section{Lo spostamento della GSR e il ciclo di ricerca}\label{sec:bm-ricerca}

Sapere calcolare $L'(i)$ non basta ancora a far girare l'algoritmo: manca il caso in
cui $L'(i)=0$, cioè quando il suffisso $\alpha$ già riconosciuto non ha alcuna
ulteriore occorrenza utile dentro $P$. In quel caso non si può buttare via tutto:
bisogna allineare il più lungo prefisso di $P$ che sia anche suffisso di $\alpha$.

\begin{definizione}[$\ell'(i)$]\label{def:bm-lminuscolo}
Per $i\in\set{2,\dots,n+1}$, $\ell'(i)$ è la lunghezza del più lungo suffisso di
$\sub{P}{i}{n}$ che è anche un prefisso di $P$, e $\ell'(i)=0$ se un tale suffisso
non esiste. In particolare $\ell'(i)\le n-i+1$.
\end{definizione}

Anche $L'$ e $\ell'$ si leggono direttamente dal vettore $N$:
\begin{align*}
  L'(i) &= \max\set{\,j \;:\; 1\le j\le n-1,\ N_j[P]=n-i+1\,},\\
  \ell'(i) &= \max\set{\,j \;:\; 1\le j\le n-i+1,\ N_j[P]=j\,},
\end{align*}
dove il massimo su un insieme vuoto vale $0$; entrambi i vettori si riempiono con
una singola scansione di $N$, quindi in tempo $\Th(n)$. La seconda formula si
giustifica così: $N_j[P]=j$ significa che $\pre{P}{j}$ è anche un suffisso di $P$, e
la limitazione $j\le n-i+1$ impone che quel suffisso stia dentro $\sub{P}{i}{n}$.

Lo spostamento prodotto dalla GSR è allora
\[
  \sigma(i)\defeq
  \begin{cases}
    n-\ell'(2) & \text{se } i=1 \quad (\text{occorrenza completa}),\\[3pt]
    n-L'(i)  & \text{se } 2\le i\le n \ \text{e}\ L'(i)>0,\\[3pt]
    n-\ell'(i) & \text{se } 2\le i\le n \ \text{e}\ L'(i)=0,\\[3pt]
    1        & \text{se } i=n+1 \quad (\alpha=\vuota).
  \end{cases}
\]
L'ultimo caso è obbligato: se il primo confronto fallisce, nessun suffisso di $P$ è
stato riconosciuto e la GSR non porta alcuna informazione. Negli altri casi
$\sigma(i)\ge 1$, perché $L'(i)\le n-1$ e $\ell'(i)\le n-i+1\le n-1$.

\begin{algorithm}[H]
\caption{Boyer--Moore: ciclo di ricerca}
\begin{algorithmic}[1]
\Require testo $T$ di lunghezza $m$, pattern $P$ di lunghezza $n$
\Ensure le posizioni iniziali di tutte le occorrenze di $P$ in $T$
\State calcola $R_j(x)$, $L'$ e $\ell'$ con il preprocessing di $P$
\State $k \gets n$ \Comment{$P[n]$ è allineato con $T[k]$}
\While{$k \le m$}
  \State $i \gets n+1$; \quad $h \gets k+1$
  \While{$i>1$ \textbf{ e } $P[i-1]=T[h-1]$}
    \State $i \gets i-1$; \quad $h \gets h-1$
  \EndWhile
  \If{$i=1$}
    \State riporta l'occorrenza che inizia in $k-n+1$
    \State $k \gets k+\sigma(1)$
  \Else
    \State $x \gets T[h-1]$ \Comment{il carattere «cattivo»}
    \State $k \gets k+\max\set{\,(i-1)-R_{i-1}(x),\ \sigma(i)\,}$
  \EndIf
\EndWhile
\end{algorithmic}
\end{algorithm}

L'invariante del ciclo interno è $\sub{P}{i}{n}=\sub{T}{h}{k}$; all'uscita, o $i=1$
e si è trovata un'occorrenza, oppure $P[i-1]\neq T[h-1]$ e si è nella situazione
descritta nella \autoref{sec:bm-regole}.

\begin{table}[H]
\centering
\begin{tabular}{ll}
\toprule
\textbf{Fase} & \textbf{Costo} \\
\midrule
Preprocessing BCR, matrice & $\Th(n\abs{\alfa})$ tempo e spazio, interrogazioni $\Oh(1)$ \\
Preprocessing BCR, liste   & $\Th(n+\abs{\alfa})$ spazio, costo ammortizzato equivalente \\
Preprocessing GSR ($Z$ su $P^r$) & $\Th(n)$ tempo e spazio \\
Ricerca, caso migliore     & $\Th(m/n)$ caratteri ispezionati (sublineare) \\
Ricerca, caso peggiore     & $\Th(mn)$ (quadratico) \\
\bottomrule
\end{tabular}
\caption{Riepilogo dei costi di Boyer--Moore.}
\end{table}

Il caso peggiore quadratico è quello già visto di $T=a^{m}$, $P=a^{n}$: è proprio
per eliminarlo, senza rinunciare alle due euristiche, che serve la variante di
Apostolico--Giancarlo.

\chapter{Matching multi-pattern: keyword tree e Aho--Corasick}
\label{cap:aho-corasick}

\scan{12}
Generalizziamo il problema del pattern matching esatto al caso in cui i pattern da cercare
non siano uno solo, ma un'intera collezione.

\section{Il problema: exact set matching}
\label{sec:exact-set-matching}

\begin{definizione}[Exact set matching]\label{def:exact-set-matching}
Dati un testo $T \in \S^{*}$ e una collezione di pattern
\[
  \mathcal{P} \;=\; \set{P_1, \dots, P_z}, \qquad P_i \in \S^{*},
\]
trovare tutte le occorrenze in $T$ di tutti gli elementi di $\mathcal{P}$.
\end{definizione}

\paragraph{Notazione.}
Poniamo $n_i = \abs{P_i}$ per $i \in \set{1,\dots,z}$ e
\[
  n \;=\; \sum_{i=1}^{z} n_i ,
\]
e indichiamo con $k$ il numero complessivo di occorrenze da riportare, cioè la dimensione
dell'output. In questo capitolo $n$ è dunque la lunghezza \emph{complessiva dei pattern} e
non quella del testo: la lunghezza del testo la scriviamo sempre, esplicitamente,
$\abs{T}$. Lo scenario tipico è quello in cui $T$ è enorme rispetto ai pattern, per cui il
costo dominante è il numero di scansioni del testo.

\paragraph{Soluzione naive.}
Applichiamo $z$ volte il miglior algoritmo per singolo pattern di cui disponiamo (KMP,
oppure l'algoritmo $Z$), che costa $\Oh(\abs{T} + n_i)$ sul pattern $P_i$. In totale
\[
  \sum_{i=1}^{z} \Oh\bigl(\abs{T} + n_i\bigr)
  \;=\; \Oh\bigl(z\,\abs{T} + n\bigr).
\]

Il \emph{punto debole} è evidente: stiamo scandendo $T$ per $z$ volte. È il fattore $z$
davanti ad $\abs{T}$ che vogliamo eliminare.

\paragraph{La sfida.}
Riusciamo a ridurci a \emph{una sola} scansione del testo? Per farlo dobbiamo combinare
tutti i pattern di $\mathcal{P}$ in un'unica struttura dati, cioè in un albero.

\section{Il keyword tree}
\label{sec:keyword-tree}

\scan{12}
Il \emph{keyword tree} $K(\mathcal{P})$ è l'albero che si ottiene sovrapponendo i prefissi
comuni degli elementi di $\mathcal{P}$: ogni cammino dalla radice a una foglia si legge come
uno dei pattern e, viceversa, ogni pattern corrisponde al cammino verso una foglia. Il
guadagno è immediato: leggendo un solo carattere di $T$ e scendendo di un arco si avanza
simultaneamente su \emph{tutti} i pattern che condividono il prefisso già letto, cioè si
scandiscono più elementi di $\mathcal{P}$ in parallelo.

Nel disegno sono già indicati, tratteggiati, i \emph{failure link}: sono i collegamenti che
permettono di riprendere la ricerca da un altro nodo dell'albero quando la discesa si blocca
su un mismatch, senza mai arretrare sul testo. Li definiremo formalmente nella
\autoref{sec:failure-link}.

\begin{figure}[H]
\centering
% Vista panoramica del keyword tree K(P): otto cammini radice-foglia (i pattern
% P_1..P_8) e due failure link tratteggiati.  I tratti in accento sono le coppie
% di sottocammini che portano la stessa stringa (suffisso/prefisso): per ogni
% failure link <v, n_v> sono marcati il cammino root ~> n_v e il suffisso di pari
% lunghezza del cammino root ~> v.  Le due coppie sono
%   alfa = root~>A  e  H2~>C   (failure link C -> A)   lunghezza 4.611
%   beta = root~>D  e  H~>B    (failure link B -> D)   lunghezza 1.392
% I due tratti di ciascuna coppia hanno la STESSA lunghezza (e' l'informazione
% del disegno).  Il tratto di destra parte da H2, sotto D, e non da D: nel
% manoscritto fra i due tratti marcati del bordo destro c'e' un segmento a penna
% sottile che li separa, senza il quale nel PDF si fondono in un unico tratto.
\begin{tikzpicture}[
    ramo/.style    = {draw=graydark, line width=0.55pt, line cap=round, line join=round},
    marcato/.style = {draw=accent, line width=2.4pt, line cap=round, line join=round,
                      opacity=0.7},
    nodo/.style    = {circle, fill=graydark, inner sep=0pt, minimum size=3.4pt},
    foglia/.style  = {circle, fill=graydark, inner sep=0pt, minimum size=2.4pt},
    flink/.style   = {draw=black, line width=0.7pt, dash pattern=on 3pt off 2.2pt,
                      -{Triangle[length=5pt,width=4pt]}},
    et/.style      = {font=\footnotesize},
  ]

  % ---- nodi -------------------------------------------------------------
  \coordinate (r)  at ( 0.00, 5.50);   % radice
  \coordinate (m1) at (-2.00, 4.07);   % biforcazione sul cammino r -> A
  \coordinate (A)  at (-3.78, 2.86);   % arrivo del failure link lungo
  \coordinate (B)  at (-0.63, 2.53);   % partenza del failure link corto
  \coordinate (D)  at ( 1.16, 4.73);   % arrivo del failure link corto (poco profondo)
  \coordinate (F)  at ( 2.73, 2.97);   % biforcazione sul bordo destro
  \coordinate (C)  at ( 4.655, 0.935); % partenza del failure link lungo (profondo)
  \coordinate (E)  at (-1.26, 1.65);
  \coordinate (G)  at ( 2.10, 1.76);
  \coordinate (P1) at (-4.83, 0.28);
  \coordinate (P2) at (-2.63, 0.28);
  \coordinate (P3) at (-1.79, 0.44);
  \coordinate (P4) at (-0.74, 0.22);
  \coordinate (P5) at ( 0.95, 0.55);
  \coordinate (P6) at ( 1.89, 0.33);
  \coordinate (P7) at ( 3.15, 0.33);
  \coordinate (P8) at ( 5.05, 0.30);
  \coordinate (H)  at ($(m1)!0.33!(B)$);    % inizio del tratto marcato che entra in B
  \coordinate (H2) at ($(D)!0.2328!(F)$);   % inizio del tratto marcato che scende in C

  % ---- tratti marcati (sotto ai rami) -----------------------------------
  \draw[marcato] (r) -- (m1) -- (A);      % alfa del failure link C -> A
  \draw[marcato] (H2) -- (F) -- (C);      % alfa del failure link C -> A
  \draw[marcato] (r) -- (D);              % beta del failure link B -> D
  \draw[marcato] (H) -- (B);              % beta del failure link B -> D

  % ---- rami dell'albero -------------------------------------------------
  \draw[ramo] (r) -- (m1) -- (A);
  \draw[ramo] (A) -- (P1);
  \draw[ramo] (A) -- (-3.10,1.95) -- (P2);
  \draw[ramo] (m1) -- (B);
  \draw[ramo] (B) -- (E);
  \draw[ramo] (E) -- (P3);
  \draw[ramo] (E) -- (-0.95,0.95) -- (P4);
  \draw[ramo] (B) -- (0.30,1.55) -- (P5);
  \draw[ramo] (r) -- (D) -- (F) -- (C) -- (P8);
  \draw[ramo] (F) -- (G);
  \draw[ramo] (G) -- (P6);
  \draw[ramo] (G) -- (2.85,1.05) -- (P7);

  % ---- pallini ----------------------------------------------------------
  \node[nodo] at (r) {};
  \node[nodo] at (A) {};
  \node[nodo] at (B) {};
  \node[nodo] at (C) {};
  \node[nodo] at (D) {};
  \foreach \l in {P1,P2,P3,P4,P5,P6,P7,P8} \node[foglia] at (\l) {};

  % ---- failure link -----------------------------------------------------
  \draw[flink] (C) .. controls (3.30,0.62) and (1.30,1.42) .. (-0.70,2.05)
                   .. controls (-1.95,2.40) and (-2.95,2.64) .. (A);
  \draw[flink] (B) .. controls (0.55,2.95) and (1.18,3.85) .. (D);
  \node[et,align=center] at (1.50,3.12) {\emph{failure}\\[-1pt]\emph{link}};

  % ---- etichette delle foglie -------------------------------------------
  \node[et,below=2pt] at (P1) {$P_1$};
  \node[et,below=2pt] at (P2) {$P_2$};
  \node[et,below=2pt] at (P3) {$P_3$};
  \node[et,below=2pt] at (P4) {$P_4$};
  \node[et,below=2pt] at (P5) {$P_5$};
  \node[et,below=2pt] at (P6) {$P_6$};
  \node[et,below=2pt] at (P7) {$P_7$};
  \node[et,below=2pt] at (P8) {$P_8$};

  \node[et,above=2pt] at (r) {$\mathrm{root}$};
\end{tikzpicture}

\caption{Il keyword tree $K(\mathcal{P})$ a colpo d'occhio: i cammini dalla radice alle $z$
foglie sono i pattern $P_1,\dots,P_z$. Tratteggiati, due failure link: puntano sempre da un
nodo a un nodo di profondità minore. In evidenza, per ciascun failure link
$\tuplev{v, \nv{v}}$, i due sottocammini che portano la stessa stringa: il cammino dalla
radice a $\nv{v}$ e il suffisso di pari lunghezza del cammino dalla radice a $v$.}
\label{fig:keyword-tree-panoramica}
\end{figure}

\subsection{Definizione}

\scan{13}

\begin{definizione}[Keyword tree]\label{def:keyword-tree}
Dato $\mathcal{P} = \set{P_1, \dots, P_z}$ con $P_i \in \S^{*}$ per
$i \in \set{1, \dots, z}$, il \emph{keyword tree} di $\mathcal{P}$ è un albero $K$ i cui
archi sono etichettati con elementi di $\S$ e tale che:
\begin{enumerate}
  \item archi uscenti dallo stesso nodo hanno etichette distinte;
  \item ogni cammino radice--foglia corrisponde a $P_i$ per qualche $i \in \set{1,\dots,z}$,
        e viceversa ogni $P_i$ si legge lungo un cammino radice--foglia.
\end{enumerate}
Scriviamo $K(\mathcal{P})$ per il keyword tree di $\mathcal{P}$.
\end{definizione}

\begin{osservazione}
Le due condizioni dicono, rispettivamente:
\begin{itemize}
  \item la (1) che ogni stringa individua in $K$ un \emph{cammino univoco}: la discesa dalla
        radice è deterministica e due nodi distinti non possono avere la stessa etichetta di
        cammino. I nodi di $K$ sono dunque in biiezione con i prefissi degli elementi di
        $\mathcal{P}$: il keyword tree è il \emph{trie} dei pattern, cioè una
        rappresentazione \emph{non ridondante sui prefissi};
  \item la (2) che i \emph{cammini} radice--foglia sono in corrispondenza con gli
        \emph{elementi di $\mathcal{P}$}.
\end{itemize}
\end{osservazione}

\subsection{Ipotesi semplificativa}

Assumiamo che nessun $P_i$ sia \emph{fattore proprio} di $P_j$ per $i \neq j$, cioè che
nessun pattern sia sottostringa propria di un altro. Scaricheremo questa ipotesi nella
\autoref{sec:output-link}.

\begin{osservazione}
L'ipotesi garantisce che foglie di $K(\mathcal{P})$ ed elementi di $\mathcal{P}$ siano in
corrispondenza biunivoca --- in particolare le foglie sono esattamente $z$ --- e soprattutto
che durante la ricerca ogni occorrenza venga scoperta arrivando in una foglia: se infatti
$P_j$ fosse sottostringa di $P_i$, percorrendo il cammino di $P_i$ si passerebbe
\emph{sopra} un'occorrenza di $P_j$ senza mai visitarne la foglia. Per la sola biiezione
foglie--pattern basterebbe in realtà l'ipotesi più debole «nessun $P_i$ è prefisso di
$P_j$».
\end{osservazione}

\subsection{Misure di \texorpdfstring{$K(\mathcal{P})$}{K(P)}}

Le tre misure di $K(\mathcal{P})$ sono
\[
\begin{aligned}
  \text{dimensione di } K(\mathcal{P}) \;&\le\; \sum_{i=1}^{z} n_i \;=\; n, \\[2pt]
  \text{altezza di } K(\mathcal{P}) \;&=\; \max_{1 \le i \le z} \set{n_i}, \\[2pt]
  \text{grado (fan-out) di } K(\mathcal{P}) \;&\le\; \min\set{\abs{\S},\, z}.
\end{aligned}
\]

\begin{osservazione}
Per \emph{dimensione} si intende il numero di archi, che in un albero è pari al numero di
nodi meno uno: gli archi di $K(\mathcal{P})$ sono tanti quanti i prefissi non vuoti
\emph{distinti} degli elementi di $\mathcal{P}$, perché ogni pattern contribuisce al più
$n_i$ archi e i prefissi condivisi vengono contati una sola volta. Dunque
$\#\text{archi} \le n$, con uguaglianza se e solo se nessuna coppia di pattern condivide un
prefisso non vuoto; lo spazio è quindi $\Oh(n)$, e $\Th(n)$ soltanto nel caso peggiore. È
proprio questa condivisione il risparmio che il keyword tree realizza rispetto a $z$
strutture separate.

L'altezza è la lunghezza del pattern più lungo. Il grado di un nodo è $\le \abs{\S}$ per la
proprietà (1) della definizione~\ref{def:keyword-tree}, ed è $\le z$ perché ogni sottoalbero figlio
contiene almeno una foglia e le foglie sono esattamente $z$.
\end{osservazione}

\begin{osservazione}[Nota implementativa]
Rappresentando i figli di un nodo con un array indicizzato su $\S$ si ha accesso $\Oh(1)$ al
figlio dato il carattere, ma spazio $\Oh(n\,\abs{\S})$; con liste, alberi bilanciati o
tabelle hash lo spazio resta $\Oh(n)$ e l'accesso costa $\Oh(\log\abs{\S})$ oppure $\Oh(1)$
atteso. Tutte le complessità «lineari» che seguono vanno intese nel modello ad accesso
$\Oh(1)$ ai figli.
\end{osservazione}

\section{I failure link}
\label{sec:failure-link}

Per utilizzare $K(\mathcal{P})$ durante la ricerca dobbiamo completarlo con i \emph{failure
link}, pensati per consentire al puntatore sul testo di avanzare \emph{sempre}: in caso di
mismatch non si ricomincia la discesa dalla radice riportando indietro il puntatore su $T$,
ma si salta al nodo la cui etichetta di cammino è il più lungo suffisso \emph{proprio} di
quanto già riconosciuto che sia ancora etichetta di un cammino dalla radice in
$K(\mathcal{P})$ --- equivalentemente, che sia ancora prefisso di almeno un pattern.

\begin{figure}[H]
\centering
% Idea del failure link: se la discesa si blocca in v si salta a n_v, il nodo la cui
% etichetta di cammino alfa e' il piu' lungo suffisso proprio di L(v) che sia ancora
% un cammino dalla radice.  I due tratti marcati alfa hanno la stessa lunghezza.
\begin{tikzpicture}[
    cammino/.style = {draw=graydark, line width=0.6pt, line cap=round, line join=round},
    marcato/.style = {draw=accent, line width=2.6pt, line cap=round, opacity=0.7},
    nodo/.style    = {circle, fill=graydark, inner sep=0pt, minimum size=3.6pt},
    flink/.style   = {draw=black, line width=0.7pt, dash pattern=on 3.5pt off 2.5pt,
                      -{Triangle[length=5.5pt,width=4.5pt]}},
    prosegue/.style= {draw=graydark, line width=0.6pt, dotted, line cap=round},
    et/.style      = {font=\footnotesize},
  ]

  % ---- sagoma di K(P) ---------------------------------------------------
  \fill[black!5, draw=gray!55, line width=0.4pt]
      (0,4.8) .. controls (2.02,4.2) and (3.70,2.49) .. (4.4,0.145)
              .. controls (2.2,-0.30) and (-2.2,-0.30) .. (-4.4,0.145)
              .. controls (-3.70,2.49) and (-2.02,4.2) .. cycle;

  % ---- nodi -------------------------------------------------------------
  \coordinate (r)  at ( 0.00, 4.80);
  \coordinate (u)  at (-1.02, 3.07);
  \coordinate (v)  at (-2.30, 0.90);
  \coordinate (nv) at ( 1.36, 2.70);

  % ---- cammini: i due tratti che portano alfa sono marcati ---------------
  \draw[marcato] (u) -- (v);
  \draw[marcato] (r) -- (nv);
  \draw[cammino] (r) -- (u) -- (v);
  \draw[cammino] (r) -- (nv);
  \draw[prosegue] (nv) -- (1.78,2.05);

  \node[nodo] at (r)  {};
  \node[nodo] at (u)  {};
  \node[nodo] at (v)  {};
  \node[nodo] at (nv) {};

  % ---- failure link -----------------------------------------------------
  \draw[flink] (v) -- (nv)
      node[et,midway,sloped,above=1pt] {\emph{failure link}};

  % ---- etichette --------------------------------------------------------
  \node[et,above=2pt]           at (r)  {$\mathrm{root}$};
  \node[et,anchor=east,xshift=-3pt]  at (u)  {$u$};
  \node[et,anchor=east,xshift=-3pt]  at (v)  {$v$};
  \node[et,anchor=west,xshift=4pt,yshift=3pt] at (nv) {$\nv{v}$};
  \node[et] at (-2.12, 2.24) {$\a$};
  \node[et] at ( 1.08, 4.00) {$\a$};
  \node[et] at ( 2.95, 0.70) {$K(\mathcal{P})$};
\end{tikzpicture}

\caption{Idea del failure link. Se la discesa si blocca nel nodo $v$, si salta al nodo
$\nv{v}$, la cui etichetta di cammino $\a$ è il più lungo suffisso proprio dell'etichetta di
cammino di $v$ che sia ancora un cammino dalla radice in $K(\mathcal{P})$: la porzione di
testo già riconosciuta resta allineata e il puntatore su $T$ non arretra.}
\label{fig:failure-link-idea}
\end{figure}

\subsection{Formalizzazione}

\scan{14}

\begin{definizione}[Etichetta di cammino, failure node]\label{def:ac-failure-node}
Dato un nodo $v$ di $K(\mathcal{P})$ definiamo:
\begin{itemize}
  \item $\L(v)$: la concatenazione delle etichette degli archi nel cammino
        $\mathrm{root}(K(\mathcal{P})) \rightsquigarrow v$;
  \item $\nv{v}$: il nodo di $K(\mathcal{P})$ tale che $\L(\nv{v})$ è il più lungo suffisso
        proprio di $\L(v)$ fra quelli che etichettano un cammino uscente dalla radice, cioè
        che sono della forma $\L(u)$ per qualche nodo $u$ di $K(\mathcal{P})$
        (equivalentemente: che sono prefissi di almeno un pattern di $\mathcal{P}$);
  \item $\mathit{lp}(v) \defeq \abs{\L(\nv{v})}$.
\end{itemize}
\end{definizione}

\begin{osservazione}
Attenzione al pedice, perché la lettera $n$ è ormai sovraccarica: in $n_i$ il pedice è
l'indice di un pattern e $n_i$ è una \emph{lunghezza}, in $\nv{v}$ il pedice è un nodo e
$\nv{v}$ è a sua volta un \emph{nodo} di $K(\mathcal{P})$.

Ogni arco di $K(\mathcal{P})$ porta esattamente un carattere, quindi
$\abs{\L(u)} = \mathrm{depth}(u)$ per ogni nodo $u$: dunque $\mathit{lp}(v)$ è la
\emph{profondità} di $\nv{v}$. Inoltre $\L(\nv{v})$ è un suffisso \emph{proprio} di $\L(v)$,
e perciò
\[
  \mathrm{depth}(\nv{v}) \;<\; \mathrm{depth}(v),
\]
fatto da cui dipenderà la terminazione di tutto ciò che segue. Per $z = 1$ il keyword tree
degenera in un cammino e $\mathit{lp}(v)$ è la lunghezza del più lungo bordo proprio del
prefisso già letto: la mappa $v \mapsto \nv{v}$ è la generalizzazione ad albero della
funzione di fallimento di KMP.
\end{osservazione}

\begin{lemma}[Buona definizione di $\nv{v}$]\label{lem:ac-nv}
Per ogni nodo $v \in K(\mathcal{P})$ con $v \neq \mathrm{root}(K(\mathcal{P}))$ esiste ed è
unico il nodo $\nv{v}$.
\end{lemma}

\begin{dimostrazione*}
Esercizio. \emph{N.B.} $\nv{v}$ può essere la radice.
\end{dimostrazione*}

\begin{osservazione}
L'enunciato va inteso per $v \neq \mathrm{root}(K(\mathcal{P}))$: infatti
$\L(\mathrm{root}) = \vuota$ non ha suffissi propri. Per la radice si adotta la convenzione
$\nv{\mathrm{root}} \defeq \mathrm{root}$ e $\mathit{lp}(\mathrm{root}) \defeq 0$.
\end{osservazione}

\begin{definizione}[Failure link]\label{def:ac-failure-link}
La coppia $\tuplev{v, \nv{v}}$ si chiama \emph{failure link}.
\end{definizione}

\subsection{Calcolo dei failure link}
\label{sec:calcolo-failure-link}

I failure link vanno calcolati contestualmente alla costruzione di $K(\mathcal{P})$, cioè
nella stessa fase di preprocessing, prima di toccare $T$: serve dunque un \emph{algoritmo
per determinare $\nv{v}$ dato $v$}. L'idea è ricorsiva e si legge sulla
figura~\ref{fig:ac-failure-link-passo}: se $v'$ è il padre di $v$ e $x$ è l'etichetta dell'arco
$v' \to v$ (cosicché $\L(v) = \L(v')\,x$), allora, noto $\nv{v'}$, basta guardare se
$\nv{v'}$ ha un arco uscente etichettato $x$.

\begin{figure}[H]
\centering
% Passo ricorsivo del calcolo dei failure link in Aho--Corasick.
% A sinistra il cammino root ~> u ~> v' -> v; a destra il cammino root ~> n_{v'} -> n_v.
% I due tratti marcati portano la stessa stringa alfa = L(n_{v'}); i due archi x sono
% paralleli e di uguale lunghezza, cosi' la regola si legge a colpo d'occhio.
\begin{tikzpicture}[
    cammino/.style = {draw=graydark, line width=0.6pt, line cap=round, line join=round},
    marcato/.style = {draw=accent, line width=2.6pt, line cap=round, opacity=0.7},
    tacca/.style   = {draw=black, line width=1.1pt, line cap=round},
    prosegue/.style= {draw=graydark, line width=0.6pt, dotted, line cap=round},
    noto/.style    = {draw=black, line width=0.7pt, dash pattern=on 3.5pt off 2.5pt,
                      -{Triangle[length=5.5pt,width=4.5pt]}},
    cercato/.style = {draw=gray, line width=0.7pt, dash pattern=on 1.6pt off 2pt,
                      -{Triangle[length=5pt,width=4pt]}},
    et/.style      = {font=\footnotesize},
    pic/.style     = {font=\scriptsize},
  ]

  % ---- sagoma di K(P) ---------------------------------------------------
  \fill[black!5, draw=gray!55, line width=0.4pt]
      (0,5.6) .. controls (2.30,4.90) and (4.20,2.90) .. (5.0,0.15)
              .. controls (2.5,-0.35) and (-2.5,-0.35) .. (-5.0,0.15)
              .. controls (-4.20,2.90) and (-2.30,4.90) .. cycle;

  % ---- nodi -------------------------------------------------------------
  \coordinate (r)   at ( 0.000, 5.60);
  \coordinate (u)   at (-1.350, 3.55);
  \coordinate (vp)  at (-2.850, 1.55);
  \coordinate (v)   at (-3.470, 0.55);
  \coordinate (np)  at ( 1.325, 3.48);
  \coordinate (nv)  at ( 1.945, 2.48);

  % ---- tratti marcati (la stessa stringa alfa) --------------------------
  \draw[marcato] (u) -- (vp);
  \draw[marcato] (r) -- (np);

  % ---- cammino di sinistra ---------------------------------------------
  \draw[cammino] (r) -- (u) -- (vp) -- (v);
  \draw[prosegue] (v) -- (-3.80,0.02);

  % ---- cammino di destra ------------------------------------------------
  \draw[cammino] (r) -- (np) -- (nv);
  \draw[cammino] (np) -- (2.16,3.14);          % moncone dell'arco y
  \draw[prosegue] (nv) -- (2.26,1.98);

  % ---- tacche trasversali sui nodi --------------------------------------
  \draw[tacca] ($(u)!3pt!90:(vp)$)   -- ($(u)!3pt!-90:(vp)$);
  \draw[tacca] ($(vp)!4pt!90:(v)$)   -- ($(vp)!4pt!-90:(v)$);
  \draw[tacca] ($(v)!4pt!90:(vp)$)   -- ($(v)!4pt!-90:(vp)$);
  \draw[tacca] ($(np)!4pt!90:(nv)$)  -- ($(np)!4pt!-90:(nv)$);
  \draw[tacca] ($(nv)!4pt!90:(np)$)  -- ($(nv)!4pt!-90:(np)$);
  \fill[graydark] (r) circle (1.8pt);

  % ---- failure link -----------------------------------------------------
  \draw[noto]    (vp) -- (np)
      node[pic,pos=0.56,sloped,above=1pt] {\emph{failure link} (noto)};
  \draw[cercato] (v) -- (nv)
      node[pic,pos=0.50,sloped,below=1pt,text=gray] {\emph{failure link} cercato};

  % ---- etichette del cammino di sinistra --------------------------------
  \node[et,above=3pt] at (r) {$\mathrm{root}$};
  \node[et,anchor=east,xshift=-4pt] at (u)  {$u$};
  \node[et,anchor=east,xshift=-5pt,yshift=4pt] at (vp) {$v'$};
  \node[et,anchor=east,xshift=-5pt,yshift=-3pt] at (v) {$v$};
  \node[et,anchor=east] at (-3.44,1.14) {$x$};
  \node[et] at (-2.52,2.74) {$\a$};

  % ---- etichette del cammino di destra ----------------------------------
  \node[et,anchor=west] at (1.60,3.70) {$\nv{v'}$};
  \node[et,anchor=west,xshift=3pt,yshift=-2pt] at (nv) {$\nv{v}$};
  \node[et,anchor=east] at (1.52,2.84) {$x$};
  \node[et,anchor=west] at (2.24,3.06) {$y \neq x$};
  \node[et] at (1.06,4.76) {$\a$};

  \node[et] at (3.55,0.60) {$K(\mathcal{P})$};
\end{tikzpicture}

\caption{Passo ricorsivo del calcolo dei failure link. Il tratto marcato sul cammino di
sinistra (da $u$ a $v'$) e il cammino marcato di destra (da $\mathrm{root}$ a $\nv{v'}$)
portano la stessa stringa $\a$: è la relazione $\L(\nv{v'}) = \a$, il più lungo suffisso
proprio di $\L(v')$ che sia anche prefisso di un pattern. Se $\nv{v'}$ possiede un arco
uscente etichettato $x$, il nodo così raggiunto è $\nv{v}$.}
\label{fig:ac-failure-link-passo}
\end{figure}

\scan{15}
Svolgiamo la ricorsione, procedendo sulla profondità dei nodi.

\paragraph{Caso base.}
$v = \mathrm{root}$: banale, si pone $\nv{\mathrm{root}} = \mathrm{root}$. Banale è anche il
caso dei figli della radice: se $\mathrm{depth}(v) = 1$ allora $\L(v)$ è un singolo
carattere e il suo unico suffisso proprio è $\vuota = \L(\mathrm{root})$, dunque
$\nv{v} = \mathrm{root}$.

\paragraph{Passo induttivo.}
Sia ora $\mathrm{depth}(v) \ge 2$: i nodi di profondità $\le 1$ sono già coperti dal caso
base, e vanno esclusi qui perché per essi il ragionamento che segue darebbe la risposta
sbagliata (se $v' = \mathrm{root}$ allora $\nv{v'} = \mathrm{root}$, che un figlio raggiunto
dall'arco etichettato $x$ ce l'ha --- è $v$ stesso --- e si concluderebbe $\nv{v} = v$
anziché $\nv{v} = \mathrm{root}$).
Se devo determinare $\nv{v}$, sia $v' = \mathrm{parent}(v)$ e sia $x$ l'etichetta dell'arco
$\tuplev{v', v}$, cosicché $\L(v) = \L(v')\,x$. Poiché
\[
  \mathrm{depth}(v') < \mathrm{depth}(v)
  \quad\Longrightarrow\quad
  \text{ho già } \nv{v'},
\]
posso ragionare per casi:
\begin{itemize}
  \item se $\nv{v'}$ ha un figlio raggiunto da un arco etichettato $x$, allora $\nv{v}$ è
        quel figlio;
  \item altrimenti passo a $\nv{\nv{v'}}$ e determino se \emph{questo} ha un figlio
        etichettato $x$, e così via\dots
\end{itemize}
Si risale cioè la catena dei failure link
\[
  \nv{v'},\quad \nv{\nv{v'}},\quad \nv{\nv{\nv{v'}}},\quad \dots
\]
finché o si trova un nodo con un arco uscente etichettato $x$ --- e $\nv{v}$ è il figlio così
raggiunto --- oppure si arriva alla radice senza trovarlo, e in tal caso
$\nv{v} = \mathrm{root}$. La risalita termina sempre, perché lungo la catena la profondità
decresce strettamente: l'algoritmo non può andare in loop.

\paragraph{Ordine di elaborazione.}
Perché la ricorsione sia ben fondata i nodi vanno elaborati per profondità crescente, cioè
con una visita per livelli: quando si calcola $\nv{v}$ il valore $\nv{v'}$ del padre deve
essere già disponibile.

\begin{algorithm}[H]
\caption{Calcolo dei failure link di $K(\mathcal{P})$}
\label{alg:failure-link}
\begin{algorithmic}[1]
\Require il keyword tree $K(\mathcal{P})$
\Ensure $\nv{v}$ e $\mathit{lp}(v)$ per ogni nodo $v$ di $K(\mathcal{P})$
\State $\nv{\mathrm{root}} \gets \mathrm{root}$
\State $\mathit{lp}(\mathrm{root}) \gets 0$
\For{ogni figlio $u$ della radice}
  \State $\nv{u} \gets \mathrm{root}$
  \State $\mathit{lp}(u) \gets 0$
\EndFor
\For{$d \gets 2, \dots, \text{altezza di } K(\mathcal{P})$}
  \For{ogni nodo $v$ con $\mathrm{depth}(v) = d$}
    \State $v' \gets \mathrm{parent}(v)$
    \State $x \gets$ etichetta dell'arco $\tuplev{v', v}$
    \State $u \gets \nv{v'}$
    \While{$u \neq \mathrm{root}$ e $u$ non ha archi uscenti etichettati $x$}
      \State $u \gets \nv{u}$
    \EndWhile
    \If{$u$ ha un arco uscente $\tuplev{u, u'}$ etichettato $x$}
      \State $\nv{v} \gets u'$
    \Else
      \State $\nv{v} \gets \mathrm{root}$
    \EndIf
    \State $\mathit{lp}(v) \gets \mathrm{depth}(\nv{v})$
  \EndFor
\EndFor
\end{algorithmic}
\end{algorithm}

\paragraph{Complessità.}
Il numero di nodi analizzati per determinare un singolo $\nv{v}$ è
$\le \mathrm{depth}(v)$: da qui, mediante un'\textbf{analisi ammortizzata}, si ottiene una
complessità \textbf{lineare}.

\begin{osservazione}
Il conto per singolo nodo darebbe soltanto il limite pessimistico
$\Oh\bigl(n \cdot \max_i n_i\bigr)$. L'analisi ammortizzata va invece fatta lungo un intero
cammino radice--foglia, cioè lungo un pattern $P_i$: usando come potenziale la profondità
del candidato, $\mathit{lp}(\cdot) \ge 0$, si osserva che passando da un nodo al figlio
successivo essa cresce al più di $1$ (si aggiunge un solo carattere), mentre ogni iterazione
della risalita lungo la catena dei failure link la fa calare di almeno $1$. Il numero
complessivo di risalite lungo il cammino è dunque $\le n_i$: al peggio si paga il doppio dei
passi. Sommando su tutti i pattern,
\[
  T_{\text{preprocessing}} \;=\; \Oh\Bigl(\sum_{i=1}^{z} n_i\Bigr) \;=\; \Oh(n),
\]
nel modello in cui l'accesso al figlio dato il carattere costa $\Oh(1)$.
\end{osservazione}

% RICOSTRUZIONE (in fase di cucitura, non presente nel quaderno): il manoscritto non
% formalizza mai la fase di ricerca.  Le posizioni l e c, l'invariante L(v) = T[l..c-1],
% l'aggiornamento l <- c - lp(v), v <- n_v e la complessita' O(|T|+k) sono ricostruiti da
% Gusfield sec. 3.4 e da L04_sintesi sec. 12.  Senza questa sezione il capitolo definirebbe
% i failure link senza mai usarli.  Da registrare in NOTE-EDITORIALI.md (Ricostruzioni).
\section{La fase di ricerca}
\label{sec:ricerca-ac}

Il preprocessing ha prodotto $K(\mathcal{P})$ con tutti i suoi failure link; la scansione di
$T$ procede ora mantenendo un nodo corrente $v$ di $K(\mathcal{P})$, una posizione corrente
$c$ in $T$ e la posizione $\ell$ in cui inizia l'allineamento corrente, con l'invariante
$\L(v) = \sub{T}{\ell}{c-1}$. Finché da $v$ esce un arco etichettato $T[c]$ si scende
lungo quell'arco e si incrementa $c$; se si raggiunge una foglia si riporta l'occorrenza del
pattern corrispondente. In caso di mismatch si segue il failure link, ponendo nell'ordine
\[
  \ell \;\gets\; c - \mathit{lp}(v), \qquad v \;\gets\; \nv{v},
\]
senza mai decrementare $c$ (se il blocco avviene già sulla radice si incrementa
semplicemente $c$): è esattamente per questo che i failure link sono stati introdotti.
Poiché $c$ non arretra mai e a ogni salto $\ell$ cresce strettamente, la scansione costa
$\Oh(\abs{T} + k)$.

\section{Scarico dell'ipotesi: gli output link}
\label{sec:output-link}

\scan{15}

\begin{domanda}
Che succede se una stringa è \emph{fattore} di un'altra stringa?
\end{domanda}

\begin{esempio}
Si consideri
\[
  \mathcal{P} = \set{\str{acatt},\ \str{ca}},
  \qquad
  T = \str{acatg}.
\]
Qui $\str{ca}$ è sottostringa di $\str{acatt}$ (posizioni $2$--$3$): l'ipotesi semplificativa
«nessun $P_i$ è sottostringa di un altro $P_j$» è violata.
\end{esempio}

\begin{figure}[H]
\centering
% Keyword tree di P = {acatt, ca}.  La scansione di T = acatg scende lungo il ramo
% sinistro (a, c, a, t) e si blocca sulla g: la foglia di ca non viene mai raggiunta
% e l'occorrenza di ca in T[2..3] va perduta.
\begin{tikzpicture}[
    arco/.style   = {draw=graydark, line width=0.55pt, line cap=round},
    nodo/.style   = {circle, fill=black, inner sep=0pt, minimum size=3.2pt},
    car/.style    = {font=\footnotesize\ttfamily, text=graydark},
    et/.style     = {font=\footnotesize},
    lieve/.style  = {font=\scriptsize\ttfamily, text=gray},
    scan/.style   = {draw=accent, line width=2.8pt, line cap=round, line join=round,
                     opacity=0.55},
    chev/.style   = {draw=accent!65!black, line width=0.7pt, line cap=round,
                     -{Stealth[length=4.5pt,width=3.4pt]}},
  ]

  % ---- nodi -------------------------------------------------------------
  \coordinate (r)  at ( 0.0,  0.0);
  \coordinate (u1) at ( 1.2, -1.1);
  \coordinate (u2) at ( 2.2, -2.2);
  \coordinate (v1) at (-1.0, -1.1);
  \coordinate (v2) at (-1.9, -2.2);
  \coordinate (v3) at (-2.8, -3.3);
  \coordinate (v4) at (-3.7, -4.4);
  \coordinate (v5) at (-4.6, -5.5);

  % ---- cammino di scansione (evidenziato) -------------------------------
  \draw[scan, rounded corners=8pt]
      (0.15,-0.35) -- (-0.55,-1.05) -- (-1.45,-2.15) -- (-2.35,-3.25) -- (-3.45,-4.55);
  \draw[accent!65!black, line width=0.9pt, line cap=round,
        -{Triangle[length=5pt,width=4.5pt]}] (-3.50,-4.55) -- (-3.50,-5.58);
  \foreach \p in {{-0.55,-1.05},{-1.45,-2.15},{-2.35,-3.25}}
      \draw[chev] ($(\p)+(51:0.22)$) -- ($(\p)+(231:0.16)$);

  % ---- archi ------------------------------------------------------------
  \draw[arco] (r) -- (v1) -- (v2) -- (v3) -- (v4) -- (v5);
  \draw[arco] (r) -- (u1) -- (u2);

  \foreach \n in {r,u1,u2,v1,v2,v3,v4,v5} \node[nodo] at (\n) {};

  % ---- etichette degli archi -------------------------------------------
  \node[car] at (-0.74,-0.34) {a};
  \node[car] at (-1.70,-1.45) {c};
  \node[car] at (-2.60,-2.55) {a};
  \node[car] at (-3.50,-3.65) {t};
  \node[car] at (-4.40,-4.75) {t};
  \node[car] at ( 0.82,-0.32) {c};
  \node[car] at ( 1.94,-1.44) {a};

  % ---- fine dei due pattern --------------------------------------------
  \node[lieve,anchor=north] at ($(v5)+(0,-0.12)$) {acatt};
  \node[lieve,anchor=west]  at ($(u2)+(0.16,-0.06)$) {ca};

  % ---- punto di fallimento ---------------------------------------------
  \draw[accent!65!black, line width=1.5pt, line cap=round]
      (-3.74,-4.88) -- (-3.26,-5.36)  (-3.74,-5.36) -- (-3.26,-4.88);

  \node[et, anchor=north west, align=left] at (-3.08,-4.94)
      {ho perso l'occorrenza\\ di \str{ca}!};
\end{tikzpicture}

\caption{Keyword tree di $\mathcal{P} = \set{\str{acatt}, \str{ca}}$. Scandendo
$T = \str{acatg}$ la discesa segue gli archi $\str{a}, \str{c}, \str{a}, \str{t}$ e si
arresta sul carattere $\str{g}$ ($\times$): la foglia di \str{ca} non viene mai visitata e
l'occorrenza di \str{ca} in $\sub{T}{2}{3}$ va perduta.}
\label{fig:ac-output-link}
\end{figure}

Percorrendo il cammino di \str{acatt} si passa «sopra» un'occorrenza di \str{ca} senza mai
visitare la foglia corrispondente: ho perso l'occorrenza di \str{ca}!

\paragraph{Soluzione.}
Aggiungo a ogni nodo un \emph{output link} (equivalentemente, una \emph{output list}) di
occorrenze da riportare, quelle dovute alle sottostringhe. Per ogni nodo $v$ si pone
\[
  \out(v) \;=\;
  \begin{cases}
    \nv{v} & \text{se } \L(\nv{v}) \in \mathcal{P},\\[2pt]
    \out(\nv{v}) & \text{altrimenti,}
  \end{cases}
\]
con $\out(v)$ indefinito se la catena raggiunge la radice senza incontrare la fine di un
pattern. Arrivati in $v$ durante la ricerca si percorre la catena degli output link
riportando tutte le occorrenze che terminano nella posizione corrente. Poiché $\out(v)$ si
calcola in $\Oh(1)$ per nodo \emph{durante} la stessa visita per livelli che calcola i
failure link (algoritmo~\ref{alg:failure-link}), il preprocessing resta $\Oh(n)$ e la ricerca
resta $\Oh(\abs{T} + k)$. Scaricata l'ipotesi, va anche marcato come «fine di un pattern»
ogni nodo, interno o foglia, la cui etichetta di cammino appartiene a $\mathcal{P}$.

\section{Bilancio}

Mettendo insieme le due fasi, l'algoritmo di Aho--Corasick costa
\[
  T_{\text{preprocessing}} = \Oh(n),
  \qquad
  T_{\text{ricerca}} = \Oh(\abs{T} + k),
  \qquad
  T_{\text{totale}} = \Oh\bigl(n + \abs{T} + k\bigr),
\]
da confrontare con l'$\Oh(z\,\abs{T} + n)$ della soluzione naive: il fattore $z$ è sparito e
il testo viene scandito una volta sola, che era esattamente la sfida posta nella
\autoref{sec:exact-set-matching}.

\chapter{Suffix trie, suffix tree e algoritmo di Ukkonen}

\section{Dai pattern al testo: il suffix trie}

\scan{16}
Il keyword tree $K(\mathcal{P})$ introdotto per Aho--Corasick è un albero costruito a
partire da una \emph{collezione di pattern} $\mathcal{P}$. Nulla però ci obbliga a prendere
come collezione un insieme di pattern dati dall'esterno.

\begin{osservazione}
Posso usare i \emph{suffissi del testo} come collezione di pattern.
\end{osservazione}

L'albero che si ottiene in questo modo, cioè
\[
  K\bigl(\set{\suf{T}{1},\ \suf{T}{2},\ \dots,\ \suf{T}{n}}\bigr),
  \qquad n = \abs{T},
\]
è quello che in letteratura si chiama \emph{suffix trie} di $T$. Il cambio di prospettiva è
sostanziale: non si tratta più di un pre-processing del \emph{pattern} --- come in KMP, in
Boyer--Moore o nell'algoritmo $Z$ --- ma di un \textbf{pre-processing del testo}. La versione
compattata di questa struttura, il \emph{suffix tree}, arriverà subito dopo.

\paragraph{La sentinella.}
Al testo si aggiunge in coda un carattere \emph{sentinella} $\$ \notin \S$, e si lavora su
$T\$$ anziché su $T$. Non è un dettaglio decorativo: la sentinella garantisce che nessun
suffisso sia prefisso di un altro, e quindi che nel trie ogni suffisso termini in una foglia
distinta (nei disegni che seguono tutte le foglie sono infatti marcate con $\$$).

\section{La più lunga sottostringa ripetuta}

\begin{domanda}[Sottostringa ripetuta più lunga]
Data una stringa $T$, determinare la più lunga sottostringa \emph{ripetuta} in $T$, cioè la
più lunga $\b$ che occorre in $T$ a partire da (almeno) due posizioni distinte.
\end{domanda}

\begin{figure}[H]
\centering
% La piu' lunga sottostringa ripetuta: due occorrenze di beta in T, sentinella $ in coda.
\begin{tikzpicture}[x=1cm,y=1cm]

  % la stringa T
  \draw[line width=1.1pt] (0,0) -- (10,0);
  \node[anchor=south] at (9.6,0.12) {$T$};

  % sentinella, vistosamente piu' grande del testo
  \node[anchor=west,inner sep=1pt] at (10.05,0) {\Large$\$$};

  % prima occorrenza di beta
  \draw[decorate,decoration={brace,mirror,amplitude=5pt}] (1.4,-0.12) -- (3.6,-0.12);
  \node[anchor=north] at (2.5,-0.42) {$\b$};

  % seconda occorrenza di beta (stessa lunghezza)
  \draw[decorate,decoration={brace,mirror,amplitude=5pt}] (5.0,-0.12) -- (7.2,-0.12);
  \node[anchor=north] at (6.1,-0.42) {$\b$};

\end{tikzpicture}

\caption{La più lunga sottostringa ripetuta: $\b$ occorre a partire da due posizioni
distinte di $T$. La stringa è terminata dalla sentinella $\$$.}
\label{fig:sottostringa-ripetuta}
\end{figure}

Posto $n = \abs{T}$, la storia delle complessità note per questo problema è la seguente.
\begin{itemize}
  \item Soluzione naive: $\Oh(n^{2})$.
  \item Knuth: $\Oh(n \log n)$, con un algoritmo divide et impera il cui costo $C(n)$
        soddisfa la ricorrenza
        \[
          C(n) \;=\; 2\,C\!\left(\frac{n}{2}\right) + \Oh(n).
        \]
  \item $\Oh(n)$: Knuth ipotizzava che non si potesse fare, ma \emph{una volta tanto ha
        sbagliato}.
\end{itemize}

La strada per il tempo lineare passa proprio per il keyword tree dei suffissi di $T\$$.

\begin{figure}[H]
\centering
% Keyword tree dei suffissi di T$ (suffix trie): schema generico, ogni foglia porta la
% sentinella $. Nessuna etichetta di carattere sugli archi (come nel manoscritto).
\begin{tikzpicture}[
    x=1cm,y=1cm,
    arco/.style={draw=black!65,line width=0.6pt},
    interno/.style={circle,fill=black!65,inner sep=0pt,minimum size=2.4pt},
    radice/.style={circle,fill=black,inner sep=0pt,minimum size=3.4pt},
    foglia/.style={inner sep=1pt,fill=white,font=\footnotesize}]

  % ---- coordinate ------------------------------------------------------
  \coordinate (r)   at ( 0.00, 0.00);   % radice

  \coordinate (la)  at ( 0.95,-0.80);   % foglia del suffisso "$"
  \coordinate (lb)  at ( 4.30,-3.30);   % foglia del suffisso piu' lungo (tutta T$)

  \coordinate (v1)  at (-2.30,-1.50);
  \coordinate (lc1) at (-3.60,-2.60);
  \coordinate (v1p) at (-2.50,-2.70);
  \coordinate (lc2) at (-3.20,-3.75);
  \coordinate (lc3) at (-1.90,-3.75);

  \coordinate (v2)  at (-0.80,-1.70);
  \coordinate (ld1) at (-1.65,-2.90);
  \coordinate (v2p) at (-0.40,-2.80);
  \coordinate (ld2) at (-0.90,-3.85);
  \coordinate (ld3) at ( 0.30,-3.85);

  % ---- archi -----------------------------------------------------------
  % (a) arco cortissimo: il suffisso costituito dalla sola sentinella
  \draw[arco] (r) to[bend left=8] (la);

  % (b) arco lunghissimo e curvo: il suffisso piu' lungo, ovvero tutto T$
  \draw[arco] (r) .. controls (2.10,-0.35) and (4.05,-1.55) .. (lb);

  % (c) ramo sinistro
  \draw[arco] (r) to[bend right=6] (v1);
  \draw[arco] (v1) -- (lc1);
  \draw[arco] (v1) -- (v1p);
  \draw[arco] (v1p) -- (lc2);
  \draw[arco] (v1p) -- (lc3);

  % (d) ramo centrale
  \draw[arco] (r) -- (v2);
  \draw[arco] (v2) -- (ld1);
  \draw[arco] (v2) -- (v2p);
  \draw[arco] (v2p) -- (ld2);
  \draw[arco] (v2p) -- (ld3);

  % ---- nodi ------------------------------------------------------------
  \node[radice]  at (r)   {};
  \node[interno] at (v1)  {};
  \node[interno] at (v1p) {};
  \node[interno] at (v2)  {};
  \node[interno] at (v2p) {};

  % ---- foglie: tutte marcate con la sentinella -------------------------
  \foreach \f in {la,lb,lc1,lc2,lc3,ld1,ld2,ld3}
    \node[foglia] at (\f) {$\$$};

\end{tikzpicture}

\caption{Il keyword tree dei suffissi di $T\$$ (suffix trie): ogni foglia corrisponde a un
suffisso e porta la sentinella $\$$.}
\label{fig:kwtree-suffissi}
\end{figure}

Se ho il keyword tree dei suffissi a disposizione, mi basta visitarlo e cercare il nodo
branching di profondità massima. Scriviamo $\cammino{v}$ per l'etichetta del cammino dalla
radice al nodo $v$, cioè per la stringa che quel cammino compita.

\begin{proposizione}[Sottostringa ripetuta più lunga]\label{prop:ripetuta}
\scan{17}
Sia $v$ un nodo branching di profondità massima del suffix trie di $T\$$. Allora
$\cammino{v}$ è la più lunga sottostringa ripetuta di $T$.
\end{proposizione}

\begin{dimostrazione}
Se $v$ è un nodo branching, ha almeno due figli e quindi almeno due foglie nel suo
sottoalbero, cioè almeno due suffissi distinti di cui $\cammino{v}$ è prefisso comune:
$\cammino{v}$ occorre allora in (almeno) due posizioni distinte di $T$, ed è ripetuta.
Viceversa, se $\b$ occorre in due posizioni distinte $i \neq j$, allora $\b$ è prefisso sia
del suffisso che comincia in $i$ sia di quello che comincia in $j$: i due cammini che dalla
radice compitano questi due suffissi condividono i primi $\abs{\b}$ archi e prima o poi si
separano, e il nodo in cui si separano è branching e ha profondità almeno $\abs{\b}$.
Massimizzare $\abs{\b}$ equivale dunque a massimizzare la profondità di un nodo branching.
La corrispondenza «nodi branching $\leftrightarrow$ sottostringhe ripetute» è esatta proprio
perché si lavora su $T\$$: con la sentinella ogni suffisso finisce in una foglia distinta,
mentre senza di essa un suffisso può essere prefisso di un altro (si pensi a $T = aa$) e la
sottostringa ripetuta $a$ non corrisponderebbe ad alcun nodo branching.
\end{dimostrazione}

Restano aperte due domande, che occupano il resto del capitolo:
\begin{enumerate}
  \item quanto è grande il keyword tree dei suffissi di $T$?
  \item come lo costruisco?
\end{enumerate}

\section{Notazione, definizione e dimensione del suffix trie}

\paragraph{Notazione.}
Il testo è $T = t_1 t_2 \cdots t_n$, di lunghezza $n = \abs{T}$, sull'alfabeto $\S$. Per ogni
$i \in \set{1,\dots,n}$ poniamo
\begin{align*}
  T^{\,i} &= \sub{T}{1}{i} = t_1 t_2 \cdots t_i,
          & &\text{$i$-esimo \emph{prefisso} di $T$},\\
  T_{i}   &= \sub{T}{i}{n} = t_i t_{i+1} \cdots t_n,
          & &\text{$i$-esimo \emph{suffisso} di $T$},
\end{align*}
con la convenzione $T^{\,0} = \vuota$, e indichiamo con
\[
  \mathcal{S}(T) \;=\; \set{\, T_i \;:\; i \in \set{1,\dots,n} \,}
\]
l'insieme di tutti i suffissi di $T$.

\begin{definizione}[Suffix trie]\label{def:strie}
Il \emph{suffix trie} di $T$ è il keyword tree dell'insieme dei suoi suffissi:
\[
  \STrie{}(T) \defeq K\bigl(\mathcal{S}(T)\bigr).
\]
\end{definizione}

Il termine \emph{trie}, per inciso, viene dall'inglese re\emph{trie}val.

Affrontiamo il primo dei due problemi.

\begin{domanda}
È possibile trovare esempi di $T$ tali che $\abs{\STrie{}(T)} \in \Th\bigl(\abs{T}^2\bigr)$?
\end{domanda}

Sì: basta prendere $T = a^m b^m$, e dunque $\abs{T} = 2m$.

\begin{figure}[H]
\centering
% Suffix trie di T = a^m b^m: spina dorsale di m archi 'a', da ogni suo nodo pende una
% catena di m archi 'b'  ==>  Theta(m^2) nodi.
\begin{tikzpicture}[
    x=0.95cm,y=0.95cm,
    arco/.style={draw=black,line width=0.5pt},
    nd/.style={circle,fill=black,inner sep=0pt,minimum size=2.6pt},
    rt/.style={circle,fill=black,inner sep=0pt,minimum size=3.8pt},
    et/.style={font=\scriptsize,inner sep=1.2pt},
    pnt/.style={circle,fill=black,inner sep=0pt,minimum size=1.5pt}]

  % =====================================================================
  % radice
  % =====================================================================
  \coordinate (r) at (0,0);

  % =====================================================================
  % catena dei suffissi b^q che pende dalla radice (verso destra-basso)
  % =====================================================================
  \coordinate (u1) at (0.85,-0.30);
  \coordinate (u2) at (1.70,-0.60);
  \coordinate (u3) at (2.55,-0.90);
  \draw[arco] (r)  -- (u1) node[et,midway,anchor=south]{$b$};
  \draw[arco] (u1) -- (u2) node[et,midway,anchor=south]{$b$};
  \draw[arco] (u2) -- (u3) node[et,midway,anchor=south]{$b$};
  \foreach \k in {1,2,3} \node[pnt] at ($(u3)+\k*(0.28,-0.10)$) {};
  \node[nd] at (u1) {}; \node[nd] at (u2) {}; \node[nd] at (u3) {};

  % =====================================================================
  % spina dorsale degli a (quasi verticale, m archi in tutto)
  % =====================================================================
  \coordinate (v1) at (-0.22,-0.95);
  \coordinate (v2) at (-0.44,-1.90);
  \coordinate (v3) at (-0.66,-2.85);
  \coordinate (v4) at (-0.88,-3.80);
  \coordinate (vm) at (-1.15,-5.35);

  \draw[arco] (r)  -- (v1) node[et,midway,anchor=east,xshift=-1pt]{$a$};
  \draw[arco] (v1) -- (v2) node[et,midway,anchor=east,xshift=-1pt]{$a$};
  \draw[arco] (v2) -- (v3) node[et,midway,anchor=east,xshift=-1pt]{$a$};
  \draw[arco] (v3) -- (v4) node[et,midway,anchor=east,xshift=-1pt]{$a$};
  % ellissi verticale sulla spina
  \draw[arco] (v4) -- (-0.93,-4.05);
  \foreach \k in {1,2,3} \node[pnt] at ($(-0.93,-4.05)+\k*(-0.045,-0.27)$) {};
  \draw[arco] (-1.09,-5.10) -- (vm);

  \node[nd] at (v1) {}; \node[nd] at (v2) {}; \node[nd] at (v3) {};
  \node[nd] at (v4) {}; \node[nd] at (vm) {};
  \node[rt] at (r) {};

  % =====================================================================
  % catene di b pendenti da ciascun nodo della spina (a 45 gradi)
  % di ognuna si disegnano i primi due archi, poi l'ellissi
  % =====================================================================
  \foreach \v in {v1,v2,v3,v4}{%
    \coordinate (p1) at ($(\v)+(0.62,-0.62)$);
    \coordinate (p2) at ($(\v)+(1.24,-1.24)$);
    \draw[arco] (\v) -- (p1) node[et,midway,anchor=south west]{$b$};
    \draw[arco] (p1) -- (p2) node[et,midway,anchor=south west]{$b$};
    \node[nd] at (p1) {}; \node[nd] at (p2) {};
    \foreach \k in {1,2,3} \node[pnt] at ($(p2)+\k*(0.24,-0.24)$) {};
  }

  % =====================================================================
  % ultima catena, quella appesa a v_m: compita l'intera stringa a^m b^m
  % =====================================================================
  \coordinate (w1) at (-0.93,-5.97);
  \coordinate (w2) at (-0.71,-6.59);
  \coordinate (w3) at (-0.49,-7.21);
  \draw[arco] (vm) -- (w1) node[et,midway,anchor=west]{$b$};
  \draw[arco] (w1) -- (w2) node[et,midway,anchor=west]{$b$};
  \draw[arco] (w2) -- (w3) node[et,midway,anchor=west]{$b$};
  \node[nd] at (w1) {}; \node[nd] at (w2) {}; \node[nd] at (w3) {};
  \foreach \k in {1,2,3} \node[pnt] at ($(w3)+\k*(0.09,-0.26)$) {};

  % =====================================================================
  % graffe di misura
  % =====================================================================
  % la spina e' fatta di m archi 'a'
  \draw[decorate,decoration={brace,amplitude=5pt}] (-1.75,-5.25) -- (-1.75,0);
  \node[font=\small] at (-2.20,-2.62) {$m$};
  % l'ultima catena e' fatta di m archi 'b'
  \draw[decorate,decoration={brace,amplitude=5pt}] (-1.75,-7.30) -- (-1.75,-5.55);
  \node[font=\small] at (-2.20,-6.42) {$m$};

\end{tikzpicture}

\caption{Il suffix trie di $T = a^m b^m$. La «spina dorsale» degli $a$ è lunga $m$ e da
ciascuno dei suoi nodi pende una catena di $m$ archi etichettati $b$: in tutto
$\Th(m^2) = \Th(\abs{T}^2)$ nodi.}
\label{fig:suffix-trie-ambm}
\end{figure}

Le sottostringhe di $a^m b^m$ sono infatti tutte e sole le stringhe della forma $a^p b^q$ con
$0 \le p \le m$ e $0 \le q \le m$, a due a due distinte; poiché i nodi del suffix trie di $T$
sono in biiezione con le sottostringhe distinte di $T$ --- ogni nodo è individuato dalla
stringa che il cammino dalla radice compita, e tali stringhe sono esattamente i prefissi dei
suffissi di $T$, cioè le sottostringhe di $T$ (alla radice corrisponde la stringa vuota
$\vuota$) --- il suffix trie di $T$ ha esattamente
\[
  (m+1)^2 \;=\; \frac{\abs{T}^2}{4} + \abs{T} + 1 \;=\; \Th\bigl(\abs{T}^2\bigr)
\]
nodi. Nel disegno: dalla radice pende la catena dei suffissi $b^q$, lunga $m$, e da ciascuno
dei nodi $a^p$ della spina dorsale degli $a$ pende una catena di $m$ archi etichettati $b$,
la cui foglia rappresenta il suffisso $a^p b^m$.

Il suffix trie è dunque inutilizzabile come indice: lo spazio è quadratico. È questa la
motivazione della compattazione che segue.

\section{Dal suffix trie al suffix tree}

\scan{18}
Guardiamo da vicino un frammento di suffix trie: lungo i rami si formano catene di nodi con
un solo figlio, ed è la proliferazione di questi nodi a rendere quadratica la dimensione
della struttura. L'idea è contrarre ogni catena in un unico arco, etichettato dalla stringa
che la catena compita: nel disegno, \str{abb}.

\begin{figure}[H]
\centering
% Contrazione di una catena unaria in un frammento di suffix trie:
% i nodi v1, v2 (un solo figlio) spariscono e il cammino abb diventa un unico arco [k,p].
\begin{tikzpicture}[
    x=0.92cm,y=0.92cm,
    arco/.style={draw=black,line width=0.5pt},
    nd/.style={circle,fill=black,inner sep=0pt,minimum size=2.8pt},
    ndg/.style={circle,fill=black,inner sep=0pt,minimum size=3.6pt},
    vuoto/.style={circle,draw=black,fill=white,line width=0.5pt,
                  inner sep=0pt,minimum size=4pt},
    et/.style={font=\footnotesize,inner sep=1.5pt},
    ann/.style={draw=accent,line width=1.1pt}]

  % =====================================================================
  % LIVELLO A -- il frammento di trie
  % =====================================================================
  \coordinate (u)  at ( 0.00, 0.00);
  \coordinate (x)  at (-1.30,-1.20);   % altro figlio di u (non interessa)
  \coordinate (v1) at ( 1.15,-1.15);
  \coordinate (v2) at ( 2.30,-2.30);
  \coordinate (w)  at ( 3.45,-3.45);
  \coordinate (f1) at ( 2.85,-4.55);
  \coordinate (f2) at ( 4.15,-4.55);

  % l'albero prosegue verso la radice
  \draw[arco] (0,1.90)
      .. controls ( 0.22, 1.72) and (-0.22, 1.56) .. (0,1.40)
      .. controls ( 0.22, 1.22) and (-0.22, 1.06) .. (0,0.90)
      .. controls ( 0.22, 0.72) and (-0.22, 0.56) .. (0,0.40)
      -- (u);

  \draw[arco] (u)  -- (x);
  \draw[arco] (u)  -- (v1) node[et,midway,anchor=south west]{$a$};
  \draw[arco] (v1) -- (v2) node[et,midway,anchor=south west]{$b$};
  \draw[arco] (v2) -- (w)  node[et,midway,anchor=south west]{$b$};
  \draw[arco] (w)  -- (f1);
  \draw[arco] (w)  -- (f2);

  \node[ndg] at (u)  {};
  \node[vuoto] at (x) {};
  \node[nd] at (v1) {};
  \node[nd] at (v2) {};
  \node[ndg] at (w)  {};
  \node[nd] at (f1) {};
  \node[nd] at (f2) {};

  % =====================================================================
  % LIVELLO B -- l'annotazione della compattazione
  % =====================================================================
  % cancellatura dei due nodi con un solo figlio
  \foreach \v in {v1,v2}{%
    \draw[ann] ($(\v)+(-0.38,-0.13)$) -- ($(\v)+(0.38,0.13)$);
    \draw[ann] ($(\v)+(-0.20,0.32)$)  -- ($(\v)+(0.20,-0.32)$);
  }

  % nuovo arco compattato u --> w, che scavalca la catena passando a sinistra
  \draw[ann] (u) .. controls (-0.90,-1.80) and (1.20,-4.30) .. (w);

  % etichetta dell'arco compattato
  \node[accent,align=center,anchor=east,font=\footnotesize] at (-0.35,-2.55)
        {\str{abb}\\[1pt]$[k,p]$};

\end{tikzpicture}

\caption{Contrazione di una catena unaria in un frammento di suffix trie: i due nodi
con un solo figlio vengono cancellati e il cammino \str{abb} diventa un unico arco, che verrà
etichettato con la coppia di indici $[k,p]$.}
\label{fig:compattazione-abb}
\end{figure}

Per garantire la linearità della struttura dati modifichiamo dunque il suffix trie in
\emph{suffix tree}. Servono due operazioni.

\paragraph{1) Eliminazione dei cammini unari.}
Eliminiamo i cammini massimali di nodi con un solo figlio: se
\[
  x_1 \xrightarrow{\;a_1\;} x_2 \xrightarrow{\;a_2\;} \cdots
       \xrightarrow{\;a_{h-1}\;} x_h
\]
è un cammino i cui nodi interni $x_2,\dots,x_{h-1}$ hanno esattamente un figlio, ed è
massimale (né $x_1$ né $x_h$ sono nodi unari), lo sostituiamo con il solo arco $(x_1,x_h)$.

\begin{figure}[H]
\centering
\begin{tikzpicture}[
    x=1cm, y=1cm, font=\footnotesize,
    vuoto/.style={circle,draw=black,line width=0.5pt,inner sep=0pt,minimum size=4.5pt},
    pieno/.style={circle,fill=black,inner sep=0pt,minimum size=3.6pt},
    arco/.style={-{Latex[length=4.5pt,width=3.2pt]},line width=0.6pt,
                 shorten >=2pt,shorten <=2pt}]

  %% ---------- pannello sinistro: la catena unaria x_1 ... x_h ----------
  \node[vuoto] (x1) at (0    , 0    ) {};
  \node[pieno] (x2) at (1    ,-0.95 ) {};
  \node[pieno] (x3) at (2    ,-1.90 ) {};
  \node[pieno] (xm) at (3.2  ,-3.04 ) {};
  \node[vuoto] (xh) at (4.2  ,-3.99 ) {};

  \draw[arco] (x1) -- node[left=1pt] {$a_1$}     (x2);
  \draw[arco] (x2) -- node[left=1pt] {$a_2$}     (x3);
  \draw[arco] (xm) -- node[left=1pt] {$a_{h-1}$} (xh);

  % i nodi omessi: tre puntini allineati sulla stessa diagonale
  \fill (2.36,-2.242) circle (0.7pt);
  \fill (2.60,-2.470) circle (0.7pt);
  \fill (2.84,-2.698) circle (0.7pt);

  \node[above right=-2pt] at (x1) {$x_1$};
  \node[right=2pt]        at (x2) {$x_2$};
  \node[right=2pt]        at (x3) {$x_3$};
  \node[right=2pt]        at (xm) {$x_{h-1}$};
  \node[right=2pt]        at (xh) {$x_h$};

  %% ---------- separatore ----------
  \draw[line width=0.45pt,double distance=1.5pt,-{Implies[length=4pt,width=5pt]}]
        (5.35,-2.0) -- (6.75,-2.0);

  %% ---------- pannello destro: l'unico arco (x_1,x_h) ----------
  \node[vuoto] (y1) at (7.6 ,-1.10) {};
  \node[vuoto] (yh) at (8.8 ,-2.24) {};
  \draw[arco] (y1) -- (yh);

  \node[above right=-2pt] at (y1) {$x_1$};
  \node[right=2pt]        at (yh) {$x_h$};

\end{tikzpicture}

\caption{Operazione 1: la catena massimale di nodi con un solo figlio $x_1,\dots,x_h$ viene
sostituita dall'unico arco $(x_1,x_h)$.}
\label{fig:compattazione-catena}
\end{figure}

\paragraph{2) Etichette come coppie di indici.}
Etichetto l'arco $(x_1,x_h)$ così introdotto con una coppia $[k,p]$ di indici nel testo tale
che
\[
  \sub{T}{k}{p} \;=\; a_1 a_2 \cdots a_{h-1}.
\]

\begin{osservazione}
Di coppie $[k,p]$ che individuano quella sottostringa ce ne può essere più d'una --- tante
quante le occorrenze di $a_1 a_2 \cdots a_{h-1}$ in $T$ --- e la scelta fra esse è del tutto
indifferente: quel che conta è la stringa compitata, non la posizione da cui la si legge.
\end{osservazione}

\begin{osservazione}
Affinché 2) funzioni dobbiamo tenere in memoria anche il testo: l'etichetta non contiene più
i caratteri, ma soltanto un puntatore alla porzione di $T$ che li contiene. La struttura dati
finale è dunque il trie compattato \emph{più} il testo $T$.
\end{osservazione}

\subsection{Linearità}

\begin{lemma}[Alberi con soli nodi branching]\label{lem:branching}
Sia $\T$ un albero finito con $\ell$ foglie e $I$ nodi interni, in cui ogni nodo interno è
\emph{branching}, cioè ha almeno due figli. Allora
\[
  I \;\le\; \ell - 1,
  \qquad\text{e dunque}\qquad
  I + \ell \;\le\; 2\ell - 1 ,
\]
cioè $\T$ ha al più $2\ell-1$ nodi. Vale l'uguaglianza se e solo se ogni nodo interno ha
esattamente due figli, cioè se e solo se $\T$ è binario.
\end{lemma}

\begin{dimostrazione}
In un albero gli archi sono $I + \ell - 1$ (uno per nodo, tranne la radice); d'altra parte
ogni arco esce da esattamente un nodo interno e ogni nodo interno ha almeno due archi
uscenti, dunque
\[
  I + \ell - 1 \;=\; \sum_{u \text{ interno}} \deg^{+}(u) \;\ge\; 2I ,
  \qquad\text{cioè}\qquad I \;\le\; \ell - 1 .
\]
Sommando le foglie, il numero di nodi è $I + \ell \le 2\ell - 1$. L'uguaglianza vale
esattamente quando $\deg^{+}(u) = 2$ per ogni nodo interno $u$.
\end{dimostrazione}

\begin{figure}[H]
\centering
% Il lemma "ogni nodo interno e' branching" a colpo d'occhio.
% Stesso albero (4 foglie, 3 nodi interni) disegnato due volte:
%  a sinistra l'addebito nodo interno -> foglia, a destra la contrazione di una ciliegia.
\begin{tikzpicture}[
    x=1cm, y=1cm, font=\footnotesize,
    arco/.style={draw=black,line width=0.5pt},
    nodo/.style={circle,draw=black,fill=white,line width=0.5pt,
                 inner sep=0pt,minimum size=4.5pt},
    addebito/.style={draw=graydark,line width=0.5pt,densely dashed,
                     -{Latex[length=4pt,width=3pt]}}]

  %% ===================================================================
  %% PANNELLO SINISTRO -- ogni nodo interno addebitato a una foglia
  %% ===================================================================
  \begin{scope}[shift={(0,0)}]
    \coordinate (r)  at (1.50, 2.00);
    \coordinate (u)  at (0.75, 1.00);
    \coordinate (w)  at (2.25, 1.00);
    \coordinate (f1) at (0.25, 0.00);
    \coordinate (f2) at (1.25, 0.00);
    \coordinate (f3) at (1.75, 0.00);
    \coordinate (f4) at (2.75, 0.00);

    \draw[arco] (r) -- (u);  \draw[arco] (r) -- (w);
    \draw[arco] (u) -- (f1); \draw[arco] (u) -- (f2);
    \draw[arco] (w) -- (f3); \draw[arco] (w) -- (f4);

    % l'addebito: da ogni nodo interno si scende a una foglia del suo sottoalbero
    \draw[addebito] (1.50,1.82) -- (1.50,-0.30);
    \draw[addebito] (0.75,0.82) -- (0.75,-0.30);
    \draw[addebito] (2.25,0.82) -- (2.25,-0.30);

    \foreach \n in {r,u,w,f1,f2,f3,f4} \node[nodo] at (\n) {};

    \node[anchor=north,align=center] at (1.50,-0.62)
          {addebito nodo interno $\to$ foglia\\[1pt] $I \le \ell$};
  \end{scope}

  %% ===================================================================
  %% PANNELLO DESTRO -- contrazione di una ciliegia
  %% ===================================================================
  \begin{scope}[shift={(5.60,0)}]
    \coordinate (R)  at (1.50, 2.00);
    \coordinate (U)  at (0.75, 1.00);
    \coordinate (W)  at (2.25, 1.00);
    \coordinate (g1) at (0.25, 0.00);
    \coordinate (g2) at (1.25, 0.00);
    \coordinate (g3) at (1.75, 0.00);
    \coordinate (g4) at (2.75, 0.00);

    \draw[arco] (R) -- (U);  \draw[arco] (R) -- (W);
    \draw[arco] (U) -- (g1); \draw[arco] (U) -- (g2);
    \draw[arco] (W) -- (g3); \draw[arco] (W) -- (g4);

    % la ciliegia: nodo branching i cui figli sono tutte foglie
    \draw[draw=accent,line width=1.1pt]
      plot[smooth cycle,tension=0.75] coordinates
      {(0.75,1.40) (1.36,0.52) (1.42,-0.34) (0.06,-0.42) (-0.20,0.45)};

    \foreach \n in {R,U,W,g1,g2,g3,g4} \node[nodo] at (\n) {};

    \node[anchor=north,align=center] at (1.50,-0.62)
          {contrazione di una ciliegia\\[1pt] $(I,\ell) \to (I-1,\ell-1)$};
  \end{scope}

\end{tikzpicture}

\caption{Il Lemma a colpo d'occhio, su un albero con $4$ foglie e $3$ nodi interni
($7 = 2\cdot 4 - 1$ nodi). A sinistra: ogni nodo interno viene addebitato a una foglia
distinta del proprio sottoalbero, da cui $I \le \ell$. A destra: la contrazione di una
\emph{ciliegia} (nodo branching i cui figli sono tutti foglie, cerchiata) toglie un nodo
interno e una foglia, e permette di concludere per induzione.}
\label{fig:lemma-branching}
\end{figure}

Il Lemma, insieme alle operazioni 1) e 2), garantisce la linearità del suffix tree. Dopo la
compattazione, infatti, ogni nodo interno è branching e le foglie sono una per suffisso,
cioè $n+1$ lavorando su $T\$$ e comunque $\Oh(n)$; quindi
\[
  \#\text{nodi} \;\le\; 2(n+1) - 1 \;=\; \Oh(n),
\]
gli archi sono uno di meno, e ciascuno di essi porta due soli interi. Sommando il testo, che
va comunque tenuto in memoria, lo spazio complessivo è $\Oh(n)$.

\section{A che cosa serve il suffix tree}

\scan{19}
Una volta costruito l'albero dei suffissi del testo, la ricerca di un pattern $P$ si riduce a
compitare $P$ lungo l'unico cammino che parte dalla radice: se il cammino esiste e termina
nel nodo (eventualmente implicito) $v$, le foglie del sottoalbero $\albero{v}$ radicato in
$v$ sono esattamente le occorrenze di $P$ in $T$.

\begin{figure}[H]
\centering
% Ricerca esatta col suffix tree: si compita P dalla radice, si arriva al nodo v,
% e le foglie del sottoalbero T(v) sono le occorrenze di P in T.
% La "montagna" e' il suffix tree dell'intero testo, disegnato come sagoma.
\begin{tikzpicture}[
    x=0.86cm, y=0.86cm, font=\footnotesize,
    profilo/.style={draw=black,line width=0.7pt,line join=round},
    sotto/.style={draw=black,line width=0.4pt,line join=round},
    barra/.style={draw=black,line width=1.1pt,line cap=round},
    foglia/.style={circle,fill=black,inner sep=0pt,minimum size=3.2pt}]

  %% ===================================================================
  %% 1) promemoria: P e' corto, T e' lungo
  %% ===================================================================
  \draw[barra] (0.90,2.55) -- (2.10,2.55);
  \draw[barra] (0.90,1.80) -- (4.10,1.80);
  \node[anchor=south west,inner sep=1pt] at (0.90,2.61) {$P$};
  \node[anchor=south west,inner sep=1pt] at (0.90,1.86) {$T$};

  %% ===================================================================
  %% 2) la montagna: l'intero suffix tree di T
  %% ===================================================================
  \draw[profilo]                                  % versante sinistro
     (7.50, 3.40) -- (7.15, 2.95) -- (7.05, 2.60) -- (6.58, 2.10)
  -- (6.45, 1.70) -- (5.99, 1.20) -- (5.85, 0.80) -- (5.40, 0.30)
  -- (5.25,-0.10) -- (4.81,-0.60) -- (4.65,-1.00) -- (4.32,-1.35)
  -- (4.20,-1.60);
  \draw[profilo]                                  % versante destro
     (7.50, 3.40) -- (7.85, 2.95) -- (7.95, 2.60) -- (8.42, 2.10)
  -- (8.55, 1.70) -- (9.01, 1.20) -- (9.15, 0.80) -- (9.60, 0.30)
  -- (9.75,-0.10) -- (10.19,-0.60) -- (10.35,-1.00) -- (10.68,-1.35)
  -- (10.80,-1.60);
  \draw[draw=graylight,line width=0.4pt] (4.20,-1.60) -- (10.80,-1.60);

  % la radice, in cima
  \node[foglia] at (7.50,3.40) {};
  \node[anchor=south,inner sep=2pt] at (7.50,3.44) {radice};

  %% ===================================================================
  %% 3) il cammino che compita P, dalla radice fino a v
  %% ===================================================================
  \draw[draw=black,line width=1.1pt,line join=round]
     (7.50,3.40) -- (7.44,2.90) -- (7.53,2.45) -- (7.38,1.95)
  -- (7.48,1.45) -- (7.40,0.90);
  \node[anchor=west,inner sep=2pt,font=\normalsize] at (7.62,2.20) {$P$};

  %% ===================================================================
  %% 4) il nodo v e il sottoalbero T(v) che vi e' appeso
  %% ===================================================================
  \draw[sotto]
     (7.40,0.90) -- (7.15,0.35) -- (7.22,0.00) -- (6.72,-0.60)
  -- (6.62,-1.00) -- (6.00,-1.40);
  \draw[sotto]
     (7.40,0.90) -- (7.68,0.40) -- (7.62,0.00) -- (8.12,-0.55)
  -- (8.22,-0.95) -- (8.90,-1.40);
  % base del sottoalbero: linea ondulata, su cui stanno le foglie
  \draw[sotto,decorate,decoration={snake,amplitude=0.6mm,segment length=5mm}]
     (6.00,-1.40) -- (8.90,-1.40);

  \node[foglia] at (6.55,-1.40) {};
  \node[foglia] at (8.30,-1.40) {};

  \node[circle,fill=black,inner sep=0pt,minimum size=4pt] at (7.40,0.90) {};
  \node[anchor=east,inner sep=3pt] at (7.40,0.90) {$v$};

  \node[anchor=west,inner sep=1pt] (etv) at (9.10,-0.62) {$\albero{v}$};
  \draw[-{Latex[length=3.5pt,width=2.5pt]},line width=0.4pt]
        (etv.west) -- (8.45,-0.88);

  \node[anchor=north,text=graydark] at (7.45,-1.95)
        {occorrenze di $P$ in $T$};

\end{tikzpicture}

\caption{Uso del suffix tree per la ricerca esatta: il cammino che compita $P$ dalla radice
termina nel nodo $v$; le foglie del sottoalbero $\albero{v}$ appeso a $v$ sono le occorrenze
di $P$ in $T$.}
\label{fig:suffix-tree-ricerca-montagna}
\end{figure}

\begin{osservazione}
A lato del disegno gli appunti richiamano Boyer--Moore («e anche Boyer--Moore»), senza
ulteriori dettagli.
\end{osservazione}

\section{L'algoritmo di Ukkonen}

\scan{19}
\begin{domanda}
Come costruisco il suffix tree?
\end{domanda}

\begin{osservazione}
Soluzioni al problema esistono già dagli anni Settanta --- Weiner (1973) e McCreight (1976)
--- ma quella di Ukkonen è molto più chiara ed è quella che tutti citano.
\end{osservazione}

L'algoritmo di Ukkonen (1995), presentato nell'articolo \emph{On-line construction of suffix
trees}, è \emph{on-line} sul testo $T = t_1 t_2 \cdots t_n$: ne legge un carattere alla volta,
da sinistra a destra, e mantiene in ogni istante la struttura relativa al prefisso già letto.
Il passo generico trasforma la struttura del prefisso $T^{\,i-1}$ in quella del prefisso
$T^{\,i}$ leggendo il carattere $t_i$:
\[
  \STrie{}(T^{\,i-1}) \;\xrightarrow{\ t_i\ }\; \STrie{}(T^{\,i}).
\]

\begin{figure}[H]
\centering
% Un passo di Ukkonen: dal prefisso T^{i-1} (gia' processato) a T^i, leggendo t_i.
% La riga superiore ricorda un generico suffisso T_j che attraversa il passo.
\begin{tikzpicture}[
    x=0.84cm, y=0.84cm, font=\small,
    testo/.style={draw=black,line width=1.1pt,line cap=round},
    tacca/.style={draw=black,line width=0.5pt,line cap=round}]

  % la cella del carattere corrente t_i, sullo sfondo
  \fill[black!14] (5.40,-0.10) rectangle (5.90,0.55);

  %% ---------------- le due barre --------------------------------------
  \draw[testo] (0,0)    -- (12,0);       % il testo T
  \draw[testo] (3,0.45) -- (12,0.45);    % il suffisso T_j
  \node[anchor=west,inner sep=4pt] at (12,0) {$T$};

  %% ---------------- tacche --------------------------------------------
  \draw[tacca] (0,-0.20) -- (0,0.25);                    % inizio del testo
  \draw[tacca] (3,-0.28) -- (3,0.66);                    % posizione j
  \draw[tacca] (5.40,-0.30) -- (5.40,1.00);              % coppia che
  \draw[tacca] (5.90,-0.30) -- (5.90,1.00);              % isola t_i

  \node[anchor=south,inner sep=2pt] at (3,0.68) {$T_j$};
  \node[anchor=north,inner sep=2pt] at (3,-0.30) {$j$};

  \node[anchor=west,inner sep=1pt] (lti) at (6.55,1.28) {$t_i$};
  \draw[-{Latex[length=4pt,width=3pt]},line width=0.5pt]
        (lti.west) -- (5.66,1.06);

  %% ---------------- graffa sul prefisso gia' letto ---------------------
  \draw[line width=0.5pt,decorate,decoration={brace,mirror,amplitude=5pt}]
        (0,-0.95) -- (5.40,-0.95);

  % collegamento fra la tacca in i e la riga in basso
  \draw[draw=graylight,line width=0.4pt,densely dashed]
        (5.40,-0.34) -- (5.40,-1.72);

  %% ---------------- la riga del passo ---------------------------------
  \node[anchor=center] at (2.70,-1.95) {$\STrie{}(T^{\,i-1})$};
  \draw[-{Latex[length=4.5pt,width=3.2pt]},line width=0.5pt]
        (5.20,-1.95) -- (6.40,-1.95);
  \node[anchor=south,inner sep=2pt] at (5.80,-1.88) {$t_i$};
  \node[anchor=west] at (6.60,-1.95) {$\STrie{}(T^{\,i})$};

\end{tikzpicture}

\caption{Un passo dell'algoritmo di Ukkonen. Il prefisso $T^{\,i-1}$ è già stato processato;
leggendo il carattere $t_i$ si passa a $\STrie{}(T^{\,i})$. La riga superiore evidenzia un
generico suffisso $T_j$, che nel passo corrente contribuisce con la sua parte
$t_j \cdots t_{i-1}$.}
\label{fig:ukkonen-passo-prefissi}
\end{figure}

La strategia viene esposta in due tempi:
\begin{itemize}
  \item[(a)] si illustra la strategia nel caso dei \emph{suffix trie};
  \item[(b)] si modifica la strategia e la si adatta ai \emph{suffix tree}.
\end{itemize}
In entrambi i casi la struttura su cui si lavora è la versione \emph{aumentata}
(\emph{augmented}) della struttura dati --- l'automa a stati finiti con lo stato ausiliario
$\bot$ e la funzione dei suffix link --- che definiamo subito.

Costruiremo dunque $\STrie{}(T^{\,i})$ per $i = 1,\dots,n$ --- \emph{in quest'ordine!} --- a
partire da $\STrie{}(T^{\,i-1})$, appendendo $t_i$ sotto ogni suffisso di $T^{\,i-1}$; il caso
base è $\STrie{}(T^{\,0})$, il trie della stringa vuota, cioè il solo nodo radice.

\subsection{Il suffix trie come automa aumentato}

\scan{20}
\begin{definizione}[Suffix trie come automa aumentato]\label{def:strie-automa}
Dato il testo $T$, il \emph{suffix trie} di $T$ è la quintupla
\[
  \STrie{}(T)\;=\;\bigl(\,Q\cup\set{\bot},\;\mathit{root},\;F,\;g,\;f\,\bigr),
\]
dove, scrivendo $\bar{x}$ per lo stato associato alla stringa $x$,
\begin{align*}
  Q             &= \set{\bar{x} \;:\; x \text{ sottostringa di } T}
                   &&\text{insieme degli stati,}\\
  \mathit{root} &= \bar{\vuota} &&\text{stato iniziale,}\\
  F             &= \set{\bar{x} \;:\; x \text{ suffisso di } T} &&\text{stati finali,}\\
  g(\bar{x},a)  &= \overline{xa} \ \text{ se } xa \text{ è sottostringa di } T
                   \ \text{(altrimenti indefinita)}, &&\text{funzione di transizione}\\
                &\phantom{{}={}} \qquad g(\bot,a)=\mathit{root}
                   \ \text{ per ogni } a\in\S &&\text{(parziale),}\\
  f(\bar{x})    &= \bar{y} \ \text{ se } x = a\,y \ (a\in\S),
                   \qquad f(\mathit{root}) = \bot &&\text{suffix link.}
\end{align*}
Gli stati sono dunque in corrispondenza biunivoca ($1{:}1$) con le sottostringhe di $T$ (alla
stringa vuota corrisponde $\mathit{root}$), mentre $\bot$ è uno stato «dummy» ausiliario.
\end{definizione}

Le prime quattro componenti, $\bigl(Q\cup\set{\bot},\mathit{root},F,g\bigr)$, sono quelle di
un automa a stati finiti deterministico (DFA), con la sola differenza che $g$ è
\emph{parziale}: da $\bar{x}$ non esce alcun arco etichettato $a$ quando $xa$ non compare in
$T$. È l'ultima componente $f$ a rendere la struttura un \emph{automa aumentato}
(\emph{augmented automaton}). I suffix tree che
vogliamo costruire sono, allo stesso modo, automi deterministici a stati finiti
\emph{aumentati} che accettano tutti e soli i suffissi di un testo.

\begin{osservazione}
Accanto alla riga di $F$ gli appunti annotano, con una freccia, la parola «necessario»: è $F$
a fare di $\STrie{}(T)$ un automa che riconosce \emph{esattamente} i suffissi di $T$, e non
semplicemente le sue sottostringhe.
\end{osservazione}

\scan{21}
Il suffix link $f$ porta dallo stato associato alla stringa $x$ allo stato associato alla
stringa che si ottiene da $x$ cancellandone il primo carattere; la clausola
$f(\mathit{root}) = \bot$ lo rende totale.

\subsection{Il boundary path}

\begin{osservazione}
La funzione $f$ mi consente di costruire un cammino che visita \emph{tutti} i suffissi di
$T$. Partendo dallo stato $\bar{x}$ associato al suffisso più lungo, $x = t_1 t_2 \cdots t_n
= T$, e applicando ripetutamente $f$ si ottiene
\[
  \overline{t_1\cdots t_n}
  \;\xrightarrow{\ f\ }\; \overline{t_2\cdots t_n}
  \;\xrightarrow{\ f\ }\; \cdots
  \;\xrightarrow{\ f\ }\; \overline{t_n}
  \;\xrightarrow{\ f\ }\; \mathit{root}
  \;\xrightarrow{\ f\ }\; \bot .
\]
Questo cammino prende il nome di \textbf{boundary path}.
\end{osservazione}

\begin{figure}[H]
\centering
% Boundary path di STrie(T).
% La "montagna" schematizza il trie: il vertice e' la radice, la linea ondulata
% in basso e' la frontiera degli stati piu' profondi.  Su di essa stanno gli
% stati dei suffissi di T, collegati fra loro dai suffix link f.
\begin{tikzpicture}[
    scale=1.15,
    >={Stealth[length=4.5pt]},
    font=\footnotesize,
    stato/.style={circle,fill,inner sep=0pt,minimum size=4.4pt}]

  % ---- contorno schematico dello STrie -------------------------------------
  \draw[gray]
    (0,3.0) -- (-0.45,2.55) -- (-0.85,2.42) -- (-1.35,1.62)
         -- (-1.70,1.42) -- (-2.05,0.62) -- (-2.30,0.10);
  \draw[gray]
    (0,3.0) -- (0.40,2.50) -- (0.85,2.40) -- (1.45,1.78)
         -- (1.90,1.70) -- (2.45,1.15) -- (2.80,0.85);
  % frontiera degli stati piu' profondi
  \draw[gray] plot[smooth,tension=0.85] coordinates
    {(-2.30,0.10) (-1.90,-0.05) (-1.50,0.18) (-1.05,-0.02) (-0.60,0.22)
     (-0.15,0.30) (0.30,0.58) (0.80,0.36) (1.20,0.44) (1.60,0.62)
     (2.05,0.86) (2.45,0.70) (2.80,0.85)};

  % ---- stati dei suffissi ---------------------------------------------------
  \node[stato] (x) at (-1.90,-0.05) {};
  \node[stato] (y) at (-0.15, 0.30) {};
  \node[stato] (z) at ( 1.60, 0.62) {};
  \node[below left=1pt and 0pt] at (x) {$\bar{x}$};
  \node[above left=-2pt and -4pt] at (y) {$\bar{y}$};

  % ---- suffix link ----------------------------------------------------------
  \draw[->,thick] (x) to[out=-55,in=-125]
        node[pos=0.28,below=-1pt] {$f$} (y);
  \draw[->,thick] (y) to[out=70,in=125] (z);

  % ---- etichette esterne ----------------------------------------------------
  \node[anchor=east,align=right] at (-3.55,0.80)
       {suffisso\\[1pt] $t_1\cdots t_n = x$};
  \draw[->,gray,decorate,
        decoration={snake,amplitude=.5mm,segment length=3.2mm,post length=2mm}]
       (-3.45,0.95) -- (-2.38,0.55);

  \node[anchor=north] at (0.75,-1.05) {suffisso $t_2\cdots t_n$};
  \draw[->,gray] (0.62,-0.95) to[out=140,in=-35] (y);

  \node[anchor=west] at (2.95,0.08) {suffisso $t_3\cdots t_n$};
  \draw[->,gray] (2.88,0.20) to[out=165,in=-20] (z);
\end{tikzpicture}

\caption{Il \emph{boundary path} di $\STrie{}(T)$: applicando ripetutamente il suffix link $f$
a partire dallo stato $\bar{x}$ del suffisso più lungo $x = t_1\cdots t_n$ si incontrano,
nell'ordine, gli stati di tutti i suffissi di $T$.}
\label{fig:boundary-path}
\end{figure}

\scan{22}
Il boundary path è ciò che definisce i nodi finali dello \STrie{}. Durante il passo che legge
$t_i$ il cammino che interessa è quello di $\STrie{}(T^{\,i-1})$: percorrendolo si incontrano,
nell'ordine, gli stati
\[
  s_j \;=\; \overline{t_j t_{j+1} \cdots t_{i-1}}
        \;=\; f^{\,j-1}\bigl(\overline{t_1\cdots t_{i-1}}\bigr),
  \qquad j = 1,\dots,i,
\]
cioè $s_1 = \mathit{top} = \overline{t_1\cdots t_{i-1}}$, poi $s_2 = \overline{t_2\cdots
t_{i-1}}$, e così via fino a $s_i = \overline{\vuota} = \mathit{root}$, seguito da $\bot$:
sono esattamente gli stati che rappresentano i suffissi di $T^{\,i-1}$, e per costruzione
$s_{j+1} = f(s_j)$.

\subsection{Costruzione del suffix trie}

\scan{21}
Passerò da $\STrie{}(T^{\,i-1})$ a $\STrie{}(T^{\,i})$ visitando il boundary path e
aggiungendo una transizione, quando necessario, ad ogni passo: a ogni nodo che vi si incontra
--- cioè a ogni nodo legato a un suffisso di $T^{\,i-1}$ --- appendo un nuovo nodo con un arco
etichettato $t_i$, \textbf{se} da quel nodo non escono già archi etichettati con $t_i$.

Il motivo per cui questo basta è che l'insieme degli stati cresce solo lungo i suffissi:
\[
  Q\bigl(T^{\,i}\bigr) \;=\; Q\bigl(T^{\,i-1}\bigr) \;\cup\;
  \set{\, \overline{t_j\cdots t_i} \;:\; 1 \le j \le i \,} .
\]
Infatti ogni sottostringa di $T^{\,i}$ che non sia già sottostringa di $T^{\,i-1}$ deve
terminare con il carattere $t_i$, cioè è un suffisso di $T^{\,i}$; e i suffissi di $T^{\,i}$
sono esattamente gli stati del boundary path di $\STrie{}(T^{\,i-1})$ estesi con $t_i$.

\begin{algorithm}[H]
\caption{Costruzione di $\STrie{}(T)$ per $T = t_1 t_2 \cdots t_n$ (Ukkonen, Algoritmo 1)}
\label{alg:ukkonen-trie}
\begin{algorithmic}[1]
  \State crea gli stati $\mathit{root}$ e $\bot$
  \State $g(\bot, a) \gets \mathit{root}$ per ogni $a \in \S$;\quad
         $f(\mathit{root}) \gets \bot$
  \State $\mathit{top} \gets \mathit{root}$
  \For{$i \gets 1$ \textbf{to} $n$}
    \State $r \gets \mathit{top}$
    \While{$g(r, t_i)$ è indefinita}
      \State crea un nuovo stato $r'$ e poni $g(r, t_i) \gets r'$
      \If{$r \neq \mathit{top}$}
        \State $f(\mathit{oldr}') \gets r'$
        \Comment{$\mathit{oldr}'$ appartiene ancora al passo precedente}
      \EndIf
      \State $\mathit{oldr}' \gets r'$
      \State $r \gets f(r)$
    \EndWhile
    \State $f(\mathit{oldr}') \gets g(r, t_i)$
    \State $\mathit{top} \gets g(\mathit{top}, t_i)$
  \EndFor
\end{algorithmic}
\end{algorithm}

\begin{osservazione}
Il ciclo \textbf{while} esegue sempre almeno un'iterazione, perché $\mathit{top}$ è lo stato
più profondo del trie e non ha archi uscenti; e termina sempre, perché risalendo il boundary
path con $r \gets f(r)$ si arriva al più a $\bot$, dove $g(\bot, t_i) = \mathit{root}$ è
definita per ogni carattere. La guardia $r \neq \mathit{top}$ serve a non sovrascrivere il
suffix link dell'ultimo stato creato al passo precedente: i suffix link vengono assegnati
all'indietro (allo stato creato subito dopo), e il primo stato creato nel passo $i$ --- il
figlio di $\mathit{top}$ --- non ha, in quel passo, alcun predecessore cui attaccarsi.
La variabile $\mathit{top}$ mantiene l'invariante
$\mathit{top} = \overline{t_1\cdots t_i}$ alla fine del passo $i$-esimo, cioè è sempre il capo
del boundary path corrente.
\end{osservazione}

\paragraph{Costo.}
Per costruire $\STrie{}(T)$ eseguo $n$ volte l'operazione
$\STrie{}(T^{\,i-1}) \longrightarrow \STrie{}(T^{\,i})$, e ognuna costa $\Oh(i)$: il boundary
path di $\STrie{}(T^{\,i-1})$ ha infatti $i$ stati (i suffissi di $T^{\,i-1}$, compresa
$\vuota$, cioè la radice) più lo stato $\bot$. Sommando,
\[
  \sum_{i=1}^{n} \Oh(i) \;=\; \Oh\bigl(\abs{T}^{2}\bigr),
\]
cioè un \textbf{algoritmo quadratico}. Ed è \emph{ottimo}: la sola dimensione di
$\STrie{}(T)$ può essere $\Th(\abs{T}^{2})$ --- si pensi a $T = a^m b^m$ --- quindi nessun
algoritmo che costruisca esplicitamente il suffix trie può costare meno.

\subsection[Il boundary path si può accorciare]{Non serve percorrere tutto il boundary path}

\scan{22}
\begin{osservazione}
\textbf{Non} è sempre necessario visitarlo tutto!
\end{osservazione}

Gli stati del boundary path rappresentano infatti suffissi via via più corti: se $\bar s$
possiede già una $t_i$-transizione, allora $s\,t_i$ occorre in $T^{\,i-1}$, e quindi anche
$s'\,t_i$ occorre in $T^{\,i-1}$ per ogni suffisso $s'$ di $s$. Tutti gli stati che seguono
$\bar s$ sul boundary path hanno perciò già la loro $t_i$-transizione, e il ciclo
\textbf{while} dell'Algoritmo~\ref{alg:ukkonen-trie} può fermarsi al primo nodo $r$ per cui
$g(r,t_i)$ è definita --- come in effetti fa.

Questo non cambia la complessità asintotica finché si costruisce il trie: l'algoritmo resta
quadratico per la sola ragione che già la dimensione di $\STrie{}(T)$ può essere
$\Th(\abs{T}^{2})$. È però l'osservazione su cui poggia tutto il guadagno del passaggio al
suffix tree.

\subsection[Dallo STrie allo STree: i due ingredienti]%
{Dallo \texorpdfstring{\STrie{}}{STrie} allo \texorpdfstring{\STree{}}{STree}: i due ingredienti}

Per passare dallo \STrie{} allo \STree{} servono \emph{due ingredienti}:
\begin{enumerate}
  \item[(i)] tenere i \emph{soli nodi branching};
  \item[(ii)] rimpiazzare le etichette-stringa con coppie di interi,
        $w \longrightarrow [k,p]$.
\end{enumerate}
Entrambi servono a limitare a $\Oh(\abs{T})$ la dimensione dello \STree{}.

\begin{domanda}
Come uso (i) e (ii) nella definizione e nella costruzione dello \STree{}?
\end{domanda}

\textbf{(i)} non è un problema: mi limiterò ai nodi branching e definirò le transizioni
\emph{su stringhe},
\[
  g'(\bar x,\,w) \;=\; \overline{xw},
\]
dove il secondo argomento $w$ non è più un singolo carattere ma una stringa, rappresentata
dalla coppia di interi $[k,p]$, cioè $w = \sub{T}{k}{p}$.

\subsection[Stati espliciti e impliciti]{Stati espliciti, stati impliciti e coppie di riferimento}

\scan{23}
\textbf{(ii)} Per costruire $\STree{}(T^{\,i})$ a partire da $\STree{}(T^{\,i-1})$ dovrò
\emph{reintrodurre stati}. Introduciamo perciò due tipi di stato:
\begin{itemize}
  \item stati \emph{espliciti} (\emph{tree}): i nodi branching --- e le foglie --- che
        sopravvivono nell'albero compattato;
  \item stati \emph{impliciti} (\emph{trie-only}): quelli che esistono nel trie ma non
        nell'albero, perché cadono all'interno di un arco compattato.
\end{itemize}

\begin{definizione}[Il suffix tree come automa]\label{def:stree-automa}
\[
  \STree{}(T) \;=\; \bigl(\, Q' \cup \set{\bot},\ \mathit{root},\ g',\ f' \,\bigr),
\]
dove $Q'$ è l'insieme degli stati espliciti, $g'$ la funzione di transizione su stringhe
appena definita e $f'$ la funzione dei suffix link ristretta a $Q'$.

Gli stati --- impliciti ed espliciti --- sono dati da \emph{reference pair} (coppie di
riferimento): coppie $(s,w)$ con $s$ stato esplicito e $w$ stringa, cioè, rappresentando $w$
con la coppia di interi, $\bigl(s,[k,p]\bigr)$.
\end{definizione}

\begin{osservazione}
Di $w$ interessa soltanto il \emph{primo carattere}: è quello che decide quale arco uscente da
$s$ imboccare. Se $t_k = a$ chiamiamo la transizione una \emph{$a$-transizione}.
\end{osservazione}

\begin{figure}[H]
\centering
% Due esempi di reference pair (s,w): stesso s, stringhe w di lunghezza diversa,
% quindi due stati impliciti distinti dello stesso arco.
% Il nodo pieno e' lo stato esplicito (branching) s; lungo l'arco lungo stanno
% gli stati impliciti (cerchietti vuoti); la graffa misura la stringa w.
\begin{tikzpicture}[
    scale=1.2,
    font=\footnotesize,
    esplicito/.style={circle,fill,inner sep=0pt,minimum size=5pt},
    implicito/.style={circle,draw,fill=white,inner sep=0pt,minimum size=4pt},
    graffa/.style={decorate,decoration={brace,amplitude=4pt},gray}]

  % il secondo schizzo e' identico al primo: cambia solo dove finisce la graffa
  \foreach \dx/\tfin/\lx/\ly in {0/0.42/0.87/-0.25, 3.6/0.82/1.13/-0.77}{%
    \begin{scope}[xshift=\dx cm]
      % l'altro figlio di s: s e' un branching point
      \draw (0,0) -- (-0.85,-0.85);
      % arco lungo, con gli stati impliciti che il trie ha e l'albero no
      \draw (0,0) -- (1.30,-2.60);
      \foreach \t in {0.22,0.42,0.62,0.82}
        \node[implicito] at ({1.30*\t},{-2.60*\t}) {};
      % stato esplicito
      \node[esplicito] at (0,0) {};
      \node[above left=-3pt and -3pt] at (0,0) {$s$};
      % la stringa w della reference pair
      \draw[graffa] (0.34,0.17) -- ({1.30*\tfin+0.34},{-2.60*\tfin+0.17});
      \node at (\lx,\ly) {$w$};
      % entrambe le coppie sono rappresentazioni lecite
      \node at (1.85,-1.35) {$\checkmark$};
    \end{scope}}
\end{tikzpicture}

\caption{Due esempi di reference pair $(s,w)$: dallo stato esplicito $s$ (nodo pieno) si legge
la stringa $w$ scendendo lungo l'arco, i cui stati intermedi (i cerchietti vuoti) sono
impliciti. Le due coppie individuano due stati impliciti diversi dello stesso arco e sono
entrambe rappresentazioni lecite; quale sia la forma \emph{canonica} è la questione discussa
subito dopo.}
\label{fig:reference-pair}
\end{figure}

\noindent\lacunainl{i due schizzi del quaderno sono troncati dal bordo inferiore del foglio:
la figura è una ricostruzione}

\scan{24}
Nella coppia di riferimento $(s,w)$ lo stato esplicito $s$ non è determinato univocamente: uno
stesso punto dell'albero è individuato da \emph{tutti} gli stati espliciti che lo precedono
sul cammino dalla radice. Fra queste rappresentazioni ne scegliamo una sola.

\begin{definizione}[Coppia di riferimento canonica]\label{def:coppia-canonica}
Una coppia di riferimento $(s,w)$ è \emph{canonica} quando, leggendo $w$ a partire da $s$, non
si incontra alcun branching point; equivalentemente, quando $s$ è il branching point
\emph{più profondo possibile} fra quelli dai quali il punto in questione si raggiunge leggendo
una stringa (e dunque $w$ è la più corta possibile).
\end{definizione}

\begin{osservazione}
Se $s$ è branching, $(s,\vuota)$ è la rappresentazione canonica dello stato $s$.
\end{osservazione}

\begin{figure}[H]
\centering
% Coppie di riferimento canoniche e non canoniche.
% Catena di stati impliciti (cerchietti vuoti, archi tratteggiati) che scende da
% s, attraversa il branching point s' e prosegue: (s,abcd) e' canonica,
% (s,abcdeabc) no, perche' scavalca s'.
\begin{tikzpicture}[
    scale=1.05,
    >={Stealth[length=4.5pt]},
    font=\footnotesize,
    esplicito/.style={circle,draw,fill=white,inner sep=0pt,minimum size=4.6mm},
    implicito/.style={circle,draw,fill=white,inner sep=0pt,minimum size=3.4pt},
    catena/.style={dashed,dash pattern=on 2pt off 1.6pt},
    ann/.style={black!55}]

  % ---- nodi -----------------------------------------------------------------
  \node[esplicito,double,double distance=0.7pt] (s)  at (0.00, 0.00) {};
  \node[implicito] (i1) at (0.65,-0.52) {};
  \node[implicito] (i2) at (1.30,-1.04) {};
  \node[implicito] (i3) at (1.95,-1.56) {};
  \node[implicito] (i4) at (2.60,-2.08) {};
  \node[esplicito] (sp) at (3.25,-2.60) {};
  \node[implicito] (j1) at (3.90,-3.12) {};
  \node[implicito] (j2) at (4.55,-3.64) {};
  \node[implicito] (fin) at (5.20,-4.16) {};

  \node at ($(s)+(0.34,0.30)$)  {$s$};
  \node at ($(sp)+(0.44,0.34)$) {$s'$};

  % ---- archi ----------------------------------------------------------------
  \draw (-0.62,0.50) -- (s);                 % arco entrante dal padre
  \draw (sp) -- (2.62,-3.16);                % s' e' un branching point
  \foreach \da/\db/\lab in {s/i1/a, i1/i2/b, i2/i3/c, i3/i4/d, i4/sp/e,
                            sp/j1/a, j1/j2/b, j2/fin/c}
    \draw[catena] (\da) -- node[pos=0.5,above right=-4pt and -4pt] {$\lab$} (\db);

  % ---- annotazioni ----------------------------------------------------------
  % coppia canonica: da s si legge abcd senza incontrare branching point
  \fill (i4) circle (1.7pt);
  \node[ann,anchor=east] at (0.35,-2.00) {$(s,abcd)$};
  \draw[ann,->] (0.45,-2.02) to[out=-12,in=180] (i4);

  % coppia non canonica: il cammino attraversa il branching point s'
  % (il punto indicato e' implicito: l'unico branching point della catena e' s')
  \fill (fin) circle (1.7pt);
  \node[ann,anchor=west] at (3.30,-0.30) {$(s,abcdeabc)$};
  \draw[ann,->] (4.20,-0.60) to[out=-80,in=25] (fin);
\end{tikzpicture}

\caption{Coppie di riferimento canoniche e non canoniche. Da $s$ parte una catena di stati
impliciti etichettata $abcde$ che termina nel branching point $s'$; da $s'$ prosegue una
seconda catena etichettata $abc$. La coppia $(s,abcd)$ è canonica, perché fra $s$ e il punto
indicato non si incontra alcun branching point; la coppia $(s,abcdeabc)$ non lo è, perché
attraversa $s'$: la sua forma canonica è $(s',abc)$.}
\label{fig:coppia-canonica}
\end{figure}

\subsection{Active point ed end point}

\scan{23}
Riprendiamo il boundary path di $\STree{}(T^{\,i-1})$, cioè la successione di stati
$s_1,\dots,s_i,s_{i+1}$ con $s_j = \overline{t_j\cdots t_{i-1}}$ e $s_{j+1} = f(s_j)$, che
dopo $s_i = \mathit{root}$ termina nello stato ausiliario $s_{i+1} = \bot$. Passare da
$\STree{}(T^{\,i-1})$ a $\STree{}(T^{\,i})$ significa garantire che ciascuno di questi stati
abbia una $t_i$-transizione; il punto è che \emph{non è necessario percorrere tutto} il
boundary path, e che il cammino si spezza in tre tratti.

\begin{figure}[H]
\centering
% Il boundary path di STrie(T^{i-1}) percorso quando si estende con t_i, e la
% sua tripartizione (Ukkonen).  Nel quaderno il cammino e' una spezzata a dente
% di sega che sale in diagonale: qui e' raddrizzato in una catena verticale,
% percorsa dal basso (top, lo stato piu' profondo) verso l'alto (perp).
\begin{tikzpicture}[
    scale=0.95,
    >={Stealth[length=4.5pt]},
    font=\footnotesize,
    stato/.style={circle,draw,fill=white,inner sep=0pt,minimum size=4pt},
    estremo/.style={circle,fill,inner sep=0pt,minimum size=4.6pt},
    tacca/.style={line width=1.1pt},
    bolla/.style={draw,ellipse,align=center,inner sep=1.5pt,font=\scriptsize},
    ann/.style={black!55}]

  % ---- la catena del boundary path: gli stati s_j legati dai suffix link f ---
  \node[estremo] (n0) at (0,0.0) {};
  \foreach \i/\y in {1/0.7, 2/1.4, 3/2.1, 4/2.8, 5/3.5, 6/4.2, 7/4.9}
    \node[stato] (n\i) at (0,\y) {};
  \node[estremo] (n8) at (0,5.6) {};

  \foreach \a/\b in {0/1, 2/3, 3/4, 4/5, 5/6, 6/7}
    \draw[->] (n\a) -- (n\b);
  % tratti in cui molti stati sono omessi
  \draw[->,dotted,thick] (n1) -- (n2);
  \draw[->,dotted,thick] (n7) -- (n8);

  \node[anchor=west] at (0.16,2.45) {$f$};

  % ---- estremi ---------------------------------------------------------------
  \node[anchor=east] at (-0.22,0) {top};
  \node[anchor=east,font=\scriptsize] at (-0.92,0) {$(=T^{\,i-1})$};
  \node[anchor=south,font=\large] at (0,5.72) {$\bot$};

  % ---- etichetta del cammino -------------------------------------------------
  \node[anchor=west,align=left] at (1.45,6.15) {boundary path\\(bp)};
  \draw[ann] (1.38,6.02) to[out=195,in=25] (0.13,5.32);

  % ---- tacche: active point ed end point ------------------------------------
  \draw[tacca] (-0.17,1.68) -- (0.17,1.92);
  \draw[tacca] (-0.17,3.78) -- (0.17,4.02);

  \node[bolla] (ap) at (2.30,1.80) {active\\point};
  \draw[ann,->] (ap.west) to[out=180,in=-25] (0.21,1.90);
  \node[anchor=west,align=left] at (3.35,1.80) {primo stato\\non foglia};

  \node[bolla] (ep) at (2.30,3.90) {end\\point};
  \draw[ann,->] (ep.west) to[out=180,in=-25] (0.21,4.00);

  % ---- le tre zone -----------------------------------------------------------
  \draw[ann,decorate,decoration={brace,amplitude=4pt,mirror}]
       (0.52,-0.10) -- (0.52,1.62);
  \node[anchor=west] at (0.90,0.30) {appendiamo a foglia};

  \draw[ann,decorate,decoration={brace,amplitude=4pt}]
       (-0.52,1.98) -- (-0.52,3.72);
  \node[anchor=east,align=right] at (-0.88,2.85) {appendiamo\\nuove foglie};

  \draw[ann,decorate,decoration={brace,amplitude=4pt,mirror}]
       (0.52,4.08) -- (0.52,5.50);
  \node[anchor=west,align=left] at (0.90,4.60) {non serve appendere\\nulla};

  % ---- dettaglio: l'arco verso una foglia ha etichetta aperta ---------------
  \fill (-1.45,-0.80) circle (1.6pt);
  \draw (-1.45,-0.80) -- (-2.70,-1.85);
  \draw[line width=1.1pt] (-2.70,-1.85) -- (-3.25,-2.32);
  \fill (-2.70,-1.85) circle (2.1pt);
  \fill (-3.25,-2.32) circle (2.1pt);
  \node[anchor=east] at (-2.50,-1.02) {$[k,\infty]$};
  \node[anchor=east] at (-3.02,-1.94) {bp};
\end{tikzpicture}

\caption{Il \emph{boundary path} percorso quando si estende l'albero con il carattere $t_i$:
parte dallo stato $\mathit{top}$, il più profondo (quello che rappresenta tutto il prefisso
$T^{\,i-1}$), e risale di suffix link in suffix link fino allo stato ausiliario $\bot$.
Finché gli stati incontrati sono foglie basta \emph{appendere a foglia}; dall'\emph{active
point} fino all'\emph{end point} si \emph{appendono nuove foglie}; dall'\emph{end point} in
poi \emph{non serve appendere nulla}, perché la $t_i$-transizione c'è già.}
\label{fig:boundary-path-tripartizione}
\end{figure}

\scan{24}
\begin{definizione}[Active point ed end point]\label{def:active-end-point}
Sul boundary path di $\STree{}(T^{\,i-1})$:
\begin{itemize}
  \item l'\emph{active point} è $s_{j'}$, dove $j'$ è il minimo indice tale che $s_{j'}$
        \textbf{non è una foglia};
  \item l'\emph{end point} è $s_{j''}$, dove $j''$ è il minimo indice tale che $s_{j''}$
        \textbf{ha già un arco uscente etichettato con $t_i$}.
\end{itemize}
Entrambi esistono sempre: la radice non è una foglia, e l'end point al più coincide con
$\bot$, dove $g(\bot,t_i) = \mathit{root}$ è definita per ogni $t_i \in \S$ --- come accade,
per esempio, quando $t_i$ non occorre affatto in $T^{\,i-1}$.
\end{definizione}

\begin{figure}[H]
\centering
% Boundary path di STree(T^{i-1}): active point s_{j'} ed end point s_{j''}.
% Il cammino si percorre da SINISTRA a DESTRA seguendo i suffix link f
% (la profondita' degli stati diminuisce spostandosi a destra).
\begin{tikzpicture}[
    font=\footnotesize,
    >={Stealth[length=2mm]},
    stato/.style   = {circle,fill=black,inner sep=0pt,minimum size=3.4pt},
    figlio/.style  = {circle,fill=black,inner sep=0pt,minimum size=2.4pt},
    arco/.style    = {line width=0.5pt},
    slink/.style   = {->,line width=0.6pt,blue!50!black},
    richiamo/.style= {line width=0.4pt,black!45,densely dashed},
    nota/.style    = {black!70,line width=0.4pt,->},
    bolla/.style   = {ellipse,draw,line width=0.5pt,inner sep=1pt,
                      minimum width=10mm,minimum height=6.5mm},
  ]

  %% ---- active point: primo stato che non e' una foglia --------------
  \node[stato] (A) at (0,0) {};
  \draw[arco] (A) -- (-0.75,-0.90) node[figlio]{};
  \draw[arco] (A) -- (-0.10,-0.95) node[figlio]{};
  \draw[arco] (A) -- ( 0.65,-0.85) node[figlio]{};

  % suffix link in arrivo dallo stato precedente del boundary path
  \draw[slink] (-1.90,-0.62) .. controls (-1.15,-0.42) and (-0.65,-0.16) .. (A);
  \node at (-1.12,0.02) {$f$};

  % richiamo verso l'etichetta
  \draw[richiamo] (0.15,0.16) -- (0.50,0.95) -- (0.72,1.50) -- (0.98,2.05);
  \node[bolla] (E1) at (1.15,2.42) {$s_{j'}$};
  \node at (1.15,3.28) {\emph{active point}};
  \node[anchor=west,text width=2.9cm,align=left] (N1) at (2.30,2.72)
        {primo stato del boundary path che non è una foglia};
  \draw[nota] (N1.west) -- (E1.east);

  %% ---- stati intermedi, elisi --------------------------------------
  \draw[arco,->] (0.22,0.14) -- (1.72,1.00);
  \fill (2.05,1.15) circle (0.8pt);
  \fill (2.40,1.29) circle (0.8pt);
  \fill (2.75,1.43) circle (0.8pt);
  \fill (3.10,1.57) circle (0.8pt);
  \fill (3.45,1.71) circle (0.8pt);

  %% ---- end point: primo stato con una t_i-transizione --------------
  \node[stato] (B) at (5.10,1.95) {};
  \draw[slink] (0.24,-0.08) .. controls (1.90,-0.55) and (4.00,1.25) .. (4.97,1.88);

  \draw[richiamo] (5.22,2.14) -- (5.45,2.55) -- (5.68,2.85) -- (5.85,3.20);
  \node[bolla] (E2) at (5.95,3.60) {$s_{j''}$};
  \node at (5.95,4.35) {\emph{end point}};
  \node[anchor=north west,text width=2.9cm,align=left] (N2) at (6.85,3.15)
        {primo stato del boundary path che ha già una \mbox{$t_i$-transizione}};
  \draw[nota] (N2.north west) -- (E2.south east);

  %% ---- regione di lavoro ------------------------------------------
  \draw[blue!50!black,line width=0.6pt,decorate,
        decoration={brace,amplitude=5pt,mirror}] (0.20,-1.40) -- (5.10,-1.40);
  \node[anchor=north] at (2.65,-1.78) {regione in cui devo «lavorare»};

\end{tikzpicture}

\caption{Il boundary path di $\STree{}(T^{\,i-1})$, percorso da sinistra a destra seguendo i
suffix link $f$. L'\emph{active point} $s_{j'}$ è il primo stato che non è una foglia,
l'\emph{end point} $s_{j''}$ è il primo stato che possiede già una $t_i$-transizione: fra i
due sta la regione in cui si deve effettivamente lavorare.}
\label{fig:active-end-point}
\end{figure}

Gli stati $s_1,\dots,s_{j'-1}$ sono foglie e si aggiornano da soli: l'arco che porta a una
foglia è etichettato con una coppia \emph{aperta} $[k,\infty]$, dove $\infty$ sta per «fine
del prefisso corrente», e passando da $T^{\,i-1}$ a $T^{\,i}$ l'etichetta non va nemmeno
toccata. Da $s_{j''}$ in poi non c'è invece nulla da fare. Resta solo il tratto
$s_{j'},\dots,s_{j''-1}$, dove occorre creare le nuove foglie etichettate $t_i$ e, se il punto
è implicito, spezzare l'arco introducendo un nuovo branching point.

\begin{osservazione}\label{oss:propagazione-ti}
Che da $s_{j''}$ in poi non ci sia più nulla da fare segue dal fatto che la proprietà «avere
una $t_i$-transizione» si propaga lungo i suffix link: se $s_j$ ha un arco uscente etichettato
con $t_i$, allora $t_j\cdots t_{i-1}t_i$ è sottostringa di $T^{\,i-1}$, quindi lo è anche
$t_{j+1}\cdots t_{i-1}t_i$ e dunque anche $s_{j+1}=f(s_j)$ ha una $t_i$-transizione.
\end{osservazione}

\begin{figure}[H]
\centering
% Le t_i-transizioni si propagano lungo i suffix link: tre stati consecutivi
% del boundary path (a profondita' decrescente, quindi disegnati sempre piu' in alto),
% ciascuno con la propria t_i-transizione.
\begin{tikzpicture}[
    font=\footnotesize,
    >={Stealth[length=2mm]},
    stato/.style  = {circle,fill=black,inner sep=0pt,minimum size=3.4pt},
    figlio/.style = {circle,fill=black,inner sep=0pt,minimum size=2.6pt},
    arco/.style   = {line width=0.5pt},
    slink/.style  = {->,line width=0.6pt,blue!50!black},
  ]

  \node[stato]  (s1) at (0.00, 0.00) {};
  \node[stato]  (s2) at (1.70, 0.70) {};
  \node[stato]  (s3) at (3.40, 1.40) {};

  \node[figlio] (c1) at (0.00,-1.10) {};
  \node[figlio] (c2) at (1.70,-0.40) {};
  \node[figlio] (c3) at (3.40, 0.30) {};

  % prosecuzione del cammino verso la radice
  \draw[arco,black!45] (s1) -- (-0.70,0.72);

  \draw[arco] (s1) -- (c1) node[midway,right,inner sep=2pt] {$t_i$};
  \draw[arco] (s2) -- (c2) node[midway,right,inner sep=2pt] {$t_i$};
  \draw[arco] (s3) -- (c3) node[midway,right,inner sep=2pt] {$t_i$};

  \draw[slink] (s1) to[bend left=32] node[midway,above,inner sep=2pt] {$f$} (s2);
  \draw[slink] (s2) to[bend left=32] node[midway,above,inner sep=2pt] {$f$} (s3);

\end{tikzpicture}

\caption{Le $t_i$-transizioni si propagano lungo i suffix link: una volta comparse sul
boundary path non spariscono più.}
\label{fig:propagazione-ti}
\end{figure}

Modifico quindi la costruzione descritta per $\STrie{}(T^{\,i})$, adattandola a
$\STree{}(T^{\,i})$:
\begin{enumerate}
  \item[1)] $\STree{}(T^{\,i})$ è un automa (senza insieme degli stati finali $F$) i cui stati
        sono sempre branching, oppure foglie: in questo modo si riduce il numero di stati e
        diventa possibile un algoritmo lineare;
  \item[2)] uso active point ed end point per limitare il lavoro alla regione
        $\text{active} \longrightarrow \text{end}$.
\end{enumerate}

\subsection[Dettagli: end point e active point]%
{I dettagli: l'end point di un passo è l'active point del successivo}

\scan{25}
\textbf{3) Dettagli.}
\begin{enumerate}
  \item[i)] Definisco gli stati \emph{impliciti} ed \emph{espliciti} e li identifico con una
    \emph{reference pair}
    \[
      (s,w)\;=\;\bigl(s,[k,p]\bigr),
    \]
    dove $s$ è uno stato esplicito (branching) e $w = t_k\cdots t_p = \sub{T}{k}{p}$ è la
    stringa che si compita scendendo da $s$ fino allo stato rappresentato. La stringa $w$ non
    viene mai memorizzata: bastano i due indici $[k,p]$, perché il testo $T$ è comunque tenuto
    in memoria. Fra le molte coppie che rappresentano lo stesso stato si usa sempre quella
    \emph{canonica}, cioè quella con $s$ il più profondo possibile.

  \item[ii)] L'\emph{end point} di $\STree{}(T^{\,i-1})$ è l'\emph{active point} di
    $\STree{}(T^{\,i})$.
\end{enumerate}

Il punto ii) va letto così: l'aggiornamento che porta da $\STree{}(T^{\,i-1})$ a
$\STree{}(T^{\,i})$ percorre il boundary path e si arresta nell'end point, cioè nel primo
stato che possiede già una $t_i$-transizione; se tale end point è la coppia
$\bigl(s,[k,i-1]\bigr)$, allora l'active point di $\STree{}(T^{\,i})$ è
\[
  \bigl(s,[k,i]\bigr),
\]
cioè lo stesso stato esteso del carattere $t_i$ appena letto. Non occorre dunque cercarlo: lo
si ha già in mano alla fine del passo precedente.

La conseguenza è che la scansione del boundary path non ricomincia mai da capo: il passo
$i+1$ riparte esattamente da dove si era fermato il passo $i$, l'indice del suffisso da cui si
riparte non decresce mai e i costi delle singole visite telescopizzano. Il numero complessivo
di stati visitati lungo tutti i boundary path resta così $\Oh(n)$.

Il lavoro da fare in ciascuno stato visitato è sempre lo stesso: se lo stato non ha una
$t_i$-transizione, lo si rende esplicito (spezzando l'arco su cui giace, se è uno stato
implicito) e vi si innesta un nuovo arco etichettato $t_i$ verso una nuova foglia; poi si
passa allo stato successivo del boundary path seguendo il suffix link $f$ e si ripete.

\begin{figure}[H]
\centering
% Due passi consecutivi lungo il boundary path.
% Convenzioni: cerchio vuoto = nodo esplicito / foglia; tratto ondulato = cammino
% di lunghezza arbitraria (stringa w); pallino pieno dentro il quadratino = stato
% corrente, reso esplicito spezzando l'arco; freccia f = suffix link.
\begin{tikzpicture}[
    font=\footnotesize,
    >={Stealth[length=2mm]},
    esplicito/.style = {circle,draw,line width=0.5pt,inner sep=0pt,minimum size=3.6mm},
    split/.style     = {rectangle,draw,line width=0.4pt,inner sep=0pt,minimum size=4.6mm},
    stato/.style     = {circle,fill=black,inner sep=0pt,minimum size=3.2pt},
    arco/.style      = {line width=0.5pt},
    onda/.style      = {line width=0.5pt,decorate,
                        decoration={snake,amplitude=1.4pt,segment length=5.5pt,
                                    pre length=1pt,post length=1pt}},
    slink/.style     = {->,line width=0.6pt,blue!50!black},
  ]

  %% =============== frammento sinistro (stato s_j) =====================
  \node[esplicito] (Lc)   at (0.00, 0.00) {};
  \node[split]     (Ls)   at (0.60,-1.90) {};
  \node[stato]            at (Ls)         {};
  \node[esplicito] (Ll)   at (0.40,-3.50) {};
  \node[esplicito] (Lnew) at (1.95,-3.30) {};

  \draw[onda] (Lc) -- (Ls);
  \node at (0.02,-0.98) {$w$};

  \draw[arco] (Ls) -- (Ll);
  \draw[arco] (Ls) -- (Lnew) node[midway,above right,inner sep=1pt] {$t_i$};

  % suffix link in arrivo dal passo precedente
  \draw[slink] (-2.20,-0.90) .. controls (-1.40,-0.55) and (-0.85,-0.18) .. (Lc);
  \node at (-1.55,-0.20) {$f$};

  %% =============== frammento destro (stato s_{j+1}) ===================
  \node[esplicito] (Rc)   at (5.50, 1.15) {};
  \node[split]     (Rs)   at (6.00,-0.75) {};
  \node[stato]            at (Rs)         {};
  \node[esplicito] (Rl)   at (5.80,-2.35) {};
  \node[esplicito] (Rnew) at (7.35,-2.15) {};

  \draw[onda] (Rc) -- (Rs);
  \node at (6.02,0.17) {$w$};

  \draw[arco] (Rs) -- (Rl);
  \draw[arco] (Rs) -- (Rnew) node[midway,above right,inner sep=1pt] {$t_i$};

  %% =============== passaggio da un frammento all'altro ================
  \draw[slink] (Ls) .. controls (2.30,-1.35) and (3.90,0.85) ..
        node[pos=0.62,above,inner sep=2pt] {$f$} (Rc);

\end{tikzpicture}

\caption{Due passi consecutivi lungo il boundary path. In ciascuno dei due frammenti si scende
dal nodo esplicito (cerchio) lungo la stringa $w$ fino allo stato corrente (pallino pieno); lì
l'arco viene spezzato (quadratino) e si innesta il nuovo arco etichettato $t_i$ verso una
nuova foglia. Il suffix link $f$ porta allo stato successivo, dove si ripete la stessa
operazione con la stessa $w$.}
\label{fig:ukkonen-boundary-step}
\end{figure}

Il suffix link è però definito soltanto sugli stati espliciti: sulle coppie di riferimento $f$
agisce sulla prima componente,
\[
  f\bigl(r,[k,p]\bigr)\;=\;\bigl(f(r),[k,p]\bigr),
\]
lasciando invariata la coppia di indici; il risultato va poi ri-canonizzato, scendendo da
$f(r)$ finché non si raggiunge il nodo esplicito più profondo. Nella figura seguente $s$
denota proprio il nodo esplicito ottenuto dopo la ri-canonizzazione e $[k,p]$ la coppia
aggiornata: lo stato $s'$ in cui si atterra è descritto dalla coppia $\bigl(s,[k,p]\bigr)$,
dove $t_k$ è il primo carattere dell'arco uscente da $s$ (è il solo carattere che serve per
scegliere la transizione) e $[k,p]$ la porzione di testo ancora da percorrere.
Anche in $s'$ si innesta il nuovo arco etichettato $t_i$.

\begin{figure}[H]
\centering
% Il suffix link letto sulle coppie di riferimento:
% f porta nello stato s', rappresentato dalla coppia (s,[k,p]).
% Convenzioni: tratto ondulato = cammino di lunghezza arbitraria (sottoalbero
% o cammino eliso); pallino pieno = stato.
\begin{tikzpicture}[
    font=\footnotesize,
    >={Stealth[length=2mm]},
    stato/.style   = {circle,fill=black,inner sep=0pt,minimum size=3.2pt},
    arco/.style    = {line width=0.5pt},
    onda/.style    = {line width=0.5pt,decorate,
                      decoration={snake,amplitude=1.4pt,segment length=5.5pt,
                                  pre length=1pt,post length=1pt}},
    slink/.style   = {->,line width=0.6pt,blue!50!black},
    richiamo/.style= {->,line width=0.5pt,black!45,densely dashed},
  ]

  %% ---- parte simbolica ---------------------------------------------
  \node at (0.10,2.72) {$f$};
  \draw[slink] (-0.40,1.65) .. controls (0.55,2.30) and (1.10,2.55) .. (1.90,2.66);
  \node[anchor=west] (rp) at (2.05,2.66) {$\bigl(s,[k,p]\bigr)$};

  %% ---- frammento di albero -----------------------------------------
  \node[stato] (s)  at (4.40, 1.70) {};
  \node at (4.40,2.06) {$s$};

  % sottoalbero eliso che parte da s
  \draw[onda] (s) -- (6.30,1.48);

  % arco uscente da s, il cui primo carattere e' t_k
  \draw[onda] (s) -- (4.40,0.00);
  \draw[line width=1.3pt] (4.22,1.28) -- (4.58,1.32);
  \node[anchor=east] at (4.12,1.30) {$t_k$};

  % lo stato raggiunto
  \node[stato] (sp) at (4.40,0.00) {};
  \node[anchor=east] at (4.20,0.20) {$s'$};
  \node[anchor=west] at (4.58,0.06) {$[k,p]$};

  % la coppia di riferimento denota proprio s'
  \draw[richiamo] (2.55,2.34) .. controls (2.25,0.85) and (3.05,-0.50) .. (4.20,-0.06);

  % sottoalbero gia' esistente, eliso
  \draw[onda] (sp) -- (4.05,-1.85);

  % nuovo arco etichettato t_i verso la nuova foglia
  \node[stato] (nf) at (5.60,-1.50) {};
  \draw[arco] (sp) -- (nf) node[midway,anchor=west,inner sep=3pt] {$t_i$};

\end{tikzpicture}

\caption{Il suffix link letto sulle coppie di riferimento: lo stato $s'$ raggiunto applicando
$f$ è rappresentato dalla coppia $\bigl(s,[k,p]\bigr)$, con $t_k$ primo carattere dell'arco
uscente da $s$; da $s'$ pende il nuovo arco etichettato $t_i$.}
\label{fig:ukkonen-refpair-f}
\end{figure}

\begin{osservazione}
Gli appunti si fermano qui: i tre ingredienti (soli nodi branching, active/end point,
reference pair canoniche) sono enunciati, ma le procedure di Ukkonen che li realizzano
(\emph{update}, \emph{canonize}, \emph{test-and-split}) e la dimostrazione formale del tempo
$\Oh(n)$ non vengono scritte.
\end{osservazione}


\chapter{Karp--Rabin}
\label{cap:karp-rabin}

\section{Coda: le applicazioni del suffix tree}

\scan{26}
Un'ultima annotazione chiude il blocco sui suffix tree. Il suffix tree è la struttura che
realizza il \emph{pre-processing del testo}; Gusfield vi dedica due capitoli di applicazioni,
che si dividono in due famiglie:
\begin{itemize}
  \item applicazioni \emph{semplici};
  \item applicazioni \emph{non semplici}, che nascono dalla combinazione
        ``suffix tree $+$ Harel--Tarjan''.
\end{itemize}

Il motivo della seconda voce è che, una volta costruito il suffix tree di $T$, moltissime
domande su $T$ si riducono a calcolare il $\lca$ di due foglie. Il teorema di Harel--Tarjan
garantisce che, dopo un pre-processing lineare nella dimensione dell'albero, ogni query di
lowest common ancestor viene risolta in tempo $\Oh(1)$: è questo che rende praticabili --- e
al tempo stesso non banali --- le applicazioni della seconda famiglia. Al teorema, e alla sua
dimostrazione, è dedicato un capitolo a sé più avanti.

\begin{figure}[H]\centering% stree-lca-montagna -- pag. 26
% Il suffix tree disegnato "a montagna": radice in cima, foglie sulla linea di base,
% il nodo lca(l1,l2) a meta' altezza raggiunto in O(1) grazie a Harel--Tarjan.
\begin{tikzpicture}[
    >=Latex,
    ramo/.style={draw=black!75,line width=0.5pt},
    nodo/.style={circle,fill=black!80,inner sep=0pt},
    et/.style={font=\footnotesize}
  ]

  % ---- intestazione dello schizzo -------------------------------------------
  \node[font=\small] at (0.1,4.15) {\STree $+$ \underline{Harel--Tarjan}};

  % ---- sagoma "a montagna" del suffix tree ----------------------------------
  \draw[line width=0.9pt,draw=black!55,fill=black!5]
      (0.40,3.50) -- (0.10,2.85) -- (-0.15,2.60) -- (-0.45,2.35)
      -- (-0.85,1.95) -- (-1.05,1.45) -- (-1.55,1.05) -- (-1.85,0.60)
      -- (-2.45,0.02) -- (2.45,0.02)
      -- (2.00,0.55) -- (1.80,1.05) -- (1.35,1.45) -- (1.20,2.05)
      -- (0.85,2.35) -- (0.70,2.80) -- cycle;

  % ---- "terreno": e' il livello delle foglie --------------------------------
  \draw[line width=0.9pt,draw=black!55]
      plot[smooth,tension=0.8] coordinates
      {(-2.90,0.12) (-2.30,-0.02) (-1.60,0.07) (-0.90,-0.02)
       (-0.10,0.07) (0.70,-0.02) (1.50,0.07) (2.20,-0.02) (2.90,0.12)};

  % ---- cammino radice -> lca (nel quaderno e' ripassato due volte) ----------
  \draw[line width=2.4pt,black!20] (0.40,3.50) -- (0.10,1.85);
  \draw[ramo,dashed] (0.40,3.50) -- (0.10,1.85);

  % ---- archi lca -> foglie --------------------------------------------------
  \draw[ramo] (0.10,1.85) -- (-0.55,1.05) -- (-0.95,0.02);
  \draw[ramo] (0.10,1.85) -- (0.35,0.02);

  % ---- nodi -----------------------------------------------------------------
  \node[nodo,minimum size=3pt] (rad) at (0.40,3.50) {};
  \node[nodo,minimum size=6pt] (a)   at (0.10,1.85) {};
  \node[nodo,minimum size=4pt] (b)   at (-0.55,1.05) {};
  \node[nodo,minimum size=4pt] (l1)  at (-0.95,0.02) {};
  \node[nodo,minimum size=4pt] (l2)  at (0.35,0.02) {};

  % ---- etichette ------------------------------------------------------------
  \node[et,right=3pt] at (rad.east) {radice};
  \node[font=\small\itshape] at (-1.30,0.62) {\STree};

  \draw[black!45,line width=0.4pt] (a) -- (2.10,2.45);
  \node[et,anchor=west] at (2.10,2.45) {$\lca(\ell_1,\ell_2)$};

  \node[et,below=2pt] at (l1) {$\ell_1$};
  \node[et,below=2pt] at (l2) {$\ell_2$};

\end{tikzpicture}

\caption{Applicazioni ``non semplici'': dentro il suffix tree di $T$ si individua il lowest
common ancestor di due foglie, che con il pre-processing di Harel--Tarjan costa $\Oh(1)$.}
\label{fig:stree-lca-montagna}
\end{figure}

\section{Impostazione: dalle stringhe ai numeri}

Siano $T,P \in \S^{*}$ e poniamo $m = \abs{P}$, $n = \abs{T}$. Negli appunti l'alfabeto di
riferimento è quello delle cifre decimali,
\[
  \S = \set{0,1,\dots,9}, \qquad d = \abs{\S} = 10,
\]
ma la scelta è puramente espositiva: l'unica ipotesi che conta davvero è che l'alfabeto abbia
taglia costante, $d = \abs{\S} \in \Oh(1)$. In generale identifichiamo ogni carattere con la
cifra corrispondente in $\set{0,1,\dots,d-1}$, così che scrivere $P[i]$ o $T[j]$ dentro una
formula abbia senso.

Il problema è il consueto problema di pattern matching: determinare tutti gli \emph{shift
validi}, cioè gli $s \in \set{0,1,\dots,n-m}$ tali che $P = \sub{T}{s+1}{s+m}$.

\begin{figure}[H]\centering% kr-finestra-testo -- pag. 26
% Il testo T con la finestra corrente T[s+1..s+m] e la finestra successiva,
% traslata di una sola posizione; il pattern P ha la stessa lunghezza m.
\begin{tikzpicture}[
    >=Latex,
    et/.style={font=\footnotesize},
    graffa/.style={decorate,decoration={brace,amplitude=4pt,raise=1pt},
                   draw=black!70,line width=0.5pt}
  ]

  % ---- il testo T -----------------------------------------------------------
  % la finestra corrente, evidenziata sul testo
  \draw[line width=3.2pt,black!18] (1.00,0) -- (3.20,0);
  \draw[line width=0.9pt,black!80] (0,0) -- (8.40,0);
  \draw[line width=1.6pt,black]    (7.95,0) -- (8.40,0);
  \node[et,right=3pt] at (8.40,0) {$T$};

  \draw[black!60,line width=0.4pt] (1.00,-0.11) -- (1.00,0.11);
  \draw[black!60,line width=0.4pt] (3.20,-0.11) -- (3.20,0.11);
  \node[et,below=1pt] at (1.00,-0.11) {$s+1$};
  \node[et,below=1pt] at (3.20,-0.11) {$s+m$};

  % ---- finestra corrente  t_s = T[s+1..s+m] --------------------------------
  \draw[graffa] (1.00,0.10) -- (3.20,0.10);
  \node[et] at (2.10,0.60) {$\sub{T}{s+1}{s+m}$};

  % ---- finestra successiva: la stessa, traslata di una posizione ------------
  \draw[graffa] (1.45,1.40) -- (3.65,1.40);
  \node[et] at (2.55,1.90) {$\sub{T}{s+2}{s+m+1}$};

  % ---- lo scorrimento di una sola posizione --------------------------------
  \draw[->,black!70,line width=0.5pt] (3.20,1.05) -- (3.65,1.05);
  \node[et,right=3pt,black!70] at (3.65,1.05) {shift di una posizione};

  % ---- il pattern P, allineato sotto la finestra corrente ------------------
  \draw[black!35,dashed,line width=0.4pt] (1.00,-0.55) -- (1.00,-1.15);
  \draw[black!35,dashed,line width=0.4pt] (3.20,-0.55) -- (3.20,-1.15);
  \draw[line width=0.9pt,black!80] (1.00,-1.15) -- (3.20,-1.15);
  \draw[line width=1.6pt,black]    (2.85,-1.15) -- (3.20,-1.15);
  \node[et,right=3pt] at (3.20,-1.15) {$P$};

  % ---- la lunghezza comune e' m --------------------------------------------
  \draw[decorate,decoration={brace,mirror,amplitude=4pt,raise=1pt},
        draw=black!70,line width=0.5pt] (1.00,-1.25) -- (3.20,-1.25);
  \node[et] at (2.10,-1.72) {$m$};

\end{tikzpicture}

\caption{Il testo $T$, la finestra corrente $\sub{T}{s+1}{s+m}$ e la finestra successiva
(traslata di una posizione); il pattern $P$ ha la stessa lunghezza $m$ della finestra.}
\label{fig:kr-finestra-testo}
\end{figure}

\begin{osservazione}[Idea di partenza]
$\sub{P}{1}{m}$ può essere visto come un \emph{numero} scritto in base $d = \abs{\S}$: ogni
carattere è una cifra e la stringa non è altro che la sua rappresentazione posizionale. La
stessa lettura si applica, evidentemente, a ogni finestra di $T$ lunga $m$.
\end{osservazione}

Poniamo dunque:
\begin{itemize}
  \item $p$ è il valore di $\sub{P}{1}{m}$ letto in base $d$, cioè
    \[
      p \;=\; \sum_{i=1}^{m} P[i]\, d^{\,m-i}
        \;=\; P[1]d^{\,m-1} + P[2]d^{\,m-2} + \dots + P[m-1]d + P[m];
    \]
  \item $t_s$ è il valore di $\sub{T}{s+1}{s+m}$ letto in base $d$, cioè
    \[
      t_s \;=\; \sum_{i=1}^{m} T[s+i]\, d^{\,m-i}.
    \]
\end{itemize}

Poiché la scrittura posizionale in base $d$ di stringhe di lunghezza \emph{fissa} $m$ è
iniettiva, vale
\[
  s \text{ è uno shift valido} \iff p = t_s,
\]
e il confronto fra stringhe si è trasformato in un confronto fra interi.

\subsection{Lo schema di partenza}

Lo schema di algoritmo che ne discende è il seguente.
\begin{enumerate}
  \item Inizializzo $p$ e $t_0$.
  \item Calcolo $t_{s+1}$ a partire da $t_s$ con uno \emph{shift} e una \emph{somma}: dalla
    finestra esce il carattere $T[s+1]$ ed entra $T[s+m+1]$, quindi basta togliere la cifra
    più significativa, moltiplicare per $d$ (lo shift) e sommare la nuova cifra meno
    significativa.
  \item A ogni passo confronto $p$ con $t_s$: se sono uguali segnalo un'occorrenza,
    altrimenti eseguo lo shift e passo a $s+1$.
\end{enumerate}

\begin{figure}[H]\centering% kr-schema-loop -- pag. 26
% Lo schema "ingenuo" di Karp--Rabin: un unico test p = t_s, ripetuto per tutti
% gli shift, con aggiornamento incrementale della finestra; l'annotazione in basso
% e' il punto debole (il test non costa O(1)).
\begin{tikzpicture}[
    >=Latex,
    et/.style={font=\footnotesize},
    filo/.style={draw=black!75,line width=0.5pt},
    blocco/.style={draw=black!75,line width=0.7pt,minimum width=3.6cm,
                   minimum height=1.1cm,inner sep=4pt,font=\small}
  ]

  % ---- il blocco di test ----------------------------------------------------
  \node[blocco] (test) at (0,0) {test\quad $p = t_s$};

  % ---- uscite a ventaglio dal lato destro -----------------------------------
  \coordinate (fork) at (1.80,0);
  \node[et,anchor=west] (occ)   at (3.35, 0.80) {occorrenza};
  \node[et,anchor=west] (shift) at (3.35,-0.80) {shift e somma};

  \draw[filo,->] (fork) -- (3.25, 0.80);
  \draw[filo,->] (fork) -- (3.25,-0.80);
  \node[et] at (2.30, 0.62) {s\`i};
  \node[et] at (2.30,-0.62) {no};

  % ---- segnalata l'occorrenza, la scansione prosegue comunque ---------------
  \draw[filo,->,black!45,densely dashed] (occ.south) -- (shift.north);

  % ---- retroazione: il ciclo su tutti gli shift ------------------------------
  \draw[filo,->]
      (shift.south) to[out=-90,in=0] (1.10,-2.95)
                    to[out=180,in=-90] (-3.45,-0.55)
                    to[out=90,in=180] (test.west);
  \node[et,below=2pt,black!70] at (1.10,-2.95) {ciclo su $s = 0,1,\dots,n-m$};

  % ---- annotazione: il test non e' a costo unitario -------------------------
  \draw[filo,black!60] (-0.50,-0.55) -- (-1.30,-1.05);
  \node[anchor=north,align=center,font=\small] at (0.75,-1.05)
    {{\large\bfseries NON} prende tempo $\Oh(1)$!\\[2pt]
     \footnotesize Non si pu\`o usare il criterio di costo unitario};

\end{tikzpicture}

\caption{Lo schema di Karp--Rabin: un unico test $p = t_s$ ripetuto per tutti gli shift, con
aggiornamento incrementale della finestra.}
\label{fig:kr-schema-loop}
\end{figure}

\begin{osservazione}[Il punto debole]
Il test $p = t_s$ \textbf{non} prende tempo $\Oh(1)$: su questi dati non si può usare il
criterio di costo unitario. Infatti $p$ e i $t_s$ sono numeri di $m$ cifre in base $d$, cioè
richiedono $\Th(m \log d) = \Th(m)$ bit; non appena $m$ supera l'ordine di $\log n$ non stanno
più in una parola di macchina, e allora ogni confronto (come ogni aggiornamento) costa
$\Th(m)$ anziché $\Oh(1)$. Ripetendolo per tutti gli $n-m+1$ shift si ritorna a
$\Th\bigl((n-m+1)\,m\bigr)$: nessun guadagno rispetto all'algoritmo naive.
\end{osservazione}

\section{Il rimedio: lavorare modulo \texorpdfstring{$q$}{q}}

\scan{27}
Il problema, insomma, è che \emph{$p$ e i $t_s$ sono troppo grandi}. La via d'uscita è non
manipolarli mai, ma manipolare soltanto i loro resti
\[
  p \mod q, \qquad t_s \mod q,
\]
dove $q$ è un parametro scelto in modo da garantire che $p \mod q$ (e analogamente ogni
$t_s \mod q$) sia manipolabile in tempo $\Oh(1)$.

\begin{proposizione}[Il test modulare è un filtro a un lato]\label{prop:kr-filtro}
Per ogni intero $q \ge 1$ e per ogni shift $s$ vale l'implicazione
\[
  p \mod q \neq t_s \mod q \implies p \neq t_s,
\]
e dunque $s$ non è uno shift valido. L'implicazione inversa, invece, \emph{non} vale in
generale: non appena $q < d^{\,m}$ esistono un pattern e una finestra di testo con
$p \neq t_s$ e $p \mod q = t_s \mod q$, cioè
\[
  p \mod q = t_s \mod q \;\nRightarrow\; p = t_s .
\]
\end{proposizione}

\begin{dimostrazione}
La prima implicazione è la contronominale del fatto che interi uguali hanno lo stesso resto:
se $p = t_s$ allora $p \mod q = t_s \mod q$. Dunque, se i due resti differiscono,
necessariamente $p \neq t_s$ e quindi --- essendo la scrittura in base $d$ di una stringa di
lunghezza fissa $m$ iniettiva --- lo shift $s$ non è valido.

La seconda non vale: $p$ e $t_s$ variano in $\set{0,1,\dots,d^{\,m}-1}$, che ha $d^{\,m}$
elementi, mentre i resti modulo $q$ sono soltanto $q$; se $q < d^{\,m}$, per il principio dei
cassetti due valori distinti finiscono sullo stesso resto. Basta dunque considerare il caso in
cui $p \neq t_s$ ma
\[
  \abs{p - t_s} \mod q = 0,
\]
cioè $q \mid p - t_s$ con $p \neq t_s$: allora $p$ e $t_s$ hanno lo stesso resto modulo $q$
pur essendo diversi, e il test modulare non li distingue. (Se $q \ge d^{\,m}$ il test è
esatto, ma è proprio il caso che si vuole evitare: i resti sarebbero grandi quanto $p$ e
$t_s$, e si tornerebbe al problema di partenza.)
\end{dimostrazione}

\begin{osservazione}
Il confronto fra resti è quindi un filtro a un solo lato: quando i resti differiscono lo
shift si può scartare con certezza (mai un falso negativo), ma quando coincidono non si
conclude nulla, perché $p \mod q = t_s \mod q$ è solo condizione \emph{necessaria} perché $s$
sia uno shift valido. I falsi positivi sono esattamente le posizioni con $p \neq t_s$ e
$q \mid p - t_s$; ognuna di esse va smascherata con una verifica esplicita.
\end{osservazione}

\paragraph{Problemi.}
\begin{enumerate}
  \item[i)] come scelgo $q$?
  \item[ii)] come calcolo e aggiorno $t_s \mod q$?
  \item[iii)] cosa faccio quando $p \mod q = t_s \mod q$?
\end{enumerate}

\begin{dubbio}[Annotazione a matita, in parte cancellata]
A metà pagina compare, a matita e quasi del tutto cancellata, la fattorizzazione
\[
  q = d^{\,n}\cdot a
  \qquad\longleftrightarrow\qquad
  \begin{cases}
    q' \mod d^{\,n} & \dfrac{1}{d^{\,n}}\\[2ex]
    q'' \mod a      & \dfrac{1}{a}
  \end{cases}
  \qquad\qquad \frac{1}{d^{\,n}}\cdot\frac{1}{a},
\]
accanto a uno schizzo (anch'esso a matita e cancellato) di un segmento con alcune tacche e un
cappio, privo di etichette. L'esponente, letto $n$, non ha nulla a che vedere con
$n = \abs{T}$, e la lettera $a$ è qui un cofattore, non il carattere entrante dello
pseudocodice.

La lettura più plausibile è che si tratti di una verifica della probabilità $1/q$ di
collisione: se $q$ si fattorizza come $d^{\,n}\cdot a$ con $\gcd(d^{\,n},a)=1$, per il
teorema cinese del resto la condizione $p \equiv t_s \pmod q$ equivale alla congiunzione di
$p \equiv t_s \pmod{d^{\,n}}$ e $p \equiv t_s \pmod a$, eventi ``indipendenti'' di
probabilità $1/d^{\,n}$ e $1/a$, il cui prodotto è appunto $1/q$. Si noti però che un $q$ di
quella forma sarebbe una pessima scelta: la riduzione modulo $d^{\,n}$ (con $d = \abs{\S}$ la
base) legge soltanto gli ultimi $n$ caratteri della finestra. È anche per evitare queste
degenerazioni che si prenderà $q$ primo.
\end{dubbio}

\section{L'algoritmo}

Rimandata la scelta di $q$, l'algoritmo prende la forma seguente. Da qui in avanti le
variabili che si chiamano $p$ e $t_s$ contengono i \emph{resti} $p \mod q$ e $t_s \mod q$,
non più i due interi grandi: è l'abuso di notazione consueto, e non crea ambiguità perché
$p$ e $t_s$ nel loro significato originario non vengono mai calcolati.

\begin{algorithm}[H]
\caption{\textsc{Karp--Rabin}$(T,P,d,q)$}
\label{alg:karp-rabin}
\begin{algorithmic}[1]
  \State $m \gets P.\mathit{length}$
  \State $n \gets T.\mathit{length}$
  \State $h \gets d^{\,m-1} \mod q$
  \State $p \gets 0$
  \State $t_0 \gets 0$
  \For{$i \gets 1$ \textbf{to} $m$} \Comment{regola di Horner, modulo $q$}
    \State $p   \gets (d\cdot p   + P[i]) \mod q$
    \State $t_0 \gets (d\cdot t_0 + T[i]) \mod q$
  \EndFor
  \For{$s \gets 0$ \textbf{to} $n-m$}
    \If{$p = t_s$} \Comment{filtro}
      \If{$\sub{P}{1}{m} = \sub{T}{s+1}{s+m}$} \Comment{verifica}
        \State \textbf{match} allo shift $s$
      \EndIf
    \EndIf
    \If{$s < n-m$} \Comment{aggiornamento della finestra}
      \State $r \gets (T[s+1]\cdot h) \mod q$
      \State $a \gets T[s+m+1]$
      \State $t_{s+1} \gets \bigl(d\,(t_s - r) + a\bigr) \mod q$
    \EndIf
  \EndFor
\end{algorithmic}
\end{algorithm}

Negli appunti una graffa raccoglie le righe $1$--$9$ e le etichetta $\Oh(m)$: è tutto il
pre-processing, cioè le due valutazioni con la regola di Horner, una moltiplicazione e una
somma per carattere, eseguite direttamente modulo $q$. Anche la riga $3$ rientra nel budget:
$h = d^{\,m-1} \mod q$ si calcola con l'esponenziazione veloce in $\Oh(\log m)$ operazioni, o
banalmente in $\Oh(m)$.

Le tre righe interne al secondo ciclo sono l'aggiornamento incrementale della finestra
descritto al passo $2$ dello schema: $r$ è il contributo della cifra che esce, $a$ è la cifra
che entra, e in forma chiusa
\[
  t_{s+1} = \bigl(d\,(t_s - T[s+1]\cdot h) + T[s+m+1]\bigr) \mod q,
  \qquad h = d^{\,m-1} \mod q,
\]
cioè $\Oh(1)$ operazioni per shift; la guardia $s < n-m$ serve solo a non leggere $T[s+m+1]$
oltre la fine del testo. Che tutto questo si possa fare direttamente sui resti dipende dal
fatto che la riduzione modulo $q$ commuta con somma, differenza e prodotto.

Il confronto $p = t_s$ è il filtro della Proposizione~\ref{prop:kr-filtro}, e il confronto
carattere per carattere $\sub{P}{1}{m} = \sub{T}{s+1}{s+m}$ è la verifica che esso rende
necessaria --- la risposta al problema iii): al più $m$ confronti, ed è il punto in cui si
concentrerà l'analisi di complessità.

\section{La scelta di \texorpdfstring{$q$}{q}}

\scan{28}
Resta aperto il problema i). Il parametro $q$ va preso \emph{primo}, e più grande è, meglio
è: il che è ovvio, dato che la probabilità di un falso positivo sarà $1/q$. Dall'altro lato
$q$ non può crescere quanto si vuole, perché tutto l'impianto si regge sull'ipotesi che le
operazioni modulo $q$ costino $\Oh(1)$: $q$ (e con esso $d\cdot q$) deve stare in una parola
di macchina, cioè occupare $\Oh(\log n)$ bit. I due vincoli sono compatibili --- l'analisi
della prossima sezione mostrerà che basta $q$ dell'ordine di $m$, e $m \le n$ --- e resta da
capire se un primo della taglia voluta esista davvero e se sia facile trovarlo: per questo
useremo il teorema sulla distribuzione dei primi.

Il modulo mappa gli interi in $\set{0,\dots,q-1}$. In modo uniforme? Sì --- o, più
onestamente, \emph{lo assumiamo}: è esattamente questa ipotesi che permette di dire che per
uno shift non valido la congruenza $p \equiv t_s \pmod{q}$ si presenta con probabilità $1/q$.

\section{Analisi di complessità}

Prima dell'analisi vera e propria, negli appunti si apre a margine una parentesi
terminologica.

\begin{osservazione}[Smoothed complexity]
La \emph{smoothed complexity} ``liscia'' la superficie di complessità. Ci sono algoritmi, come
QuickSort, che benché su carta siano $\Oh(n^{2})$, nella pratica sono solitamente più veloci:
l'analisi smoothed serve proprio a rendere conto di questo scarto, facendo pesare di meno i
casi peggiori isolati --- gli \emph{outlier}.
\end{osservazione}

\begin{dubbio}[Ultima riga della parentesi]
L'ultima riga della parentesi si legge ``nella smoothed complexity si \dots{} gruppi maggiori
ridotti dai outlier'', con due verbi separati da una barra che non è stato possibile
decifrare. La lettura qui adottata --- si valuta il costo su \emph{gruppi} (intorni) di input
anziché sul singolo input, così che il contributo degli outlier risulti ridotto --- è
congetturale.
\end{dubbio}

Veniamo all'analisi. Siano
\begin{itemize}
  \item $v$ il numero di posizioni in cui $p \equiv t_s \pmod{q}$ \emph{e} $P$ occorre
        davvero (occorrenze \emph{valide});
  \item $k$ il numero di posizioni in cui $p \equiv t_s \pmod{q}$ \emph{ma} $P$ non occorre
        (occorrenze \emph{spurie}).
\end{itemize}
In entrambi i casi scatta il confronto carattere per carattere, che costa $\Oh(m)$.

Da qui due conseguenze immediate:
\begin{itemize}
  \item se $v$ è alto, Karp--Rabin costa $\Oh(m\cdot n)$, esattamente quanto l'algoritmo
        naive;
  \item assumo dunque che $v$ sia $\Oh(1)$.
\end{itemize}

\begin{osservazione}[Annotazione a matita]
Accanto all'ipotesi $v = \Oh(1)$ il quaderno annota, a matita: ``non è vero per testi
completamente random''. L'osservazione va però capovolta: se $T$ è davvero uniforme e casuale
su un alfabeto di taglia $d$, il numero atteso di occorrenze di $P$ è $(n-m+1)/d^{\,m}$, che
è praticamente nullo non appena $d^{\,m} \gg n$, cioè per $m \gtrsim \log_d n$: in quel regime
l'ipotesi $v = \Oh(1)$ è ampiamente soddisfatta. (Per $m$ molto piccolo il valore atteso di
$v$ può essere $\Th(n)$ anche su testo casuale, ma allora $m\,\E[v]$ resta $\Oh(n)$ e la
complessità non ne risente.) È sui testi \emph{ripetitivi} che l'ipotesi cade: con $P = a^{m}$
e $T = a^{n}$ si ha $v = n-m+1$, e in quel caso nessun algoritmo che verifichi tutte le
occorrenze può scendere sotto $\Th\bigl((n-m+1)\,m\bigr)$.
\end{osservazione}

Quanto a $k$, sono costretto a spendere $\Oh(m)$ ogni volta che
\[
  p - t_s \equiv 0 \pmod{q},
\]
cosa che --- sotto l'ipotesi di uniformità di poco fa --- avviene con probabilità $1/q$;
dunque
\[
  \E[k] \;\le\; \frac{n-m+1}{q} \;=\; \Oh\!\left(\frac{n}{q}\right).
\]

Mettendo tutto insieme (il primo addendo è la scansione del testo con l'aggiornamento
incrementale della finestra, il secondo sono le verifiche esplicite) si ottiene, in valore
atteso,
\begin{align*}
  \Oh(n) + \Oh\bigl(m\,(v+k)\bigr)
    &= \Oh(n) + \underbrace{\Oh(m\,v)}_{v \,=\, \Oh(1)} + \Oh(m\,k) \\[2pt]
    &= \Oh(n) + \Oh(m) + \Oh\!\left(\frac{m\cdot n}{q}\right),
\end{align*}
e con $q \gtrsim m$ il tutto è \emph{lineare}: $\Oh(n+m) = \Oh(n)$.

\begin{osservazione}[Che cosa è davvero aleatorio]
Il conto appena fatto è un conto \emph{in valore atteso}, e l'aspettazione è presa rispetto
all'ipotesi di uniformità dei resti. Con $q$ fissato una volta per tutte quell'ipotesi è una
comodità di analisi e non un teorema: un testo costruito ad arte può far collidere tutti gli
shift e riportare il costo a $\Th\bigl((n-m+1)\,m\bigr)$. La versione rigorosa sceglie $q$ primo \emph{a caso}
in un intervallo opportuno, ed è allora la scelta di $q$ --- non il testo --- a fornire
l'aleatorietà su cui si prende il valore atteso.
\end{osservazione}

\chapter{Shift--And}
\label{cap:shift-and}

\scan{29}
Karp--Rabin ci insegna che, guardando con attenzione il \emph{modello di calcolo} --- e
scegliendo bene il numero primo $q$ ---, si possono disegnare algoritmi nuovi: là la ricetta era
\emph{codifico in numeri e manipolo numeri}. L'algoritmo \emph{Shift--And} nasce dalla stessa
idea, con una codifica diversa: \emph{codifico in stringhe di bit e manipolo stringhe di bit}.
Il guadagno verrà dal parallelismo che la parola di macchina regala sulle operazioni bit a bit.

Fissiamo le notazioni per tutto il capitolo: il pattern è $P=\sub{P}{1}{n}$ e il testo è
$T=\sub{T}{1}{m}$, entrambi sull'alfabeto $\alfa$; dunque $\abs{P}=n$ e $\abs{T}=m$. In tutte le
matrici che seguono il pattern indicizza le \emph{righe} e il testo le \emph{colonne}.


\section{La matrice di bit (dot-plot)}

\begin{definizione}[Matrice di bit, o \emph{dot-plot}]
\label{def:shift-and-dotplot}
La \emph{matrice di bit} di $P$ contro $T$ è la matrice $D$ con $n$ righe --- una per ogni
posizione del pattern --- e $m$ colonne --- una per ogni posizione del testo --- definita da
\[
  D[i,j]\;=\;
  \begin{cases}
    1 & \text{se } P[i]=T[j] \quad \text{(\emph{match})},\\[2pt]
    0 & \text{se } P[i]\neq T[j] \quad \text{(\emph{mismatch})},
  \end{cases}
  \qquad 1\le i\le n,\quad 1\le j\le m .
\]
\end{definizione}

\begin{figure}[H]
\centering
% Dot-plot: la matrice di bit D di P (righe) contro T (colonne).
% La cella (i,j) e' individuata dall'incrocio fra la verticale punteggiata
% sotto T[j] e l'orizzontale tratteggiata a destra di P[i].
\begin{tikzpicture}[
    x=1cm, y=1cm,
    font=\footnotesize,
    gl/.style={font=\scriptsize, text=black!60, inner sep=1pt},
    et/.style={inner sep=2pt},
  ]

  % --- cornice appena accennata della matrice (nessuna griglia interna) ---
  \draw[gray!45, dashed, line width=0.3pt] (0,0) rectangle (6,4.2);

  % --- colonne: il testo ---
  \node[et, anchor=south] at (0.50,4.28) {$T[1]$};
  \node[et, anchor=south] at (1.35,4.28) {$T[2]$};
  \node[et, anchor=south] at (2.15,4.28) {$\cdots$};
  \node[et, anchor=south] at (3.60,4.28) {$T[j]$};
  \node[et, anchor=south] at (4.45,4.28) {$\cdots$};
  \node[et, anchor=south] at (5.60,4.28) {$T[m]$};

  % --- righe: il pattern ---
  \node[et, anchor=east] at (-0.18,3.78) {$P[1]$};
  \node[et, anchor=east] at (-0.18,3.22) {$P[2]$};
  \node[et, anchor=east] at (-0.18,2.68) {$\vdots$};
  \node[et, anchor=east] at (-0.18,1.89) {$P[i]$};
  \node[et, anchor=east] at (-0.18,1.12) {$\vdots$};
  \node[et, anchor=east] at (-0.18,0.42) {$P[n]$};

  % --- le due linee che individuano la cella (i,j) ---
  \draw[black!55, densely dotted] (3.60,4.20) -- (3.60,0);
  \draw[black!55, dashed]        (-0.05,1.89) -- (6.75,1.89);

  % --- tratto diagonale discendente: allineamento diretto ---
  \draw[line width=0.7pt]
       (0.35,3.92) -- (0.58,3.74) -- (0.72,3.66) -- (0.98,3.42)
    -- (1.12,3.34) -- (1.38,3.08) -- (1.58,2.90);
  \node[gl, anchor=south west] at (1.00,3.50) {\emph{match}};

  % --- tratto antidiagonale ascendente: allineamento invertito ---
  \draw[line width=0.7pt]
       (2.00,1.55) -- (2.22,1.76) -- (2.36,1.84) -- (2.62,2.10)
    -- (2.76,2.18) -- (3.02,2.44) -- (3.22,2.64);
  \node[gl, anchor=south east, align=center] at (3.32,2.70)
       {\emph{match}\\\emph{inverso}};

  % --- valore della cella e biforcazione 0/1 ---
  \node[et, anchor=west] at (6.82,1.89) {$0/1$};
  \node[gl, anchor=south] at (7.08,2.20) {$D[i,j]$};

  \draw[black!80] (7.48,1.99) -- (8.14,2.84);
  \draw[black!80] (7.48,1.79) -- (8.14,0.96);

  \node[et, anchor=west] at (8.16,2.90) {$0$};
  \node[et, anchor=west] at (8.16,0.90) {$1$};

  \node[et, anchor=west, align=left] at (8.60,2.90)
       {$P[i]\neq T[j]$\\[-2pt] {\scriptsize\itshape (mismatch)}};
  \node[et, anchor=west, align=left] at (8.60,0.90)
       {$P[i]=T[j]$\\[-2pt] {\scriptsize\itshape (match)}};

\end{tikzpicture}

\caption{La matrice di bit (\emph{dot-plot}) di $P$ contro $T$: righe indicizzate dal pattern,
colonne dal testo. La cella $(i,j)$ vale $1$ se e solo se $P[i]=T[j]$.}
\label{fig:shift-and-dotplot}
\end{figure}

\begin{osservazione}
Il dot-plot è l'oggetto con cui, in bioinformatica, si confrontano a occhio due sequenze: un
tratto \emph{diagonale} di $1$ --- che scende verso destra --- segnala un match fra un fattore di
$P$ e un fattore di $T$, mentre un tratto \emph{antidiagonale} --- che sale verso destra ---
segnala un match \emph{invertito}, cioè un match fra un fattore di $P$ e un fattore di $T$ letto
al contrario.
\end{osservazione}

Costruire la matrice per intero costa
\[
  \text{tempo per costruirla}\;\longrightarrow\;\Oh(n\cdot m),
  \qquad
  \text{spazio occupato}\;\longrightarrow\;\Oh(n\cdot m)\ \text{bit}.
\]

\begin{dubbio}
Accanto alla parola «bit», sottolineata, il quaderno riporta un'annotazione marginale di una o
due parole che non è stato possibile decifrare.
\end{dubbio}

E allora: questa matrice serve davvero a fare \emph{pattern matching}? Sì, ma non così com'è. Il
punto è che non occorre tenere in memoria un'intera matrice: basta muoversi \emph{per colonne},
un carattere di testo alla volta.

\begin{figure}[H]
\centering
% Il vettore di stato letto per colonne: dall'array al tempo j a quello al tempo j+1
% bastano uno SHIFT (gli 1 scendono di una posizione) e un AND con la maschera.
% ATTENZIONE: le colonne disegnate sono quelle della matrice M dello Shift--And,
% non del dot-plot D (per D la colonna j e' semplicemente la maschera U_{T[j]}).
\begin{tikzpicture}[
    x=1cm, y=1cm,
    font=\footnotesize,
    gl/.style={font=\scriptsize, text=black!60, inner sep=1pt},
    fr/.style={-{Stealth[length=4pt,width=3pt]}, black!70, line width=0.5pt},
  ]

  % --- la matrice, che non viene mai costruita per intero ---
  \draw[black!75] (0,0) rectangle (8.6,5);

  % --- etichette di orientamento (aggiunte per leggibilita') ---
  \node[gl, anchor=south west] at (0.05,5.05) {$T[1]$};
  \node[gl, anchor=south east] at (8.55,5.05) {$T[m]$};
  \node[gl, anchor=east] at (-0.10,4.72) {$P[1]$};
  \node[gl, anchor=east] at (-0.10,0.28) {$P[n]$};

  % --- le due colonne adiacenti j e j+1 ---
  \draw[black!75] (3.30,0) -- (3.30,5);
  \draw[black!75] (3.85,0) -- (3.85,5);
  \draw[black!75] (4.40,0) -- (4.40,5);

  % --- frecce dall'alto: chi sono le due strisce ---
  \draw[fr] (2.55,5.82) to[out=0,in=115] (3.55,5.02);
  \draw[fr] (5.62,5.82) to[out=180,in=65] (4.15,5.02);
  \node[anchor=east, align=center] at (2.45,5.88) {array al\\tempo $j$};
  \node[anchor=west, align=center] at (5.72,5.88) {array al\\tempo $j+1$};

  % --- l'1 scende di una posizione passando da j a j+1 ---
  \node[inner sep=1pt] at (3.575,2.80) {$1$};
  \draw[black!60, densely dotted, -{Stealth[length=3.5pt,width=2.5pt]}]
        (3.575,2.58) -- (3.575,2.28);
  \node[inner sep=1pt] at (4.125,2.42) {$1$};

  % --- un 1 sull'ultima riga: occorrenza completa di P ---
  \node[inner sep=1pt] at (6.00,0.32) {$1$};
  \node[gl, anchor=west] at (6.25,0.32) {occorrenza di $P$};

  % --- le due operazioni ---
  \draw[fr] (2.98,-0.70) to[out=0,in=250] (3.52,-0.05);
  \draw[fr] (4.78,-0.70) to[out=180,in=290] (4.20,-0.05);
  \node[anchor=east] at (2.90,-0.76) {\textsc{shift}};
  \node[anchor=west] at (4.86,-0.76) {\textsc{and}};

\end{tikzpicture}

\caption{Dall'array al tempo $j$ all'array al tempo $j+1$: uno \textsc{shift} verso il basso, che
fa scendere di una posizione gli $1$ già accesi, seguito da un \textsc{and}. L'$1$ disegnato
sull'ultima riga corrisponde a un'occorrenza del pattern. Le colonne disegnate sono quelle della
matrice $M$ della Sezione~\ref{sec:shift-and-matrice}, non del dot-plot $D$.}
\label{fig:shift-and-colonne}
\end{figure}

Ogni colonna del dot-plot è un array di $n$ bit, uno per ciascuna posizione del pattern, e non
costa nulla ottenerla: la colonna $j$ dipende dal solo carattere $T[j]$. Non è però questa la
colonna che conviene far evolvere --- in $D$ un $1$ dice soltanto che due singoli caratteri
coincidono, non che un'occorrenza del pattern è arrivata fin lì. Quello che l'algoritmo tiene in
vita è un unico array di $n$ bit che riassume lo stato dell'intera ricerca: la colonna della
matrice $M$ che definiremo nella Sezione~\ref{sec:shift-and-matrice}. Per passare dall'array al
tempo $j$ a quello al tempo $j+1$ bastano allora due operazioni bit a bit: uno \textsc{shift},
che fa scendere di una posizione tutti gli $1$ già accesi, e un \textsc{and} con una maschera che
dipende dal carattere di testo appena letto --- la maschera della
Definizione~\ref{def:shift-and-maschere}. Un $1$ che arrivi fino all'ultima riga segnalerà
un'occorrenza di $P$ in $T$.


\section{Le maschere del pattern}

\scan{30}
Il pattern è «assimilabile» a una sequenza di $n$ bit, e questo suggerisce di precalcolare, una
volta per tutte, dove ciascun carattere dell'alfabeto compare in $P$.

\begin{definizione}[Maschere del pattern]
\label{def:shift-and-maschere}
Per ogni carattere $x\in\alfa$ si definisce il vettore di bit $U_x\in\set{0,1}^n$ ponendo
\[
  U_x[i]\;=\;
  \begin{cases}
    1 & \text{se } P[i]=x,\\[2pt]
    0 & \text{altrimenti,}
  \end{cases}
  \qquad i=1,\dots,n .
\]
\end{definizione}

La tabella $\set{U_x}_{x\in\alfa}$ mi dice \emph{quale} carattere occorre in $P[i]$: fissata la
posizione $i$, c'è un solo $x$ per cui $U_x[i]=1$, ed è $x=P[i]$. Le maschere si costruiscono una
volta per tutte in preprocessing, a partire dal solo pattern, e occupano $\abs{\alfa}\cdot n$ bit.

Uso allora $\vec{P}$ (il pattern visto come vettore di bit), $T[j]$ e le maschere $U_x[i]$ per
determinare tutte le occorrenze di $P$ in $T$ \emph{senza} costruire il dot-plot: la $j$-esima
colonna del dot-plot è esattamente $U_{T[j]}$, quindi è già precalcolata e la matrice $D$ non va
mai materializzata.


\section{La matrice dello Shift--And}
\label{sec:shift-and-matrice}

Consideriamo ora la matrice $M$ dello Shift--And, matrice che \emph{non} costruiamo per intero ma
che ci permette di risolvere il problema:
\[
  M[i,j]=1 \iff \pre{P}{i}\ \text{è un suffisso di}\ \pre{T}{j},
  \qquad 0\le i\le n,\quad 0\le j\le m ,
\]
con la convenzione $\pre{P}{0}=\vuota$: la riga $i=0$ è allora identicamente $1$ --- la stringa
vuota è suffisso di ogni $\pre{T}{j}$ --- e la colonna $j=0$ è nulla dalla riga $1$ in giù.
Detto altrimenti, $M[i,j]=1$ se e solo se c'è un'occorrenza del prefisso $\pre{P}{i}$ che termina
esattamente in posizione $j$; in particolare gli $1$ dell'ultima riga, quelli con $i=n$, sono
tutte e sole le occorrenze complete di $P$ in $T$. Attenzione a non confondere $M$ con il
dot-plot $D$: dove entrambe sono definite, cioè per $i,j\ge 1$, da $M[i,j]=1$ segue in
particolare $P[i]=T[j]$, quindi $M\le D$ componente per componente, ma non vale il viceversa.

Le condizioni per scrivere correttamente un $1$ in $M[i,j]$ sono due:
\begin{itemize}
  \item[a)] devo avere un $1$ in $M[i-1,\,j-1]$ \quad$\longrightarrow$\quad \textsc{shift};
  \item[b)] deve essere $P[i]=T[j]$ \quad$\longrightarrow$\quad \textsc{and}.
\end{itemize}
Le due condizioni non sono che la lettura ricorsiva della definizione: una stringa di lunghezza
$i$ è suffisso di $\pre{T}{j}$ se e solo se il suo ultimo carattere è $T[j]$ e il suo prefisso di
lunghezza $i-1$ è suffisso di $\pre{T}{j-1}$.

Lette colonna per colonna --- sia $M_j\in\set{0,1}^n$ la $j$-esima colonna di $M$ ristretta alle
righe $1,\dots,n$, cosicché $M_j[i]=M[i,j]$ per $i=1,\dots,n$ --- le due condizioni diventano due
sole operazioni su vettori di bit:
\begin{itemize}
  \item la (a) chiede, per ogni $i$, il bit di riga $i-1$ della colonna precedente: basta far
    scorrere $M_{j-1}$ di una posizione verso il basso, immettendo in testa (\emph{pad}) un $1$,
    \[ \operatorname{shift}(b_1b_2\cdots b_n)\;=\;1\,b_1b_2\cdots b_{n-1}\,; \]
  \item la (b) chiede, per ogni $i$, che sia $P[i]=T[j]$: è esattamente l'\textsc{and} bit a bit
    con la maschera $U_{T[j]}$.
\end{itemize}
In forma vettoriale la ricorrenza è dunque
\[
  M_j\;=\;\operatorname{shift}\bigl(M_{j-1}\bigr)\ \mathbin{\textsc{and}}\ U_{T[j]},
  \qquad j=1,\dots,m ,
\]
ed è da qui che l'algoritmo prende il nome: \emph{uno} shift e \emph{un} and per ogni carattere
del testo.

\begin{osservazione}
Il $1$ immesso in testa dallo \emph{shift} non è un dettaglio implementativo: è la riga $i=0$ di
$M$, quella che i vettori $M_j$ non contengono e che vale costantemente $1$, cioè il fatto che in
\emph{ogni} posizione del testo può iniziare una nuova occorrenza. Con uno shift che immette $0$
la colonna diventerebbe identicamente nulla e l'algoritmo non troverebbe mai nulla. Per la ragione simmetrica la colonna iniziale è nulla, $M_0=0^n$: nessun prefisso non
vuoto di $P$ è suffisso di $\pre{T}{0}=\vuota$.
\end{osservazione}


\section{L'algoritmo}

\begin{algorithm}[H]
\caption{\textsc{Shift-And}}\label{alg:shift-and}
\begin{algorithmic}[1]
\Require $P[1\,..\,n]$, $T[1\,..\,m]$
\Ensure le posizioni di inizio delle occorrenze di $P$ in $T$
\ForAll{$x\in\alfa$}
  \State $U_x \gets 0^n$
\EndFor
\For{$i \gets 1$ \textbf{to} $n$}
  \State $U_{P[i]}[i] \gets 1$ \Comment{$U_x[i]=1$ sse $P[i]=x$}
\EndFor
\State $M_0 \gets 0^n$
\For{$j \gets 1$ \textbf{to} $m$}
  \State $M_j \gets \operatorname{shift}\bigl(M_{j-1}\bigr) \mathbin{\textsc{and}} U_{T[j]}$
  \If{$M_j[n]=1$}
    \State \textbf{output} $j-n+1$
  \EndIf
\EndFor
\end{algorithmic}
\end{algorithm}

\begin{esempio}
Siano $P=\str{aab}$ (quindi $n=3$) e $T=\str{baab}$ (quindi $m=4$). Scrivendo i vettori di bit
dall'alto in basso, da $P[1]$ a $P[3]$, le maschere sono
\[
  U_{\str{a}}=110,\qquad U_{\str{b}}=001,
\]
e la matrice $M$ è quella della Figura~\ref{fig:shift-and-esempio}. Si parte da $M_0=000$ e:
\[
  \begin{aligned}
    M_1 &= \operatorname{shift}(000)\mathbin{\textsc{and}}U_{\str{b}} = 100\mathbin{\textsc{and}}001 = 000, \\
    M_2 &= \operatorname{shift}(000)\mathbin{\textsc{and}}U_{\str{a}} = 100\mathbin{\textsc{and}}110 = 100, \\
    M_3 &= \operatorname{shift}(100)\mathbin{\textsc{and}}U_{\str{a}} = 110\mathbin{\textsc{and}}110 = 110, \\
    M_4 &= \operatorname{shift}(110)\mathbin{\textsc{and}}U_{\str{b}} = 111\mathbin{\textsc{and}}001 = 001 .
  \end{aligned}
\]
L'unico $1$ dell'ultima riga sta in colonna $j=4$: c'è un'occorrenza di $P$ che termina in $4$ e
comincia in $j-n+1=2$, e infatti $\sub{T}{2}{4}=\str{aab}=P$.
\end{esempio}

\begin{figure}[H]
\centering
% La matrice M dello Shift--And per P = aab e T = baab.
% M[i,j] = 1  sse  P[1..i] e' suffisso di T[1..j].
\begin{tikzpicture}[
    x=1cm, y=1cm,
    font=\small,
    gl/.style={font=\scriptsize, text=black!60, inner sep=1pt},
    dia/.style={black!55, dashed, -{Stealth[length=4pt,width=3pt]}, line width=0.5pt},
  ]

  % --- fondo tenue sull'ultima riga (i = n) ---
  \fill[black!8] (0,-1.5) rectangle (3.75,-2.25);

  % --- griglia 3 x 5 (celle di 0.75 cm) ---
  \draw[gray!55, line width=0.4pt, step=0.75] (0,-2.25) grid (3.75,0);
  % separatore fra la colonna di inizializzazione j=0 e le colonne j>=1
  \draw[black!70, line width=1pt] (0.75,0.02) -- (0.75,-2.27);

  % --- intestazioni di colonna: indice j e carattere T[j] ---
  \node[gl, anchor=south] at (0.375,0.44) {$0$};
  \node[gl, anchor=south] at (1.125,0.44) {$1$};
  \node[gl, anchor=south] at (1.875,0.44) {$2$};
  \node[gl, anchor=south] at (2.625,0.44) {$3$};
  \node[gl, anchor=south] at (3.375,0.44) {$4$};

  \node[gl, anchor=south] at (0.375,0.06) {$M_0$};
  \node[anchor=south, inner sep=1pt] at (1.125,0.06) {\texttt{b}};
  \node[anchor=south, inner sep=1pt] at (1.875,0.06) {\texttt{a}};
  \node[anchor=south, inner sep=1pt] at (2.625,0.06) {\texttt{a}};
  \node[anchor=south, inner sep=1pt] at (3.375,0.06) {\texttt{b}};

  % --- intestazioni delle due colonnine di sinistra ---
  \node[gl, anchor=south] at (-1.05,0.06) {$i$};
  \node[gl, anchor=south] at (-0.42,0.06) {$P[i]$};

  % --- etichette di riga: indice i e carattere P[i] ---
  \node[gl] at (-1.05,-0.375) {$1$};
  \node[gl] at (-1.05,-1.125) {$2$};
  \node[gl] at (-1.05,-1.875) {$3$};
  \node[inner sep=1pt] at (-0.42,-0.375) {\texttt{a}};
  \node[inner sep=1pt] at (-0.42,-1.125) {\texttt{a}};
  \node[inner sep=1pt] at (-0.42,-1.875) {\texttt{b}};

  % --- contenuto della matrice ---
  % riga i=1
  \node at (0.375,-0.375) {$0$};
  \node at (1.125,-0.375) {$0$};
  \node at (1.875,-0.375) {$1$};
  \node at (2.625,-0.375) {$1$};
  \node at (3.375,-0.375) {$0$};
  % riga i=2
  \node at (0.375,-1.125) {$0$};
  \node at (1.125,-1.125) {$0$};
  \node at (1.875,-1.125) {$0$};
  \node at (2.625,-1.125) {$1$};
  \node at (3.375,-1.125) {$0$};
  % riga i=3
  \node at (0.375,-1.875) {$0$};
  \node at (1.125,-1.875) {$0$};
  \node at (1.875,-1.875) {$0$};
  \node at (2.625,-1.875) {$0$};
  \node at (3.375,-1.875) {$1$};

  % --- propagazione in diagonale del prefisso che cresce (lo shift) ---
  \draw[dia] (2.05,-0.55) -- (2.45,-0.95)
     node[pos=0.5, sloped, above=3pt, fill=white, inner sep=1pt,
          font=\scriptsize\itshape, text=black!60] {shift};
  \draw[dia] (2.80,-1.30) -- (3.13,-1.63);

  % --- l'ultima riga: gli 1 qui sono occorrenze ---
  \draw[decorate, decoration={brace, amplitude=4pt}, black!60]
        (3.85,-1.5) -- (3.85,-2.25);
  \node[gl, anchor=west, align=left] at (4.18,-1.875)
       {riga $n$: gli $1$ qui\\sono occorrenze di $P$};

  % --- l'occorrenza trovata ---
  \draw[accentdark, line width=0.9pt] (3.375,-1.875) circle (0.29);
  \draw[accentdark, -{Stealth[length=4pt,width=3pt]}, line width=0.5pt]
        (3.55,-2.14) to[out=300,in=170] (4.30,-2.82);
  \node[anchor=west, align=left, font=\scriptsize, text=accentdark] at (4.38,-2.82)
       {occorrenza: $\sub{T}{2}{4}=\str{aab}=P$,\\inizio in $j-n+1=2$};

  % --- maschere applicate colonna per colonna ---
  \node[gl, anchor=north] at (1.125,-2.36) {$U_{\str{b}}$};
  \node[gl, anchor=north] at (1.875,-2.36) {$U_{\str{a}}$};
  \node[gl, anchor=north] at (2.625,-2.36) {$U_{\str{a}}$};
  \node[gl, anchor=north] at (3.375,-2.36) {$U_{\str{b}}$};

  % --- promemoria delle maschere ---
  \node[gl, anchor=north west, align=left] at (-1.20,-2.50)
       {$U_{\str{a}}=110$\\$U_{\str{b}}=001$};

\end{tikzpicture}

\caption{La matrice $M$ dello Shift--And per $P=\str{aab}$ e $T=\str{baab}$. Ogni colonna si
ottiene dalla precedente con uno \emph{shift} (che immette un $1$ in testa) seguito da un
\textsc{and} con la maschera del carattere letto; l'unico $1$ dell'ultima riga segnala
l'occorrenza $\sub{T}{2}{4}=P$.}
\label{fig:shift-and-esempio}
\end{figure}


\section{Complessità}

\[
  \Oh(n\cdot m)\ \text{tempo},\qquad \Oh(n)\ \text{bit di spazio (\emph{working space})} .
\]

\begin{osservazione}
La stima $\Oh(n\cdot m)$ è una \emph{sovrastima}: conta le celle di $M$, che però non viene mai
costruita. Il ciclo esegue $m$ iterazioni, ciascuna con \emph{uno} shift e \emph{un} and su
vettori di $n$ bit; nel modello word-RAM con parola di $w$ bit il costo della ricerca è
\[
  \Th\!\left(m\left\lceil \frac{n}{w}\right\rceil\right),
\]
che vale $\Th(m)$ non appena $n\le w$, cioè quando il pattern ha «una lunghezza ragionevole»
(a margine, a matita, una nota in gran parte illeggibile richiama proprio questo:
\lacunainl{\dots\ \str{/64}\ \dots}, la parola di macchina da $64$ bit).
Analogamente, lo spazio $\Oh(n)$ bit è quello della sola colonna corrente $M_j$: va aggiunta la
tabella delle maschere, $\abs{\alfa}\cdot n$ bit, trascurabile per alfabeti di dimensione
costante. Il testo viene letto in \emph{streaming}, un carattere alla volta.
\end{osservazione}


\section{Variazioni sul tema}

Lo Shift--And è la prima delle due variazioni sul tema del pattern matching annunciate dopo
Karp--Rabin: quella che sfrutta il parallelismo della parola di macchina, impacchettando in una
stringa di bit lo stato dell'intera ricerca e facendolo avanzare con una sola istruzione di shift
e una sola di and. La seconda variazione è la \emph{ricerca approssimata}, cioè il pattern con
errori: non viene sviluppata qui e sarà ripresa più avanti, con le nozioni di distanza e di
allineamento. L'estensione bit-parallela naturale in quella direzione consiste nel mantenere
$k+1$ vettori invece di uno --- uno per ogni numero di errori ammessi da $0$ a $k$ --- ed è la
famiglia Shift-Add / Wu--Manber.

\chapter{Lowest Common Ancestor in tempo costante}

\scan{31}

Questo capitolo costruisce, seguendo Harel--Tarjan, una struttura dati che dopo un
pre-processing lineare risponde a ogni query ``qual è il più basso antenato comune di
due nodi?'' in tempo costante. Il percorso è: prima si motiva il problema a
partire dalle applicazioni dei suffix tree, poi lo si risolve sul caso facile
dell'albero binario completo, infine si riporta il caso generale a quello facile
tramite una mappa opportuna.

\section{Perché serve il \texorpdfstring{$\lca$}{lca}: applicazioni dei suffix tree}

I suffix tree hanno \emph{tante} applicazioni: il capitolo~7 di Gusfield ne raccoglie
diciassette, da APL1 ad APL17. Per alcune di esse serve saper calcolare il $\lca$: è
questo il motivo per cui si studia l'algoritmo di Harel--Tarjan.

La più semplice è la ricerca esatta. Dati il testo $T$ e il pattern $P$, si costruisce
$\STree(T)$ e si segue dalla radice il cammino etichettato $P$: il sottoalbero appeso
al punto in cui tale cammino termina ha una dimensione pari al numero di occorrenze di
$P$ in $T$ (le sue foglie \emph{sono} le occorrenze).

\begin{figure}[H]
  \centering
  % Ricerca esatta di P nello STree di T (Gusfield, APL1).
% Il cammino etichettato P scende dalla radice; il sottoalbero appeso al punto
% in cui il cammino termina ha per foglie le occorrenze di P in T.
\begin{tikzpicture}[>=Latex,line join=round]
  \footnotesize

  % ---- sagoma irregolare dello STree(T) ------------------------------------
  \draw[line width=0.7pt,rounded corners=1.2pt]
       (0.40,3.20) -- (0.15,2.78) -- (0.05,2.42) -- (-0.36,2.16) -- (-0.50,1.80)
    -- (-0.96,1.56) -- (-1.14,1.16) -- (-1.70,0.96) -- (-1.94,0.60)
    -- (-2.50,0.44) -- (-2.82,0.16) -- (-3.20,0.00);
  \draw[line width=0.7pt,rounded corners=1.2pt]
       (0.40,3.20) -- (0.72,2.90) -- (0.80,2.50) -- (1.16,2.34) -- (1.30,1.94)
    -- (1.72,1.78) -- (1.86,1.34) -- (2.30,1.14) -- (2.50,0.74)
    -- (2.96,0.50) -- (3.12,0.20) -- (3.30,0.00);

  % ---- livello delle foglie ------------------------------------------------
  \draw[line width=1.1pt] (-3.75,0) -- (3.75,0);
  \node[anchor=east] at (-3.85,0.04) {foglie};

  % ---- sottoalbero delle occorrenze (tratteggiato) -------------------------
  \begin{scope}
    \clip (0.35,1.75) -- (-1.30,0) -- (0.95,0) -- cycle;
    \foreach \k in {-3.4,-3.25,...,0.75}{%
      \draw[gray,line width=0.35pt] (\k,-0.30) -- (\k+2.60,2.30);}
  \end{scope}
  \draw[line width=0.7pt] (-1.30,0) -- (0.35,1.75) -- (0.95,0);

  % ---- cammino etichettato P ----------------------------------------------
  \draw[line width=0.9pt]
       (0.40,3.20) .. controls (0.60,2.92) and (0.16,2.68) .. (0.36,2.42)
                   .. controls (0.56,2.18) and (0.18,1.98) .. (0.35,1.75);
  \node[anchor=west] at (0.50,2.14) {$P$};

  % ---- fine del cammino ----------------------------------------------------
  \fill (0.35,1.75) circle (2.2pt);

  % ---- etichetta della sagoma ---------------------------------------------
  \draw[gray,line width=0.4pt] (-2.12,1.74) -- (-1.08,1.40);
  \node[anchor=east] at (-2.18,1.74) {$\STree(T)$};

  % ---- puntatore alla nota -------------------------------------------------
  \draw[line width=0.6pt,->]
       (-0.20,-0.08) .. controls (0.05,-0.42) and (-0.25,-0.56) .. (0.12,-0.86)
                     .. controls (0.44,-1.10) and (0.70,-0.94) .. (1.00,-1.14);
  \node[anchor=west,align=left] at (1.06,-1.16)
       {tante foglie quante le\\ occorrenze di $P$ in $T$};
\end{tikzpicture}

  \caption{Ricerca di $P$ nello \STree di $T$: il cammino etichettato $P$ scende dalla
    radice fino a un punto dell'albero; il sottoalbero (tratteggiato) appeso a quel
    punto ha tante foglie quante sono le occorrenze di $P$ in $T$.}
  \label{fig:stree-occorrenze}
\end{figure}

\begin{osservazione}
  Il testo viene scandito una volta sola: dopo la costruzione dello \STree, il tempo
  per i pattern successivi è proporzionale alla dimensione del pattern più il numero di
  occorrenze, cioè $\Th(\abs{P}+k)$ per un pattern $P$ con $k$ occorrenze.
\end{osservazione}

Con due testi si procede allo stesso modo, concatenandoli con due terminatori distinti
e non appartenenti all'alfabeto:
\[
  T_1\,\text{\$}\,T_2\,\text{\texteuro}
  \qquad\longrightarrow\qquad
  \STree\bigl(T_1\,\text{\$}\,T_2\,\text{\texteuro}\bigr).
\]
In preprocessing si etichetta ogni foglia a seconda che il suo suffisso contenga o no
il primo separatore: se lo contiene --- forma $\dots\text{\$}\dots\text{\texteuro}$ ---
il suffisso comincia dentro $T_1$; se non lo contiene --- forma
$\dots\text{\texteuro}$, senza \$ --- comincia dentro $T_2$.

\begin{figure}[H]
  \centering
  % Suffix tree generalizzato: STree(T_1 $ T_2 EURO).
% Si scende fino al nodo alpha e si leggono le foglie del suo sottoalbero;
% l'etichetta apposta in preprocessing dice se il suffisso comincia in T_1
% oppure in T_2 (Gusfield, APL4: piu' lunga sottostringa comune).
\begin{tikzpicture}[>=Latex,line join=round]
  \footnotesize

  % ---- sagoma dello STree --------------------------------------------------
  \draw[line width=0.7pt] (0,3.00) -- (-2.20,0) (0,3.00) -- (2.20,0);
  \node[anchor=west] at (0.55,2.72)
       {$\STree\bigl(T_1\,\text{\$}\,T_2\,\text{\texteuro}\bigr)$};
  \draw[gray,line width=0.4pt] (0.50,2.70) -- (0.16,2.62);

  % ---- livello delle foglie, percorso da sinistra a destra -----------------
  \draw[line width=1.0pt,->] (-2.55,0) -- (3.05,0);

  % ---- cammino dalla radice fino ad alpha ----------------------------------
  \draw[line width=0.9pt]
       (0,3.00) .. controls (-0.18,2.78) and (0.04,2.62) .. (-0.14,2.42)
                .. controls (-0.32,2.24) and (-0.12,2.20) .. (-0.30,2.05);
  \fill (-0.30,2.05) circle (2.0pt);
  \node[anchor=west] at (-0.18,2.14) {$\a$};

  % ---- sottoalbero radicato in alpha ---------------------------------------
  \draw[line width=0.7pt] (-0.95,0) -- (-0.30,2.05) -- (0.85,0);

  % ---- foglie del sottoalbero, etichettate in preprocessing ----------------
  \foreach \x/\lab in {-0.78/T_1,-0.26/T_2,0.26/T_2,0.70/T_1}{%
    \fill (\x,0) circle (1.4pt);
    \node[anchor=north,inner sep=2pt,font=\scriptsize] at (\x,-0.06) {$\lab$};}

  \draw[line width=0.5pt,decorate,decoration={brace,amplitude=4pt}]
       (0.90,-0.52) -- (-1.00,-0.52);
  \node[anchor=north] at (-0.05,-0.70) {foglie di $\albero{\a}$};

  % ---- nota ----------------------------------------------------------------
  \node[anchor=west,align=left] at (3.15,0)
       {lettura delle foglie:\\
        suffisso con \$ (forma $\dots\text{\$}\dots\text{\texteuro}$)
        $\Rightarrow$ dentro $T_1$\\
        suffisso senza \$ (forma $\dots\text{\texteuro}$)
        $\Rightarrow$ dentro $T_2$};
\end{tikzpicture}

  \caption{Nello \STree di $T_1\,\text{\$}\,T_2\,\text{\texteuro}$ si scende fino a un nodo $\a$ e si
    leggono le foglie del sottoalbero radicato in $\a$: le etichette apposte in
    preprocessing dicono se la stringa $\a$ occorre in $T_1$, in $T_2$ o in entrambi.}
  \label{fig:stree-generalizzato}
\end{figure}

\section{Il problema e il piano di lavoro}

\begin{definizione}[Lowest common ancestor]\label{def:lca}
  Dato un \emph{albero} $\T$ e due suoi \emph{nodi} $x,y$, si vuole trovare
  $\lca(x,y)$: il nodo $w$ più profondo tale che
  \[
    x,y \in \albero{w},
  \]
  dove $\albero{w}$ denota il sottoalbero di $\T$ radicato in $w$. In particolare
  $x,y\in\albero{\lca(x,y)}$.
\end{definizione}

L'obiettivo è rispondere a ogni query in tempo $\Oh(1)$, dopo un pre-processing che
resti lineare.

\scan{32}

\begin{esempio}
  Lavoreremo con alberi \emph{qualsiasi}, cioè non necessariamente binari né
  bilanciati: si veda l'albero della Figura~\ref{fig:lca-albero-dfs}.
\end{esempio}

\begin{figure}[H]
  \centering
  % Albero generico T con i nodi numerati in dfs-numbering (preordine).
% Accanto ai nodi 1, 2, 5 e' indicato l'intervallo dei numeri DFS del
% sottoalbero radicato in quel nodo: e' sempre un intervallo contiguo.
\begin{tikzpicture}[
    nodo/.style={circle,fill,inner sep=0pt,minimum size=3.6pt},
    arco/.style={line width=0.7pt},
    num/.style={inner sep=1.5pt},
    ran/.style={inner sep=1.5pt,text=gray,font=\scriptsize}]
  \footnotesize

  % ---- nodi ----------------------------------------------------------------
  \node[nodo] (n1)  at (3.00, 0.00) {};
  \node[nodo] (n2)  at (1.30,-1.05) {};
  \node[nodo] (n5)  at (4.40,-0.95) {};
  \node[nodo] (n3)  at (0.10,-2.05) {};
  \node[nodo] (n4)  at (1.40,-2.10) {};
  \node[nodo] (n6)  at (3.35,-2.10) {};
  \node[nodo] (n7)  at (4.15,-2.00) {};
  \node[nodo] (n8)  at (5.65,-1.90) {};
  \node[nodo] (n9)  at (5.15,-2.90) {};
  \node[nodo] (n10) at (5.95,-2.90) {};

  % ---- archi ---------------------------------------------------------------
  \draw[arco] (n1) -- (n2)  (n1) -- (n5)
              (n2) -- (n3)  (n2) -- (n4)
              (n5) -- (n6)  (n5) -- (n7)  (n5) -- (n8)
              (n8) -- (n9)  (n8) -- (n10);

  % ---- numeri DFS ----------------------------------------------------------
  \node[num,anchor=south east] at ([shift={(-0.04,0.06)}]n1.center)  {$1$};
  \node[num,anchor=south east] at ([shift={(-0.04,0.06)}]n2.center)  {$2$};
  \node[num,anchor=south west] at ([shift={( 0.06,0.02)}]n5.center)  {$5$};
  \node[num,anchor=north east] at ([shift={(-0.04,-0.06)}]n3.center) {$3$};
  \node[num,anchor=north west] at ([shift={( 0.06,-0.06)}]n4.center) {$4$};
  \node[num,anchor=north east] at ([shift={(-0.04,-0.06)}]n6.center) {$6$};
  \node[num,anchor=north]      at ([shift={( 0.00,-0.08)}]n7.center) {$7$};
  \node[num,anchor=west]       at ([shift={( 0.10,0.02)}]n8.center)  {$8$};
  \node[num,anchor=north]      at ([shift={( 0.00,-0.08)}]n9.center) {$9$};
  \node[num,anchor=north west] at ([shift={( 0.02,-0.08)}]n10.center){$10$};

  % ---- intervalli dei sottoalberi ------------------------------------------
  \node[ran,anchor=south west] at ([shift={( 0.08,0.10)}]n1.center) {$[1,10]$};
  \node[ran,anchor=south east] at ([shift={(-0.32,0.10)}]n2.center) {$[2,4]$};
  \node[ran,anchor=south west] at ([shift={( 0.30,0.06)}]n5.center) {$[5,10]$};
\end{tikzpicture}

  \caption{Un albero generico $\T$ con i nodi numerati secondo la visita in profondità
    (preordine). Accanto ai nodi $1$, $2$ e $5$ è riportato l'intervallo dei numeri DFS
    del rispettivo sottoalbero.}
  \label{fig:lca-albero-dfs}
\end{figure}

Numereremo i nodi in \emph{dfs-numbering}: a ogni nodo si assegna il suo numero
d'ordine nella visita in profondità (preordine). L'assegnazione costa \emph{tempo
lineare} e si può fare mentre si legge l'input. D'ora in avanti identifichiamo ogni
nodo con il proprio numero DFS.

\begin{osservazione}[Proprietà del dfs-numbering]
  I discendenti del nodo $v$ costituiscono un \emph{intervallo}
  \[
    [\,v+1,\ j\,],
  \]
  dove $j-v$ è il numero di tali discendenti. Equivalentemente, i numeri DFS dei nodi
  di $\albero{v}$ formano l'intervallo contiguo
  \[
    [\,v,\ v+\abs{\albero{v}}-1\,],
  \]
  cioè sono assegnati \emph{senza buchi}; in particolare
  $j = v + \abs{\albero{v}} - 1$.
\end{osservazione}

Il problema da risolvere è dunque il calcolo di
\[
  \lca(x,y) \qquad \text{in tempo } \Oh(1),
\]
dove $x$ e $y$ sono i numeri DFS di due nodi di $\T$. Lo strumento tecnico saranno le
\emph{bitmap}: rappresentando l'informazione con stringhe di bit si rendono efficienti
le operazioni di base.

\paragraph{Piano.}
\begin{enumerate}
  \item Risolvo il problema su $\B$, albero binario \emph{completo} (caso facile).
  \item Definisco una mappa
  \[
    I \colon \T \longrightarrow \B, \qquad v \longmapsto I(v),
  \]
  in generale \emph{non} iniettiva.
  \item Uso la mappa e la soluzione in $\Oh(1)$ su $\B$ per risolvere il problema su
  $\T$.
\end{enumerate}

\section{Un modello di calcolo a parole logaritmiche}\label{sec:lca-modello}

\scan{33}

Prima di attaccare il problema devo mettere in discussione il modello di calcolo, cioè
dichiarare esplicitamente quali operazioni considero a costo unitario. In effetti non
lo sto davvero modificando: quelle che seguono sono le assunzioni standard del modello
RAM con parole di lunghezza logaritmica.

Gli assegnamenti
\[
  i \leftarrow i+j, \qquad
  i \leftarrow k, \qquad
  i \leftarrow i/k, \qquad \dots
\]
costano $\Oh(1)$ \emph{se} gli interi coinvolti sono esprimibili mediante $c\log n$
bit, dove $n$ è la dimensione dell'input (per noi: il numero di nodi di $\T$).

Assumo inoltre di poter eseguire in tempo $\Oh(1)$ le seguenti operazioni su array di
$c\log n$ bit, che chiameremo \emph{parole}:
\begin{itemize}
  \item \textsc{and}, \textsc{or}, \textsc{xor} (bit a bit);
  \item \textsc{shift} (\emph{left}/\emph{right}); scriveremo $a \ll s$ e $a \gg s$ per
        lo scorrimento di $a$ di $s$ posizioni verso sinistra e verso destra;
  \item costruzione di una \textsc{mask};
  \item individuazione della posizione dell'1 significativo più a sinistra
        (\emph{leftmost}) o più a destra (\emph{rightmost}).
\end{itemize}

Le posizioni dei bit di una parola sono sempre numerate \emph{da destra} a partire
da~$1$: la posizione $1$ è quella del bit meno significativo. Nella
Sezione~\ref{sec:lca-primitive}, a capitolo concluso, verificheremo che tutte queste
primitive si realizzano davvero in $\Oh(1)$ a partire dalle sole operazioni
aritmetiche, al prezzo di un pre-processing $\Oh(n)$.

\section{Passo 1: il \texorpdfstring{$\lca$}{lca} in un albero binario completo}
\label{sec:lca-passo1}

È il caso facile del piano. Se $\B$ ha $p$ foglie, i suoi livelli contengono
rispettivamente $1, 2, 4, 8, \dots$ nodi e sono dunque in numero logaritmico.

\begin{figure}[H]
  \centering
  % L'albero binario completo B: p foglie, livelli di 1, 2, 4, 8, ... nodi,
% profondita' d = log_2 p.  Sono esplicitati solo i primi tre livelli; i due
% rami esterni (= lati del triangolo) proseguono fino alla linea ondulata
% delle foglie, a indicare che l'albero continua.
\begin{tikzpicture}[
    nodo/.style={circle,draw,fill=white,line width=0.6pt,
                 inner sep=0pt,minimum size=4.4pt},
    arco/.style={line width=0.7pt}]
  \footnotesize

  % ---- rami esterni: coincidono con i lati del triangolo -------------------
  \draw[arco] (0,2.80) -- (-2.70,0);
  \draw[arco] (0,2.80) -- ( 2.70,0);

  % ---- archi interni dei primi tre livelli ---------------------------------
  \draw[arco] (-0.72,2.05) -- (-0.42,1.26);
  \draw[arco] ( 0.72,2.05) -- ( 0.42,1.26);

  % ---- nodi dei primi tre livelli ------------------------------------------
  \node[nodo] at ( 0.00,2.80) {};
  \node[nodo] at (-0.72,2.05) {};
  \node[nodo] at ( 0.72,2.05) {};
  \node[nodo] at (-1.44,1.30) {};
  \node[nodo] at (-0.42,1.26) {};
  \node[nodo] at ( 0.42,1.26) {};
  \node[nodo] at ( 1.44,1.30) {};

  % ---- livello delle foglie: linea ondulata --------------------------------
  \draw[line width=0.9pt,smooth,samples=160,domain=-2.70:2.70]
        plot (\x,{0.075*sin(\x*235)+0.030*sin(\x*470)});
  \node[anchor=north] at (0,-0.28) {$p$};

  \node at (-1.55,0.50) {$\B$};

  % ---- taglia dei livelli e profondita' ------------------------------------
  \node at (3.30,2.80) {$1$};
  \node at (3.30,2.05) {$2$};
  \node at (3.30,1.30) {$4$};
  \node at (3.30,0.55) {$8$};
  \node at (3.30,0.00) {$\vdots$};
  \draw[line width=0.5pt,decorate,decoration={brace,amplitude=5pt}]
       (3.72,2.98) -- (3.72,-0.22);
  \node[anchor=west] at (4.02,1.38) {$d=\log_2 p$};
\end{tikzpicture}

  \caption{L'albero binario completo $\B$: $p$ foglie, livelli di $1,2,4,8,\dots$ nodi,
    profondità $d=\log_2 p$.}
  \label{fig:b-albero-completo}
\end{figure}

\scan{34}

Riassumendo:
\[
  \B:\qquad
  p \ \text{foglie},\qquad
  2p-1 \ \text{nodi}\ \ (\text{di cui } p-1 \ \text{interni}),\qquad
  d=\log_2 p \ \text{di profondità}.
\]

Ogni nodo $v$ di $\B$ è individuato dal cammino che lo raggiunge a partire dalla
radice: un bit per arco, con la convenzione $0=\text{scendo a sinistra}$,
$1=\text{scendo a destra}$. Per avere codifiche tutte della stessa lunghezza, la
stringa del cammino si completa con un $1$ e con tanti $0$ quanti ne servono: un
generico nodo di $\B$ si codifica così con $d+1$ bit.

\begin{figure}[H]
  \centering
  % path-number-alpha -- codifica di un nodo v di B: i bit del cammino radice->v
% (la stringa alpha), poi il terminatore 1, poi gli 0 di riempimento (quaderno, pag. 33).
\begin{tikzpicture}[
    >={Stealth[length=5pt,width=4pt]},
    font=\small,
    sagoma/.style={semithick,draw=black!65},
    nodo/.style={circle,fill=black,inner sep=0pt,minimum size=3pt},
    parola/.style={thick},
    sep/.style={semithick},
  ]

  %% ------------------------------------------------------------------
  %% Pannello sinistro: la sagoma dell'albero binario completo B
  %% ------------------------------------------------------------------
  \draw[sagoma]
      (-2.20,0.00)
        to[out=68,in=228]  (-1.55,1.05)
        to[out=48,in=252]  (-1.00,1.95)
        to[out=72,in=232]  (-0.50,2.90)   % apice = radice
        to[out=-42,in=152] (-0.02,2.25)
        to[out=-28,in=172] ( 0.52,1.95)
        to[out=-14,in=142] ( 1.20,1.00)
        to[out=-44,in=118] ( 1.90,0.00)
      % base ondulata
        to[out=186,in=-6]  ( 0.75,-0.07)
        to[out=174,in=-6]  (-0.55, 0.06)
        to[out=174,in=-12] (-2.20, 0.00);

  % cammino radice -> v (linea leggermente ondulata)
  \draw[semithick,->]
      (-0.50,2.82)
        to[out=-86,in=94] (-0.42,2.45)
        to[out=-86,in=94] (-0.58,2.10)
        to[out=-86,in=94] (-0.42,1.78)
        to[out=-86,in=98] (-0.50,1.52);
  \node[nodo] (v) at (-0.50,1.44) {};
  \node[right=2pt] at (v) {$v$};
  \node[anchor=west] at (-0.34,2.06) {$\a$};

  %% ------------------------------------------------------------------
  %% Pannello destro: il path number di v come parola di d+1 bit
  %% ------------------------------------------------------------------
  \def\ay{1.10}          % quota della parola
  \def\xa{2.60}          % delimitatore sinistro
  \def\xb{5.40}          % fine della zona alpha
  \def\xc{5.85}          % fine della cella del terminatore
  \def\xd{8.20}          % delimitatore destro

  \node[anchor=west,font=\footnotesize] at (\xa,\ay+0.78) {path number di $v$};

  \draw[parola] (\xa,\ay) -- (\xd,\ay);
  \draw[parola] (\xa,\ay-0.16) -- (\xa,\ay+0.16);
  \draw[parola] (\xd,\ay-0.16) -- (\xd,\ay+0.16);
  \draw[sep]    (\xb,\ay-0.11) -- (\xb,\ay+0.32);
  \draw[sep]    (\xc,\ay-0.11) -- (\xc,\ay+0.32);

  \node[above] at ({0.5*(\xa+\xb)},\ay+0.04) {$\a$};
  \node[above] at ({0.5*(\xb+\xc)},\ay+0.04) {$1$};
  \foreach \i in {0,...,6}{%
    \node[above] at ({\xc+(\i+0.5)*(\xd-\xc)/7},\ay+0.04) {$0$};}

  \draw[decorate,decoration={brace,amplitude=5pt,mirror,raise=2pt},draw=black!65]
      (\xa,\ay-0.16) -- (\xd,\ay-0.16)
      node[midway,below=8pt,font=\footnotesize]{$d+1$ bit};

\end{tikzpicture}

  \caption{Codifica di un nodo $v$ di $\B$: i bit $\a$ del cammino radice--$v$, poi
    un $1$, poi tutti $0$.}
  \label{fig:path-number-alpha}
\end{figure}

\begin{definizione}[Path number]\label{def:path-number}
  Sia $v$ un nodo di $\B$ a profondità $t$ (la radice ha $t=0$) e sia
  $\a=b_1b_2\cdots b_t$ la sequenza di bit che descrive il cammino dalla radice a $v$.
  Il \emph{path number} di $v$ è la parola di $d+1$ bit
  \[
    \underbrace{\ \underbrace{b_1b_2\cdots b_t}_{\a}\ \underbrace{1}_{\text{terminatore}}\ \underbrace{0\cdots 0}_{d-t}\ }_{d+1\ \text{bit}}
  \]
\end{definizione}

\begin{figure}[H]
  \centering
  % lca-b-path-number -- albero binario completo con p=4 foglie e d=2: ogni nodo porta
% il proprio path number in binario e, sotto, in base dieci (quaderno, pag. 34).
\begin{tikzpicture}[
    >={Stealth[length=5pt,width=4pt]},
    font=\small,
    nodo/.style={circle,fill=black,inner sep=0pt,minimum size=3.4pt},
    pn/.style={align=center,inner sep=1.5pt,font=\footnotesize},
    arco/.style={semithick},
  ]

  %% ---- nodi ---------------------------------------------------------
  \node[nodo] (r)  at ( 0.00,2.60) {};   % radice
  \node[nodo] (s)  at (-1.70,1.30) {};   % figlio sinistro
  \node[nodo] (d)  at ( 1.70,1.30) {};   % figlio destro
  \node[nodo] (f1) at (-2.55,0.00) {};
  \node[nodo] (f2) at (-0.85,0.00) {};
  \node[nodo] (f3) at ( 0.85,0.00) {};
  \node[nodo] (f4) at ( 2.55,0.00) {};

  %% ---- archi --------------------------------------------------------
  \draw[arco] (r) -- (s);  \draw[arco] (r) -- (d);
  \draw[arco] (s) -- (f1); \draw[arco] (s) -- (f2);
  \draw[arco] (d) -- (f3); \draw[arco] (d) -- (f4);

  %% ---- path number: binario sopra, valore decimale sotto ------------
  \node[pn,above=3pt]      at (r)  {$100$\\$4$};
  \node[pn,left=4pt]       at (s)  {$010$\\$2$};
  \node[pn,right=4pt]      at (d)  {$110$\\$6$};
  \node[pn,below=3pt]      at (f1) {$001$\\$1$};
  \node[pn,below=3pt]      at (f2) {$011$\\$3$};
  \node[pn,below=3pt]      at (f3) {$101$\\$5$};
  \node[pn,below=3pt]      at (f4) {$111$\\$7$};

  %% ---- annotazione "d+1 bit" sulla codifica della radice ------------
  \node[anchor=west,font=\footnotesize] (bit) at (0.95,3.55) {$d+1$ bit};
  \draw[->,semithick,shorten >=2pt] (bit.west) to[out=180,in=25] (0.28,3.32);

  %% ---- graffa: i tre livelli, profondita' d=2 -----------------------
  \draw[decorate,decoration={brace,amplitude=6pt,raise=2pt},draw=black!65]
      (3.45,2.60) -- (3.45,0.00)
      node[midway,right=9pt,font=\footnotesize]{$d=2$};

\end{tikzpicture}

  \caption{Albero binario completo con $p=4$ foglie, $d=2$: ogni nodo è etichettato con
    il proprio path number (in binario e, sotto, in base dieci).}
  \label{fig:lca-b-path-number}
\end{figure}

Dato $\B$, assegno ad ogni nodo il suo path number e lo manipolo come array di bit: è
un pre-processing lineare. Da qui in avanti identifichiamo ogni nodo di $\B$ con il
proprio path number.

\begin{domanda}
  Come si calcola $\lca_{\B}(x,y)$ in $\B$?
\end{domanda}

\begin{figure}[H]
  \centering
  % lca-b-triangolo -- z = lca_B(x,y): il cammino radice->z e' gamma, e da z i due
% cammini verso x e verso y si separano subito (quaderno, pag. 34).
\begin{tikzpicture}[
    font=\small,
    nodo/.style={circle,fill=black,inner sep=0pt,minimum size=3.4pt},
    livello/.style={semithick,draw=black!55},
    cammino/.style={semithick,rounded corners=2pt},
  ]

  %% ---- la sagoma triangolare di B -----------------------------------
  \coordinate (apice) at (0,4.40);
  \draw[semithick] (-3.00,0) -- (apice) -- (3.00,0) -- cycle;
  \node[font=\footnotesize,text=black!70] at (-2.05,0.55) {$\B$};

  %% ---- i due livelli: quello della radice e quello di z -------------
  \draw[livello] (-1.10,4.40) -- (2.50,4.40);
  \draw[livello] (-1.60,3.35) -- (2.50,3.35);

  % misura della profondita' di z, cioe' della lunghezza di gamma
  \draw[<->,>={Stealth[length=4.5pt,width=3.5pt]},semithick]
      (2.15,4.36) -- (2.15,3.39)
      node[midway,right=2pt,font=\footnotesize]{$\abs{\y}$};

  %% ---- il cammino radice -> z ---------------------------------------
  \node[nodo] (z) at (0,3.35) {};
  \node[right=3pt] at (z) {$z$};
  \draw[line width=1.1pt,-{Stealth[length=6pt,width=5pt]}]
      (0,3.44) -- (0,4.32) node[pos=0.15,left=2pt]{$\y$};

  %% ---- da z i due cammini divergono subito --------------------------
  \draw[cammino]
      (0,3.28) -- (-0.42,2.94) -- (-0.24,2.60) -- (-0.76,2.26)
             -- (-0.58,1.96) -- (-1.20,1.78) -- (-1.55,1.60);
  \draw[cammino]
      (0,3.28) -- ( 0.36,2.94) -- ( 0.18,2.60) -- ( 0.66,2.26)
             -- ( 0.48,1.96) -- ( 1.00,1.78) -- ( 1.25,1.60);

  \node[nodo] (x) at (-1.55,1.60) {};
  \node[nodo] (y) at ( 1.25,1.60) {};
  \node[below left=0pt and 1pt]  at (x) {$x$};
  \node[below right=0pt and 1pt] at (y) {$y$};

\end{tikzpicture}

  \caption{$z=\lca_{\B}(x,y)$ in $\B$: il cammino dalla radice a $z$ è $\y$, e da $z$ i
    due cammini verso $x$ e verso $y$ si separano subito.}
  \label{fig:lca-b-triangolo}
\end{figure}

Siano $x$ e $y$ due nodi distinti di $\B$, nessuno dei due antenato dell'altro, e sia
$z=\lca_{\B}(x,y)$. Detto $\y$ il cammino dalla radice a $z$, i cammini che portano a
$x$ e a $y$ hanno entrambi $\y$ come prefisso e subito dopo divergono: uno prosegue a
sinistra (bit $0$) e l'altro a destra (bit $1$). Se $x<y$, sui path number questo si
legge così:
\[
  \begin{aligned}
    x &: \ \y\ 0\ \ast\cdots\ast\ 1\ 0\cdots 0,\\
    y &: \ \y\ 1\ \ast\cdots\ast\ 1\ 0\cdots 0,\\
    z &: \ \y\ 1\ 0\cdots\cdots\cdots\cdots 0,
  \end{aligned}
\]
dove $\ast\cdots\ast$ sono blocchi di bit arbitrari e ciascun $1$ evidenziato è il
terminatore del rispettivo nodo.

\begin{figure}[H]
  \centering
  % lca-b-array-bit -- i path number di x, y e z=lca_B(x,y) come array di d+1 bit:
% la prima colonna in cui x e y differiscono e' quella del terminatore di z
% (quaderno, pag. 34).
\begin{tikzpicture}[
    >={Stealth[length=4.5pt,width=3.5pt]},
    font=\small,
    parola/.style={thick},
    tacca/.style={semithick},
    etichetta/.style={anchor=east,font=\footnotesize},
    guida/.style={dashed,draw=black!55,thin},
  ]

  %% ---- geometria comune ---------------------------------------------
  \def\xa{0.00}   % delimitatore sinistro
  \def\xb{2.80}   % fine del prefisso comune
  \def\xc{3.30}   % fine della colonna di divergenza
  \def\xd{5.20}   % fine del blocco di bit arbitrari
  \def\xe{5.70}   % fine del terminatore di x e di y
  \def\xf{8.40}   % delimitatore destro
  \def\ra{3.60}   % riga di x
  \def\rb{2.10}   % riga di y
  \def\rc{0.30}   % riga di z

  % la colonna in cui x e y divergono, evidenziata su tutte e tre le righe
  \fill[black!6] (\xb,\rc-0.45) rectangle (\xc,\ra+0.90);

  % una parola: linea spessa con i due delimitatori agli estremi
  \newcommand{\parolabit}[1]{%
    \draw[parola] (\xa,#1) -- (\xf,#1);
    \draw[parola] (\xa,#1-0.16) -- (\xa,#1+0.16);
    \draw[parola] (\xf,#1-0.16) -- (\xf,#1+0.16);}
  % una tacca interna di separazione fra celle
  \newcommand{\tacca}[2]{\draw[tacca] (#1,#2-0.10) -- (#1,#2+0.34);}
  % contenuto di una cella, scritto sopra il segmento
  \newcommand{\cella}[4]{\node[above] at ({0.5*(#1+#2)},#3+0.04) {#4};}

  %% ---- riga 1: path number di x -------------------------------------
  \parolabit{\ra}
  \tacca{\xb}{\ra} \tacca{\xc}{\ra} \tacca{\xd}{\ra} \tacca{\xe}{\ra}
  \cella{\xa}{\xb}{\ra}{$\a$}
  \cella{\xb}{\xc}{\ra}{$0$}
  \cella{\xc}{\xd}{\ra}{$*\cdots*$}
  \cella{\xd}{\xe}{\ra}{$1$}
  \cella{\xe}{\xf}{\ra}{$0\cdots 0$}
  \node[etichetta] at (\xa-0.25,\ra) {p.n.\ di $x$};

  %% ---- riga 2: path number di y -------------------------------------
  \parolabit{\rb}
  \tacca{\xb}{\rb} \tacca{\xc}{\rb} \tacca{\xd}{\rb} \tacca{\xe}{\rb}
  \cella{\xa}{\xb}{\rb}{$\b$}
  \cella{\xb}{\xc}{\rb}{$1$}
  \cella{\xc}{\xd}{\rb}{$*\cdots*$}
  \cella{\xd}{\xe}{\rb}{$1$}
  \cella{\xe}{\xf}{\rb}{$0\cdots 0$}
  \node[etichetta] at (\xa-0.25,\rb) {p.n.\ di $y$};

  %% ---- riga 3: path number di z = lca(x,y) --------------------------
  \parolabit{\rc}
  \tacca{\xb}{\rc} \tacca{\xc}{\rc}
  \cella{\xa}{\xb}{\rc}{$\y$}
  \cella{\xb}{\xc}{\rc}{$1$}
  \cella{\xc}{\xf}{\rc}{$0\cdots\cdots\cdots 0$}
  \node[etichetta] at (\xa-0.25,\rc) {p.n.\ di $z$};

  %% ---- la colonna di divergenza attraversa le tre righe -------------
  \draw[guida] (\xb,\rc-0.45) -- (\xb,\ra+0.90);
  \draw[guida] (\xc,\rc-0.45) -- (\xc,\ra+0.90);
  \node[font=\footnotesize] at ({0.5*(\xb+\xc)},\ra+1.12) {colonna $\ell$};

  %% ---- il prefisso comune di x e y e' esattamente gamma -------------
  \draw[decorate,decoration={brace,amplitude=5pt,mirror,raise=2pt},draw=black!65]
      (\xa,\rb-0.16) -- (\xb,\rb-0.16)
      node[midway,below=8pt,font=\footnotesize] (gm) {$\a=\b=\y$};
  \coordinate (divg) at ({0.5*(\xb+\xc)},\rb-0.32);
  \draw[->,semithick,draw=black!65] (gm.east) to[out=0,in=250] (divg);

\end{tikzpicture}

  \caption{I path number di $x$, $y$ e $z=\lca_{\B}(x,y)$ visti come array di bit. La
    colonna evidenziata è la prima in cui $x$ e $y$ differiscono: lì cade esattamente
    il terminatore di $z$.}
  \label{fig:lca-b-array-bit}
\end{figure}

La prima posizione (da sinistra) in cui i due path number differiscono è dunque
esattamente la posizione del terminatore di $z$: il path number di $z$ si ottiene
ricopiando i bit su cui $x$ e $y$ concordano, mettendo $1$ in quella posizione e $0$ in
tutte le posizioni successive. Sia $\ell$ tale posizione, contata da destra: è la
posizione dell'$1$ più significativo di $x \oplus y$.

\begin{algorithm}[H]
  \caption{$\lca_{\B}$ in un albero binario completo (nessuno dei due nodi antenato
    dell'altro)}
  \label{alg:lca-b}
  \begin{algorithmic}[1]
    \Require i path number $x,y$ di due nodi di $\B$, ciascuno su $d+1$ bit, nessuno
      dei due antenato proprio dell'altro
    \Ensure il path number di $\lca_{\B}(x,y)$
    \State $s \gets x \oplus y$ \Comment{1 \textsc{xor}}
    \If{$s = 0$} \Comment{$x$ e $y$ sono lo stesso nodo}
      \State \Return $x$
    \EndIf
    \State $\ell \gets \textsc{Leftmost1}(s)$ \Comment{1 \emph{leftmost} $1$}
    \State $m \gets \max\{x,y\}$ \Comment{$m$ ha il bit $1$ in posizione $\ell$}
    \State \Return $\bigl(m \gg (\ell-1)\bigr) \ll (\ell-1)$ \Comment{2 \emph{shift}}
  \end{algorithmic}
\end{algorithm}

In forma simmetrica (senza il confronto) lo stesso risultato si scrive
$\bigl((x \gg \ell)\ll \ell\bigr)\ \textsc{or}\ 2^{\ell-1}$. Il conto delle operazioni
è
\[
  1\ \textsc{xor} \ +\ 1\ \textit{leftmost } 1 \ +\ 2\ \textit{shift}
  \qquad\Longrightarrow\qquad \Oh(1) \ \text{nel nuovo modello di calcolo.}
\]

\begin{esercizio}
  Risolvere il caso in cui $x$ è antenato o discendente di $y$.
\end{esercizio}

\begin{osservazione}
  L'esercizio non è un caso di scuola: l'algoritmo di query della
  Sezione~\ref{sec:lca-query} invocherà $\lca_{\B}$ su $I(x)$ e $I(y)$, e nulla
  impedisce che quei due path number siano uno antenato dell'altro in $\B$; accade
  anzi spesso che coincidano, cioè che $I(x)=I(y)$, ed è per questo che
  l'Algoritmo~\ref{alg:lca-b} tratta a parte il caso $s=0$. La versione generale resta
  comunque $\Oh(1)$: si verifica che la posizione da cui ripartire non è $\ell$, ma il
  massimo fra $\ell$ e le posizioni dell'$1$ \emph{meno} significativo di $x$ e di $y$.
\end{osservazione}

\section{Passo 2: la mappa \texorpdfstring{$I\colon\T\to\B$}{I : T -> B}}
\label{sec:lca-passo2}

\scan{35}

Il secondo passo consiste nel mandare l'albero generico $\T$ dentro l'albero binario
completo $\B$, per poter riusare l'algoritmo del passo~1. La mappa che costruiremo,
\[
  I \colon \T \longrightarrow \B, \qquad v \longmapsto I(v),
\]
\emph{non} sarà iniettiva.

\begin{osservazione}
  $\B$ non verrà davvero costruito: resta un albero \emph{implicito}. I suoi nodi
  \emph{sono} i loro path number, e ogni operazione su di essi è un'operazione su
  parole di bit.
\end{osservazione}

Identifico i nodi di $\T$ con i \emph{dfs-numeri} loro assegnati in una fase di
pre-processing, cioè con i numeri prodotti da una visita in profondità: una fase che
resta sempre lineare.

\begin{figure}[H]
  \centering
  % L'albero generico T con i dfs-numeri e la loro codifica binaria su 4 bit.
\begin{tikzpicture}[
    x=1.22cm, y=1.20cm,
    nodo/.style={circle, fill=black, inner sep=0pt, minimum size=3.4pt},
    ramo/.style={line width=0.5pt, black!70},
    et/.style={align=center, inner sep=1.2pt, font=\footnotesize}
  ]

  % --- posizioni dei nodi (dfs-numbering in preordine) ---
  \coordinate (v1)  at ( 0.00,  0.00);
  \coordinate (v2)  at (-1.50, -1.20);
  \coordinate (v5)  at ( 2.40, -1.00);
  \coordinate (v3)  at (-2.50, -2.40);
  \coordinate (v4)  at (-1.10, -2.40);
  \coordinate (v6)  at ( 1.20, -2.40);
  \coordinate (v7)  at ( 2.45, -2.40);
  \coordinate (v8)  at ( 3.90, -2.20);
  \coordinate (v9)  at ( 3.45, -3.45);
  \coordinate (v10) at ( 5.00, -3.20);

  % --- archi ---
  \draw[ramo] (v1) -- (v2) (v1) -- (v5)
              (v2) -- (v3) (v2) -- (v4)
              (v5) -- (v6) (v5) -- (v7) (v5) -- (v8)
              (v8) -- (v9) (v8) -- (v10);

  % --- nodi ---
  \foreach \n in {1,2,3,4,5,6,7,8,9,10} { \node[nodo] at (v\n) {}; }

  % --- etichette: dfs-numero (interno) e codifica binaria (esterna) ---
  \node[et, anchor=south]      at ($(v1)+(0,0.10)$)      {$0001$\\$1$};
  \node[et, anchor=south east] at ($(v2)+(-0.08,0.02)$)  {$0010$\\$2$};
  \node[et, anchor=south west] at ($(v5)+(0.10,0.02)$)   {$0101$\\$5$};
  \node[et, anchor=north]      at ($(v3)+(0,-0.10)$)     {$3$\\$0011$};
  \node[et, anchor=north]      at ($(v4)+(0,-0.10)$)     {$4$\\$0100$};
  \node[et, anchor=north]      at ($(v6)+(0,-0.10)$)     {$6$\\$0110$};
  \node[et, anchor=north]      at ($(v7)+(0,-0.10)$)     {$7$\\$0111$};
  \node[et, anchor=south west] at ($(v8)+(0.08,0.10)$)   {$8$\\$1000$};
  \node[et, anchor=north]      at ($(v9)+(0,-0.10)$)     {$9$\\$1001$};
  \node[et, anchor=north]      at ($(v10)+(0,-0.10)$)    {$10$\\$1010$};

\end{tikzpicture}

  \caption{L'albero generico $\T$ con i dfs-numeri e la loro codifica binaria su
    $\lfloor\log_2 n\rfloor+1 = 4$ bit ($n=10$ nodi).}
  \label{fig:lca-albero-I}
\end{figure}

\begin{definizione}[la funzione $h$]\label{def:h}
  Per ogni numero $k$, $h(k)$ è la posizione del suo \emph{rightmost} $1$, cioè la
  posizione, contata da destra a partire da $1$, del bit $1$ meno significativo nella
  rappresentazione binaria di $k$.
\end{definizione}

\begin{osservazione}
  $h(k)$ è l'\emph{altezza di $k$ in $\B$}: è l'altezza del nodo di $\B$ che ha $k$
  come path number, dove l'altezza di un nodo è il numero di nodi sul cammino che lo
  congiunge a una foglia. Infatti un nodo a profondità $t$ ha path number
  $b_1\cdots b_t\,1\,0^{\,d-t}$, il cui $1$ meno significativo sta in posizione
  $d-t+1$, che è esattamente il numero di nodi sul cammino da esso a una foglia.
  Inoltre $h(k)$ si calcola in $\Oh(1)$ con la primitiva \textsc{rightmost}.
\end{osservazione}

\begin{definizione}[la mappa $I$]\label{def:mappa-I}
  Dato $v \in \T$, indico con $I(v)$ il nodo $w \in \albero{v}$ tale che $h(w)$ è
  massimo.
\end{definizione}

\begin{osservazione}
  Due avvertenze sulla definizione. Primo: $h(w)$ è la funzione aritmetica appena
  definita, applicata al \emph{dfs-numero} di $w$; non è l'altezza di $w$ dentro $\T$.
  Secondo: $I(v)$ è, letteralmente, un nodo di $\T$; lo si identifica con il nodo di
  $\B$ che ha $I(v)$ come path number, ed è in questo senso che $I$ è una mappa da
  $\T$ a $\B$.
\end{osservazione}

\begin{esempio}\label{ex:I-esempio}
  Nell'albero della Figura~\ref{fig:lca-albero-I}
  \[
    I(1) = I(5) = I(8) = 8 .
  \]
  Infatti $h(8) = h(1000) = 4$ è il massimo valore di $h$ su tutto l'albero, e il nodo
  $8$ appartiene ai sottoalberi radicati in $1$, in $5$ e in $8$.
\end{esempio}

\begin{osservazione}
  $I(v) = v$ per ogni foglia $v$: il sottoalbero di una foglia è la foglia stessa.
\end{osservazione}

\begin{osservazione}
  $h(I(v))$ non decresce muovendosi verso la radice: se $u$ è un antenato di $v$ allora
  $\albero{v} \subseteq \albero{u}$, e quindi il massimo di $h$ su $\albero{u}$ è
  almeno quello su $\albero{v}$, cioè $h(I(u)) \ge h(I(v))$. Detto altrimenti,
  l'immagine di $v$ sale in $\B$ man mano che $v$ sale in $\T$.
\end{osservazione}

La mappa $I \colon \T \to \B$ definisce implicitamente una relazione di equivalenza sui
nodi di $\T$,
\[
  u \sim v \quad\iff\quad I(u) = I(v),
\]
le cui classi sono i gruppi di nodi evidenziati nella
Figura~\ref{fig:lca-run-map}.

\begin{domanda}
  La funzione $I$ è ben definita?
\end{domanda}

\subsection{Buona definizione e run}

\scan{36}

Perché la mappa $I$ sia ben definita occorre che per ogni elemento del dominio
$v \in \T$ ci sia \emph{uno e un solo} elemento $w \in \B$ tale che $I(v) = w$. Un tale
$w$ esiste sempre --- $\albero{v}$ è finito e non vuoto, dunque il massimo di $h$ su
$\albero{v}$ viene certamente raggiunto: c'è sempre almeno un elemento del codominio.
Resta perciò da provare che non ne possono essere \emph{due}.

\begin{proposizione}[buona definizione di $I$]\label{prop:I-ben-definita}
  Per ogni nodo $v \in \T$ esiste un \emph{unico} nodo $w \in \albero{v}$ tale che
  \[
    h(w) \;=\; \max_{q \in \albero{v}} h(q).
  \]
\end{proposizione}

\begin{dimostrazione}
  Per assurdo, assumiamo che esistano $u, w \in \albero{v}$ con $u \neq w$ e
  \[
    h(u) \;=\; h(w) \;\geq\; h(q) \qquad \text{per ogni } q \in \albero{v}.
  \]
  Poniamo $i \defeq h(u) = h(w)$ e supponiamo, senza perdita di generalità, $w < u$ (il
  confronto è fra i numeri DFS). Sia $k$ la posizione del bit più significativo in cui
  $u$ e $w$ differiscono: poiché $w < u$, in posizione $k$ il numero $u$ ha $1$ e il
  numero $w$ ha $0$. Inoltre $k > i$: infatti $u$ e $w$ hanno entrambi l'$1$ meno
  significativo in posizione $i$ e tutti $0$ nelle posizioni $< i$, quindi non possono
  differire in una posizione $\leq i$.

  Costruiamo allora il numero $N$ che coincide con $u$ e con $w$ in tutte le posizioni
  $> k$, vale $1$ in posizione $k$ ed è nullo in tutte le posizioni $< k$:
  \[
    \begin{array}{c||c|c|c|c|c}
      & \text{pos.} > k & k & k-1,\dots,i+1 & i & i-1,\dots,1\\
      \hline\hline
      u & \a & 1 & \cdots      & 1 & 0\cdots 0\\
      \hline
      w & \a & 0 & \cdots      & 1 & 0\cdots 0\\
      \hline
      N & \a & 1 & 0\cdots 0   & 0 & 0\cdots 0
    \end{array}
  \]
  dove $\a$ denota il prefisso comune (i bit dalla posizione più significativa fino
  alla posizione $k+1$, uguali in $u$, $w$ e $N$) e i puntini nelle righe di $u$ e $w$
  stanno per bit arbitrari.

  È ovvio che
  \[
    w \;<\; N \;<\; u :
  \]
  $N$ e $w$ coincidono sopra la posizione $k$, dove però $N$ ha $1$ e $w$ ha $0$, da
  cui $N > w$; $N$ e $u$ coincidono da sinistra fino a $k$ incluso, ma sotto $k$ il
  numero $N$ ha solo zeri mentre $u$ ha un $1$ in posizione $i < k$, da cui $N < u$.

  Per la proprietà del dfs-numbering (i numeri dei nodi di
  $\albero{v}$ formano un intervallo contiguo, cioè sono assegnati \emph{senza buchi}),
  da $w, u \in \albero{v}$ e $w < N < u$ segue
  \[
    N \in \albero{v}.
  \]
  Ma per costruzione l'$1$ meno significativo di $N$ sta in posizione $k$, cioè
  $h(N) = k > i$, e quindi
  \[
    h(N) \;>\; h(u) \;=\; h(w),
  \]
  contro la massimalità di $h(u) = h(w)$ su $\albero{v}$: assurdo.
\end{dimostrazione}

\begin{figure}[H]
  \centering
  % A sinistra l'albero generico T con i sette run; a destra l'albero binario
% completo B con i sette nodi immagine marcati.
\begin{tikzpicture}[
    x=1cm, y=1cm,
    nodo/.style={circle, fill=black, inner sep=0pt, minimum size=3.2pt},
    nodog/.style={circle, fill=black!45, inner sep=0pt, minimum size=2.6pt},
    anello/.style={circle, draw=black, line width=0.5pt, inner sep=0pt,
                   minimum size=7pt},
    ramo/.style={line width=0.5pt, black!70},
    ramog/.style={line width=0.45pt, black!45},
    run/.style={dash pattern=on 2pt off 1.6pt, line width=0.5pt, black!60},
    et/.style={align=center, inner sep=1pt, font=\scriptsize},
    etg/.style={et, text=black!50}
  ]

  %% ================= PANNELLO SINISTRO: albero T =================
  \coordinate (t1)  at (1.90,  0.00);
  \coordinate (t2)  at (0.85, -0.90);
  \coordinate (t5)  at (2.85, -0.90);
  \coordinate (t3)  at (-0.05, -1.80);
  \coordinate (t4)  at (0.85, -1.80);
  \coordinate (t6)  at (2.10, -1.80);
  \coordinate (t7)  at (2.95, -1.80);
  \coordinate (t8)  at (3.80, -1.80);
  \coordinate (t9)  at (3.58, -2.70);
  \coordinate (t10) at (4.40, -2.70);

  \draw[ramo] (t1) -- (t2) (t1) -- (t5)
              (t2) -- (t3) (t2) -- (t4)
              (t5) -- (t6) (t5) -- (t7) (t5) -- (t8)
              (t8) -- (t9) (t8) -- (t10);

  \foreach \n in {1,2,3,4,5,6,7,8,9,10} { \node[nodo] at (t\n) {}; }

  % --- i sette run ---
  \begin{scope}[rotate around={-43.45:(2.85,-0.90)}]
    \draw[run] (2.85,-0.90) ellipse [x radius=1.57, y radius=0.27];
  \end{scope}
  \draw[run] (0.85,-1.35) ellipse [x radius=0.25, y radius=0.70];
  \foreach \n in {3,6,7,9,10} { \draw[run] (t\n) circle [radius=0.22]; }

  % --- etichette: dfs-numero e path number su 4 bit ---
  \node[et, anchor=south]      at (1.90,  0.28) {$0001$\\$1$};
  \node[et, anchor=south east] at (0.62, -0.72) {$0010$\\$2$};
  \node[et, anchor=south west] at (3.02, -0.55) {$0101$\\$5$};
  \node[et, anchor=north]      at (-0.05, -2.10) {$3$\\$0011$};
  \node[et, anchor=north]      at (0.85, -2.10) {$4$\\$0100$};
  \node[et, anchor=north]      at (2.10, -2.10) {$6$\\$0110$};
  \node[et, anchor=north]      at (2.95, -2.10) {$7$\\$0111$};
  \node[et, anchor=south west] at (4.00, -1.50) {$8=1000$};
  \node[et, anchor=north]      at (3.58, -2.96) {$9$\\$1001$};
  \node[et, anchor=north]      at (4.40, -2.96) {$10$\\$1010$};

  \node[anchor=west, font=\small] at (-0.15, 0.40) {$\T$};

  %% ================= PANNELLO DESTRO: albero binario completo B =========
  \coordinate (b1000) at (8.40,  0.00);
  \coordinate (b0100) at (7.00, -0.85);
  \coordinate (b1100) at (9.80, -0.85);
  \coordinate (b0010) at (6.30, -1.70);
  \coordinate (b0110) at (7.70, -1.70);
  \coordinate (b1010) at (9.10, -1.70);
  \coordinate (b1110) at (10.50,-1.70);
  \coordinate (b0001) at (5.95, -2.55);
  \coordinate (b0011) at (6.65, -2.55);
  \coordinate (b0101) at (7.35, -2.55);
  \coordinate (b0111) at (8.05, -2.55);
  \coordinate (b1001) at (8.75, -2.55);
  \coordinate (b1011) at (9.45, -2.55);
  \coordinate (b1101) at (10.15,-2.55);
  \coordinate (b1111) at (10.85,-2.55);

  \draw[ramog] (b1000) -- (b0100) (b1000) -- (b1100)
               (b0100) -- (b0010) (b0100) -- (b0110)
               (b1100) -- (b1010) (b1100) -- (b1110)
               (b0010) -- (b0001) (b0010) -- (b0011)
               (b0110) -- (b0101) (b0110) -- (b0111)
               (b1010) -- (b1001) (b1010) -- (b1011)
               (b1110) -- (b1101) (b1110) -- (b1111);

  % nodi non marcati
  \foreach \n in {0010,0101,1100,1110,0001,1011,1101,1111}
    { \node[nodog] at (b\n) {}; }
  % nodi marcati: immagine di un run di T
  \foreach \n in {1000,0100,0110,1010,0011,0111,1001}
    { \node[anello] at (b\n) {}; \node[nodo] at (b\n) {}; }

  \node[et, anchor=south]      at (8.40,  0.14) {$1000$};
  \node[et, anchor=south east] at (6.88, -0.74) {$0100$};
  \node[etg,anchor=south west] at (9.92, -0.74) {$1100$};
  \node[et, anchor=north east] at (7.52, -1.82) {$0110$};
  \node[et, anchor=south east] at (8.96, -1.60) {$1010$};
  \node[et, anchor=north]      at (6.65, -2.70) {$0011$};
  \node[et, anchor=north]      at (8.05, -2.70) {$0111$};
  \node[et, anchor=north]      at (8.75, -2.70) {$1001$};

  \node[anchor=west, font=\small] at (5.85, 0.40) {$\B$};

  %% ================= la mappa I =================
  \draw[-{Latex[length=4pt,width=3.2pt]}, dashed, line width=0.6pt, black!60]
        (5.08,-2.00) -- (5.86,-2.00)
        node[midway, above=1pt, font=\scriptsize, text=black] {$I$};

\end{tikzpicture}

  \caption{A sinistra l'albero generico $\T$ con la numerazione DFS (in binario su $4$
    bit) e i sette \emph{run} evidenziati; a destra l'albero binario completo $\B$, in
    cui sono marcati i nodi che sono immagine di qualche nodo di $\T$. La mappa $I$
    collassa ogni run in un unico nodo di $\B$.}
  \label{fig:lca-run-map}
\end{figure}

La mappa $I$ definisce dei \emph{run} di nodi in $\T$.

\begin{definizione}[run]\label{def:run}
  Un \emph{run} è una collezione massimale di nodi di $\T$ mappati nello stesso $I(v)$;
  equivalentemente, è una classe di equivalenza della relazione
  $u \sim v \iff I(u) = I(v)$.
\end{definizione}

\begin{osservazione}
  Ogni run è un \emph{cammino verticale} di $\T$. Anzitutto due nodi dello stesso run
  sono confrontabili: se $I(u)=I(v)$ allora il nodo $I(u)$ appartiene sia ad
  $\albero{u}$ sia ad $\albero{v}$, e due sottoalberi che hanno un nodo in comune sono
  confrontabili per inclusione, cioè uno fra $u$ e $v$ è antenato dell'altro. Il run è
  inoltre \emph{contiguo}: se $u$ è antenato di $v$ e $I(u)=I(v)=w$, ogni nodo $m$ del
  cammino da $u$ a $v$ soddisfa $\albero{v}\subseteq\albero{m}\subseteq\albero{u}$,
  quindi il massimo di $h$ su $\albero{m}$ è compreso fra quello su $\albero{v}$ e
  quello su $\albero{u}$, che valgono entrambi $h(w)$; poiché $w\in\albero{v}\subseteq
  \albero{m}$, quel massimo è realizzato proprio da $w$ e dunque $I(m)=w$. Una catena
  contigua di nodi è un cammino verticale, e lungo un cammino verticale «più vicino
  alla radice» equivale a «numero DFS più piccolo».
\end{osservazione}

\section{Passo 3: il pre-processing}\label{sec:lca-preprocessing}

\scan{37}

Eseguiamo un pre-processing lineare in cui, per ogni nodo $v$ di $\T$, calcoliamo
\[
  \begin{cases}
    I(v),\\[2pt]
    L(v) = \text{il \emph{leader} del suo run},\\[2pt]
    A_v = \text{un array ausiliario}.
  \end{cases}
\]

\begin{osservazione}
  Il nodo di $\albero{v}$ su cui $h$ è massimo --- cioè $I(v)$ --- è il nodo più
  \emph{basso} del run: è il \emph{bottom} del run.
\end{osservazione}

\begin{definizione}[leader di un run]\label{def:leader}
  Definisco \emph{leader} (o \emph{head}) di un run il nodo più \emph{alto} del run,
  cioè quello più vicino alla radice di $\T$.
\end{definizione}

\paragraph{Pre-processing.}
\begin{enumerate}
  \item \textsc{dfs} di $\T$ $\longrightarrow$ computo il \emph{dfs-numbering} e $h(v)$
        per ogni $v$.

  \item Usando $h(\cdot)$ determino $I(\cdot)$ per ogni $v$ mediante una visita
        \emph{bottom-up}. Nella stessa visita determino il leader di ogni run (la
        tabella dei leader è indicizzata dai valori di $I$: $L(v)=L(I(v))$). Ora, per
        ogni $v$, leggo $I(v)$ e $L(v)$ in tempo $\Oh(1)$.

  \item Costruisco $\B$, l'albero binario completo di profondità $d$, dove $d+1$ è il
        numero di bit necessari a scrivere i numeri DFS di $\T$, cioè
        $d+1=\lfloor\log_2 n\rfloor+1$. I path number dei nodi di $\B$ sono allora
        tutte e sole le parole non nulle di $d+1$ bit, ossia gli interi da $1$ a
        $2^{d+1}-1\ge n$: le due numerazioni si possono identificare. (Se $n$ non è una
        potenza di $2$ si ha $d+1=\lceil \log_2 n\rceil$; se lo è, serve un bit in più.)

  \item Per ogni $v$ costruisco $A_v$, array di $d+1$ bit definito da
        \[
          A_v[i] = 1 \iff \exists\, w \text{ antenato di } v \text{ in } \T
          \text{ tale che } h\bigl(I(w)\bigr) = i ,
        \]
        con la convenzione che $v$ è antenato di se stesso.
\end{enumerate}

\begin{figure}[H]
  \centering
  % Il vettore di bit A_v: un array di d+1 = floor(log_2 n)+1 bit; i bit sono numerati da
% destra a partire da 1.
\begin{tikzpicture}[
    x=1cm, y=1cm,
    bordo/.style={line width=0.7pt},
    sep/.style={line width=0.5pt},
    onda/.style={decorate, decoration={snake, amplitude=2pt, segment length=22pt},
                 line width=0.5pt}
  ]

  % --- la parola ---
  \def\hh{0.72}
  \draw[bordo] (0,0) rectangle (7,\hh);

  % separatori fra le celle
  \foreach \x in {3.5, 4.2, 4.9, 5.6, 6.3} { \draw[sep] (\x,0) -- (\x,\hh); }

  % contenuto delle celle: ~ = bit qualsiasi
  \draw[onda] (0.30,0.36) -- (3.20,0.36);
  \draw[onda, decoration={snake, amplitude=1.6pt, segment length=13pt}]
        (4.31,0.36) -- (4.79,0.36);
  \node[font=\small] at (3.85,0.36) {$1$};
  \node[font=\small] at (5.25,0.36) {$1$};
  \node[font=\small] at (5.95,0.36) {$0$};
  \node[font=\small] at (6.65,0.36) {$0$};

  % --- etichetta del vettore ---
  \node[anchor=east, font=\small] at (-0.45,0.36) {$A_v$};

  % --- la posizione puntata ---
  \node[anchor=south, font=\footnotesize] at (3.85,\hh+0.06) {$i$};

  \draw[-{Latex[length=4pt,width=3.4pt]}, line width=0.6pt]
        (3.85,-0.06) -- (3.85,-0.80);

  \node[anchor=north, align=center, font=\small] at (3.85,-0.92)
       {significa che esiste $w$ antenato di $v$\\
        che mappa ad altezza $i$ in $\B$};

\end{tikzpicture}

  \caption{Il vettore di bit $A_v$: un $1$ in posizione $i$ significa che esiste un
    antenato $w$ di $v$ che viene mappato ad altezza $i$ in $\B$.}
  \label{fig:av-bitvector}
\end{figure}

\begin{osservazione}
  Tutti e quattro i passi costano $\Oh(n)$: in particolare $A_v$ si calcola con una
  visita top-down ponendo
  \[
    A_v = A_{\mathrm{padre}(v)} \ \textsc{or}\ 2^{\,h(I(v))-1},
  \]
  cioè con un'operazione $\Oh(1)$ per nodo. Il pre-processing complessivo è dunque
  lineare.
\end{osservazione}

\begin{osservazione}
  Il bit $1$ più a destra di $A_v$ sta in posizione $h(I(v))$: per
  la monotonia di $h\circ I$ ogni antenato $u$ di $v$ ha $h(I(u)) \ge h(I(v))$, e la
  posizione $h(I(v))$ è occupata da $v$ stesso.
\end{osservazione}

\section{Rispondere alle query in tempo costante}\label{sec:lca-query}

\scan{38}

Restano due cose da fare: mostrare che \emph{$I$ mantiene la relazione
antenato/discendente}, e scrivere un algoritmo che, usando $I$, $L$ e gli $A_v$,
risponda a $\lca(x,y)$ in $\Oh(1)$ nel modello di calcolo fissato.

\begin{lemma}[la mappa $I$ rispetta la relazione antenato/discendente]\label{lem:lca-I-ant}
  Siano $x$ e $z$ nodi di $\T$. Allora
  \[
    x\in\albero{z}
    \quad\Longrightarrow\quad
    I(x)\in\B\bigl(I(z)\bigr),
  \]
  dove $\albero{z}$ è il sottoalbero di $\T$ radicato in $z$ e $\B(w)$ è il
  sottoalbero di $\B$ radicato in $w$.
\end{lemma}

\begin{dimostrazione*}
  Distinguiamo due casi.
  \begin{itemize}
    \item Se $I(x)=I(z)$, allora $I(x)=I(z)\in\B\bigl(I(z)\bigr)$ e la tesi è
          immediata: $x$ e $z$ stanno sullo stesso run.

    \item Altrimenti $h\bigl(I(z)\bigr)>h\bigl(I(x)\bigr)$: da
          $\albero{x}\subseteq\albero{z}$ segue
          $h\bigl(I(x)\bigr)\le h\bigl(I(z)\bigr)$ per massimalità, e l'uguaglianza è
          esclusa dall'unicità del nodo che realizza il massimo di $h$ su $\albero{z}$
          (Proposizione~\ref{prop:I-ben-definita}).

          Sia $h\bigl(I(z)\bigr)=i$: i bit a sinistra della posizione $i$ in $I(x)$
          devono essere uguali a quelli a sinistra di $i$ in $I(z)$.

          Da qui segue che $I(x)$ è discendente di $I(z)$ in $\B$. Il sottoalbero di
          $\B$ radicato in un nodo $w$ è infatti esattamente l'insieme dei path number
          che coincidono con $w$ in tutti i bit a sinistra della posizione $h(w)$
          \emph{e} hanno l'$1$ meno significativo in posizione $\le h(w)$: senza
          quest'ultima clausola vi rientrerebbero anche gli antenati di $w$ (per
          esempio, con $d=2$, la radice $100$ coincide con $w=110$ nell'unico bit a
          sinistra di $h(w)=2$). Per $I(x)$ valgono entrambe le condizioni: la prima
          per quanto appena osservato, la seconda perché in questo caso
          $h\bigl(I(x)\bigr)<i$.
  \end{itemize}
\end{dimostrazione*}

\begin{esercizio}[il passaggio lasciato in sospeso]\label{ex:lca-bit-sx}
  Nelle ipotesi del Lemma~\ref{lem:lca-I-ant} e posto $i=h\bigl(I(z)\bigr)$, dimostrare
  che i bit di $I(x)$ a sinistra della posizione $i$ coincidono con quelli di $I(z)$.
\end{esercizio}

\emph{Traccia.} $I(x)$ e $I(z)$ sono numeri DFS di nodi di $\albero{z}$, che per la
proprietà DFS formano un intervallo contiguo. Se i due numeri differissero in qualche
posizione $k>i$, si prenda la più significativa di tali posizioni e si ponga
\[
  N=\bigl(\text{bit comuni a sinistra di }k\bigr)\,\Vert\,1\,\Vert\,0\cdots 0 ,
\]
dove $\Vert$ denota la giustapposizione di stringhe di bit. Allora $\min\{I(x),I(z)\}<N<\max\{I(x),I(z)\}$ e $h(N)=k>i$: dunque $N$ sarebbe il
numero di un nodo di $\albero{z}$ con $h$ maggiore di $h\bigl(I(z)\bigr)$, contro la
definizione di $I(z)$.

\bigskip

I due teoremi che seguono sono il cuore dell'algoritmo di query: il primo dice che
basta conoscere \emph{l'altezza in $\B$} dell'immagine del $\lca$ per ricostruire il
$\lca$ stesso, il secondo dice come si calcola quell'altezza.

\begin{teorema}[determinare $z$ noto $h(I(z))$]\label{thm:lca-teo1}
  Sia $z=\lca(x,y)$. Se conosco $h\bigl(I(z)\bigr)$, allora determino $z$ in $\Oh(1)$.
\end{teorema}

\begin{teorema}[determinare $h(I(z))$ da $A_x$ e $A_y$]\label{thm:lca-teo2}
  Sia $z=\lca(x,y)$ e sia
  \[
    j=\min\Bigl\{\,k\ :\ k\ \ge\ h\bigl(\lca_{\B}(I(x),I(y))\bigr)
    \ \wedge\ A_x[k]=A_y[k]=1\,\Bigr\}.
  \]
  Allora $j=h\bigl(I(z)\bigr)$.
\end{teorema}

\emph{Traccia della dimostrazione del Teorema~\ref{thm:lca-teo2}.} Poniamo
$b=\lca_{\B}(I(x),I(y))$. Da un lato $z$ è antenato comune di $x$ e $y$, quindi
$A_x\bigl[h(I(z))\bigr]=A_y\bigl[h(I(z))\bigr]=1$; inoltre, per il
Lemma~\ref{lem:lca-I-ant}, $I(z)$ è antenato in $\B$ sia di $I(x)$ sia di $I(y)$, dunque
è antenato di $b$ e $h(I(z))\ge h(b)$: il valore $h(I(z))$ è quindi uno dei candidati
su cui si prende il minimo, da cui $j\le h(I(z))$. Dall'altro lato, sia $k\ge h(b)$ un
candidato: esistono un antenato $u$ di $x$ e un antenato $w$ di $y$ con
$h(I(u))=h(I(w))=k$. Poiché $I(u)$ è antenato di $I(x)$, $I(w)$ è antenato di $I(y)$ ed
entrambi hanno altezza $k\ge h(b)$, i due sono lo stesso nodo di $\B$ (in un albero
l'antenato di altezza fissata è unico, e sopra $b$ i cammini verso $I(x)$ e verso
$I(y)$ coincidono): $I(u)=I(w)$, cioè $u$ e $w$ stanno sullo stesso run. Essendo il run
un cammino verticale, il più alto fra $u$ e $w$ è antenato
comune di $x$ e $y$, quindi antenato di $z$, e per
la monotonia di $h\circ I$ lungo i cammini radice--foglia $k\ge h(I(z))$. Dunque $j\ge h(I(z))$, e le due
disuguaglianze danno la tesi.

\begin{dimostrazione*}[del Teorema~\ref{thm:lca-teo1}]
  Sia $j=h\bigl(I(z)\bigr)$ noto (lo fornisce il Teorema~\ref{thm:lca-teo2}). Siano
  $\bar x$ e $\bar y$ i primi nodi del run di $z$ che si incontrano lungo il cammino
  che risale da $x$, rispettivamente da $y$, verso la radice. Allora $z$ è uno tra
  $\bar x$ e $\bar y$: precisamente quello con dfs-numbering più piccolo, cioè quello
  più in alto (i due nodi stanno sullo stesso run, che è un cammino verticale, quindi
  sono confrontabili).

  \begin{figure}[H]
    \centering
    % Frammento di T: il run di z (catena a zig-zag in diagonale) con xbar e ybar.
% Convenzione grafica del manoscritto: i cammini di lunghezza arbitraria sono
% resi con linee a zig-zag.
\begin{tikzpicture}[
    scale=1.25,
    nodo/.style={circle, fill=black, inner sep=0pt, minimum size=3.6pt},
    zig/.style={decorate, decoration={zigzag, segment length=7pt, amplitude=1.9pt},
                line width=0.55pt, line cap=round},
    runz/.style={zig, line width=1pt},
    et/.style={font=\small, inner sep=1.5pt}
  ]

  \coordinate (TOP) at ( 0.00, 4.40);   % estremo libero: il run prosegue verso il leader
  \coordinate (XB)  at ( 0.70, 3.50);
  \coordinate (YB)  at ( 2.05, 2.00);
  \coordinate (BOT) at ( 3.22, 0.60);   % fondo del run
  \coordinate (X)   at (-0.90, 1.75);
  \coordinate (Y)   at ( 0.75, 0.00);

  % --- il run di z ---
  \draw[runz] (TOP) -- (XB) -- (YB) -- (BOT);

  % --- cammino che risale da x fino a xbar (tratto liscio, leggermente arcuato) ---
  \draw[line width=0.55pt] (XB) to[out=214, in=50] (X);

  % --- cammino che risale da y fino a ybar ---
  \draw[zig] (YB) -- (Y);

  % --- nodi ---
  \node[nodo] at (XB)  {};
  \node[nodo] at (YB)  {};
  \node[nodo] at (BOT) {};
  \node[nodo] at (X)   {};
  \node[nodo] at (Y)   {};

  % --- etichette ---
  \node[et, anchor=south west] at ($(XB)+(0.10,0.08)$)  {$\bar x\,\textcolor{black!55}{=z}$};
  \node[et, anchor=south west] at ($(YB)+(0.17,0.16)$)  {$\bar y$};
  \node[et, anchor=north east] at ($(X)+(-0.06,-0.06)$) {$x$};
  \node[et, anchor=north east] at ($(Y)+(-0.06,-0.06)$) {$y$};
  \node[et, anchor=north west, text=black!55]
       at ($(BOT)+(0.10,0.02)$) {$I(z)$};

\end{tikzpicture}

    \caption{Il run di $z$ (linea a zig-zag che scende in diagonale): $\bar x$ e
      $\bar y$ sono i primi nodi del run incontrati risalendo da $x$ e da $y$. Poiché
      $\bar x$ è il più alto dei due, $z=\bar x$.}
    \label{fig:lca-run-xbar-ybar}
  \end{figure}

  \scan{39}

  Resta da determinare $\bar x$ in tempo costante; per $\bar y$ si procede in modo
  simmetrico.
  \begin{itemize}
    \item Se $h(I(x))=j$, allora $x$ sta già sul run di $z$, cioè $x=\bar x$: ho
          finito.
    \item Altrimenti, sia $w$ l'ultimo nodo del cammino $x\rightsquigarrow\bar x$ che
          viene mappato in un nodo di altezza $<j$ (tutti i nodi del cammino che stanno
          strettamente sotto $\bar x$ hanno infatti $h(I(\cdot))<j$).
  \end{itemize}

  Il nodo $w$ è figlio di $\bar x$ ed è il \emph{leader} del proprio run; dunque, se
  determino $w$ ho finito, perché $\bar x=\mathrm{padre}(w)$.

  \begin{figure}[H]
    \centering
    % lca-xbar-w -- il caso h(I(x)) != j.
% Convenzione grafica (la stessa di lca-run-xbar-ybar):
%   * un RUN e' una spezzata a scalini, spessa e grigia;
%   * un cammino generico di T e' una linea liscia e sottile.
\begin{tikzpicture}[
    >=Latex,
    nodo/.style     = {circle, fill=black, inner sep=0pt, minimum size=3.4pt},
    run/.style      = {line width=1.2pt, black!62, line cap=round, line join=round},
    runoltre/.style = {line width=1.2pt, black!45, line join=round,
                       dash pattern=on 0pt off 3pt, line cap=round},
    cammino/.style  = {line width=0.5pt, black!85},
    arco/.style     = {line width=1pt, black},
    ann/.style      = {font=\footnotesize\itshape, text=black!55, inner sep=1.5pt},
    guida/.style    = {line width=0.4pt, black!45},
    lbl/.style      = {font=\small, inner sep=1.5pt},
  ]

  \coordinate (x)    at (0,0);
  \coordinate (w)    at (3.0,2.3);
  \coordinate (xbar) at (3.65,2.85);

  % ---- run di z: scalinata che scende da sinistra in alto e termina in xbar.
  %      Prosegue oltre il bordo verso l'alto (li' stanno z e il leader del run).
  \draw[runoltre] (2.15,4.45) -- ++(0,0.30) -- ++(-0.30,0) -- ++(0,0.30)
                              -- ++(-0.30,0) -- ++(0,0.30);
  \draw[run] (2.15,4.45)
    \foreach \i in {1,...,5} { -- ++(0.30,0) -- ++(0,-0.32) };

  % ---- run di w: parte esattamente da w (che ne e' il leader) e scende a destra.
  \draw[run] (w)
    \foreach \i in {1,...,6} { -- ++(0,-0.28) -- ++(0.46,0) }
    -- ++(0,-0.28) coordinate (fondow);

  % ---- l'arco di T fra xbar e w: xbar = padre(w)
  \draw[arco] (xbar) -- (w);

  % ---- cammino (liscio) che da w scende fino a x
  \draw[cammino] (w) to[out=228,in=44] (x);

  % ---- nodi
  \node[nodo] at (x)      {};
  \node[nodo] at (w)      {};
  \node[nodo] at (xbar)   {};
  \node[nodo] at (fondow) {};

  % ---- etichette
  \node[lbl] at ($(xbar)+(0.26,0.14)$) {$\bar x$};
  \node[lbl] at ($(w)+(-0.24,0.16)$)   {$w$};
  \node[lbl] at ($(x)+(-0.26,0.02)$)   {$x$};

  % ---- annotazioni
  \node[ann, anchor=east] at (1.42,3.85) {run di $z$};
  \draw[guida] (1.50,3.87) -- (2.45,4.22);

  \node[ann, anchor=west] at (4.62,2.12) {leader del run};
  \draw[guida] (4.56,2.12) .. controls (3.90,2.30) .. (3.14,2.26);

  \node[ann, anchor=west] at (4.70,1.60) {run di $w$};
\end{tikzpicture}

    \caption{Il caso $h(I(x))\neq j$: $\bar x$ appartiene al run di $z$, mentre $w$ è
      il figlio di $\bar x$ sul cammino che scende verso $x$ ed è il leader del proprio
      run.}
    \label{fig:lca-xbar-w}
  \end{figure}

  Il nodo $w$ si determina con quattro osservazioni, ciascuna eseguibile in $\Oh(1)$:
  \begin{itemize}
    \item $h(I(w))$ è la posizione del primo $1$ di $A_x$ a destra di $j$, cioè la più
          grande posizione $<j$ in cui $A_x$ vale $1$;
    \item $I(w)$ coincide con $I(x)$ in tutti i bit a sinistra della posizione
          $h(I(w))$; quindi
          \[
            I(w)\;=\;\bigl(\text{bit di } I(x) \text{ a sinistra di } h(I(w))\bigr)
            \;\Vert\;1\;\Vert\;0\cdots0 ;
          \]
    \item $w$ è il leader del run individuato da $I(w)$, cioè $w=L(I(w))$;
    \item $\bar x$ è il padre di $w$.
  \end{itemize}
\end{dimostrazione*}

\begin{osservazione}
  Il secondo punto si appoggia al Lemma~\ref{lem:lca-I-ant}: se
  $u\in\albero{v}$ allora $I(u)\in\B(I(v))$, cioè $I(u)$ e $I(v)$ coincidono in tutti i
  bit a sinistra della posizione $h(I(v))$. Dentro $\albero{v}$, insomma, i bit a
  sinistra della posizione $h(I(v))$ restano fissi e possono cambiare solo quelli dalla
  posizione $h(I(v))$ in giù. È esattamente questo che permette di ricostruire $I(w)$ a
  partire da $I(x)$.
\end{osservazione}

\subsection{L'algoritmo}

\begin{osservazione}
  Tutte le informazioni che l'algoritmo usa devono essere state rese disponibili in
  fase di pre-processing.
\end{osservazione}

\begin{algorithm}[H]
  \caption{Query $\lca(x,y)$ in tempo costante}
  \label{alg:lca-query}
  \begin{algorithmic}[1]
    \Require due nodi $x,y$ di $\T$ (identificati con i loro numeri DFS) e le tabelle
      $I$, $L$, $A$ prodotte dal pre-processing
    \Ensure $z=\lca(x,y)$
    \State $b \gets \lca_{\B}\bigl(I(x),I(y)\bigr)$
      \Comment{Algoritmo~\ref{alg:lca-b}}
    \State $j \gets \min\bigl\{\,k \ge h(b)\ :\ A_x[k]=A_y[k]=1\,\bigr\}$
      \Comment{per il Teorema~\ref{thm:lca-teo2}, $j=h(I(z))$}
    \ForAll{$u \in \{x,y\}$} \Comment{passi 3 e 4: calcolo $\bar x$ e $\bar y$}
      \State $\ell \gets \min\{\,k : A_u[k]=1\,\}$
        \Comment{$\ell=h(I(u))$, primitiva \emph{rightmost} $1$}
      \If{$\ell = j$}
        \State $\bar u \gets u$
      \Else
        \State $k \gets \max\{\,t < j : A_u[t]=1\,\}$ \Comment{$k=h(I(w))$}
        \State $I_w \gets \bigl((I(u) \gg k)\ll k\bigr)\ \textsc{or}\ 2^{\,k-1}$
          \Comment{$I_w=I(w)$}
        \State $\bar u \gets \mathrm{padre}\bigl(L(I_w)\bigr)$
      \EndIf
    \EndFor
    \State \Return $z \gets \min(\bar x,\bar y)$
      \Comment{minimo rispetto al dfs-numbering}
  \end{algorithmic}
\end{algorithm}

\scan{40}

\begin{osservazione}
  Il confronto fra i due numeri DFS all'ultima riga basta perché $\bar x$ e $\bar y$
  stanno entrambi sul run di $z=\lca(x,y)$, e un run è un cammino verticale di $\T$
  (si veda l'osservazione dopo la Definizione~\ref{def:run}): lungo un cammino
  verticale «più vicino alla radice»
  equivale a «numero DFS più piccolo». Poiché $z$ è il più alto fra $\bar x$ e
  $\bar y$, il minimo dei due numeri DFS è esattamente $z$.
\end{osservazione}

\section{Realizzazione delle primitive}\label{sec:lca-primitive}

Torniamo sul \emph{modello di calcolo} della Sezione~\ref{sec:lca-modello}: posso
implementare, usando le sole operazioni standard, tutte le operazioni sulle
\emph{parole} che mi servono. Poiché una parola di $c\log n$ bit è la giustapposizione
di $\Oh(1)$ blocchi da $\log n$ bit, e tutte le operazioni in gioco si calcolano blocco
per blocco (ricombinando poi i risultati parziali con $\Oh(1)$ shift e \textsc{or}),
basta trattare il caso di un array di $\log n$ bit.

Gli shift si ottengono dall'aritmetica:
\[
  \text{shift di } i \text{ posizioni}
  \;\rightsquigarrow\;
  \text{moltiplico / divido per } 2^{i}.
\]
Essendo $i\le\log n$, il fattore $2^{i}$ è $\Oh(n)$, sta cioè in una parola: prodotto e
quoziente sono operazioni elementari.

Le operazioni che invece non si riducono all'aritmetica --- tipicamente quelle bit a
bit --- si realizzano per \textbf{tabulazione}: costruisco tabelle di output per tutti i
possibili
\[
  2^{\log n} \;=\; n
\]
input, una per ognuna delle operazioni. Si ottiene così un array per ogni operazione
\underline{unaria}: spazio $\Oh(n)$ per tabella, accesso in $\Oh(1)$.

\begin{esempio}
  Sia $j = 1011010_{2}$. Allora $j-1 = 1011001_{2}$ e
  \[
    j \mathbin{\textsc{and}} (j-1) \;=\; 1011000_{2},
  \]
  cioè l'\textsc{and} con il predecessore azzera l'$1$ più a destra di $j$, quello in
  posizione $h(j)=2$. È il tipico esempio di operazione bit a bit che non si ottiene
  dalle sole operazioni aritmetiche e che conviene quindi precalcolare in tabella.
\end{esempio}

Per le operazioni \emph{binarie} --- \textsc{and}, \textsc{xor}, \dots\ --- la
tabulazione ingenua non funziona: indicizzare la tabella con la coppia di operandi
richiederebbe
\[
  2^{\log n} \times 2^{\log n} \;=\; n \times n \;=\; n^{2}
\]
entrate. Il trucco è spezzare ogni parola in due metà da $\tfrac{\log n}{2}$ bit e
tabulare l'operazione sulle mezze parole: la tabella, indicizzata da una coppia di
mezze parole, ha
\[
  2^{\frac{\log n}{2}} \times 2^{\frac{\log n}{2}}
  \;=\; \sqrt{n}\times\sqrt{n} \;=\; n
\]
entrate, cioè spazio $\Oh(n)$. I due risultati parziali (metà alta e metà bassa) si
ricompongono con uno shift e un \textsc{or}, dunque ancora in $\Oh(1)$.

\begin{figure}[H]
  \centering
  % tabella-op-binarie -- tabulazione di un'operazione binaria sulle mezze parole.
% In alto: la tabella (quadrato vuoto) indicizzata da una coppia di mezze parole.
% In basso: i due operandi di log n bit, ciascuno spezzato a meta'.
\begin{tikzpicture}[
    >=Latex,
    bordo/.style   = {line width=0.9pt, black},
    parola/.style  = {line width=0.7pt, black},
    graffa/.style  = {decorate, decoration={brace, amplitude=5pt, raise=3pt},
                      line width=0.5pt, black!75},
    graffina/.style= {decorate, decoration={brace, amplitude=3.5pt, raise=2.5pt},
                      line width=0.4pt, black!65},
    et/.style      = {font=\small, inner sep=1.5pt},
    ann/.style     = {font=\footnotesize, text=black!55, inner sep=1.5pt},
  ]

  %====================== la tabella ======================
  \draw[bordo] (0,0) rectangle (4.2,4.2);

  % graffa superiore: numero di colonne
  \draw[graffa] (0,4.2) -- (4.2,4.2);
  \node[et, anchor=south] at (2.1,4.58) {$2^{\frac{\log n}{2}}$};

  % graffa sinistra: numero di righe
  \draw[graffa] (0,0) -- (0,4.2);
  \node[et, anchor=east] at (-0.40,2.1) {$2^{\frac{\log n}{2}}$};

  % contenuto (in grigio chiaro): l'entrata indicizzata da una coppia di mezze parole
  \node[ann] at (2.1,2.1) {$T[a_i,b_i]=a_i \mathbin{\mathrm{op}} b_i$};

  % annotazione a destra, con freccia verso il bordo destro della tabella
  \node[et, anchor=west] at (4.95,4.20) {spazio $\sqrt n\times\sqrt n=n$};
  \draw[-{Latex[length=3mm,width=2.2mm]}, line width=0.6pt]
        (5.15,3.92) -- (4.32,3.42);

  %====================== i due operandi ======================
  % operando a
  \draw[parola] (0.6,-1.70) rectangle (4.6,-1.25);
  \draw[parola] (2.6,-1.70) -- (2.6,-1.25);
  \node[ann] at (1.6,-1.475) {$a_1$};
  \node[ann] at (3.6,-1.475) {$a_2$};
  \node[et, anchor=east] at (0.44,-1.475) {$a$};

  % operando b
  \draw[parola] (0.6,-2.60) rectangle (4.6,-2.15);
  \draw[parola] (2.6,-2.60) -- (2.6,-2.15);
  \node[ann] at (1.6,-2.375) {$b_1$};
  \node[ann] at (3.6,-2.375) {$b_2$};
  \node[et, anchor=east] at (0.44,-2.375) {$b$};

  % graffa: la parola intera e' lunga log n bit
  \draw[graffa] (0.6,-1.25) -- (4.6,-1.25);
  \node[et, anchor=south] at (2.6,-0.97) {$\log n$ bit};

  % graffe: le due meta'
  \draw[graffina] (2.6,-2.60) -- (0.6,-2.60);
  \draw[graffina] (4.6,-2.60) -- (2.6,-2.60);
  \node[et, anchor=north] at (1.6,-2.77) {$\frac{\log n}{2}$};
  \node[et, anchor=north] at (3.6,-2.77) {$\frac{\log n}{2}$};

\end{tikzpicture}

  \caption{Tabella di un'operazione binaria indicizzata da una coppia di \emph{mezze}
    parole: $2^{(\log n)/2}\times 2^{(\log n)/2}=\sqrt n\times\sqrt n=n$ entrate, cioè
    spazio $\Oh(n)$. In basso, i due operandi di $\log n$ bit, ciascuno spezzato nelle
    due metà che indicizzano la tabella.}
  \label{fig:tabella-op-binarie}
\end{figure}

\begin{osservazione}
  Tutte le primitive usate nell'algoritmo di query --- $\lca_{\B}$, \textsc{and},
  \textsc{xor}, posizione dell'$1$ più a sinistra e dell'$1$ più a destra, maschere ---
  rientrano in questo schema: sono quindi $\Oh(1)$, e le tabelle costano
  complessivamente $\Oh(n)$ di pre-processing. L'algoritmo per il $\lca$ ha dunque
  pre-processing $\Oh(n)$ e query in tempo $\Oh(1)$.
\end{osservazione}


\part{String matching inesatto}
\chapter{Pattern matching inesatto: distanze e allineamento}

\scan{41}%
Il pattern matching \emph{esatto}, di cui ci siamo occupati finora, ha due volti
opposti: da un lato è \emph{matematicamente elegante}, dall'altro è \emph{poco
realistico per le applicazioni}. Il dominio applicativo di riferimento è la
bioinformatica (cfr.\ i corsi di \emph{Bioinformatica I} e \emph{Bioinformatica
II}), dove i dati sono intrinsecamente rumorosi: \emph{voglio poter gestire
errori}.

\begin{domanda}
  Cosa significa fare pattern matching \emph{con errori}?
\end{domanda}

Significa, prima di tutto, introdurre una (o più) nozioni di \emph{distanza}
$d(\a,\b)$ fra stringhe. Le due nozioni principali sono la \emph{distanza di
Levenshtein} e la \emph{distanza di Hamming}.


\section{Distanze fra stringhe}

\subsection{Distanza di Hamming}

Le due stringhe vengono sovrapposte e confrontate \emph{posizione per
posizione}: la distanza è il numero di posizioni in cui i caratteri
corrispondenti differiscono.

\begin{figure}[H]
  \centering
  % hamming-posizioni -- pag. 41
% Due nastri paralleli di uguale lunghezza (alpha e beta) attraversati dalle
% stesse verticali: ogni verticale e' una posizione del confronto, cioe' la
% coppia (alpha[i], beta[i]). Glossa a destra: sull'alfabeto binario la
% distanza e' il numero di 1 nella XOR delle due stringhe.
\begin{tikzpicture}[
    font=\footnotesize,
    barra/.style={line width=0.9pt},
    colonna/.style={line width=0.45pt,black!70},
  ]

  \def\L{5.4}   % lunghezza dei due nastri
  \def\N{10}    % posizioni disegnate

  % ---- i due nastri --------------------------------------------------------
  \draw[barra] (0,0)  -- (\L,0);
  \draw[barra] (0,-1) -- (\L,-1);

  % barrette di chiusura agli estremi (i nastri sono "chiusi")
  \foreach \x in {0,\L}{%
    \draw[barra] (\x,0)  -- (\x,0.26);
    \draw[barra] (\x,-1) -- (\x,-0.74);
  }

  % ---- le colonne del confronto, una per posizione --------------------------
  \foreach \k in {1,...,\N}{%
    \draw[colonna] ({\L*\k/(\N+1)},0.34) -- ({\L*\k/(\N+1)},-1.34);
  }

  % ---- etichette delle due stringhe ----------------------------------------
  \node[left,inner sep=5pt] at (0,0)  {$\a$};
  \node[left,inner sep=5pt] at (0,-1) {$\b$};

  % ---- le stringhe proseguono ----------------------------------------------
  \node[right,inner sep=4pt] at (\L,-0.5) {$\dots$};

  % ---- glossa: lettura binaria della distanza ------------------------------
  \node[right,inner sep=0pt,black!70,align=left]
    at ({\L+0.7},-0.5)
    {su alfabeto binario:\\[2pt] numero di $1$ nella $\mathrm{XOR}(\a,\b)$};

\end{tikzpicture}

  \caption{Confronto posizione per posizione fra $\a$ e $\b$: le due stringhe
  hanno la stessa lunghezza e ogni segmento verticale accoppia la $i$-esima
  posizione di $\a$ con la $i$-esima posizione di $\b$.}
  \label{fig:hamming-posizioni}
\end{figure}

\begin{definizione}[Distanza di Hamming]
  \label{def:hamming}
  Siano $\a$ e $\b$ due stringhe della \emph{stessa} lunghezza $n$. Allora
  \[
    \dham(\a,\b)\;=\;\#\set{\, i \in \set{1,\dots,n} \;:\; \a[i] \neq \b[i] \,}.
  \]
\end{definizione}

Su alfabeto binario la definizione si riscrive in modo particolarmente
compatto: la distanza di Hamming è il \emph{numero di $1$ nella}
$\mathrm{XOR}(\a,\b)$, perché lo XOR bit a bit vale $1$ esattamente nelle
posizioni in cui le due stringhe non concordano.

\begin{osservazione}
  È il punto in cui si salda il discorso sul modello di calcolo del capitolo
  precedente: lo XOR e il conteggio degli $1$ di una parola di $\log n$ bit
  costano $\Oh(1)$ grazie alle tabelle precalcolate, quindi la distanza di
  Hamming fra due stringhe binarie di $n$ bit si calcola in tempo $\Oh(n/\log
  n)$ anziché $\Oh(n)$.
\end{osservazione}

\subsection{Distanza di Levenshtein}

Cerchiamo ora una nozione che consenta di catturare la \emph{similarità} tra
stringhe: accanto a una distanza voglio anche una \emph{misura di similarità}.
È questo il ruolo della distanza di Levenshtein, che a differenza di quella di
Hamming non richiede stringhe di uguale lunghezza.

\scan{42}%
La introduciamo misurando il \emph{lavoro} da fare per ottenere una stringa
dall'altra.

\begin{esempio}
  Due sequenze di DNA:
  \[
    \sigma_1 = \str{acgtcatca},
    \qquad
    \sigma_2 = \str{taagtgtca}.
  \]
\end{esempio}

Le operazioni elementari sono quattro; a ciascuna associamo un costo:
\[
  \begin{array}{r@{\quad}c@{\qquad}l}
    \text{costo } 0 & \mathbf{m} & \textit{match}\\[8pt]
    \text{costo } 1 &
      \left\{
        \begin{array}{@{}c@{}}
          \mathbf{i}\\[3pt] \mathbf{d}\\[3pt] \mathbf{s}
        \end{array}
      \right. &
      \begin{array}{@{}l@{}}
        \textit{insert}\\[3pt] \textit{delete}\\[3pt] \textit{substitute}
      \end{array}
  \end{array}
\]
cioè, rispettivamente, corrispondenza, inserzione, cancellazione e
sostituzione.

Qual è il costo minimo di una sequenza di operazioni in
$\set{\mathbf{m},\mathbf{i},\mathbf{d},\mathbf{s}}$ che permetta di ottenere
$\sigma_2$ da $\sigma_1$?

\begin{osservazione}
  Rispondendo a questa domanda stiamo in realtà definendo il \emph{costo
  dell'allineamento} di $\sigma_1$ e $\sigma_2$.
\end{osservazione}

\begin{definizione}[Allineamento]
  \label{def:allineamento}
  Un \emph{allineamento} (o \emph{edit string}) di $\sigma_1$ e $\sigma_2$ è una
  stringa
  \[
    \eta \;\in\; \set{\mathbf{m},\mathbf{i},\mathbf{d},\mathbf{s}}^{*}
  \]
  che, letta come programma, converte $\sigma_1$ in $\sigma_2$. Il suo
  \emph{costo} è la somma dei costi delle operazioni che la compongono.
\end{definizione}

\begin{figure}[H]
  \centering
  % allineamento-esempio -- pag. 42
% Allineamento di sigma_1 = acgtcatca con sigma_2 = taagtgtca.
% Tre righe incolonnate: sopra i caratteri di sigma_1, in mezzo (riquadrata)
% l'edit string eta = i m s m m s d m m m, sotto i caratteri di sigma_2.
% Colonna 1: solo sigma_2  -> inserzione (i);  colonna 7: solo sigma_1 -> cancellazione (d).
\begin{tikzpicture}[
    font=\small,
    car/.style={inner sep=1pt},
  ]

  \def\dx{0.72}   % passo fra le colonne
  \def\dy{0.66}   % distanza fra le righe

  % ---- riga 1: sigma_1 (la colonna 1 e' vuota: e' l'inserzione) -------------
  \foreach \c [count=\k from 2] in {a,c,g,t,c,a,t,c,a}
    \node[car] at ({\k*\dx},\dy) {\texttt{\c}};

  % ---- riga 2: l'allineamento eta ------------------------------------------
  \foreach \c [count=\k from 1] in {i,m,s,m,m,s,d,m,m,m}
    \node[car] at ({\k*\dx},0) {$\mathbf{\c}$};

  % ---- riga 3: sigma_2 (la colonna 7 e' vuota: e' la cancellazione) ---------
  \foreach \c [count=\k from 1] in {t,a,a,g,t,g}
    \node[car] at ({\k*\dx},-\dy) {\texttt{\c}};
  \foreach \c [count=\k from 8] in {t,c,a}
    \node[car] at ({\k*\dx},-\dy) {\texttt{\c}};

  % ---- il riquadro attorno alla sola edit string ----------------------------
  \draw[line width=0.7pt]
    ({0.5*\dx-0.06},-0.28) rectangle ({10.5*\dx+0.06},0.28);

  % ---- etichette a sinistra -------------------------------------------------
  \node[left,inner sep=2pt] at ({0.5*\dx-0.20},\dy)  {$\sigma_1 =$};
  \node[left,inner sep=2pt] at ({0.5*\dx-0.20},-\dy) {$\sigma_2 =$};
  \node[left,inner sep=2pt] at ({0.5*\dx-0.20},0)    {$\eta \;\Rightarrow$};

\end{tikzpicture}

  \caption{Un allineamento di $\sigma_1=\str{acgtcatca}$ e
  $\sigma_2=\str{taagtgtca}$: la riga centrale, riquadrata, è l'allineamento
  $\eta=\mathbf{i}\,\mathbf{m}\,\mathbf{s}\,\mathbf{m}\,\mathbf{m}\,\mathbf{s}\,\mathbf{d}\,\mathbf{m}\,\mathbf{m}\,\mathbf{m}$;
  sopra e sotto stanno i caratteri di $\sigma_1$ e di $\sigma_2$ che ciascuna
  operazione mette in corrispondenza. È il migliore possibile?}
  \label{fig:allineamento-esempio}
\end{figure}

\begin{definizione}[Distanza di Levenshtein]
  \label{def:levenshtein}
  La \emph{distanza di Levenshtein} (o \emph{distanza di edit}) fra $\sigma_1$ e
  $\sigma_2$ è il costo minimo di un allineamento che converte $\sigma_1$ in
  $\sigma_2$:
  \[
    \dedit(\sigma_1,\sigma_2)\;=\;\min\set{\,\mathrm{costo}(\eta)\;:\;
      \eta \text{ converte } \sigma_1 \text{ in } \sigma_2\,}.
  \]
\end{definizione}

\begin{domanda}[Problema]
  Come determino $\eta$, date $\sigma_1$ e $\sigma_2$, in modo che il costo sia
  \emph{minimo} (e che $\eta$ converta $\sigma_1$ in $\sigma_2$)?
\end{domanda}

Le due distanze non sono indipendenti. Un allineamento che non usi mai
$\mathbf{i}$ né $\mathbf{d}$ è esattamente un confronto posizione per posizione,
cioè un allineamento di tipo Hamming; poiché la distanza di Levenshtein è il
minimo del costo su \emph{tutti} gli allineamenti, si ha
\[
  \dham(\sigma_1,\sigma_2) \;\ge\; \dedit(\sigma_1,\sigma_2)
\]
per ogni coppia di stringhe di uguale lunghezza (l'unico caso in cui $\dham$ è
definita).


\section{Costi e punteggi}

\scan{43}%
Un allineamento si può valutare in due modi \emph{duali}: gli si può associare
un \emph{costo}, che vogliamo minimizzare, oppure un \emph{punteggio}, che
vogliamo massimizzare.
\[
  \underset{\displaystyle\downarrow}{\text{costo}}
  \qquad\text{---}\qquad
  \underset{\displaystyle\uparrow}{\text{punteggio}}
\]

\subsection{Matrici di costo e di punteggio}

Non tutte le sostituzioni devono pesare allo stesso modo: posso assegnare costi
(o punteggi) diversi alle diverse sostituzioni, raccogliendoli in una matrice
indicizzata dalle coppie di caratteri dell'alfabeto.

\begin{esempio}
  Nel DNA le due purine $(a,g)$ sono simili tra loro, e così le due pirimidine
  $(c,t)$: una sostituzione interna a una delle due classi va quindi penalizzata
  meno di una sostituzione fra classi diverse. Una possibile matrice di
  \emph{punteggio} è allora
  \[
    \begin{array}{c|rrrr}
         &  a &  c &  g &  t \\ \hline
       a &  2 & -1 &  1 & -1 \\
       c & -1 &  2 & -1 &  1 \\
       g &  1 & -1 &  2 & -1 \\
       t & -1 &  1 & -1 &  2
    \end{array}
  \]
\end{esempio}

\begin{esempio}
  Posso penalizzare in modo specifico l'\emph{apertura} di un $\gap$: nel DNA è
  infatti molto probabile che un gap, una volta aperto, si estenda.
\end{esempio}

\subsection{Funzioni di costo per i gap}

\scan{47}%
In un allineamento chiamiamo \emph{gap} un blocco massimale di spazi
consecutivi e indichiamo con $\gamma(n)$ il costo di un gap di lunghezza $n$.
I modelli di costo che si usano sono tre.

\begin{description}
  \item[Lineare:] un gap di lunghezza $n$ ha costo
    \[
      \gamma(n)\;\equiv\;n\cdot g .
    \]
  \item[Convesso:] un gap di lunghezza $n$ ha costo $\gamma(n)$ tale che
    \[
      \gamma(i+1)-\gamma(i)\;\le\;\gamma(i)-\gamma(i-1)
      \qquad \forall\, i\ge 1
    \]
    (gli incrementi non crescono: allungare un gap già lungo costa
    marginalmente meno).
  \item[Affine:]
    \[
      \gamma(n)\;\equiv\;g_o+(n-1)\cdot g_e ,
    \]
    dove $g_o$ è il costo di \emph{apertura} del gap e $g_e$ quello di
    \emph{estensione}.
\end{description}

Nel seguito, salvo avviso contrario, useremo lo schema di costi \emph{unitari}
introdotto sopra --- $0$ per il match, $1$ per ciascuna delle altre tre
operazioni --- che è il caso lineare con $g=1$: né la matrice di punteggio né la
penalizzazione dell'apertura dei gap intervengono negli algoritmi che seguono.

\subsection{Quale distanza usare?}

\scan{43}%
Come regola di massima:
\begin{align*}
  \text{stringhe lunghe} &\;\longleftrightarrow\; \text{Levenshtein},\\
  \text{stringhe corte}  &\;\longleftrightarrow\; \text{Hamming}.
\end{align*}


\section{Programmazione dinamica}

\subsection{La ricorrenza}

Poniamo $n_1=\abs{\sigma_1}$ e $n_2=\abs{\sigma_2}$ e indichiamo con
\[
  \Dm{i}{j}\;=\;\dedit\bigl(\pre{\sigma_1}{i},\ \pre{\sigma_2}{j}\bigr)
\]
la distanza fra i due prefissi di lunghezza $i$ e $j$. Un allineamento ottimo
dei due prefissi termina necessariamente con una delle tre mosse possibili
sull'ultima colonna, e questo dà la ricorrenza.

\begin{figure}[H]
  \centering
  % dp-stencil -- pag. 43
% Stencil "in avanti" della ricorsione: dalla cella (i,j) si puo' passare a
% (i,j+1) (avanza solo sigma_2: inserzione), a (i+1,j) (avanza solo sigma_1:
% cancellazione) oppure a (i+1,j+1) (avanzano entrambe: match o sostituzione).
% NB: nel manoscritto gli assi sono scambiati (i+1 sull'orizzontale e j+1 sulla
% verticale); qui si adotta la convenzione usata in tutto il capitolo, cioe'
% i cresce verso il basso (righe, su sigma_1) e j verso destra (colonne, su sigma_2).
\begin{tikzpicture}[
    font=\footnotesize,
    >={Stealth[length=4.5pt,width=3.5pt]},
    mossa/.style={->,line width=0.6pt},
    glossa/.style={black!60,font=\scriptsize},
  ]

  \def\h{0.55}   % semilato della cella

  % ---- la cella (i,j) ------------------------------------------------------
  \draw[line width=0.6pt] (-\h,-\h) rectangle (\h,\h);
  \node[above left,inner sep=2pt] at (-\h,\h) {$(i,j)$};

  % ---- avanza solo sigma_2: inserzione -> (i, j+1) --------------------------
  \draw[mossa] (\h,0) -- (3.30,0);
  \node[above,inner sep=3pt] at (3.10,0) {$j+1$};
  \node[glossa,below,inner sep=4pt] at (1.95,0) {inserzione};

  % ---- avanza solo sigma_1: cancellazione -> (i+1, j) -----------------------
  \draw[mossa] (0,-\h) -- (0,-3.00);
  \node[below,inner sep=3pt] at (0,-2.95) {$i+1$};
  \node[glossa,left,inner sep=4pt] at (0,-1.75) {cancellazione};

  % ---- avanzano entrambe: match o sostituzione -> (i+1, j+1) ----------------
  \draw[mossa] (\h,-\h) -- (2.70,-2.70)
        node[glossa,pos=0.45,sloped,above,inner sep=2pt]
          {match / sostituzione};
  \node[right,inner sep=4pt,align=center] at (2.70,-2.70) {$i+1$\\$j+1$};

\end{tikzpicture}

  \caption{Le tre transizioni elementari della ricorsione: dalla coppia di
  indici $(i,j)$ si passa a $(i+1,j)$ (avanza solo $\sigma_1$: cancellazione), a
  $(i,j+1)$ (avanza solo $\sigma_2$: inserzione) oppure a $(i+1,j+1)$ (avanzano
  entrambe: match o sostituzione).}
  \label{fig:dp-stencil}
\end{figure}

Posto $\mathrm{cost}(x,y)=0$ se $x=y$ e $\mathrm{cost}(x,y)=1$ altrimenti, si
ha dunque
\[
  \Dm{i}{0}=i \quad \forall i, \qquad \Dm{0}{j}=j \quad \forall j,
\]
\[
  \Dm{i}{j}=\min\bigl\{\,\Dm{i-1}{j-1}+\mathrm{cost}(\sigma_1[i],\sigma_2[j]),\;
  \Dm{i-1}{j}+1,\; \Dm{i}{j-1}+1 \,\bigr\},
\]
e la distanza cercata è $\Dm{n_1}{n_2}$. I casi base dicono semplicemente che
per allineare un prefisso con la stringa vuota $\vuota$ non si può fare altro
che cancellarne (o inserirne) tutti i caratteri.

\subsection{Algoritmo ricorsivo esponenziale}

\scan{44}%
La ricorrenza si traduce direttamente in un algoritmo ricorsivo: per allineare
il prefisso $\pre{\sigma_1}{i}$ con il prefisso $\pre{\sigma_2}{j}$ basta
provare le tre mosse possibili sull'ultima colonna dell'allineamento e tenere
la migliore.

\begin{algorithm}[H]
\caption{\textsc{allineamento\_exp}: allineamento globale, versione ricorsiva}
\begin{algorithmic}[1]
\Function{allineamento\_exp}{$\sigma_1,\sigma_2,i,j$}
  \If{$i=0$ \textbf{oppure} $j=0$}
    \State \Return $i+j$ \Comment{casi base}
  \EndIf
  \If{$\sigma_1[i]=\sigma_2[j]$}
    \State $c_S \gets \Call{allineamento\_exp}{\sigma_1,\sigma_2,i-1,j-1}$
           \Comment{match: costo $0$}
  \Else
    \State $c_S \gets 1+\Call{allineamento\_exp}{\sigma_1,\sigma_2,i-1,j-1}$
           \Comment{sostituzione}
  \EndIf
  \State $c_D \gets 1+\Call{allineamento\_exp}{\sigma_1,\sigma_2,i-1,j}$
         \Comment{cancellazione di $\sigma_1[i]$}
  \State $c_I \gets 1+\Call{allineamento\_exp}{\sigma_1,\sigma_2,i,j-1}$
         \Comment{inserzione di $\sigma_2[j]$}
  \State \Return $\min(c_S,c_D,c_I)$
\EndFunction
\end{algorithmic}
\end{algorithm}

Le tre variabili raccolgono i costi delle tre mosse: $c_S$ per la mossa
diagonale (match, gratuito, oppure sostituzione, di costo $1$), $c_D$ per la
cancellazione di $\sigma_1[i]$ --- si consuma un carattere di $\sigma_1$ e si
apre un gap su $\sigma_2$ --- e $c_I$ per l'inserzione di $\sigma_2[j]$. I casi
base sono corretti: se $i=0$ resta da allineare $\vuota$ con
$\pre{\sigma_2}{j}$ e servono esattamente $j$ inserzioni, e simmetricamente per
$j=0$; da cui \textbf{return} $i+j$.

\medskip
\noindent$\leadsto$\; \textbf{Il costo è esponenziale!}

\begin{osservazione}
  Sia $C(n)$ il massimo numero di chiamate su un'istanza con $i+j=n$. Le tre
  chiamate ricorsive hanno taglia $n-2$ (diagonale), $n-1$ e $n-1$, quindi
  \[
    C(n) \;\le\; 2\,C(n-1) + C(n-2) + 1, \qquad C(0)=C(1)=1 .
  \]
  È una maggiorazione e non un'uguaglianza, perché i casi base scattano non
  appena \emph{uno} dei due indici si annulla, e non quando $n$ diventa piccolo.
  L'equazione caratteristica $x^2=2x+1$ dà $x=1\pm\sqrt2$, dunque
  \[
    C(n)\;=\;\Oh\!\big((1+\sqrt{2})^{\,n}\big),
    \qquad n=n_1+n_2 .
  \]
  Una stima più grossolana ($3$ chiamate di taglia $n-1$) dà $\Oh(3^{\,n})$: in
  entrambi i casi il costo è esponenziale nella somma delle lunghezze.
\end{osservazione}

\subsection{Tabulazione e lower bound}

Ma i possibili valori della coppia $(i,j)$ sono solo $n_1\times n_2$: i
sottoproblemi distinti sono pochi e la ricorsione li ricalcola un numero
esponenziale di volte. Conviene allora \emph{tabulare} i risultati $\leadsto$
una matrice di dimensione quadratica. Da qui il \emph{lower bound}
\[
  \Om\big(n_1\times n_2\big).
\]

\begin{osservazione}
  Due precisazioni. (i) Le coppie $(i,j)$ ammissibili sono in realtà
  $(n_1+1)(n_2+1)$, perché gli indici partono da $0$ (prefisso vuoto): la stima
  $n_1\times n_2$ è corretta come ordine di grandezza. (ii) Il \emph{lower
  bound} $\Om(n_1\cdot n_2)$ è un limite inferiore \emph{del metodo}, non del
  problema: riempire tutta la matrice costa almeno quanto il numero delle sue
  celle. Che non si possa fare sostanzialmente meglio per la distanza di edit è
  noto solo come risultato \emph{condizionale} (sotto l'ipotesi SETH).
\end{osservazione}

Tabulando i valori $\Dm{i}{j}$ in una matrice $(n_1+1)\times(n_2+1)$ e
riempiendola per righe (o per colonne) crescenti, ogni cella si calcola in
tempo costante dalle tre celle già note: la distanza cercata è $\Dm{n_1}{n_2}$,
e l'allineamento ottimo $\eta$ si recupera con un \emph{backtracking} dalla
cella finale verso $(0,0)$, seguendo a ritroso le mosse che hanno realizzato il
minimo.

\subsection{Che cosa dà la matrice}

\scan{45}%
Per calcolare la cella $(i,j)$ servono soltanto le tre celle $(i-1,j-1)$,
$(i-1,j)$ e $(i,j-1)$: stanno tutte nella colonna corrente o in quella
immediatamente precedente. Mi basta dunque \emph{spazio lineare}! \dots\ ma solo
per la \emph{distanza} (senza la stringa di edit).

\begin{figure}[H]
  \centering
  % dp-stencil-lineare -- pag. 45
% Stencil "all'indietro": le tre celle da cui dipende (i,j) sono (i-1,j-1),
% (i-1,j) e (i,j-1); stanno tutte nelle colonne j-1 e j, quindi per la sola
% distanza bastano due colonne in memoria (l'ovale racchiude "tutto cio' che serve").
\begin{tikzpicture}[
    font=\footnotesize,
    >={Stealth[length=4.5pt,width=3.5pt]},
    cella/.style={circle,fill=black,inner sep=0pt,minimum size=3.4pt},
    obiettivo/.style={circle,draw,line width=0.8pt,fill=black!12,
                      inner sep=0pt,minimum size=6.5pt},
    griglia/.style={black!45,line width=0.4pt},
    dipende/.style={->,line width=0.7pt},
    et/.style={font=\scriptsize,inner sep=0pt},
  ]

  \def\d{1.35}   % lato del blocco 2x2

  % ---- le quattro celle ----------------------------------------------------
  \coordinate (dd) at (0,0);        % (i-1,j-1)
  \coordinate (du) at (\d,0);       % (i-1,j)
  \coordinate (sd) at (0,-\d);      % (i,j-1)
  \coordinate (ij) at (\d,-\d);     % (i,j)   <- cella da calcolare

  % ---- residuo della griglia della matrice ---------------------------------
  \draw[griglia] (dd) -- (du);
  \draw[griglia] (dd) -- (sd);

  % ---- le tre dipendenze, tutte convergenti su (i,j) -----------------------
  \draw[dipende] (dd) -- ($(ij)+(-0.12,0.12)$);
  \draw[dipende] (du) -- ($(ij)+(0,0.17)$);
  \draw[dipende] (sd) -- ($(ij)+(-0.17,0)$);

  % ---- i nodi --------------------------------------------------------------
  \node[cella] at (dd) {};
  \node[cella] at (du) {};
  \node[cella] at (sd) {};
  \node[obiettivo] at (ij) {};

  % ---- etichette delle celle -----------------------------------------------
  \node[et,anchor=south east] at ($(dd)+(-0.18,0.16)$) {$(i-1,j-1)$};
  \node[et,anchor=south west] at ($(du)+(0.18,0.16)$)  {$(i-1,j)$};
  \node[et,anchor=north east] at ($(sd)+(-0.18,-0.16)$) {$(i,j-1)$};
  \node[et,anchor=west]       at ($(ij)+(0.22,-0.02)$)  {$(i,j)$};

  % ---- l'ovale: "tutto cio' che serve sta qui dentro" ----------------------
  \draw[black!55,dashed,line width=0.5pt,rotate around={-20:(0.5*\d,-0.5*\d)}]
        (0.5*\d,-0.5*\d) ellipse [x radius=1.02, y radius=1.40];

  % ---- le due sole colonne coinvolte ---------------------------------------
  \node[et,black!60] at (0,-\d-1.02)  {$j-1$};
  \node[et,black!60] at (\d,-\d-1.02) {$j$};
  \draw[decorate,decoration={brace,mirror,amplitude=4pt,raise=1pt},
        draw=black!55,line width=0.5pt]
        (-0.28,-\d-1.22) -- (\d+0.28,-\d-1.22);
  \node[et,black!60,anchor=north] at (0.5*\d,-\d-1.64) {due colonne};

\end{tikzpicture}

  \caption{Le tre celle da cui dipende $(i,j)$ cadono tutte fra la colonna $j-1$
  e la colonna $j$: per calcolare la sola distanza basta tenere in memoria due
  colonne.}
  \label{fig:dp-stencil-lineare}
\end{figure}

\begin{figure}[H]
  \centering
  % dp-due-colonne -- pag. 45
% La matrice DP letta "a due colonne": la colonna j-1 e' completa, la colonna j
% e' riempita fino alla riga i-1; le tre celle da cui dipende (i,j) stanno tutte
% li' dentro. Ridisegno dello schizzo: nell'originale i due segmenti verticali
% sono sfalsati (destra: dall'alto a meta'; sinistra: da meta' in giu'), qui
% sono due colonne intere con l'ombreggiatura a indicare cio' che sta in memoria.
\begin{tikzpicture}[
    x=1cm, y=1cm,
    font=\footnotesize,
    et/.style={font=\footnotesize, inner sep=1.5pt},
    dip/.style={-{Stealth[length=4.5pt,width=3.2pt]}, black!65, line width=0.55pt,
                shorten <=4pt, shorten >=5pt},
  ]

  % ---- geometria: 8 colonne (0.8) x 8 righe (0.65) --------------------------
  % colonna j-1 = [3.2,4.0], colonna j = [4.0,4.8]
  % riga i-1    = [-2.60,-3.25], riga i = [-3.25,-3.90]

  % ---- cio' che sta in memoria ---------------------------------------------
  \fill[black!16] (3.2,0) rectangle (4.0,-5.2);      % colonna j-1: completa
  \fill[black!8]  (4.0,0) rectangle (4.8,-3.25);     % colonna j: fino alla riga i-1

  % ---- griglia leggera ------------------------------------------------------
  \foreach \c in {1,...,7} \draw[black!20,line width=0.3pt] (0.8*\c,0) -- (0.8*\c,-5.2);
  \foreach \r in {1,...,7} \draw[black!20,line width=0.3pt] (0,-0.65*\r) -- (6.4,-0.65*\r);

  % ---- confini delle due colonne conservate --------------------------------
  \draw[black!50,line width=0.5pt] (3.2,0) -- (3.2,-5.2);
  \draw[black!50,line width=0.5pt] (4.0,0) -- (4.0,-5.2);
  \draw[black!50,line width=0.5pt] (4.8,0) -- (4.8,-5.2);
  \draw[black!50,line width=0.5pt] (4.0,-3.25) -- (4.8,-3.25);  % fronte di calcolo

  % ---- bordo della matrice --------------------------------------------------
  \draw[black!80, line width=0.7pt] (0,0) rectangle (6.4,-5.2);

  % ---- la cella da calcolare ------------------------------------------------
  \draw[black, line width=0.9pt, fill=white] (4.0,-3.25) rectangle (4.8,-3.90);

  % ---- le tre celle da cui dipende ------------------------------------------
  \fill[black] (3.6,-2.925) circle (2.2pt);   % (i-1,j-1)
  \fill[black] (3.6,-3.575) circle (2.2pt);   % (i,j-1)
  \fill[black] (4.4,-2.925) circle (2.2pt);   % (i-1,j)

  \draw[dip] (3.6,-2.925) -- (4.4,-3.575);
  \draw[dip] (3.6,-3.575) -- (4.4,-3.575);
  \draw[dip] (4.4,-2.925) -- (4.4,-3.575);

  \node[et, anchor=west, inner sep=2pt, fill=white] at (4.95,-3.575) {$(i,j)$};

  % ---- etichette di riga e di colonna --------------------------------------
  \node[et, anchor=north] at (3.6,-5.32) {$j-1$};
  \node[et, anchor=north] at (4.4,-5.32) {$j$};
  \node[et, anchor=east]  at (-0.10,-2.925) {$i-1$};
  \node[et, anchor=east]  at (-0.10,-3.575) {$i$};

  % ---- assi -----------------------------------------------------------------
  \node[et, anchor=south east] at (6.4,0.10) {$\sigma_2$};
  \node[et, anchor=east]       at (-0.10,-0.32) {$\sigma_1$};

  % ---- che cosa e' ombreggiato ---------------------------------------------
  \draw[black!55, line width=0.4pt, -{Stealth[length=4pt,width=3pt]}]
        (6.85,-1.30) to[out=180,in=25] (4.55,-0.70);
  \node[et, anchor=north west, text=black!55, align=left] at (6.78,-1.12)
        {le sole due colonne\\tenute in memoria};

\end{tikzpicture}

  \caption{La stessa osservazione letta sulla matrice: si conservano la colonna
  $j-1$ (già completata) e la colonna $j$ riempita fino alla riga $i-1$; le tre
  celle marcate sono esattamente quelle da cui dipende $(i,j)$.}
  \label{fig:dp-due-colonne}
\end{figure}

\begin{itemize}
  \item La matrice consente di calcolare la distanza di Levenshtein fra
        $\sigma_1$ e $\sigma_2$, cioè $\dedit(\sigma_1,\sigma_2)$, in
        $\Oh(n_1\times n_2)$ tempo \emph{e} spazio.
  \item La matrice mi permette di calcolare \emph{tutti} gli allineamenti di
        costo $\dedit(\sigma_1,\sigma_2)$.
  \item Se riduco la matrice \dubbioinl{all'ultima colonna calcolata} riduco lo
        spazio a $\Oh(\min(n_1,n_2))$ --- facendo scorrere la matrice lungo il
        lato più lungo --- \emph{ma non sono più in grado di produrre $\eta$}.
\end{itemize}

\subsection{Algoritmo di Hirschberg}

Per ottenere \emph{anche} $\eta$ restando in spazio lineare: \textbf{algoritmo
di Hirschberg}.

\begin{osservazione}
  La matrice non mi dà soltanto $\dedit(\sigma_1,\sigma_2)$: mi consente di
  avere
  \[
    \dedit\bigl(\sub{\sigma_1}{i}{j},\ \sub{\sigma_2}{i'}{j'}\bigr),
  \]
  cioè la distanza fra due \emph{sottostringhe} qualsiasi. Ogni sottorettangolo
  della matrice è a sua volta un problema di allineamento fra due sottostringhe:
  è questa la proprietà che permette di spezzare il problema.
\end{osservazione}

\begin{figure}[H]
  \centering
  % hirschberg-divide -- pag. 45
% Schema divide et impera di Hirschberg: la riga mediana i* taglia la matrice a
% meta'; il cammino ottimo la attraversa in una colonna j*; si ricorre sui due
% sottorettangoli [0,i*]x[0,j*] e [i*,n1]x[j*,n2].
% Nell'originale la colonna j* e' disegnata come DUE verticali (solo per
% separare visivamente i due sottorettangoli): qui una sola verticale, con le
% due regioni escluse dalla ricorsione lasciate in bianco. Attenzione: a questo
% livello l'INTERO rettangolo viene riempito (in avanti sopra i*, all'indietro
% sotto i*) per individuare j*; sono le CHIAMATE RICORSIVE a limitarsi ai due
% sottorettangoli ombreggiati.
\begin{tikzpicture}[
    x=1cm, y=1cm,
    font=\footnotesize,
    et/.style={font=\footnotesize, inner sep=1.5pt},
    cam/.style={black, line width=0.9pt, line join=round},
  ]

  % ---- geometria ------------------------------------------------------------
  % matrice (0,0) -- (9.6,-5.6);  riga mediana y=-2.8;  colonna j* in x=3.6

  % ---- i due sottorettangoli su cui si ricorre ------------------------------
  \fill[black!9] (0,0)      rectangle (3.6,-2.8);
  \fill[black!9] (3.6,-2.8) rectangle (9.6,-5.6);

  % ---- bordo e linee di taglio ---------------------------------------------
  \draw[black!80, line width=0.7pt] (0,0) rectangle (9.6,-5.6);
  \draw[black!70, line width=0.6pt] (0,-2.8) -- (9.6,-2.8);                 % riga i*
  \draw[black!70, line width=0.6pt, densely dashed] (3.6,0) -- (3.6,-5.6);  % colonna j*

  % ---- le due regioni escluse dalle chiamate ricorsive ---------------------
  \node[et, text=black!50, align=center] at (6.60,-1.35) {esclusa dalla\\ricorsione};
  \node[et, text=black!50, align=center] at (1.80,-4.20) {esclusa dalla\\ricorsione};

  % ---- il cammino ottimo ----------------------------------------------------
  % meta' superiore: da (0,0) alla riga mediana, nella colonna j*
  \draw[cam] (0,0) -- (0.40,-0.40) -- (0.40,-0.80) -- (0.85,-1.25) -- (1.30,-1.25)
                   -- (1.30,-1.65) -- (1.75,-2.10) -- (2.20,-2.10) -- (2.60,-2.50)
                   -- (3.10,-2.50) -- (3.60,-2.80);
  % meta' inferiore: da (i*,j*) fino a (n1,n2)
  \draw[cam] (3.60,-2.80) -- (3.95,-3.15) -- (4.45,-3.15) -- (4.85,-3.55)
                   -- (5.35,-3.55) -- (5.35,-3.95) -- (5.85,-4.30) -- (6.45,-4.30)
                   -- (6.85,-4.65) -- (7.45,-4.65) -- (7.85,-5.00) -- (8.45,-5.00)
                   -- (8.90,-5.30) -- (9.20,-5.30) -- (9.60,-5.60);

  \fill[black] (3.6,-2.8) circle (2.4pt);

  % ---- etichette ------------------------------------------------------------
  \node[et, anchor=south east] at (-0.02,0.06) {$(0,0)$};
  \node[et, anchor=north west] at (9.62,-5.66) {$(n_1,n_2)$};
  \node[et, anchor=east]       at (-0.12,-2.8) {$i^{*}$};
  \node[et, anchor=north]      at (3.60,-5.74) {$j^{*}$};

  \node[et, anchor=south] at (4.8,0.06) {$\sigma_2$};
  \node[et, rotate=90]    at (-0.58,-1.40) {$\sigma_1$};

\end{tikzpicture}

  \caption{Schema divide et impera di Hirschberg. La riga mediana $i^{*}$ taglia
  la matrice in due metà; il cammino ottimo la attraversa in una colonna
  $j^{*}$, individuata riempiendo l'\emph{intero} rettangolo --- in avanti sopra
  la riga $i^{*}$, all'indietro sotto di essa. Fatto questo, il problema si
  spezza nei due sottorettangoli $[0,i^{*}]\times[0,j^{*}]$ e
  $[i^{*},n_1]\times[j^{*},n_2]$, sui quali si ricorre; le altre due regioni,
  attraversate una sola volta a questo livello, non entrano nelle chiamate
  ricorsive.}
  \label{fig:hirschberg-divide}
\end{figure}

Poiché a ogni livello della ricorsione l'area complessiva da riempire si
dimezza, il tempo resta $\Th(n_1 n_2)$ mentre lo spazio scende a
$\Th(\min(n_1,n_2))$, questa volta \emph{senza} perdere l'allineamento $\eta$.


\section{Allineamento locale}

\scan{44}%
Invece di allineare tutto $\sigma_1$ con tutto $\sigma_2$, cerco le porzioni più
piccole del pattern che si allineano in modo ottimo con il testo.

\begin{osservazione}
  È il passaggio dal problema \emph{globale} (prefisso contro prefisso) a quello
  \emph{locale}. Nella versione che useremo i ruoli delle due stringhe si
  separano: la prima stringa è il \emph{pattern} e la chiamiamo $\sigma$, la
  seconda è il \emph{testo} e la chiamiamo $T$; si vuole allineare \emph{tutto}
  $\sigma$ con una \emph{sottostringa} di $T$, cioè con un suffisso di un
  prefisso di $T$.
\end{osservazione}

\scan{46}%
\begin{figure}[H]
  \centering
  % allineamento-locale-prima-riga -- pag. 46
% La matrice sigma x T dell'allineamento locale: la riga 0 vale 0 in OGNI
% colonna (il gap iniziale sul testo e' gratuito). Il backtracking risale da una
% cella dell'ultima riga fino a una cella qualunque della riga 0: le due
% verticali delimitano la sottostringa di T allineata contro tutto sigma.
% Nel manoscritto le etichette sono sigma_1 / sigma_2: qui sigma / T, come nel
% resto della sezione sull'allineamento locale.
\begin{tikzpicture}[
    x=1cm, y=1cm,
    font=\footnotesize,
    et/.style={font=\footnotesize, inner sep=1.5pt},
    graffa/.style={decorate, decoration={brace, amplitude=4pt, raise=2pt},
                   draw=black!65, line width=0.5pt},
  ]

  % ---- geometria: colonne 0..13 (0.6) x righe 0..3 (0.5) -------------------
  % colonna k: centro 0.6k+0.3;  riga 0: y=-0.25;  righe 1,2,3: -0.75,-1.25,-1.75
  % la colonna 0 e la riga 0 sono le condizioni al contorno

  % ---- la banda di colonne di T allineate contro sigma ---------------------
  \fill[black!7] (4.8,0) rectangle (7.2,-2.0);

  % ---- bordi della matrice (a destra resta aperta: T prosegue) -------------
  \draw[black!80, line width=0.7pt] (0,0)    -- (8.4,0);      % bordo superiore
  \draw[black!80, line width=0.7pt] (0,-2.0) -- (8.4,-2.0);   % ultima riga
  \draw[black!80, line width=0.7pt] (0,0)    -- (0,-2.0);     % bordo sinistro
  \node[et, anchor=west] at (8.55,0)     {$T$};
  \node[et, anchor=east] at (-0.20,-1.0) {$\sigma$};

  % ---- condizioni al contorno: riga 0 tutta a 0, colonna 0 uguale a i ------
  \foreach \c in {0,...,13} \node[et] at (0.6*\c+0.3,-0.25) {$0$};
  \foreach \r in {1,2,3}    \node[et] at (0.3,-0.25-0.5*\r) {$\r$};
  % la cella della riga 0 in cui termina il backtracking
  \node[circle, draw=black!75, line width=0.6pt, inner sep=1.2pt, fill=white,
        font=\footnotesize] at (4.5,-0.25) {$0$};

  % ---- le due verticali che delimitano la sottostringa ---------------------
  \draw[black!70, line width=0.6pt] (4.8,0.14) -- (4.8,-2.0);
  \draw[black!70, line width=0.6pt] (7.2,0.14) -- (7.2,-2.0);

  \draw[graffa] (4.8,0.20) -- (7.2,0.20);
  \node[et, anchor=south] at (6.0,0.46)
        {sottostringa di $T$ allineata contro tutto $\sigma$};

  % ---- il cammino di backtracking: dall'ultima riga fino alla riga 0 -------
  \draw[black, line width=0.9pt, line join=round,
        -{Stealth[length=4.5pt,width=3.2pt]}]
        (6.9,-1.75) -- (5.7,-1.75) -- (5.7,-1.25) -- (5.1,-1.25)
                    -- (5.1,-0.75) -- (4.5,-0.75) -- (4.5,-0.48);
  \fill[black] (6.9,-1.75) circle (2.1pt);

  % ---- il prefisso di T saltato a costo 0 ----------------------------------
  \draw[decorate, decoration={brace, mirror, amplitude=4pt, raise=2pt},
        draw=black!65, line width=0.5pt] (0.6,-2.05) -- (4.8,-2.05);
  \node[et, anchor=north, align=center, text=black!70] at (2.7,-2.34)
        {prefisso di $T$ saltato\\(gap iniziale a costo $0$)};

\end{tikzpicture}

  \caption{Matrice dell'allineamento locale: la prima riga vale $0$ in
  \emph{ogni} colonna. Un allineamento locale è un cammino di backtracking che
  risale fino a una cella qualsiasi della riga $0$; le due verticali delimitano
  la porzione di $T$ effettivamente allineata contro tutto $\sigma$.}
  \label{fig:all-locale-prima-riga}
\end{figure}

Assegno costo $0$ al gap iniziale e uso lo stesso algoritmo!

\begin{definizione}[Allineamento locale]
  \label{def:allineamento-locale}
  L'\emph{allineamento locale} (a distanza $\le d$) del pattern $\sigma$ nel
  testo $T$ è l'insieme delle coppie $\langle i,\eta_i\rangle$ in cui $\eta_i$ è
  l'edit string di costo minimo --- e tale costo è $\le d$ --- che allinea tutto
  $\sigma$ con un suffisso del prefisso $\pre{T}{i}$, cioè con una sottostringa
  di $T$ che termina in posizione $i$.
\end{definizione}

\begin{osservazione}
  Basta quindi compilare la stessa matrice di programmazione dinamica
  dell'allineamento globale cambiando una sola condizione al contorno. Ne
  indichiamo ancora le celle con $\Dm{i}{j}$, che d'ora in poi denota la matrice
  \emph{locale} così reinizializzata:
  \[
    \Dm{0}{j}=0 \quad \forall j, \qquad \Dm{i}{0}=i \quad \forall i,
  \]
  \[
    \Dm{i}{j}=\min\bigl\{\,\Dm{i-1}{j-1}+\mathrm{cost}(\sigma[i],T[j]),\;
    \Dm{i-1}{j}+1,\; \Dm{i}{j-1}+1 \,\bigr\}.
  \]
  Con questa inizializzazione ogni cella soddisfa l'invariante
  \[
    \Dm{i}{j}=\min_{0\le k\le j}\ \dedit\bigl(\pre{\sigma}{i},\ \sub{T}{k+1}{j}\bigr),
  \]
  cioè contiene la distanza fra il prefisso $\pre{\sigma}{i}$ e il
  \emph{migliore} suffisso del prefisso $\pre{T}{j}$: è esattamente il gap
  iniziale gratuito.
\end{osservazione}

Costo dell'allineamento locale:
\[
  \abs{\sigma}\times\abs{T},
\]
proporzionale all'area della matrice.

\begin{figure}[H]
  \centering
  % allineamento-locale-ultima-riga -- pag. 46
% La matrice sigma x T dell'allineamento locale, con l'ultima riga (indice
% |sigma|) messa in evidenza: le tacche marcano le colonne di costo minimo,
% cioe' le posizioni finali delle occorrenze approssimate di sigma in T.
\begin{tikzpicture}[
    x=1cm, y=1cm,
    font=\footnotesize,
    et/.style={font=\footnotesize, inner sep=1.5pt},
  ]

  % ---- geometria: matrice (0,0)--(8.4,-2.2); ultima riga: striscia [-1.75,-2.2]

  \fill[black!8] (0,-1.75) rectangle (8.4,-2.2);

  \draw[black!80, line width=0.7pt] (0,0) rectangle (8.4,-2.2);
  \draw[black!80, line width=0.6pt] (0,-1.75) -- (8.4,-1.75);

  % ---- etichette degli assi -------------------------------------------------
  \node[et, anchor=south] at (4.2,0.08)   {$T$};
  \node[et, anchor=east]  at (-0.18,-0.85) {$\sigma$};
  \node[et, anchor=west]  at (8.58,-1.975) {ultima riga: $i=\abs{\sigma}$};

  % ---- alcuni cammini di allineamento: partono dalla riga 0 (tutta a 0) e
  %      terminano in una colonna di costo minimo dell'ultima riga -----------
  \draw[black!35, line width=0.6pt, line join=round]
        (1.20,0) -- (1.20,-0.45) -- (1.45,-0.45) -- (1.45,-0.90)
                 -- (1.75,-0.90) -- (1.75,-1.35) -- (2.00,-1.35) -- (2.00,-1.75);
  \draw[black!35, line width=0.6pt, line join=round]
        (3.70,0) -- (3.70,-0.50) -- (4.00,-0.50) -- (4.00,-1.00)
                 -- (4.30,-1.00) -- (4.30,-1.40) -- (4.60,-1.40) -- (4.60,-1.75);
  \draw[black!35, line width=0.6pt, line join=round]
        (6.10,0) -- (6.10,-0.45) -- (6.40,-0.45) -- (6.40,-0.95)
                 -- (6.70,-0.95) -- (6.70,-1.40) -- (7.00,-1.40) -- (7.00,-1.75);

  % ---- le tacche dei minimi -------------------------------------------------
  \foreach \x in {2.00, 2.25, 4.60, 7.00, 7.25}
    \draw[black, line width=1pt] (\x,-1.85) -- (\x,-2.10);

  % ---- richiamo ai minimi ---------------------------------------------------
  \foreach \x in {2.125, 4.60, 7.125}
    \draw[black!60, line width=0.5pt, -{Stealth[length=4pt,width=3pt]}]
          (\x,-2.60) -- (\x,-2.30);
  \node[et, anchor=north, text=black!60] at (4.2,-2.62)
        {colonne di costo minimo};

\end{tikzpicture}

  \caption{La matrice $\sigma\times T$ dell'allineamento locale. I minimi
  dell'ultima riga (le tacche in basso) individuano le posizioni finali delle
  occorrenze approssimate di $\sigma$ in $T$.}
  \label{fig:all-locale-ultima-riga}
\end{figure}

\[
  \text{cerco i minimi costi}\ \iff\ \text{allineamento locale \emph{ottimo}}
\]

\begin{domanda}
  Cambia qualcosa se mi accontento della risposta alla domanda «esistono
  allineamenti locali a costo $\le d$?»
\end{domanda}

Sì: se il numero $d$ di differenze ammesse è una costante --- $d=\Oh(1)$ ---
si può fare molto meglio che riempire tutta la matrice.


\section{Verso Landau--Vishkin}

\scan{47}%
Limitandosi agli allineamenti con al più $d$ differenze, infatti, il cammino di
allineamento non può allontanarsi di più di $d$ diagonali da quella di
partenza: basta quindi calcolare una banda di $\Oh(d)$ diagonali.

\begin{figure}[H]
  \centering
  % Landau--Vishkin: la banda di O(d) diagonali nella matrice sigma x T.
% Coordinate normalizzate: (0,0) = angolo in alto a sinistra, (1,1) = angolo in
% basso a destra; l'unita' y e' negativa cosi' y cresce verso il basso.
\begin{tikzpicture}[
  x=10.8cm, y=-4.0cm,
  font=\footnotesize,
  cammino/.style={draw=blue!55!black, line width=0.7pt, line join=round},
  freccia/.style={draw=blue!55!black, line width=0.7pt,
                  -{Stealth[length=1.8mm,width=1.4mm]}},
]

  % --- banda di O(d) diagonali (le due rette hanno la stessa pendenza) ---
  \fill[black!5]
    (0.20,0) -- (0.50,0) -- (0.74,1) -- (0.44,1) -- cycle;
  \draw[gray!70, line width=0.4pt] (0.20,0) -- (0.44,1);
  \draw[gray!70, line width=0.4pt] (0.50,0) -- (0.74,1);

  % --- bordo della matrice di programmazione dinamica ---
  \draw[black, line width=0.8pt] (0,0) rectangle (1,1);

  % --- etichette degli assi ---
  \node[left=2pt]  at (0,0.45)  {$\sigma$};
  \node[above=2pt] at (0.93,0)  {$T$};

  % --- cammino di allineamento (a scaletta, dentro la banda) ---
  \draw[cammino]
    (0.32,0)       -- (0.32,0.1333)  -- (0.36,0.1333)   -- (0.36,0.2667) --
    (0.40,0.2667)  -- (0.40,0.40)    -- (0.45,0.45)     -- (0.45,0.5875) --
    (0.49333,0.5875)--(0.49333,0.725)-- (0.53667,0.725) -- (0.53667,0.8625)--
    (0.58,0.8625)  -- (0.58,1);

  % --- cella generica e scostamenti consentiti ---
  \fill[blue!55!black] (0.42,0.42) circle (1.5pt);

  \draw[freccia] (0.42,0.42) -- (0.33,0.42);
  \node[above=1pt, blue!55!black] at (0.375,0.42) {$d$};

  \draw[freccia] (0.42,0.42) -- (0.51,0.42);
  \node[above=1pt, blue!55!black] at (0.465,0.42) {$d$};

  \draw[freccia] (0.42,0.42) -- (0.42,0.70);
  \node[left=1pt, blue!55!black] at (0.42,0.56) {$d$};

\end{tikzpicture}

  \caption{Landau--Vishkin: ogni mossa non diagonale sposta il cammino di
  esattamente una diagonale, quindi con al più $d$ differenze il cammino non si
  scosta di più di $d$ diagonali da quella di partenza e resta confinato nella
  banda di $2d+1$ diagonali delimitata dalle due rette grigie. Le tre frecce
  marcate $d$ misurano questo scostamento massimo, contato in colonne (mosse
  orizzontali) o in righe (mosse verticali).}
  \label{fig:landau-vishkin-banda}
\end{figure}

Il costo complessivo dell'algoritmo di Landau--Vishkin (1985--1989) è
\[
  \Oh\bigl(d\,(\abs{\sigma}+\abs{T})\bigr),
\]
cioè $\Oh(\abs{\sigma}+\abs{T})$ --- lineare --- quando il numero $d$ di
differenze ammesse è costante, invece dei $\abs{\sigma}\times\abs{T}$ della
programmazione dinamica. La ricetta è:
\[
  \text{allineamento locale} \;+\; \text{suffix tree} \;+\; \text{Harel--Tarjan},
\]
ossia $\lca$ in tempo costante. L'algoritmo è ripreso e sviluppato nel capitolo
seguente.


\section{Riepilogo}

\scan{47}%
\begin{itemize}
  \item $\Oh(m+n)$ --- pattern matching \textbf{esatto}:
    \begin{itemize}
      \item KMP;
      \item Boyer--Moore;
      \item Aho--Corasick;
      \item Ukkonen $\Rightarrow$ \STree.
    \end{itemize}
  \item $\Oh(m\times n)$ --- pattern matching \textbf{approssimato}:
    \begin{itemize}
      \item algoritmo di \dubbioinl{Miller};
      \item distanze: Hamming e Levenshtein, quest'ultima calcolata per
        programmazione dinamica.
    \end{itemize}
\end{itemize}

\noindent
dove, come di consueto, $n=\abs{P}$ è la lunghezza del pattern e $m=\abs{T}$
quella del testo.

\chapter{Landau--Vishkin}

\scan{48}
L'algoritmo di \textbf{Landau--Vishkin} è un esempio di \emph{programmazione
dinamica ibrida}: mescola la tabella di programmazione dinamica
dell'allineamento con le strutture dati per stringhe (i suffix tree), e arriva a
\[
  \Oh(m+n)
\]
per risolvere l'allineamento a distanza $\le k$, purché $k$ sia
\emph{costante}. Nella notazione asintotica restano infatti nascoste delle
costanti, e una di esse è proprio $k$: il costo effettivo è
$\Oh\bigl(k\,(m+n)\bigr)$, che diventa lineare solo perché $k=\Oh(1)$.

Come di consueto $T$ è il testo e $P$ è il pattern, con $m=\abs{T}$ e
$n=\abs{P}$, e nelle applicazioni $\abs{P}\ll\abs{T}$. Il problema che vogliamo
risolvere è il \emph{matching con al più $k$ differenze}: individuare tutte le
posizioni di $T$ in cui termina un'occorrenza di $P$ a distanza di edit al
più $k$.

\begin{figure}[H]
  \centering
  % lv-testo-pattern -- il testo T (lungo) e il pattern P (corto), con |P| << |T|
% (quaderno, pag. 48).  Le graffe con m ed n non sono nel manoscritto: servono a
% fissare la notazione usata subito prima nel testo.
\begin{tikzpicture}[font=\footnotesize]

  % --- il testo ---
  \draw[line width=1.1pt] (0,1.45) -- (7,1.45);
  \node[right=4pt,font=\normalsize] at (7,1.45) {$T$};
  \draw[decorate,decoration={brace,mirror,amplitude=4pt},draw=black!55,
        line width=0.5pt] (0,1.30) -- (7,1.30);
  \node[text=black!55] at (3.5,0.88) {$m=\abs{T}$};

  % --- il pattern ---
  \draw[line width=1.1pt] (2.2,0) -- (3.6,0);
  \node[right=4pt,font=\normalsize] at (3.6,0) {$P$};
  \draw[decorate,decoration={brace,mirror,amplitude=4pt},draw=black!55,
        line width=0.5pt] (2.2,-0.15) -- (3.6,-0.15);
  \node[text=black!55] at (2.9,-0.55) {$n=\abs{P}$};

\end{tikzpicture}

  \caption{Il testo $T$ e il pattern $P$: si cercano le occorrenze di $P$ in $T$
  a distanza di edit $\le k$.}
  \label{fig:lv-testo-pattern}
\end{figure}

\section{L'idea: restare dentro una banda di diagonali}

La tabella di programmazione dinamica ha quindi $n+1$ righe (indicizzate dai
caratteri di $P$) e $m+1$ colonne (indicizzate dai caratteri di $T$). Se
guardassi tutta la matrice il costo sarebbe $\Oh(m\cdot n)$: l'idea è invece di
non uscire mai da una \emph{banda} di diagonali.

\begin{figure}[H]
  \centering
  % lv-banda-diagonali -- la tabella di programmazione dinamica (n+1 righe, m+1
% colonne, con m >> n: rettangolo largo e basso) e la banda di 2k+1 diagonali
% attorno alla diagonale di riferimento (quaderno, pag. 48).
\begin{tikzpicture}[
    font=\footnotesize,
    bordo/.style={line width=0.9pt, line join=round},
    onda/.style={draw=black!70, line width=0.6pt},
    diag/.style={line width=1.4pt, line cap=round},
    fre/.style={line width=0.6pt, -{Stealth[length=2.2mm,width=1.7mm]}},
  ]

  \def\Hh{2.6}

  % --- bordo della tabella: le righe sono indicizzate da P, le colonne da T ---
  \draw[bordo] (0,0) rectangle (11,\Hh);
  \node[left=4pt,font=\normalsize]  at (0,1.30)     {$P$};
  \node[above=3pt,font=\normalsize] at (10.75,\Hh) {$T$};

  % --- le due curve che delimitano la banda -------------------------------
  \draw[onda] plot[domain=0:1,samples=80,smooth,variable=\x]
      ({3.0+2.3*\x+0.13*sin(720*\x)},{\Hh*(1-\x)});
  \draw[onda] plot[domain=0:1,samples=80,smooth,variable=\x]
      ({5.8+2.3*\x+0.13*sin(720*\x)},{\Hh*(1-\x)});

  % --- la diagonale di riferimento: il tratto di match "gratis" -----------
  \draw[diag] (4.4,\Hh) -- (6.7,0);

  % --- scostamento massimo dalla diagonale --------------------------------
  \coordinate (C) at (5.55,1.30);
  \draw[fre] (C) -- (5.55,2.30);
  \draw[fre] (C) -- (5.55,0.32);
  \draw[fre] (C) -- (4.28,1.30);
  \draw[fre] (C) -- (6.82,1.30);
  \node at (5.88,0.55) {$k$};

\end{tikzpicture}

  \caption{La banda di diagonali. Il cammino di allineamento parte lungo una
  diagonale e, avendo a disposizione $k$ errori, può scostarsene di al più $k$
  posizioni in ciascuna direzione: resta perciò confinato fra le due curve che
  delimitano la banda.}
  \label{fig:lv-banda-diagonali}
\end{figure}

Ogni mossa non diagonale --- orizzontale (spazio in $P$) o verticale (spazio in
$T$) --- sposta il cammino di esattamente una diagonale, mentre le mosse
diagonali, di match o di sostituzione, lasciano invariato il numero di
diagonale. Disponendo di al più $k$ errori, il cammino resta dunque confinato in
una banda di $2k+1$ diagonali attorno a quella di partenza, e il lavoro
complessivo scende da $\Th(m\cdot n)$ a $\Oh\bigl(k\,(m+n)\bigr)$.

Restano però da attraversare rapidamente i lunghi tratti di \emph{match} lungo
una stessa diagonale: è esattamente a questo che serve l'intermezzo seguente.

\section{Intermezzo: longest common extension}

È uno dei \emph{by-product} del suffix tree.

\begin{definizione}[Longest common extension]
\label{def:lce}
  Siano $S_1$ e $S_2$ due stringhe e siano
  \[
    i\in\set{1,\dots,\abs{S_1}}, \qquad j\in\set{1,\dots,\abs{S_2}}.
  \]
  Dati $S_1$, $S_2$, $i$ e $j$, si cerca il massimo $k$ tale che
  \[
    \sub{S_1}{i}{i+k-1} \;=\; \sub{S_2}{j}{j+k-1},
  \]
  cioè la lunghezza del più lungo prefisso comune dei due suffissi
  $\suf{S_1}{i}$ e $\suf{S_2}{j}$. La indicheremo con $\lcey(S_1,S_2,i,j)$.
\end{definizione}

\begin{osservazione}
  Il $k$ della definizione di longest common extension è un semplice riuso del
  nome: non ha nulla a che vedere con il numero $k$ di errori ammessi
  nell'allineamento.
\end{osservazione}

Lo strumento per rispondere è il suffix tree generalizzato delle due stringhe,
concatenate per mezzo di due terminatori distinti:
\[
  \STree\bigl(S_1\,\text{\$}\,S_2\,\text{\texteuro}\bigr),
\]
la cui costruzione richiede tempo lineare.

\begin{figure}[H]
  \centering
  % lv-stree-lce -- lo STree generalizzato di S_1 $ S_2 € visto come triangolo:
% la stringa beta letta dalla radice a un nodo interno v e' l'estensione comune
% ai due suffissi che scendono dai rami sotto v (quaderno, pag. 48).
\begin{tikzpicture}[
    x=0.92cm, y=0.92cm,
    font=\footnotesize,
    nodo/.style={circle,fill=black,inner sep=0pt,minimum size=3.6pt},
    ramo/.style={line width=0.7pt, rounded corners=2pt, line join=round},
  ]

  %% ---- la sagoma dell'albero ---------------------------------------
  \coordinate (R) at (0,4);
  \draw[line width=0.8pt, line join=round] (-3.4,0) -- (R) -- (3.4,0) -- cycle;

  %% ---- il cammino radice -> v, etichettato beta ---------------------
  \node[nodo] at (R) {};
  \draw[ramo,line width=1.0pt]
      (0,4) -- (0.12,3.55) -- (-0.06,3.10) -- (0.14,2.65) -- (-0.04,2.25)
           -- (0.05,1.90);
  \node[font=\normalsize] at (0.45,2.85) {$\beta$};

  \node[nodo,minimum size=4.4pt] (v) at (0.05,1.90) {};
  \node[left=4pt] at (v) {$v$};

  %% ---- i due rami che scendono da v ---------------------------------
  \draw[ramo]
      (0.05,1.90) -- (-0.30,1.55) -- (-0.40,1.15) -- (-0.74,0.85)
                  -- (-0.86,0.50) -- (-1.22,0.26) -- (-1.45,0);
  \draw[ramo]
      (0.05,1.90) -- ( 0.42,1.55) -- ( 0.54,1.18) -- ( 0.82,0.86)
                  -- ( 0.94,0.48) -- ( 1.26,0.24) -- ( 1.45,0);

  \node[nodo] (fi) at (-1.45,0) {};
  \node[nodo] (fj) at ( 1.45,0) {};
  \node[below=3pt] at (fi) {$\suf{S_1}{i}$};
  \node[below=3pt] at (fj) {$\suf{S_2}{j}$};

\end{tikzpicture}

  \caption{Nello \STree di $S_1\,\text{\$}\,S_2\,\text{\texteuro}$, la stringa
  $\beta$ letta dalla radice fino a un nodo interno è l'estensione comune ai due
  suffissi che scendono dai due rami sottostanti.}
  \label{fig:lv-stree-lce}
\end{figure}

Ogni nodo dell'albero corrisponde infatti alla stringa $\beta$ che si legge
lungo il cammino dalla radice al nodo, e ogni tale $\beta$ è un \emph{prefisso
di un suffisso} di $S_1\,\text{\$}\,S_2\,\text{\texteuro}$. Posso allora
etichettare l'albero indicando quali di questi prefissi di suffissi occorrono
sia in $S_1$ sia in $S_2$: sono esattamente i candidati a essere estensioni
comuni alle due stringhe.

\scan{49}
In definitiva, la longest common extension si risolve in tempo lineare
procedendo in tre passi.

\begin{enumerate}
  \item Si costruisce il suffix tree generalizzato
        $\STree\bigl(S_1\,\text{\$}\,S_2\,\text{\texteuro}\bigr)$: usando i due
        terminatori distinti \$ e \texteuro{} ogni suffisso è associato a una
        foglia (nessun suffisso è prefisso proprio di un altro, quindi la
        corrispondenza fra suffissi e foglie è biunivoca).
  \item Con una visita \emph{bottom-up} si identificano i nodi che corrispondono
        a prefissi di suffissi sia di $S_1$ sia di $S_2$; nella stessa visita si
        calcolano le profondità in caratteri di tutti i nodi.
  \item Il valore cercato è la profondità massima di un nodo identificato da $i$
        e da $j$: è cioè la profondità, misurata in caratteri, del più profondo
        antenato comune delle due foglie corrispondenti.
\end{enumerate}

\begin{osservazione}
  I primi due passi sono un preprocessing lineare, da eseguire una volta sola.
  Aggiungendovi il preprocessing per le query di $\lca$ su
  $\STree\bigl(S_1\,\text{\$}\,S_2\,\text{\texteuro}\bigr)$ --- anch'esso
  lineare, per Harel--Tarjan --- il terzo passo diventa una singola query di
  $\lca$: si risponde quindi a $\lcey(S_1,S_2,i,j)$ in tempo $\Oh(1)$.
\end{osservazione}

\begin{figure}[H]
  \centering
  % lce-lca -- la query lce(S_1,S_2,i,j) si riduce a una query di lca sullo STree
% generalizzato di S_1 $ S_2 €: le posizioni i e j individuano due foglie, e la
% profondita' in caratteri del loro lca vale k (quaderno, pag. 49).
\begin{tikzpicture}[
    font=\footnotesize,
    nodo/.style={circle,fill=black,inner sep=0pt,minimum size=3.4pt},
    ramo/.style={line width=0.7pt, rounded corners=2pt, line join=round},
    link/.style={draw=black!65, line width=0.6pt,
                 -{Stealth[length=2.2mm,width=1.7mm]}},
  ]

  %% ================= la stringa concatenata =========================
  \draw[line width=0.7pt] (0,4.20) rectangle (6.20,4.75);
  \draw[draw=black!45,line width=0.4pt] (2.90,4.20) -- (2.90,4.75);
  \node[above=2pt] at (3.10,4.75) {$S_1\,\text{\$}\,S_2\,\text{\texteuro}$};

  % le due posizioni interrogate
  \draw[line width=0.6pt,-{Stealth[length=2mm,width=1.6mm]}]
        (1.86,3.92) -- (1.86,4.17);
  \draw[line width=0.6pt,-{Stealth[length=2mm,width=1.6mm]}]
        (3.41,3.92) -- (3.41,4.17);
  \node[font=\normalsize] at (1.86,3.68) {$i$};
  \node[font=\normalsize] at (3.41,3.68) {$j$};

  %% ================= lo STree generalizzato =========================
  \coordinate (R)  at (9.10,4.90);
  \coordinate (bl) at (7.40,1.30);
  \coordinate (br) at (11.00,1.30);
  \draw[line width=0.8pt,line join=round] (bl) -- (R) -- (br) -- cycle;
  \node[nodo] at (R) {};

  % cammino radice -> v e sua profondita' in caratteri
  \draw[ramo,line width=1.0pt]
      (9.10,4.90) -- (9.18,4.62) -- (9.00,4.35) -- (9.14,4.14) -- (9.05,3.95);
  \node[nodo,minimum size=4.4pt] (v) at (9.05,3.95) {};
  \node at (7.72,4.24) {$v=\lca$};
  \draw[draw=black!55,line width=0.4pt] (8.40,4.19) -- (8.96,4.00);

  % la profondita' in caratteri di v: e' il k cercato
  \draw[draw=black!45,line width=0.35pt,dotted] (9.10,4.90) -- (10.42,4.90);
  \draw[draw=black!45,line width=0.35pt,dotted] (9.05,3.95) -- (10.42,3.95);
  \draw[{Stealth[length=3.4pt,width=2.8pt]}-{Stealth[length=3.4pt,width=2.8pt]},
        line width=0.5pt,draw=black!70] (10.28,4.88) -- (10.28,3.97);
  \node[font=\normalsize] at (10.58,4.42) {$k$};

  % i due rami che scendono da v fino alle foglie
  \draw[ramo]
      (9.05,3.95) -- (8.72,3.50) -- (8.88,3.05) -- (8.45,2.50)
                  -- (8.58,2.00) -- (8.15,1.30);
  \draw[ramo]
      (9.05,3.95) -- (9.42,3.55) -- (9.28,3.15) -- (9.68,2.70) -- (10.00,2.60);
  \node[nodo] at (10.00,2.60) {};
  \draw[ramo] (10.00,2.60) -- (10.16,2.20) -- (10.42,1.80) -- (10.50,1.30);

  \node[nodo] (fi) at ( 8.15,1.30) {};
  \node[nodo] (fj) at (10.50,1.30) {};
  \node[font=\normalsize] at ( 8.42,1.06) {$i$};
  \node[font=\normalsize] at (10.78,1.06) {$j$};

  %% ================= dalle posizioni alle foglie ====================
  \draw[link] (1.86,3.45) .. controls (2.50,1.30) and (6.00,0.05) ..
              (8.13,1.21);
  \draw[link] (3.41,3.45) .. controls (4.30,0.70) and (7.60,-0.90) ..
              (10.48,1.22);

\end{tikzpicture}

  \caption{La query $\lcey(S_1,S_2,i,j)$ si riduce a una query di $\lca$ sul
  suffix tree generalizzato di $S_1\,\text{\$}\,S_2\,\text{\texteuro}$: la
  profondità in caratteri del più profondo antenato comune delle foglie di $i$ e
  di $j$ vale esattamente $k$.}
  \label{fig:lce-lca}
\end{figure}

\begin{osservazione}
  Perché il meccanismo funzioni devo mantenere i collegamenti fra ogni indice
  $i$, $j$ e la corrispondente foglia dell'albero, e anche delle informazioni
  sulla profondità di ogni nodo: il costo resta comunque proprio lineare.
\end{osservazione}

Nel seguito useremo la longest common extension con $S_1=P$ e $S_2=T$: il
suffix tree generalizzato da costruire una volta per tutte è dunque quello di
$P\,\text{\$}\,T\,\text{\texteuro}$, e da lì in poi ogni query
$\lcey(P,T,v,u)$ --- il più lungo tratto di match che parte da $P[v]$ e da
$T[u]$ --- costa $\Oh(1)$.

\section{Diagonali e \texorpdfstring{$d$}{d}-path}

Torniamo alla tabella di programmazione dinamica e fissiamo una convenzione per
identificare le sue diagonali: la diagonale che parte dalla cella $(0,c)$ della
prima riga porta il numero $c\ge 0$, quella che parte dalla cella $(l,0)$ della
prima colonna porta il numero $-l$. Equivalentemente, la cella $(l,c)$
appartiene alla diagonale $c-l$; i numeri di diagonale vanno da $-n$ a $m$, dove
le righe sono numerate da $0$ a $n=\abs{P}$ e le colonne da $0$ a $m=\abs{T}$, e
la diagonale $0$ è la diagonale principale. Si tratta appunto di una pura
convenzione di nomi: le diagonali non le posso costruire davvero, impiegherei
troppo tempo.

\begin{figure}[H]
  \centering
  % Convenzione di numerazione delle diagonali della tabella di programmazione
% dinamica.  Si disegna solo l'angolo in alto a sinistra: il bordo superiore
% (riga 0) e il bordo sinistro (colonna 0); la tabella resta aperta a destra e
% in basso.  La diagonale che parte da (0,c) porta il numero c, quella che
% parte da (l,0) porta il numero -l: la cella (l,c) sta sulla diagonale c-l.
\begin{tikzpicture}[
    scale=0.72,
    font=\footnotesize,
    bordo/.style={semithick},
    diag/.style={thin,black!60},
    diagprinc/.style={semithick,black}]

  % ---- i due bordi tracciati ------------------------------------------------
  \draw[bordo] (0,0) -- (14,0);      % riga 0
  \draw[bordo] (0,0) -- (0,-7.2);    % colonna 0

  % ---- numeri delle diagonali non negative (sopra la riga 0) ----------------
  \foreach \c in {0,1,2,3}
    \node[anchor=south,inner sep=3pt] at (\c,0) {$\c$};
  \node[anchor=south,inner sep=3pt] at (4.5,0) {$\cdots$};

  % ---- numeri delle diagonali negative (a sinistra della colonna 0) ---------
  \node[anchor=east,inner sep=5pt] at (0, 0)   {$0$};
  \node[anchor=east,inner sep=5pt] at (0,-1)   {$-1$};
  \node[anchor=east,inner sep=5pt] at (0,-2)   {$-2$};
  \node[anchor=east,inner sep=5pt] at (0,-3)   {$-3$};
  \node[anchor=east,inner sep=5pt] at (0,-4.3) {$\vdots$};
  \node[anchor=east,inner sep=5pt] at (0,-5.6) {$-n$};

  % ---- le diagonali ---------------------------------------------------------
  \draw[diagprinc] (0,0) -- ++(4.6,-4.6);          % diagonale principale
  \foreach \c in {1,2,3} \draw[diag] (\c,0)  -- ++(4,-4);
  \foreach \r in {1,2,3} \draw[diag] (0,-\r) -- ++(4,-4);

\end{tikzpicture}

  \caption{Convenzione per la numerazione delle diagonali della tabella di
  programmazione dinamica: la cella $(l,c)$ sta sulla diagonale $c-l$.}
  \label{fig:diagonali-numerate}
\end{figure}

\begin{definizione}[$d$-path]
\label{def:dpath}
  Un \emph{$d$-path} è un cammino nella matrice che parte dalla riga $0$ e
  ``usa'' esattamente $d$ errori, dove per errore si intende una mossa
  \begin{itemize}
    \item orizzontale;
    \item verticale;
    \item diagonale con caratteri diversi in riga e in colonna (cioè una
          sostituzione).
  \end{itemize}
  Le mosse diagonali fra caratteri uguali sono invece gratuite e non vengono
  contate.
\end{definizione}

Il legame fra i $d$-path e il problema che vogliamo risolvere è immediato.

\scan{50}
\begin{figure}[H]
  \centering
  % Un d-path che parte dalla cella (0,c) della riga 0 e termina nella cella
% (n,j) dell'ultima riga.  I tratti a 45 gradi sono corse di match (gratuite,
% attraversate con una query di longest common extension); la mossa
% orizzontale e quella verticale sono le due mosse d'errore, evidenziate dai
% pallini.  Della tabella si tracciano solo tre lati: colonna 0 a sinistra,
% riga 0 in alto (che prosegue oltre il riquadro) e riga n in basso.
\begin{tikzpicture}[
    scale=0.85,
    font=\footnotesize,
    bordo/.style={semithick},
    guida/.style={thin,black!55},
    cammino/.style={very thick},
    nodo/.style={circle,fill,inner sep=0pt,minimum size=3.6pt}]

  % ---- cornice --------------------------------------------------------------
  \draw[bordo] (0,0) -- (0,-5);        % colonna 0
  \draw[bordo] (0,0) -- (12.4,0);      % riga 0, prolungata a destra
  % riga n: tracciata a mano libera, sporge da entrambi i lati
  \draw[bordo,decorate,
        decoration={snake,amplitude=.5mm,segment length=11mm}]
       (-0.5,-5) -- (11.2,-5);

  % ---- etichette delle due stringhe -----------------------------------------
  \node[anchor=east,inner sep=4pt] at (0,-2.5) {$P$};
  \node[anchor=south east,inner sep=2pt] at (12.4,0) {$T$};

  % ---- colonne c e j sulla riga 0 -------------------------------------------
  \node[anchor=south,inner sep=3pt] at (2.8,0) {$c$};
  \node[anchor=south,inner sep=3pt] at (7.4,0) {$j$};
  \draw[guida] (7.4,0) -- (7.4,-5);    % lettura della colonna d'arrivo

  % ---- il d-path ------------------------------------------------------------
  \draw[cammino]
      (2.8,0)      -- (4.0,-1.2)       % corsa di match
   -- (4.7,-1.2)                       % mossa orizzontale: spazio in P
   -- (6.2,-2.7)                       % corsa di match
   -- (6.2,-3.8)                       % mossa verticale: spazio in T
   -- (7.4,-5.0);                      % corsa di match fino alla riga n

  \foreach \p in {(2.8,0),(4.0,-1.2),(4.7,-1.2),(6.2,-2.7),(6.2,-3.8),(7.4,-5.0),(7.4,0)}
    \node[nodo] at \p {};

\end{tikzpicture}

  \caption{Un $d$-path che parte dalla riga $0$ nella colonna $c$ e raggiunge
  l'ultima riga nella colonna $j$. I tratti diagonali sono match gratuiti; le
  mosse orizzontali e verticali (evidenziate dai pallini) sono gli errori.}
  \label{fig:lv-dpath-riga-n}
\end{figure}

\begin{osservazione}
  Se il $d$-path termina sull'ultima riga della tabella, allora ho trovato
  un'occorrenza di $P$ in $T$. Più precisamente, se il cammino parte dalla cella
  $(0,c)$ e termina nella cella $(n,j)$, esso descrive un allineamento fra $P$ e
  la sottostringa $\sub{T}{c+1}{j}$ che impiega esattamente $d$ operazioni di
  edit, e dunque
  \[
    \dedit\bigl(P,\;\sub{T}{c+1}{j}\bigr) \;\le\; d .
  \]
  In particolare, se $d\le k$ la colonna $j$ è la posizione finale di
  un'occorrenza di $P$ in $T$ con al più $k$ differenze.
\end{osservazione}

Non serve però tenere traccia di tutti i $d$-path: su ogni diagonale ne basta
uno, il più \emph{spinto in avanti}.

\begin{definizione}[$d$-path farthest reaching]
\label{def:fr-dpath}
  Il $d$-path \emph{farthest reaching} sulla diagonale $i$, che indichiamo con
  $\mathrm{f.r.}(d,i)$, è il $d$-path che termina sulla diagonale $i$ nella
  colonna $j$ \textbf{massima} fra quelle di tutti i $d$-path che terminano su
  quella diagonale.
\end{definizione}

\begin{osservazione}
  Sulla diagonale $i$ riga e colonna differiscono per la costante $i$ (la cella
  di colonna $j$ sta nella riga $j-i$): massimizzare la colonna equivale quindi
  a massimizzare la riga, e un $d$-path farthest reaching resta individuato da
  un \emph{solo} numero, la cella in cui termina.

  Per $d=0$ la nozione si semplifica ulteriormente: un $0$-path non compie mosse
  d'errore e quindi non cambia mai diagonale, cosicché la diagonale su cui
  termina coincide con la colonna da cui parte sulla riga $0$.
\end{osservazione}

\section{L'algoritmo}

Lo schema dell'algoritmo è il seguente.

\begin{algorithm}[H]
\caption{Landau--Vishkin: matching con al più $k$ differenze}
\begin{algorithmic}[1]
\Require $P$ con $\abs{P}=n$, $T$ con $\abs{T}=m$, soglia $k$
\Statex \textit{Inizializzazione: i farthest reaching $0$-path.}
\For{$i \gets 0$ \textbf{to} $m$}
  \State calcola $\lcey(P,T,1,i+1)$: è la lunghezza del $0$-path farthest
         reaching sulla diagonale $i$ \Comment{$\Oh(1)$}
\EndFor
\Statex \textit{Passo induttivo sul numero $d$ di errori.}
\For{$d \gets 1$ \textbf{to} $k$}
  \For{$i \gets -n$ \textbf{to} $m$}
    \State dai $(d-1)$-path farthest reaching sulle diagonali $i-1$, $i$, $i+1$
           ricava i tre cammini $R_1,R_2,R_3$ che terminano sulla diagonale $i$
    \State il più lontano fra $R_1,R_2,R_3$ è $\mathrm{f.r.}(d,i)$
           \Comment{$\Oh(1)$}
  \EndFor
\EndFor
\end{algorithmic}
\end{algorithm}

Il primo ciclo è l'\emph{inizializzazione}: determina i $d$-path farthest
reaching per $d=0$. Ogni query di longest common extension costa $\Oh(1)$ grazie
al preprocessing lineare descritto sopra (suffix tree generalizzato di
$P\,\text{\$}\,T\,\text{\texteuro}$ più Harel--Tarjan per la $\lca$), quindi
l'inizializzazione costa in tutto $\Oh(m)$.

Il secondo ciclo esegue $k$ iterazioni: al passo $d$ si costruiscono i $d$-path
farthest reaching per ognuna delle diagonali, prolungando in tre modi diversi i
$(d-1)$-path farthest reaching delle diagonali $i-1$, $i$ e $i+1$ e prendendo il
più lontano dei tre cammini $R_1,R_2,R_3$ così ottenuti.

Le occorrenze si leggono strada facendo: ogni volta che $\mathrm{f.r.}(d,i)$
raggiunge la riga $n$, la colonna $j=n+i$ in cui esso termina è la posizione
finale di un'occorrenza di $P$ in $T$ con al più $d\le k$ differenze. E poiché
il farthest reaching è quello che arriva più in basso sulla diagonale, se non
raggiunge la riga $n$ nessun altro $d$-path di quella diagonale ci arriva:
limitarsi ai farthest reaching non fa perdere occorrenze.

\subsection{Il passo di estensione}

\scan{51}
Il passo centrale del ciclo interno è riassunto dalla figura seguente: dai
$(d-1)$-path farthest reaching sulle tre diagonali $i-1$, $i$, $i+1$ si
costruiscono i tre cammini $R_1$, $R_2$, $R_3$ che terminano sulla diagonale
$i$, e il più lontano dei tre è il $d$-path farthest reaching sulla diagonale
$i$.

\begin{figure}[H]
  \centering
  % Passo induttivo di Landau--Vishkin.  Dai tre f.r. (d-1)-path che terminano
% sulle diagonali i-1, i, i+1 si ricavano i cammini R1 (mossa orizzontale),
% R2 (mossa diagonale con mismatch) e R3 (mossa verticale): tutti e tre
% terminano sulla diagonale i.  Il piu' lontano -- qui R3 -- individua la
% cella di riga v e colonna u (con u-v=i), da cui si scorre gratis lungo la
% diagonale con una query di longest common extension.
%
% Vincolo geometrico rispettato ovunque: le diagonali i-1, i, i+1 partono
% dalle colonne 2.5, 3.5, 4.5 della riga 0 e ogni mossa d'errore e' lunga
% esattamente una cella, cosicche' per tutti i punti d'arrivo si ha
% colonna - riga = 3.5, cioe' stanno davvero sulla diagonale i.
\begin{tikzpicture}[
    scale=0.74,
    font=\footnotesize,
    bordo/.style={semithick},
    guida/.style={thin,black!45},
    frpath/.style={semithick,black!65,rounded corners=1.6pt},
    mossa/.style={thick,-{Stealth[length=4.5pt]}},
    lce/.style={line width=1.2pt},
    nodo/.style={circle,fill,inner sep=0pt,minimum size=3.2pt},
    nodone/.style={circle,fill,inner sep=0pt,minimum size=5.2pt},
    etich/.style={inner sep=1pt,fill=white}]

  % ---- i due bordi (la tabella resta aperta a destra e in basso) ------------
  \draw[bordo] (0,0) -- (10.3,0);     % riga 0
  \draw[bordo] (0,0) -- (0,-6.3);     % colonna 0

  % ---- diagonali di riferimento ---------------------------------------------
  \foreach \c/\lab in {2.5/{i-1}, 3.5/{i}, 4.5/{i+1}} {
    \draw[guida] (\c,0) -- ++(5.2,-5.2);
    \node[anchor=south,inner sep=2pt] at (\c,0) {$\lab$};
  }

  % ---- colonna u e riga v: si incontrano sulla diagonale i ------------------
  \draw[guida] (7.9,0) -- (7.9,-5.1);
  \node[anchor=south,inner sep=2pt] at (7.9,0) {$u$};
  \draw[guida] (0,-4.4) -- (9.7,-4.4);
  \node[anchor=east,inner sep=3pt] at (0,-4.4) {$v$};

  % ---- i tre f.r. (d-1)-path ------------------------------------------------
  % termina sulla diagonale i+1
  \draw[frpath] (0.6,0) -- (0.6,-1.1) -- (1.15,-1.1) -- (1.15,-1.85)
             -- (1.9,-1.85) -- (1.9,-2.45) -- (2.75,-2.45) -- (2.75,-2.85)
             -- (3.8,-2.85) -- (3.8,-3.1) -- (5.2,-3.1) -- (5.2,-3.25)
             -- (6.4,-3.25) -- (6.4,-3.4) -- (7.9,-3.4);
  % termina sulla diagonale i
  \draw[frpath] (1.5,0) -- (1.5,-0.85) -- (2.05,-0.85) -- (2.05,-1.45)
             -- (2.8,-1.45) -- (2.8,-1.9) -- (3.45,-1.9) -- (3.45,-2.35)
             -- (4.6,-2.35) -- (4.6,-2.6) -- (6.3,-2.8);
  % termina sulla diagonale i-1
  \draw[frpath] (2.4,0) -- (2.4,-0.75) -- (2.95,-0.75) -- (2.95,-1.25)
             -- (3.5,-1.25) -- (3.5,-1.75) -- (4.05,-1.75) -- (4.05,-2.2)
             -- (4.7,-2.2);

  \node[nodo] at (4.7,-2.2) {};   % fine del f.r. (d-1)-path su i-1
  \node[nodo] at (6.3,-2.8) {};   % fine del f.r. (d-1)-path su i
  \node[nodo] at (7.9,-3.4) {};   % fine del f.r. (d-1)-path su i+1

  \draw[guida,decorate,decoration={brace,mirror,amplitude=4pt}]
       (0.9,-3.55) -- (4.3,-3.55);
  \node[anchor=north,inner sep=3pt] at (2.6,-3.78) {f.r.\ $(d-1)$-path};

  % ---- le tre mosse d'errore, tutte verso la diagonale i --------------------
  \draw[mossa] (4.7,-2.2) -- (5.7,-2.2);   % orizzontale: da i-1 a i
  \draw[mossa] (6.3,-2.8) -- (7.3,-3.8);   % diagonale con mismatch: resta su i
  \draw[mossa] (7.9,-3.4) -- (7.9,-4.4);   % verticale: da i+1 a i
  \node[etich,anchor=south] at (5.20,-2.16) {$R_1$};
  \node[etich,anchor=east]  at (6.20,-3.05) {$R_2$};
  \node[etich,anchor=west]  at (8.02,-3.90) {$R_3$};

  \node[nodo] at (5.7,-2.2) {};   % arrivo di R1 sulla diagonale i
  \node[nodo] at (7.3,-3.8) {};   % arrivo di R2 sulla diagonale i

  % ---- la cella (v,u) e l'estensione trovata dalla query lce ---------------
  \node[nodone] at (7.9,-4.4) {};          % il piu' lontano dei tre arrivi
  \draw[lce] (7.9,-4.4) -- (9.5,-6.0);
  \node[nodo] at (9.5,-6.0) {};
  \draw[guida] (9.55,-5.95) -- (10.8,-2.9);
  \node[anchor=west,inner sep=2pt] at (10.8,-2.9) {$\lcey(P,T,v+1,u+1)$};

\end{tikzpicture}

  \caption{Costruzione del $d$-path farthest reaching sulla diagonale $i$ a
  partire dai tre $(d-1)$-path farthest reaching sulle diagonali $i-1$, $i$,
  $i+1$: una sola mossa d'errore porta sulla diagonale $i$, poi si scorre gratis
  lungo la diagonale con una query di longest common extension.}
  \label{fig:lv-fr-tre-path}
\end{figure}

Ciascuno dei tre cammini si ottiene prolungando il rispettivo $(d-1)$-path
farthest reaching con \emph{una sola} mossa d'errore, quella che lo porta sulla
diagonale $i$:
\begin{itemize}
  \item dalla diagonale $i-1$ con una mossa \emph{orizzontale} \quad ($R_1$);
  \item dalla diagonale $i$ con una mossa \emph{diagonale con mismatch} \quad
        ($R_2$);
  \item dalla diagonale $i+1$ con una mossa \emph{verticale} \quad ($R_3$).
\end{itemize}
Dopo la mossa d'errore ci si trova nella cella di riga $v$ e colonna $u$ (con
$u-v=i$) e da lì si prosegue lungo la diagonale $i$ finché i caratteri
coincidono. Quello scorrimento è esattamente una query di longest common
extension
\[
  \lcey\bigl(P,\,T,\,v+1,\,u+1\bigr),
\]
che costa $\Oh(1)$ dopo il preprocessing lineare. Confrontare i tre cammini e
scegliere il più lontano richiede a sua volta tempo costante: il corpo del ciclo
interno è quindi $\Oh(1)$.

\section{Complessità}

Il preprocessing --- costruzione del suffix tree generalizzato di
$P\,\text{\$}\,T\,\text{\texteuro}$ e preparazione delle query di $\lca$ ---
costa $\Oh(m+n)$. L'inizializzazione esamina $m+1$ diagonali con una query di
longest common extension ciascuna, quindi $\Oh(m)$. Il ciclo principale esegue
$k$ iterazioni, ognuna delle quali percorre le $n+m+1$ diagonali facendo lavoro
costante su ciascuna, quindi $\Oh\bigl(k\,(n+m)\bigr)$. In totale
\[
  \Oh\bigl(k\,(m+n)\bigr),
\]
che è lineare non appena $k=\Oh(1)$: esattamente la promessa fatta all'inizio
del capitolo.


\part{Algoritmi randomizzati}
\chapter{Algoritmi randomizzati e RandQS}
\label{cha:randqs}

\scan{51}Con questa lezione si apre la seconda parte del corso, dedicata agli
\emph{algoritmi randomizzati}. Riferimenti bibliografici:
\begin{itemize}
  \item \emph{Randomized Algorithms}, di Motwani e Raghavan --- interessante,
        con molti conti da fare;
  \item \emph{An introduction to randomized algorithms}, di Richard Karp.
\end{itemize}

\section{Che cos'è un algoritmo randomizzato}
\label{sec:def-randomizzato}

\begin{definizione}[Algoritmo randomizzato]
\label{def:algoritmo-randomizzato}
Un \emph{algoritmo randomizzato} è un algoritmo che riceve, oltre al proprio
input, un \emph{flusso} (\emph{stream}) di \emph{bit casuali} che utilizza per
compiere \emph{scelte casuali}.
\end{definizione}

\begin{osservazione}
Il flusso non deve essere casuale in senso stretto: basta che sia una sequenza
che l'utente non è in grado di distinguere da una sequenza random.
\end{osservazione}

\begin{figure}[H]
  \centering
  % rand-schema-bit -- pag. 51
% Un algoritmo randomizzato riceve, accanto al proprio input, un flusso di bit
% casuali; poiche' le scelte compiute dipendono da quei bit, la complessita'
% dell'esecuzione (e in generale anche il risultato) e' una variabile aleatoria.
\begin{tikzpicture}[
    x=1cm, y=1cm, >=Latex,
    filo/.style={draw=black!80,line width=0.7pt},
    et/.style={font=\small}
  ]

  % ---- la scatola dell'algoritmo -------------------------------------------
  \node[draw=black!80,line width=1pt,minimum width=2.2cm,minimum height=1.5cm,
        inner sep=2pt,font=\small] (alg) at (0,0) {algoritmo};

  % ---- il flusso di bit casuali --------------------------------------------
  \node[et,align=center,anchor=south] (bit) at (-3.55,1.30)
        {\texttt{010110}\\[-2pt] \emph{random}};
  \draw[filo,->] (bit.south) to[out=-90,in=180] (-1.15,0.35);

  % ---- l'input ordinario ---------------------------------------------------
  \node[et,anchor=east] (in) at (-3.00,-0.40) {input};
  \draw[filo,->] (in.east) -- (-1.15,-0.40);

  % ---- l'uscita: la complessita' e' una variabile aleatoria -----------------
  \draw[filo,->] (1.15,0.20) to[out=0,in=90] (2.90,-1.40);
  \node[et,anchor=west] at (3.05,-0.60) {complessit\`a};

\end{tikzpicture}

  \caption{Un algoritmo randomizzato riceve in ingresso, accanto all'input, un
  flusso di bit casuali; di conseguenza anche la complessità dell'esecuzione
  diventa una variabile aleatoria.}
  \label{fig:rand-schema-bit}
\end{figure}

Poiché le scelte compiute dipendono dai bit casuali, sia il risultato prodotto
sia la complessità dell'esecuzione diventano variabili aleatorie. Fissando
l'una o l'altra si ottengono le due famiglie classiche:
\begin{itemize}
  \item \MonteCarlo: risultato potenzialmente sbagliato, complessità fissa;
  \item \LasVegas: complessità variabile, risultato sempre corretto.
\end{itemize}

Il primo esempio che incontreremo, \RandQS, appartiene alla seconda famiglia:
l'output è sempre l'insieme ordinato, mentre il numero di operazioni eseguite
dipende dalle scelte casuali ed è quindi una variabile aleatoria. Il taglio
minimo, di cui ci occuperemo subito dopo, è invece l'esempio prototipico della
prima.

\section{Randomized QuickSort}
\label{sec:randqs}

\scan{52}L'idea di \RandQS è quella di QuickSort, con un'unica differenza: il pivot non
viene scelto in modo deterministico (il primo elemento, l'ultimo, la mediana di
tre, \dots), ma viene \emph{estratto uniformemente a caso} dall'insieme da
ordinare.

\begin{algorithm}[H]
\caption{\RandQS}
\begin{algorithmic}[1]
\Require un insieme di numeri $S$ (elementi distinti)
\Ensure gli elementi di $S$ in ordine crescente
\State scegli un elemento $x \in S$ \textbf{uniformemente a caso};
\State $\textsc{Partition}(S,x)$: genera $S_1=\set{y\in S \mid y<x}$ e
       $S_2=\set{y\in S \mid y>x}$;
\State ordina ricorsivamente con \RandQS{}($S_1$) e \RandQS{}($S_2$); restituisci
       $S_1$ ordinato, seguito da $x$, seguito da $S_2$ ordinato.
\end{algorithmic}
\end{algorithm}

\begin{osservazione}
QuickSort ha complessità $\Oh(n\log n)$ nel caso medio \emph{se la distribuzione
degli input è uniforme}: si tratta di un'ipotesi (forte) sui possibili input.
\RandQS, invece, \textbf{non} fa alcuna ipotesi sulla distribuzione dell'input:
la casualità non sta più nei dati, ma nell'algoritmo.
\end{osservazione}

\subsection{Il caso ideale}
\label{sec:randqs-caso-ideale}

Nel caso ideale il pivot cade «in mezzo» e la partizione produce due
sottoproblemi di $n/2$ elementi ciascuno.

\begin{figure}[H]
  \centering
  % randqs-partizione -- pag. 52
% Caso ideale: il pivot estratto uniformemente a caso spezza S in due
% sottoinsiemi di n/2 elementi; le due frecce sono le chiamate ricorsive
% RandQS(S_1) e RandQS(S_2). E' la figura che giustifica T(n)=O(n)+2T(n/2).
\begin{tikzpicture}[
    x=1cm, y=1cm, >=Latex,
    filo/.style={draw=black!80,line width=0.8pt},
    et/.style={font=\footnotesize}
  ]

  % ---- la cella del pivot ---------------------------------------------------
  \draw[filo,line width=1pt] (-0.28,1.92) rectangle (0.28,2.48);
  \node[et,anchor=south] at (0,2.58) {pivot};

  % ---- le due chiamate ricorsive -------------------------------------------
  \draw[filo,->] (-0.28,1.92) to[out=-120,in=60]  (-2.10,0.55);
  \draw[filo,->] ( 0.28,1.92) to[out=-60, in=120] ( 2.10,0.55);

  % ---- i due sottoinsiemi, di n/2 elementi ciascuno -------------------------
  \foreach \c in {-2.10, 2.10} {
    \draw[filo] (\c-1.5,-0.22) rectangle (\c+1.5,0.22);
    \draw[filo] (\c-1.5,-0.36) -- (\c-1.5,0.36);
    \draw[filo] (\c+1.5,-0.36) -- (\c+1.5,0.36);
    \node at (\c,0) {$n/2$};
  }

\end{tikzpicture}

  \caption{Partizione ideale: il pivot spezza l'insieme in due metà di $n/2$
  elementi.}
  \label{fig:randqs-partizione}
\end{figure}

In questo caso l'albero delle chiamate ricorsive è bilanciato: ha altezza
$h=\log_2 n$ e ad ogni livello il costo complessivo delle partizioni è
$\Oh(n)$, perché i sottoproblemi di uno stesso livello ripartiscono fra loro gli
$n$ elementi di partenza.

\begin{figure}[H]
  \centering
  % randqs-albero-chiamate -- pag. 52
% Albero delle chiamate ricorsive nel caso bilanciato: e' un albero binario
% completo di altezza h = log_2 n; su ogni livello i sottoproblemi si
% ripartiscono gli n elementi, quindi il costo per livello e' O(n).
\begin{tikzpicture}[
    x=1cm, y=1cm, >=Latex,
    filo/.style={draw=black!80,line width=0.7pt},
    nodo/.style={circle,draw=black!80,fill=white,line width=0.7pt,
                 inner sep=0pt,minimum size=4.5pt},
    radice/.style={circle,draw=black!80,fill=black!80,inner sep=0pt,
                   minimum size=5.5pt},
    onda/.style={filo,->,decorate,
                 decoration={snake,amplitude=0.5mm,segment length=3mm,
                             post length=2mm,pre length=1mm}}
  ]

  % ---- i nodi ---------------------------------------------------------------
  \node[radice] (r)  at ( 0.00,2.40) {};
  \node[nodo]   (a)  at (-1.15,1.20) {};
  \node[nodo]   (b)  at ( 1.15,1.20) {};
  \node[nodo]   (a1) at (-1.75,0.00) {};
  \node[nodo]   (a2) at (-0.55,0.00) {};
  \node[nodo]   (b1) at ( 0.55,0.00) {};
  \node[nodo]   (b2) at ( 1.75,0.00) {};

  % ---- gli archi ------------------------------------------------------------
  \draw[filo] (r) -- (a);   \draw[filo] (r) -- (b);
  \draw[filo] (a) -- (a1);  \draw[filo] (a) -- (a2);
  \draw[filo] (b) -- (b1);  \draw[filo] (b) -- (b2);

  % ---- l'altezza dell'albero ------------------------------------------------
  \draw[filo,decorate,decoration={brace,amplitude=6pt}]
       (-2.60,-0.18) -- (-2.60,2.58);
  \node[anchor=east] at (-2.85,1.20) {$h$};

  % ---- il costo di ciascun livello ------------------------------------------
  \draw[onda] (0.25,2.40) -- (2.80,2.40);
  \draw[onda] (1.40,1.20) -- (2.80,1.20);
  \draw[onda] (1.95,0.00) -- (2.80,0.00);
  \node[anchor=west] at (2.95,2.40) {$\Oh(n)$};
  \node[anchor=west] at (2.95,1.20) {$\Oh(n)$};
  \node[anchor=west] at (2.95,0.00) {$\Oh(n)$};

\end{tikzpicture}

  \caption{Albero delle chiamate ricorsive nel caso bilanciato: $h=\log_2 n$
  livelli, ciascuno di costo complessivo $\Oh(n)$.}
  \label{fig:randqs-albero-chiamate}
\end{figure}

\[
  T(n) \;=\; \Oh(n) \;+\; 2\,T\!\left(\frac{n}{2}\right)
  \qquad\Longrightarrow\qquad
  T(n) \in \Oh(n\log n).
\]

Nulla garantisce però che il pivot estratto a caso spezzi davvero l'insieme in
due metà: la ricorrenza qui sopra vale solo nel caso bilanciato, mentre in
generale una partizione in $p$ e $q$ elementi dà $T(n)=T(p)+T(q)+\Th(n)$, che
nel caso pessimo (pivot sempre estremo) degenera in $\Th(n^2)$. Quello che
possiamo controllare non è dunque il costo di una singola esecuzione, ma il suo
\emph{valore atteso}.

\section{Analisi di \texorpdfstring{\RandQS}{RandQS} nel caso medio}
\label{sec:randqs-analisi}

Qui «caso medio» va inteso come \emph{valore atteso} di una variabile aleatoria
--- quella che rappresenta la complessità --- dove la media è presa sulle scelte
casuali dell'algoritmo, a input fissato, e non sugli input. Non si tratta quindi
dell'analisi di QuickSort nel caso medio: non c'è alcuna ipotesi sulla
distribuzione degli input.

Il costo di \RandQS è dominato dal numero di \emph{confronti} eseguiti dalle
chiamate a \textsc{Partition}: è quello che conteremo.

\begin{teorema}[Numero atteso di confronti di \RandQS]
\label{thm:randqs-confronti}
Il numero atteso di confronti eseguiti da \RandQS su un insieme di $n$ elementi
distinti è al più $2\,n\,H_n$, dove $H_n=\sum_{k=1}^{n}1/k$ è l'$n$-esimo numero
armonico. Poiché $H_n \in \Th(\log n)$, il numero atteso di confronti è
$\Oh(n\log n)$.
\end{teorema}

Indichiamo con $x_i$ l'elemento di rango $i$, cioè l'$i$-esimo elemento di $S$
nell'ordine crescente, e introduciamo per ogni coppia $1\le i<j\le n$ la
variabile aleatoria \emph{indicatrice}
\[
  X_{ij} \;=\;
  \begin{cases}
    1 & \text{se } x_i \text{ \emph{è confrontato} con } x_j,\\[2pt]
    0 & \text{altrimenti,}
  \end{cases}
  \qquad\text{ovvero}\qquad
  X_{ij} = \indic{x_i \text{ è confrontato con } x_j}.
\]

Il numero totale di confronti è allora $\sum_{i=1}^{n}\sum_{j>i} X_{ij}$, e la
quantità da stimare è
\[
  \E\!\left[\; \sum_{i=1}^{n}\sum_{j>i} X_{ij} \;\right].
\]

\subsection{Linearità del valore atteso}
\label{sec:linearita}

\scan{53}Per linearità del valore atteso --- che non richiede alcuna ipotesi di
indipendenza fra le $X_{ij}$ --- vale
\[
  \E\!\left[\;\sum_{i=1}^{n}\sum_{j>i} X_{ij}\right]
  \;=\;
  \sum_{i=1}^{n}\sum_{j>i} \E\!\left[X_{ij}\right].
\]

Poiché $X_{ij}$ è una variabile indicatrice, posto
$p_{ij} = \Pr\!\left[X_{ij}=1\right]$ si ha
\[
  \E\!\left[X_{ij}\right] \;=\; p_{ij}\cdot 1 \;+\; (1-p_{ij})\cdot 0 \;=\; p_{ij},
\]
dove $p_{ij}$ è dunque la probabilità che $x_i$ e $x_j$ vengano confrontati.

\subsection{Calcolo di \texorpdfstring{$p_{ij}$}{p\_ij}}
\label{sec:pij}

Occorre calcolare la probabilità che $x_i$ e $x_j$ siano coinvolti in un
confronto.

\begin{figure}[H]
  \centering
  % randqs-separazione -- pag. 53
% La retta tratteggiata e' la sequenza ordinata x_1 < x_2 < ... < x_n. Se il
% pivot estratto cade strettamente fra x_i e x_j, allora x_i finisce in S_1 e
% x_j in S_2: i due non verranno mai piu' confrontati (evento complementare a
% quello che definisce X_{ij}=1).
\begin{tikzpicture}[
    x=1cm, y=1cm, >=Latex,
    filo/.style={draw=black!80,line width=0.7pt},
    zig/.style={filo,->,decorate,
                decoration={snake,amplitude=0.9mm,segment length=3.6mm,
                            post length=2.5mm}}
  ]

  % ---- la sequenza ordinata degli elementi ---------------------------------
  \draw[filo,dash pattern=on 4pt off 4pt] (-5.4,0) -- (5.4,0);

  % ---- i due elementi x_i e x_j (il tratteggio si interrompe attorno) -------
  \node[fill=white,inner sep=2pt] at (-2.2,0) {$x_i$};
  \node[fill=white,inner sep=2pt] at ( 2.2,0) {$x_j$};

  % ---- il pivot estratto, strettamente compreso fra i due -------------------
  \fill[white] (0,0) circle (5pt);
  \draw[filo] (0,0) circle (2pt);

  % ---- la scelta casuale del pivot ------------------------------------------
  \draw[zig] (0,-1.75) -- (0,-0.10);
  \node[font=\footnotesize,anchor=north] at (0,-1.90) {pivot scelto};

\end{tikzpicture}

  \caption{Se il pivot estratto cade strettamente fra $x_i$ e $x_j$, i due
  elementi finiscono in sottoinsiemi diversi e non verranno mai confrontati.}
  \label{fig:randqs-separazione}
\end{figure}

\begin{osservazione}
$x_i$ e $x_j$ \emph{non} vengono confrontati se il pivot scelto è in
$\set{x_{i+1},\dots,x_{j-1}}$: un pivot interno all'intervallo separa i due
elementi, mandandone uno in $S_1$ e l'altro in $S_2$, e da quel momento in poi
non potranno più essere confrontati.
\end{osservazione}

Ogni elemento $d \in \set{x_i,\dots,x_j}$ ha la stessa probabilità di essere
scelto come pivot; quindi la probabilità di scegliere $x_i$ oppure $x_j$ ---
cioè $p_{ij}$ --- soddisfa
\[
  p_{ij} \;\le\; \frac{2}{\,j-i+1\,}
\]
(è il caso peggiore, quello in cui mi ritrovo con una partizione contenente
esclusivamente $\set{x_i,\dots,x_j}$; generalmente mi ritrovo in partizioni più
grandi, e la probabilità di pescare proprio uno dei due estremi è minore).

\begin{osservazione}
In realtà vale l'uguaglianza $p_{ij} = \dfrac{2}{j-i+1}$. Finché nessuno fra
$x_i, x_{i+1},\dots,x_j$ è stato scelto come pivot, questi $j-i+1$ elementi
stanno tutti nella stessa partizione (un pivot esterno all'intervallo non può
separarli); il primo di essi a diventare pivot è uniformemente distribuito fra i
$j-i+1$ candidati, e $x_i$ e $x_j$ vengono confrontati esattamente quando tale
primo elemento è uno dei due estremi. Per la maggiorazione che segue basta
comunque la disuguaglianza.
\end{osservazione}

Sommando su tutte le coppie, e osservando che per $i$ fissato il denominatore
$j-i+1$ percorre i valori $2,3,\dots,n-i+1$:
\[
  \sum_{i=1}^{n}\sum_{j>i} p_{ij}
  \;\le\; \sum_{i=1}^{n}\sum_{j>i} \frac{2}{\,j-i+1\,}
  \;\le\; 2\sum_{i=1}^{n} \underbrace{\sum_{k=1}^{n}\frac{1}{k}}_{H_n}
  \;=\; 2\sum_{i=1}^{n} H_n
  \;=\; 2\,n\,H_n .
\]

\subsection{Il numero armonico}
\label{sec:numero-armonico}

\[
  H_n \;=\; \sum_{k=1}^{n}\frac1k
  \;=\; 1+\frac12+\frac13+\frac14+\frac15+\frac16+\frac17+\cdots
\]
Raggruppando i termini a blocchi di lunghezza $1,2,4,8,\dots$ --- i blocchi sono
$\lfloor\log_2 n\rfloor+1$ e ogni blocco \emph{completo} ha somma compresa fra
$1/2$ e $1$, mentre l'ultimo può essere parziale --- si ottengono le due stime,
la seconda con i blocchi traslati di un posto:
\begin{align*}
  H_n &\;=\; 1
    + \underbrace{\left(\frac12+\frac13\right)}_{\le\;\frac12+\frac12\;=\;1}
    + \underbrace{\left(\frac14+\frac15+\frac16+\frac17\right)}_{\le\;4\cdot\frac14\;=\;1}
    + \cdots \;\le\; 1+\log_2 n, \\[6pt]
  H_n &\;=\; 1+\frac12
    + \underbrace{\left(\frac13+\frac14\right)}_{\ge\;\frac14+\frac14\;=\;\frac12}
    + \underbrace{\left(\frac15+\frac16+\frac17+\frac18\right)}_{\ge\;4\cdot\frac18\;=\;\frac12}
    + \cdots \;\ge\; \frac12\log_2 n .
\end{align*}

Dunque $H_n \in \Th(\log n)$ e, mettendo insieme i pezzi,
\[
  \E\!\left[\;\sum_{i=1}^{n}\sum_{j>i} X_{ij}\;\right]
  \;=\; \sum_{i=1}^{n}\sum_{j>i} p_{ij}
  \;\le\; 2\,n\,H_n \;\in\; \Th(n\log n),
\]
cioè il numero atteso di confronti di \RandQS è $\Oh(n\log n)$, come annunciato
dal Teorema~\ref{thm:randqs-confronti}.

\begin{osservazione}
Ogni scelta del pivot viene pagata in bit casuali: per estrarre uniformemente un
elemento da un insieme di $n$ elementi servono $\log_2 n$ bit letti dalla
sorgente,
\[
  \text{scelta del pivot}\;\longleftarrow\;
  \underbrace{0\,1\,0\,1\,0\;\cdots}_{\log_2 n\ \text{bit}}
\]
\end{osservazione}

\subsection{Elementi ripetuti}
\label{sec:randqs-ripetuti}

\scan{54}L'analisi svolta finora suppone che gli elementi da ordinare siano \emph{a due a
due distinti}: è l'ipotesi che rende ben definiti i ranghi $x_1<\dots<x_n$ e che
garantisce $S_1\cup\set{x}\cup S_2 = S$.

\begin{figure}[H]
  \centering
  % randqs-pivot-array -- un passo di RandQS: scelta uniforme del pivot nel segmento.
\begin{tikzpicture}[x=1cm,y=1cm,line join=round,line cap=round]
  \footnotesize
  % --- il segmento corrente: rettangolo allungato, prima e ultima cella marcate
  \draw[thick] (0,0) rectangle (8.4,0.55);
  \draw[thick] (0.6,0) -- (0.6,0.55);
  \draw[thick] (7.8,0) -- (7.8,0.55);

  % --- posizione del pivot (~45% del segmento)
  \coordinate (P) at (3.8,0.61);

  % --- scelta casuale: arco dall'inizio del segmento alla cella del pivot
  \draw[thick,-{Stealth[length=2.4mm,width=2mm]}]
    (0.35,0.75) .. controls (1.1,1.85) and (3.0,1.9) .. (P);
  \node[gray,anchor=south] at (2.1,1.62) {scelta uniforme a caso};

  % --- freccia ondulata dal basso che indica la cella del pivot
  \draw[thick,-{Stealth[length=2.2mm,width=1.8mm]}]
    plot[domain=0:1,samples=90,smooth,variable=\t]
      ({3.8-0.30*(1-\t)+0.10*(1-\t)*sin(\t*1080)},{-1.15+1.03*\t});
  \node[anchor=north east] at (3.62,-1.12) {pivot};
\end{tikzpicture}

  \caption{Un passo di \RandQS: il pivot è scelto uniformemente a caso fra gli
  elementi del segmento corrente, che si assumono a due a due distinti.}
  \label{fig:randqs-pivot-array}
\end{figure}

Se l'input contiene ripetizioni basta sostituire la sequenza $a_1,\dots,a_n$ con
la sequenza di coppie
\[
  \tuple{a_1,1},\ \tuple{a_2,2},\ \dots,\ \tuple{a_n,n},
\]
ordinate lessicograficamente: la seconda componente rende le coppie a due a due
distinte e non altera l'ordine relativo dei valori $a_i$. È una soluzione pulita
e semplice per ripulire l'analisi in presenza di copie ripetute.

\section{Che cosa viene dopo}
\label{sec:randqs-seguito}

\scan{54}Il prossimo problema è quello del \emph{taglio minimo} (\mincut): dato un grafo,
individuare un insieme di archi la cui rimozione lo sconnette e che abbia
cardinalità minima.

\begin{figure}[H]
  \centering
  % mincut-anticipazione -- un taglio di G: la spezzata separa i vertici in due parti,
% il taglio e' l'insieme degli archi che la attraversano.
\begin{tikzpicture}[x=1cm,y=1cm,line join=round,line cap=round,scale=0.92]
  \footnotesize
  % --- sagoma del grafo G
  \draw[thick] plot[smooth cycle,tension=0.7] coordinates {
    (4.44,0.00) (3.47,0.96) (2.73,1.98) (0.70,2.25) (-1.22,2.42) (-2.58,1.62)
    (-4.12,0.94) (-3.85,-0.15) (-3.85,-1.26) (-2.16,-1.81) (-0.62,-2.49)
    (0.70,-2.05) (2.48,-2.08) (3.34,-1.09)};
  \node[anchor=north east] at (-3.80,-1.30) {$G$};

  % --- frontiera del taglio: spezzata a zig-zag dal bordo superiore a quello inferiore
  \draw[thick]
    (0.70,2.25) -- (0.85,1.89) -- (0.55,1.53) -- (0.85,1.17) -- (0.55,0.81) --
    (0.85,0.45) -- (0.55,0.09) -- (0.85,-0.27) -- (0.55,-0.63) -- (0.85,-0.99) --
    (0.55,-1.35) -- (0.85,-1.71) -- (0.70,-2.05);

  % --- i quattro archi che attraversano il taglio
  \draw[thick] (0.10,1.62) -- (1.26,1.84);
  \draw[thick] (0.06,0.92) -- (1.22,0.68);
  \draw[thick] (0.12,-0.42) -- (1.28,-0.22);
  \draw[thick] (0.04,-1.28) -- (1.20,-1.10);

  % --- i due vertici separati dal taglio
  \fill (-2.00,0.35) circle (2.4pt) node[anchor=north,yshift=-2pt] {$x$};
  \fill ( 2.50,0.15) circle (2.4pt) node[anchor=north,yshift=-2pt] {$y$};

  % --- etichetta della frontiera
  \node[anchor=west] at (1.02,-1.56) {taglio};
\end{tikzpicture}

  \caption{Un taglio di $G$: la linea spezzata separa i vertici in due parti (da
  un lato $x$, dall'altro $y$); il taglio è l'insieme degli archi che la
  attraversano.}
  \label{fig:mincut-anticipazione}
\end{figure}

\begin{itemize}
  \item esistono algoritmi deterministici, ma estremamente complicati;
  \item c'è un algoritmo randomizzato di prestazioni comparabili e molto
        semplice.
\end{itemize}

Dopo ancora: il \emph{metodo probabilistico}.

\chapter{Min-Cut di Karger}

\scan{54}
Il problema del \emph{taglio minimo} (\mincut) ammette un algoritmo randomizzato molto
semplice, dovuto a Karger. È un algoritmo di tipo \MonteCarlo, e lo strumento
probabilistico su cui se ne regge l'analisi è la \emph{probabilità condizionata}.

L'input è un grafo $G=(V,E)$ \emph{non orientato} e \emph{connesso}; per tutto il capitolo
poniamo $n=\abs{V}$. La richiesta di connessione non è restrittiva: su un grafo già
sconnesso il taglio minimo è vuoto e il problema è banale.

\section{Il problema}

\scan{55}

\begin{definizione}[Taglio]\label{def:taglio}
Un \emph{taglio} (\emph{cut}) $C$ di un grafo $G=(V,E)$ è un insieme di archi tale che
\begin{itemize}
  \item $C\subseteq E$;
  \item $(V,\,E\setminus C)$ è \emph{sconnesso}.
\end{itemize}
\end{definizione}

\begin{definizione}[Taglio minimo]\label{def:min-cut}
Un \emph{taglio minimo} (\emph{min-cut}) è un taglio $C$ tale che $\abs{C}$ è minima.
\end{definizione}

Si noti che si parla di taglio in senso \emph{globale} — non ci sono due terminali $s,t$
fissati — e \emph{non pesato}: il costo di un taglio è il numero dei suoi archi.

\section{Richiami: probabilità condizionata}

La probabilità condizionata è lo strumento con cui si controlla la probabilità di errore
relativa a una \emph{sequenza} di scelte, come sarà quella dell'algoritmo che stiamo per
descrivere.

\begin{itemize}
  \item Se $\mathcal{E}_1$ ed $\mathcal{E}_2$ sono \emph{indipendenti}:
  \[
    \Pr[\mathcal{E}_1\cap\mathcal{E}_2]\;=\;\Pr[\mathcal{E}_1]\cdot\Pr[\mathcal{E}_2] .
  \]
  \item Se non sono indipendenti (caso generale):
  \[
    \Pr[\mathcal{E}_1\cap\mathcal{E}_2]\;=\;\Pr[\mathcal{E}_1\mid\mathcal{E}_2]\cdot\Pr[\mathcal{E}_2],
  \]
  dove $\Pr[\mathcal{E}_1\mid\mathcal{E}_2]$ è la \emph{probabilità condizionata}: è ciò che
  ci consente di misurare la probabilità di due eventi \emph{consecutivi}.
\end{itemize}

\begin{figure}[H]
  \centering
  % prob-condizionata-venn -- due diagrammi di Venn affiancati (quaderno, pag. 55).
% A sinistra i due eventi sono disegnati disgiunti (rappresentazione informale
% dell'indipendenza, discussa nell'osservazione che segue la figura);
% a destra il caso generale, con l'intersezione non vuota evidenziata.
\begin{tikzpicture}[
    font=\footnotesize,
    cerchio/.style={draw,line width=0.7pt},
    didascalia/.style={font=\footnotesize,text=gray}]

  % ---- pannello sinistro: eventi che "non si influenzano" -------------------
  \begin{scope}[shift={(0,0)}]
    \draw[cerchio] (0,0)   circle[radius=0.55];
    \draw[cerchio] (1.6,0) circle[radius=0.55];
    \node[anchor=south east] at (-0.28,0.36) {$\mathcal{E}_1$};
    \node[anchor=south west] at ( 1.88,0.36) {$\mathcal{E}_2$};
    \node[didascalia] at (0.8,-1.15) {eventi indipendenti};
  \end{scope}

  % ---- pannello destro: caso generale --------------------------------------
  \begin{scope}[shift={(4.7,0)}]
    % lente d'intersezione, evidenziata
    \begin{scope}
      \clip (0,0) circle[radius=0.60];
      \fill[black!18] (0.78,0) circle[radius=0.60];
    \end{scope}
    \draw[cerchio] (0,0)    circle[radius=0.60];
    \draw[cerchio] (0.78,0) circle[radius=0.60];
    \node[anchor=east] at (-0.70,0.06) {$\mathcal{E}_1$};
    \node[anchor=west] at ( 1.48,0.06) {$\mathcal{E}_2$};
    \node[didascalia] at (0.39,-1.15) {caso generale};
  \end{scope}

\end{tikzpicture}

  \caption{A sinistra due eventi «che non si influenzano»; a destra il caso generale, con
  l'intersezione $\mathcal{E}_1\cap\mathcal{E}_2$ evidenziata.}
  \label{fig:prob-condizionata-venn}
\end{figure}

\begin{osservazione}
Nel disegno di sinistra i due eventi indipendenti sono rappresentati come regioni
\emph{disgiunte}: è solo un modo per suggerire che «non si influenzano a vicenda».
Attenzione però che l'indipendenza non ha una controparte insiemistica: due eventi disgiunti
di probabilità non nulla non sono mai indipendenti, perché
$\Pr[\mathcal{E}_1\cap\mathcal{E}_2]=0\neq\Pr[\mathcal{E}_1]\Pr[\mathcal{E}_2]$.
\end{osservazione}

Se invece di due eventi ne vogliamo misurare $m$, si itera la formula precedente e si
ottiene la \emph{regola della catena} (\emph{chain rule}):
\begin{align*}
  \Pr\!\left[\bigcap_{i=1}^{m}\mathcal{E}_i\right]
  ={}& \Pr[\mathcal{E}_1]\cdot\Pr[\mathcal{E}_2\mid\mathcal{E}_1]\cdot
       \Pr[\mathcal{E}_3\mid\mathcal{E}_1\cap\mathcal{E}_2]\cdot{}\\
     & \cdots\cdot
       \Pr\!\left[\mathcal{E}_m\;\middle|\;\bigcap_{i=1}^{m-1}\mathcal{E}_i\right].
\end{align*}
È l'unica formula probabilistica che serve per tutta l'analisi che segue: le scelte
successive dell'algoritmo non sono indipendenti, quindi la regola del prodotto non basta.

\section{L'idea: contrazione di un arco}

L'algoritmo introduce le «scelte» casuali in un modo elementare: si sceglie a caso un arco
del grafo e lo si \emph{contrae}, cioè si fondono i suoi due estremi $u$ e $v$ in un unico
nodo $uv$, che eredita tutti gli archi incidenti a $u$ o a $v$. Servono due accorgimenti:

\begin{itemize}
  \item gli archi paralleli fra due nodi distinti si mantengono tutti: si lavora perciò su
        \emph{multigrafi}, cioè su grafi in cui fra due nodi possono esserci più archi;
  \item se la contrazione forma un \emph{self-loop}, lo si ignora (viene eliminato).
\end{itemize}

\begin{figure}[H]
  \centering
  % multigrafo-archi-paralleli -- schemetto isolato (quaderno, pag. 55).
% Due nodi uniti da tre archi paralleli: la molteplicita' degli archi va
% conservata, percio' la contrazione lavora su multigrafi.
\begin{tikzpicture}[
    font=\footnotesize,
    nodo/.style={circle,draw,line width=0.7pt,fill=white,
                 inner sep=0pt,minimum size=4.5pt},
    arco/.style={line width=0.7pt}]

  \node[nodo] (u) at (0,0)   {};
  \node[nodo] (v) at (2.8,0) {};

  \draw[arco] (u) to[bend left=32]  (v);
  \draw[arco] (u) --                (v);
  \draw[arco] (u) to[bend right=32] (v);

  \node[anchor=east] at ([xshift=-2pt]u.west) {$u$};
  \node[anchor=west] at ([xshift= 2pt]v.east) {$v$};

\end{tikzpicture}

  \caption{Multigrafo: fra due nodi possono esserci più archi.}
  \label{fig:multigrafo-archi-paralleli}
\end{figure}

\begin{figure}[H]
  \centering
  % karger-contrazione -- un passo di contrazione (quaderno, pag. 55).
% Sinistra: multigrafo G con l'arco scelto (u,v) evidenziato.
% Destra:  G/(u,v): u e v fondono nel super-nodo uv; gli archi u--z e v--z
%          sopravvivono entrambi (archi paralleli), l'arco contratto diventa
%          un self-loop e viene eliminato (disegnato tratteggiato).
\begin{tikzpicture}[
    font=\footnotesize,
    nodo/.style={circle,draw,line width=0.7pt,fill=white,
                 inner sep=0pt,minimum size=4.5pt},
    arco/.style={line width=0.7pt},
    scelto/.style={line width=1.4pt},
    nota/.style={align=left,font=\footnotesize},
    glossa/.style={align=center,font=\footnotesize}]

  %% ======================= grafo di partenza G ============================
  \begin{scope}[shift={(0,0)}]
    \node[nodo] (a) at (0.0, 2.6) {};
    \node[nodo] (b) at (1.5, 3.4) {};
    \node[nodo] (v) at (3.4, 2.2) {};
    \node[nodo] (u) at (0.1, 0.6) {};
    \node[nodo] (z) at (3.2,-0.5) {};

    % archi paralleli a--b (molteplicita' 2)
    \draw[arco] (a) to[bend left=28]  (b);
    \draw[arco] (a) to[bend right=28] (b);
    % resto del multigrafo
    \draw[arco] (b) -- (v);
    \draw[arco] (a) -- (u);
    \draw[arco] (v) -- (z);
    \draw[arco] (u) -- (z);
    % arco scelto
    \draw[scelto] (u) -- (v)
      node[sloped,above=2.5pt,pos=0.44]{arco scelto};

    \node[anchor=north east] at ([shift={(-0.02,-0.04)}]u.south west) {$u$};
    \node[anchor=south west] at ([shift={( 0.04, 0.02)}]v.north east) {$v$};
    \node[anchor=north]      at ([shift={( 0.00,-0.08)}]z.south)      {$z$};
  \end{scope}

  %% =========================== freccia ondulata ===========================
  \draw[line width=0.7pt,-{Stealth[length=5pt,width=4pt]}]
        plot[domain=0:1.6,samples=90,smooth]
        ({3.78+\x},{1.20+0.075*sin(\x*900)});
  \node[glossa] at (4.58,0.62) {scelgo un arco\\e lo contraggo};

  %% ================= grafo contratto  G/(u,v) =============================
  \begin{scope}[shift={(5.9,0)}]
    \node[nodo] (aa) at (0.0, 2.6) {};
    \node[nodo] (bb) at (1.2, 3.6) {};
    \node[nodo] (uv) at (1.6, 1.3) {};
    \node[nodo] (zz) at (2.9,-0.4) {};

    % la lente a--b sopravvive invariata
    \draw[arco] (aa) to[bend left=28]  (bb);
    \draw[arco] (aa) to[bend right=28] (bb);
    % eredi di b--v e di a--u
    \draw[arco] (bb) -- (uv);
    \draw[arco] (aa) -- (uv);
    % eredi di v--z e u--z: due archi paralleli
    \draw[arco] (uv) to[bend left=32]  (zz);
    \draw[arco] (uv) to[bend right=32] (zz);
    % self-loop generato dalla contrazione: eliminato
    \draw[arco,dash pattern=on 1.6pt off 1.6pt]
          (uv) .. controls +(0.95,0.75) and +(0.95,-0.75) .. (uv);

    \node[anchor=east] at ([shift={(-0.22,0.06)}]uv.west) {$uv$};
    \node[anchor=north west] at ([shift={( 0.04,-0.04)}]zz.south east) {$z$};

    \draw[line width=0.5pt] (2.30,1.56) -- (2.82,1.92);
    \node[nota,anchor=west] at (2.88,2.00)
         {se formo un\\self-loop,\\lo ignoro};
  \end{scope}

\end{tikzpicture}

  \caption{Un passo dell'algoritmo: si sceglie a caso l'arco $(u,v)$ e lo si contrae. Gli
  archi $u$--$z$ e $v$--$z$ diventano due archi paralleli fra $uv$ e $z$; l'arco contratto
  diventa un self-loop e viene eliminato.}
  \label{fig:karger-contrazione}
\end{figure}

Mantenere gli archi paralleli è essenziale: la molteplicità del fascio che unisce due
super-nodi conta esattamente quanti archi \emph{originali} di $G$ hanno un estremo in
ciascuno dei due insiemi di vertici che quei super-nodi rappresentano.

\scan{56}
Ripetendo l'operazione, dopo $n-2$ contrazioni si arriva a questa situazione:

\begin{figure}[H]
  \centering
  % karger-due-nodi-taglio -- situazione dopo n-2 contrazioni (quaderno, pag. 56).
% Restano due super-nodi a e b uniti da quattro archi paralleli; la linea
% ondulata verticale li attraversa tutti: quel fascio di archi e' un taglio
% del grafo di partenza.
\begin{tikzpicture}[
    font=\footnotesize,
    nodo/.style={circle,draw,line width=0.8pt,fill=white,
                 inner sep=0pt,minimum size=5pt},
    nodopieno/.style={circle,draw,line width=0.8pt,fill=black,
                 inner sep=0pt,minimum size=5pt},
    arco/.style={line width=0.8pt},
    taglio/.style={line width=0.7pt,black!70}]

  % ---- i due super-nodi residui --------------------------------------------
  \node[nodo]      (a) at (0,0) {};
  \node[nodopieno] (b) at (4,0) {};
  \node[anchor=north east] at ([shift={(-0.08,-0.12)}]a.south west) {$a$};
  \node[anchor=west]       at ([shift={( 0.26,-0.14)}]b.east)       {$b$};

  % ---- fascio di quattro archi paralleli -----------------------------------
  \draw[arco] (a) to[bend left=58]  (b);
  \draw[arco] (a) to[bend left=10]  (b);
  \draw[arco] (a) to[bend right=15] (b);
  \draw[arco] (a) to[bend right=58] (b);

  % ---- linea di taglio: serpentina verticale, tracciata sopra gli archi ----
  \draw[taglio] plot[domain=-1.50:1.50,samples=100,smooth]
        ({1.80+0.07*sin(\x*720)},{\x});
  \draw[taglio] (1.80,-1.50) -- (1.77,-2.30);

  % ---- annotazione ---------------------------------------------------------
  \draw[line width=0.5pt] (1.96,-2.28) -- (2.36,-2.28);
  \node[anchor=west,align=left] at (2.44,-2.44) {questo è un\\taglio};

\end{tikzpicture}

  \caption{Dopo $n-2$ contrazioni restano due soli nodi, $a$ e $b$, uniti da un fascio di
  archi paralleli: \emph{questo è un taglio} del grafo di partenza.}
  \label{fig:karger-due-nodi-taglio}
\end{figure}

\begin{algorithm}[H]
\caption{\contract: una singola esecuzione dell'algoritmo randomizzato di Karger per il \mincut}
\label{alg:contract}
\begin{algorithmic}[1]
  \Require un multigrafo non orientato e connesso $G=(V,E)$, con $n=\abs{V}$
  \Ensure un insieme di archi che è un taglio di $G$
  \For{$i \gets 1$ \textbf{to} $n-2$}
    \State scegli uniformemente a caso un arco $(u,v)\in E$
    \State cancella $u$ e $v$ da $V$ e introduci un nuovo nodo $uv$
      \Comment{contrazione: $\Oh(1)$}
    \State reindirizza su $uv$ ogni arco incidente su $u$ o su $v$
    \State elimina i self-loop che si sono generati
  \EndFor
  \State \Return gli archi fra i $2$ nodi rimasti
\end{algorithmic}
\end{algorithm}

\begin{osservazione}
L'annotazione $\Oh(1)$ per il costo di una contrazione va presa con cautela: se il
multigrafo è rappresentato esplicitamente, reindirizzare tutti gli archi incidenti su $u$ e
su $v$ costa $\Oh(n)$ e una singola esecuzione dell'algoritmo costa $\Oh(n^{2})$. Il costo
costante va riferito alla sola fusione dei due nodi, e lo si ottiene con una realizzazione
«pigra»: si tiene la lista degli archi di partenza, si fondono i nodi con una
\emph{union--find} e si scartano gli archi diventati self-loop soltanto nel momento in cui
vengono estratti. Anche così, però, il guadagno riguarda solo i grafi sparsi: il costo
complessivo di una singola esecuzione resta $\Om(\abs{E})$, e sui grafi densi
$\abs{E}=\Th(n^{2})$, quindi non si guadagna nulla rispetto alla rappresentazione
esplicita.
\end{osservazione}

L'algoritmo è \emph{corretto per costruzione}: qualunque siano le scelte casuali, ciò che
restituisce è sempre un taglio di $G$. L'unica cosa aleatoria è se sia \emph{minimo}.

\begin{lemma}[gli archi residui formano un taglio]\label{lem:taglio-residuo}
L'insieme $F$ degli archi fra i due nodi rimasti dopo $n-2$ contrazioni è un taglio di $G$.
\end{lemma}

\begin{dimostrazione*}
Siano $a$ e $b$ i due nodi rimasti dopo le $n-2$ contrazioni, e siano
$\{a_1,\dots,a_p\}$ e $\{b_1,\dots,b_q\}$ i nodi di $G$ contratti rispettivamente su $a$ e
su $b$. I due insiemi sono non vuoti e costituiscono una bipartizione di $V$.

Una contrazione non cancella mai un arco i cui estremi finiscono in parti diverse: le
uniche cancellazioni sono quelle dei self-loop, cioè degli archi i cui due estremi sono già
stati fusi nello stesso super-nodo. Ogni arco di $G$ con un estremo in $\{a_1,\dots,a_p\}$
e l'altro in $\{b_1,\dots,b_q\}$ sopravvive dunque fino alla fine, e gli archi paralleli
fra $a$ e $b$ — cioè $F$ — sono esattamente questi.

Ne segue che ogni cammino in $G$ da un $a_i$ a un $b_j$ deve attraversare un arco di $F$:
rimosso $F$, nessun $a_i$ è più collegato ad alcun $b_j$, quindi $(V,E\setminus F)$ è
sconnesso e $F$ è un taglio di $G$.
\end{dimostrazione*}

\begin{figure}[H]
  \centering
  % karger-bipartizione-taglio -- pag. 56
% La bipartizione di V indotta dalle n-2 contrazioni: a sinistra i nodi fusi nel
% super-nodo a, a destra quelli fusi in b. La linea ondulata e' il taglio F
% restituito dall'algoritmo: ogni cammino da un a_i a un b_j e' costretto ad
% attraversarla.  (Nel quaderno questo insieme e' chiamato C; nel capitolo C e'
% riservato al taglio minimo fissato per l'analisi.)
% NB: nel quaderno i pedici sono k e h; qui sono p e q, per non collidere con
% k = |C| e con l'h del numero di ripetizioni.
\begin{tikzpicture}[x=1cm,y=1cm,line join=round,line cap=round,font=\footnotesize]

  % sagoma dell'insieme V (contorno irregolare, "a mano libera")
  \def\sagoma{plot[smooth cycle,tension=0.72] coordinates {
      (2.60,0.05) (2.35,0.72) (1.65,1.18) (0.75,1.36) (-0.35,1.40) (-1.35,1.28)
      (-2.10,0.88) (-2.48,0.30) (-2.36,-0.28) (-2.46,-0.92) (-1.65,-1.30)
      (-0.55,-1.44) (0.60,-1.40) (1.75,-1.18) (2.45,-0.62)}}

  \draw[line width=0.9pt] \sagoma;

  % ---- il taglio F: frontiera ondulata fra le due classi -------------------
  % la serpentina si disegna piu' lunga della sagoma e si taglia sul bordo
  \begin{scope}
    \clip \sagoma;
    \draw[line width=0.8pt] plot[domain=-1.70:1.70,samples=140,smooth]
          ({-0.28+0.09*sin(\x*430)},{\x});
  \end{scope}
  \node[anchor=south,inner sep=1.5pt] at (-0.55,1.38) {$F$};

  % ---- classe dei nodi contratti su a (a sinistra del taglio) --------------
  \node[inner sep=1pt] at (-1.56,-0.70) {$a_{1}$};
  \fill (-1.29,-0.28) circle (0.85pt);
  \fill (-1.17,-0.14) circle (0.85pt);
  \fill (-1.05, 0.00) circle (0.85pt);
  \node[inner sep=1pt] at (-0.80, 0.42) {$a_{p}$};

  % ---- classe dei nodi contratti su b (a destra del taglio) ---------------
  \node[inner sep=1pt] at ( 0.42,-0.70) {$b_{1}$};
  \fill ( 0.60,-0.24) circle (0.85pt);
  \fill ( 0.68,-0.04) circle (0.85pt);
  \fill ( 0.76, 0.16) circle (0.85pt);
  \node[inner sep=1pt] at ( 0.94, 0.62) {$b_{q}$};

\end{tikzpicture}

  \caption{La bipartizione di $V$ indotta dalle contrazioni: a sinistra i nodi contratti su
  $a$, a destra quelli contratti su $b$. La linea ondulata è il taglio $F$, e ogni cammino
  da un $a_i$ a un $b_j$ è costretto ad attraversarla.}
  \label{fig:karger-bipartizione-taglio}
\end{figure}

\section{Analisi: quando il taglio restituito è minimo}

\scan{57}

\begin{domanda}
Qual è la probabilità che il taglio $F$ restituito dall'algoritmo sia un \emph{taglio
minimo}, e non semplicemente un taglio qualunque?
\end{domanda}

\begin{figure}[H]
  \centering
  % mincut-schema-taglio -- pag. 57
% Il grafo G con un suo taglio C (la serpentina verticale). A sinistra una zona
% rada, a destra un cluster denso (un K_4); soltanto due archi attraversano la
% frontiera, quindi |C| = 2. E' questo squilibrio fra archi interni e archi del
% taglio a rendere probabile che una contrazione casuale non distrugga C.
\begin{tikzpicture}[x=1cm,y=1cm,line join=round,line cap=round,
    font=\footnotesize,
    nodo/.style={circle,draw,line width=0.6pt,fill=white,
                 inner sep=0pt,minimum size=4.2pt},
    arco/.style={line width=0.6pt},
    contorno/.style={line width=0.9pt,black!55}]

  % sagoma del grafo G
  \def\sagoma{plot[smooth cycle,tension=0.72] coordinates {
      (3.20,-0.05) (2.90,0.75) (2.15,1.30) (1.05,1.58) (-0.30,1.62) (-1.55,1.50)
      (-2.55,1.05) (-3.10,0.35) (-3.20,-0.35) (-2.75,-1.05) (-1.80,-1.50)
      (-0.55,-1.62) (0.85,-1.58) (2.05,-1.35) (2.95,-0.80)}}

  \draw[contorno] \sagoma;
  \node[anchor=east,inner sep=2pt] at (-2.92,1.00) {$G$};

  % ---- il taglio: serpentina verticale, tagliata sul bordo della sagoma ----
  \begin{scope}
    \clip \sagoma;
    \draw[line width=0.8pt] plot[domain=-1.90:1.90,samples=150,smooth]
          ({-0.32+0.10*sin(\x*400)},{\x});
  \end{scope}
  \node[anchor=south,inner sep=1.5pt] at (-0.32,1.62) {$C$};

  % ---- parte sinistra: rada ------------------------------------------------
  \node[nodo] (A) at (-2.24, 0.16) {};
  \node[nodo] (B) at (-1.15, 0.38) {};
  \node[nodo] (D) at (-2.11,-0.96) {};
  \node[nodo] (F) at (-1.02,-0.90) {};

  \draw[arco] (A) -- (B);
  \draw[arco] (A) -- (F);   % le due diagonali
  \draw[arco] (B) -- (D);   % si incrociano
  \draw[arco] (D) -- (F);

  % ---- parte destra: cluster denso (K_4) -----------------------------------
  \draw[contorno,line width=0.6pt]
        plot[smooth cycle,tension=0.75] coordinates {
          (2.52,-0.16) (2.12,0.78) (1.18,1.14) (0.22,0.80) (-0.14,-0.16)
          (0.24,-1.10) (1.18,-1.44) (2.12,-1.08)};

  \node[nodo] (P) at (0.50, 0.38) {};
  \node[nodo] (Q) at (1.86, 0.38) {};
  \node[nodo] (R) at (1.92,-0.58) {};
  \node[nodo] (S) at (0.50,-0.70) {};

  \draw[arco] (P) -- (Q);
  \draw[arco] (Q) -- (R);
  \draw[arco] (R) -- (S);
  \draw[arco] (S) -- (P);
  \draw[arco] (P) -- (R);
  \draw[arco] (Q) -- (S);

  % ---- gli unici due archi che attraversano il taglio ----------------------
  \draw[arco] (B) -- (P);
  \draw[arco] (F) -- (S);

\end{tikzpicture}

  \caption{Il grafo $G$ e un suo taglio (la linea ondulata): la grande maggioranza degli
  archi è \emph{interna} alle due parti, pochissimi attraversano il taglio. È questo
  squilibrio a rendere probabile che una contrazione scelta a caso \emph{non} distrugga il
  taglio.}
  \label{fig:mincut-schema-taglio}
\end{figure}

\begin{dubbio}
Accanto al disegno, a matita e ormai quasi cancellata, c'è un'annotazione di quattro righe
collegata alla figura da una linea di richiamo. Se ne leggono con ragionevole sicurezza solo
l'inizio, «alta prob.\ che contra\dots», e l'ultima riga, «\dots\ da solo»; il resto non è
ricostruibile. Il senso, coerente con il disegno e con il resto della pagina, è che con alta
probabilità la contrazione cade su un arco \emph{interno} a una delle due parti e non su un
arco del taglio.
\end{dubbio}

\subsection{Gli eventi \texorpdfstring{$\mathcal{E}_i$}{E\_i} e il primo passo}

\begin{definizione}[Scelta positiva o negativa; evento $\mathcal{E}_i$]\label{def:evento-Ei}
Sia $\mathcal{E}_i$ l'evento «la scelta dell'$i$-esimo arco da contrarre è \emph{positiva}»,
dove una scelta si dice
\begin{itemize}
  \item \emph{positiva} se dopo la contrazione rimango dentro un taglio minimo, cioè se
        l'arco estratto \emph{non} appartiene al taglio minimo che voglio preservare;
  \item \emph{negativa} altrimenti.
\end{itemize}
\end{definizione}

Quello che serve calcolare è dunque la probabilità di rimanere dentro un taglio minimo al
passo $i$, cioè la probabilità dell'evento $\mathcal{E}_i$: e, come si vedrà, questa va
condizionata all'esito positivo di tutte le scelte precedenti, perché le contrazioni
successive non sono fra loro indipendenti.

\begin{figure}[H]
  \centering
  % mincut-c-in-e -- pag. 57
% Il taglio minimo C come sottoinsieme dell'insieme E degli archi: l'arco da
% contrarre e' estratto uniformemente a caso da E, quindi la probabilita' di
% colpire un arco di C e' |C|/|E|, e il disegno rende visivo che tale rapporto
% e' piccolo.
\begin{tikzpicture}[x=1cm,y=1cm,line join=round,line cap=round,font=\footnotesize]

  % ---- l'insieme E degli archi --------------------------------------------
  \draw[line width=0.9pt] plot[smooth cycle,tension=0.72] coordinates {
      (2.50,0.10) (2.30,0.85) (1.55,1.35) (0.55,1.44) (-0.20,1.34) (-0.90,1.46)
      (-1.75,1.30) (-2.35,0.75) (-2.50,0.00) (-2.30,-0.75) (-1.60,-1.30)
      (-0.55,-1.50) (0.45,-1.38) (1.10,-1.46) (1.95,-1.15)};
  \node[inner sep=1pt] at (0.55,-0.95) {$E$};

  % ---- il taglio minimo C --------------------------------------------------
  \draw[line width=0.9pt,fill=black!8] plot[smooth cycle,tension=0.75] coordinates {
      (0.12,0.22) (0.02,0.64) (-0.40,0.90) (-0.92,0.92) (-1.38,0.68)
      (-1.54,0.26) (-1.42,-0.24) (-1.10,-0.56) (-0.62,-0.60) (-0.14,-0.32)};
  \node[inner sep=1pt] at (-0.72,0.16) {$C$};

\end{tikzpicture}

  \caption{Il taglio minimo $C$ come sottoinsieme di $E$: l'arco da contrarre viene estratto
  uniformemente a caso da $E$, quindi la probabilità di colpire proprio un arco di $C$ è
  $\abs{C}/\abs{E}$.}
  \label{fig:mincut-c-in-e}
\end{figure}

Nel caso peggiore il grafo ha un \emph{unico} taglio minimo $C$: la probabilità di scegliere
proprio un arco di quel taglio è
\[
  \frac{\abs{C}}{\abs{E}} ,
\]
e di $C$ interessa soltanto la cardinalità. Se i tagli minimi fossero più d'uno, per
sbagliare occorrerebbe distruggerli tutti e la probabilità di errore sarebbe minore: il caso
di un unico taglio minimo è quindi davvero il peggiore. Fissiamo allora un taglio minimo $C$
e poniamo $k=\abs{C}$; ricordiamo che $n=\abs{V}$.

\begin{osservazione}
Per ogni vertice $v$ vale $d(v) \ge k$.

Altrimenti l'insieme degli archi incidenti a $v$ — che è già un taglio, quello che «taglia
fuori» il solo nodo $v$ — avrebbe cardinalità $d(v) < k$: assurdo, perché $k$ è la
cardinalità minima di un taglio di $G$.
\end{osservazione}

Sommando i gradi di tutti i vertici ogni arco viene contato due volte, quindi
$2\abs{E} = \sum_{v \in V} d(v) \ge n\,k$, da cui
\[
  \abs{E} \;\ge\; \frac{k \cdot n}{2}
  \qquad\Longrightarrow\qquad
  \frac{\abs{C}}{\abs{E}} \;\le\; \frac{k}{\dfrac{k \cdot n}{2}} \;=\; \frac{2}{n} .
\]

Il valore $2/n$ è dunque un maggiorante della probabilità dell'evento complementare
$\overline{\mathcal{E}_1}$, cioè della probabilità di distruggere il taglio minimo già alla
prima contrazione:
\[
  \Pr\bigl[\,\overline{\mathcal{E}_1}\,\bigr] \;\le\; \frac{2}{n} .
\]

\scan{58}
Ne segue che la probabilità che, dopo la prima contrazione, esista ancora un taglio minimo
che l'algoritmo può determinare è
\[
  \Pr[\mathcal{E}_1] \;\ge\; 1-\frac{2}{n} .
\]

\subsection{I passi successivi}

Indichiamo con $G_i$ il multigrafo ottenuto dopo $i$ contrazioni ($G_0=G$) e con $E_i$ il suo
insieme di archi.

Passiamo alla seconda contrazione. Il multigrafo $G_1$ ha $n-1$ nodi e, condizionatamente a
$\mathcal{E}_1$, ammette ancora $C$ come taglio minimo, di cardinalità $k$. Ogni nodo di
$G_1$ ha quindi grado almeno $k$, per cui
\[
  \abs{E_1}\;\ge\;\frac{k\,(n-1)}{2}
  \qquad\Longrightarrow\qquad
  \Pr\!\left[\overline{\mathcal{E}_2}\;\middle|\;\mathcal{E}_1\right]
  \;\le\;\frac{\abs{C}}{\abs{E_1}}
  \;\le\;\frac{k}{\dfrac{k\,(n-1)}{2}}
  \;=\;\frac{2}{n-1},
\]
e passando all'evento complementare
\[
  \Pr[\mathcal{E}_2\mid\mathcal{E}_1]\;\ge\;1-\frac{2}{n-1}.
\]

\begin{osservazione}
Il passaggio si ripete identico a ogni passo perché una contrazione non può \emph{diminuire}
il valore del taglio minimo: ogni taglio del multigrafo contratto $G_i$ è anche un taglio di
$G$, e quindi contiene almeno $k$ archi. Finché non si contrae un arco di $C$, il taglio
minimo di $G_i$ vale esattamente $k$, ogni nodo di $G_i$ ha grado $\ge k$ e vale di nuovo la
stima sul numero di archi.
\end{osservazione}

In generale, dopo $i-1$ contrazioni restano $n-(i-1)=n-i+1$ nodi, dunque
$\abs{E_{i-1}}\ge k\,(n-i+1)/2$ e
\[
  \Pr\!\left[\mathcal{E}_i \;\middle|\; \bigcap_{j=1}^{i-1}\mathcal{E}_j\right]
  \;\ge\; 1-\frac{2}{n-i+1}.
\]

\subsection{Probabilità di successo di una singola esecuzione}

\begin{teorema}[Karger]\label{thm:karger}
Fissato un taglio minimo $C$ di $G$, una singola esecuzione di \contract restituisce
proprio $C$ con probabilità
\[
  \Pr\!\left[\bigcap_{i=1}^{n-2}\mathcal{E}_i\right]
  \;\ge\; \frac{2}{n(n-1)} \;=\; \frac{1}{\binom{n}{2}} \;\ge\; \frac{2}{n^{2}} .
\]
\end{teorema}

\begin{dimostrazione*}
La quantità che interessa è la probabilità che la scelta dell'arco da contrarre sia
\emph{positiva} a \emph{ogni} passo, cioè $\Pr\bigl[\bigcap_{i=1}^{n-2}\mathcal{E}_i\bigr]$:
poiché le scelte successive non sono indipendenti, la si calcola con la regola della catena.
\begin{align*}
  \Pr\!\left[\bigcap_{i=1}^{n-2}\mathcal{E}_i\right]
  &= \Pr[\mathcal{E}_1]\cdot\Pr[\mathcal{E}_2\mid\mathcal{E}_1]\cdot\;\cdots\;\cdot
     \Pr\!\left[\mathcal{E}_{n-2}\;\middle|\;\bigcap_{i=1}^{n-3}\mathcal{E}_i\right]\\[4pt]
  &\ge \left(1-\frac{2}{n}\right)\left(1-\frac{2}{n-1}\right)\cdot\;\cdots\;\cdot
       \left(1-\frac{2}{\,n-(n-2)+1\,}\right)\\[4pt]
  &= \left(\frac{n-2}{n}\right)\left(\frac{n-3}{n-1}\right)\left(\frac{n-4}{n-2}\right)
     \cdots\;\frac{2}{4}\cdot\frac{1}{3}
   \;=\;\frac{2}{n(n-1)} ,
\end{align*}
dove l'ultimo passaggio segue dal fatto che il prodotto è telescopico: ogni numeratore si
semplifica con il denominatore che compare due fattori più avanti, e sopravvivono soltanto i
numeratori $2$ e $1$ degli ultimi due fattori e i denominatori $n$ e $n-1$ dei primi due.
Infine $n(n-1)\le n^{2}$ dà l'ultima disuguaglianza.
\end{dimostrazione*}

La probabilità di successo è dunque solo \emph{inversamente polinomiale} in $n$, non
esponenzialmente piccola: è precisamente questo a rendere efficace l'amplificazione per
ripetizione.

\section{Amplificazione per ripetizione}

\begin{osservazione}
La probabilità che il taglio restituito da \contract \emph{non} sia minimo è
\[
  \begin{aligned}
    \text{dopo un'esecuzione:}   &\quad \le\; 1-\frac{2}{n^{2}},\\[4pt]
    \text{dopo due esecuzioni:}  &\quad \le\; \left(1-\frac{2}{n^{2}}\right)^{\!2},\\[4pt]
    \text{dopo $h$ esecuzioni:}  &\quad \le\; \left(1-\frac{2}{n^{2}}\right)^{\!h}.
  \end{aligned}
\]
Le $h$ esecuzioni sono infatti indipendenti fra loro — a differenza dei passi \emph{interni}
a una singola esecuzione — e si restituisce il migliore degli $h$ tagli trovati: si sbaglia
soltanto se sbagliano tutte, da cui il prodotto delle probabilità di errore.
\end{osservazione}

\begin{figure}[H]
  \centering
  % karger-tradeoff -- pag. 58
% Il compromesso dell'amplificazione per ripetizione: al crescere del numero h
% di esecuzioni indipendenti sale il costo (freccia verso l'alto) e scende la
% probabilita' di errore (freccia verso il basso).
% Nel quaderno le etichette sono abbreviate («maggior complessità», «minore
% prob. di errore»): qui sono espanse.
\begin{tikzpicture}[x=1cm,y=1cm,line cap=round,font=\footnotesize,
    >={Stealth[length=5pt,width=4pt]}]

  % ---- costo: cresce -------------------------------------------------------
  \draw[->,line width=0.8pt] (0,0) -- (0,1.55);
  \node[anchor=west,inner sep=5pt] at (0.10,1.45) {maggiore complessità computazionale};

  % ---- probabilità di errore: cala ----------------------------------------
  \draw[->,line width=0.8pt] (0,-0.75) -- (0,-2.30);
  \node[anchor=west,inner sep=5pt] at (0.10,-2.20) {minore probabilità di errore};

\end{tikzpicture}

  \caption{Il compromesso fra costo e affidabilità: al crescere del numero $h$ di ripetizioni
  aumenta la complessità computazionale e diminuisce la probabilità di errore.}
  \label{fig:karger-tradeoff}
\end{figure}

\scan{59}
Resta da fissare il numero $h$ di ripetizioni indipendenti. Prendendo
\[
  h \;=\; \frac{n^{2}}{2}
\]
si ottiene
\[
  \left(1-\frac{1}{\,n^{2}/2\,}\right)^{\!n^{2}/2}
  \;=\;\left(1-\frac{2}{n^{2}}\right)^{\!n^{2}/2}
  \;<\;\frac{1}{e},
\]
cioè la probabilità che \emph{nessuna} delle $h$ esecuzioni abbia prodotto un taglio minimo
scende sotto $1/e$.

\begin{osservazione}
La disuguaglianza è il caso $m=n^{2}/2$ della stima
$\bigl(1-\tfrac{1}{m}\bigr)^{m}<e^{-1}$, valida per ogni $m>1$ e conseguenza immediata di
$1-x<e^{-x}$ — vera per ogni $x\neq 0$ — posto $x=\tfrac{1}{m}$.
\end{osservazione}

Numericamente
\[
  \frac{1}{e}\;\approx\;\frac{1}{3}\;\approx\;0{,}33
\]
(più precisamente $1/e\approx 0{,}368$): con $h=n^{2}/2$ ripetizioni la probabilità di errore
è dunque una \emph{costante} minore di $1$, indipendente da $n$. Trattandosi di una costante,
la si abbatte quanto si vuole continuando a iterare: con
$h=\frac{n^{2}}{2}\ln\frac{1}{\delta}$ esecuzioni la probabilità di errore è al più $\delta$.

\begin{osservazione}
Mettendo insieme le due stime si legge in cifre il compromesso della
\autoref{fig:karger-tradeoff}: una singola esecuzione di \contract costa $\Oh(n^{2})$ con la
rappresentazione esplicita del multigrafo, quindi $h=n^{2}/2$ ripetizioni costano
$\Oh(n^{4})$ per una probabilità di errore inferiore a $1/e$, e
$h=\Th(n^{2}\log n)$ ripetizioni costano $\Oh(n^{4}\log n)$ per una probabilità di errore
polinomialmente piccola.
\end{osservazione}

\chapter{Partizioni binarie del piano e RandAuto}

\section{Il problema della partizione binaria del piano}

\scan{60}

Dati $n$ segmenti del piano a due a due \emph{non intersecanti}, vogliamo partizionare il
piano in modo tale che in ogni regione ci sia al più un segmento, oppure una parte di esso.
Scriveremo $S=\set{s_1,\dots,s_n}$ per l'insieme dei segmenti dati, e useremo indifferentemente
$n$ e $\abs{S}$ per la loro cardinalità. Una partizione di questo tipo si chiama
\emph{partizione binaria del piano} (\emph{binary planar partition}, in breve BPP): «binaria»
perché la si costruisce tagliando ripetutamente una regione in due con una retta.

\begin{esempio}
Cinque segmenti $s_1,\dots,s_5$ e una partizione che li separa. I tagli (tratteggiati)
vengono eseguiti uno dopo l'altro, ciascuno all'interno della regione che si vuole spezzare,
e producono le sei regioni $r_1,\dots,r_6$. Uno dei tagli attraversa $s_3$ e lo spezza in due
parti, che finiscono in due regioni diverse ($r_6$ e $r_4$): per questo le regioni sono
\emph{sei} e non cinque.

\begin{figure}[H]
\centering
% Esempio di partizione binaria del piano (BPP) per cinque segmenti non
% intersecanti.  Tratto pieno = segmenti s_1,...,s_5; tratteggio = rette di
% taglio l_1,...,l_5, disegnate nell'ordine in cui la partizione le esegue
% (ogni taglio vive soltanto dentro la regione che spezza, e l_3 termina con
% una giunzione a T su l_2).  L'unico segmento attraversato da un taglio e'
% s_3: lo spezza l_2, e le due porzioni finiscono in r_4 (quella bassa) e in
% r_6 (quella alta).  Sei regioni per cinque segmenti.
%
% Ogni taglio agisce su una regione che contiene almeno due porzioni di
% segmento, come vuole la costruzione:
%   l_1 : {s_5}            | resto
%   l_2 : r_4 = {s_3 basso}| resto
%   l_3 : r_6 = {s_3 alto} | resto
%   l_4 : r_2 = {s_1}      | resto
%   l_5 : r_5 = {s_4}      | r_3 = {s_2}
\begin{tikzpicture}[
    scale=0.95,
    segmento/.style={draw=black,line width=1.3pt,line cap=round},
    taglio/.style={draw=graydark,line width=0.7pt,line cap=round,
                   dash pattern=on 3.2pt off 2.4pt},
    et/.style={font=\footnotesize},
    etg/.style={font=\footnotesize,text=graydark},
  ]

  % ---- riquadro implicito (nessuna cornice disegnata) ---------------------
  \useasboundingbox (-0.15,-0.55) rectangle (9.15,6.05);

  % ---- rette di taglio, nell'ordine di esecuzione -------------------------
  % l_1: taglia tutto il piano
  \draw[taglio] (0,4.10) -- (9,4.10);
  \node[etg] at (0.42,4.36) {$\ell_1$};
  % l_2: taglia il semipiano inferiore e spezza s_3
  \draw[taglio] (1.45,4.10) -- (3.438,0);
  \node[etg] at (3.44,-0.28) {$\ell_2$};
  % l_3: taglia la regione a destra di l_2 (giunzione a T su l_2)
  \draw[taglio] (2.75,4.10) -- (2.4683,2.00);
  \node[etg] at (3.05,3.80) {$\ell_3$};
  % l_4: taglia la regione a destra di l_3
  \draw[taglio] (5.15,4.10) -- (7.60,0);
  \node[etg] at (7.60,-0.28) {$\ell_4$};
  % l_5: taglia la regione compresa fra l_2 e l_4
  \draw[taglio] (2.856,1.20) -- (5.7296,3.13)
        node[etg,pos=0.86,below=2.5pt,sloped] {$\ell_5$};

  % ---- segmenti -----------------------------------------------------------
  \draw[segmento] (2.00,5.30) -- (2.70,4.35);           % s_5, sopra l_1
  \draw[segmento] (1.00,1.60) -- (2.15,2.95);           % s_3, spezzato da l_2
  \draw[segmento] (3.10,2.40) -- (5.00,3.75);           % s_4
  \draw[segmento] (4.60,1.55) -- (4.85,0.85);           % s_2
  \draw[segmento] (6.50,2.60) -- (8.20,2.45);           % s_1

  % punto in cui l_2 spezza s_3
  \filldraw[fill=white,draw=black,line width=0.5pt] (2.059,2.844) circle (1.7pt);

  % ---- etichette dei segmenti --------------------------------------------
  \node[et] at (3.00,4.98) {$s_5$};
  \node[et] at (0.72,2.50) {$s_3$};
  \node[et] at (3.90,2.66) {$s_4$};
  \node[et] at (5.20,1.18) {$s_2$};
  \node[et] at (7.45,2.92) {$s_1$};

  % ---- etichette delle regioni -------------------------------------------
  \node[et] at (7.50,5.20) {$r_1$};   % sopra l_1: contiene s_5
  \node[et] at (2.20,3.55) {$r_6$};   % fra l_2 e l_3, sotto l_1: porzione alta di s_3
  \node[et] at (1.30,1.35) {$r_4$};   % a sinistra di l_2: porzione bassa di s_3
  \node[et] at (2.92,2.10) {$r_5$};   % fra l_3, l_4 e sopra l_5: contiene s_4
  \node[et] at (3.90,0.60) {$r_3$};   % sotto l_5: contiene s_2
  \node[et] at (8.30,0.90) {$r_2$};   % a destra di l_4: contiene s_1
\end{tikzpicture}

\caption{Una partizione binaria del piano per cinque segmenti non intersecanti: in tratto
continuo i segmenti $s_1,\dots,s_5$, tratteggiate le rette di taglio. Ognuna delle sei
regioni $r_1,\dots,r_6$ contiene esattamente un segmento o una porzione di segmento;
$s_3$ è l'unico segmento spezzato.}
\end{figure}
\end{esempio}

\begin{domanda}
Dati $n$ segmenti, esiste sempre una partizione che genera $n$ regioni?
\end{domanda}

La risposta è \emph{no}, ed eppure i segmenti che bastano a farla cadere sono soltanto tre:
tre segmenti disposti «a girandola» — ciascuno rivolto verso il successivo, e a due a due
non intersecanti — non si lasciano separare da sole tre regioni.
Tagliando con la retta che prolunga $s_2$ si spezza $s_3$ in due parti; con altri due tagli,
le rette che prolungano $s_1$ e $s_3$, si isolano i quattro frammenti così ottenuti. Le
regioni finali sono quattro.

\begin{figure}[H]
\centering
% Tre segmenti «a girandola»: nessuna partizione binaria del piano riesce a
% separarli con tre sole regioni.  Tratto pieno = i segmenti; tratteggio = le
% rette che li prolungano, usate come tagli.  Prima si taglia con l(s_2), che
% attraversa s_3 e lo spezza; poi, nel semipiano superiore, con l(s_1) e, in
% quello inferiore, con l(s_3).  Ogni taglio e' disegnato solo dentro la
% regione che spezza, quindi l(s_1) e l(s_3) si fermano su l(s_2).
% Risultato: quattro regioni r_1,...,r_4 per tre segmenti.
\begin{tikzpicture}[
    scale=1.0,
    segmento/.style={draw=black,line width=1.3pt,line cap=round},
    taglio/.style={draw=graydark,line width=0.7pt,line cap=round,
                   dash pattern=on 3.2pt off 2.4pt},
    et/.style={font=\footnotesize},
    etg/.style={font=\footnotesize,text=graydark},
    testo/.style={font=\footnotesize},
  ]

  \useasboundingbox (-0.10,-0.10) rectangle (8.60,5.70);

  % ---- rette di taglio -----------------------------------------------------
  % l(s_2): il primo taglio, attraversa tutto il piano e spezza s_3
  \draw[taglio] (0.30,2.4042) -- (7.30,2.2875);
  \node[etg,anchor=west] at (7.38,2.30) {$\ell(s_2)$};
  % l(s_1): taglia il semipiano superiore, si ferma su l(s_2)
  \draw[taglio] (2.6191,5.30) -- (4.704,2.3311);
  \node[etg] at (2.42,5.48) {$\ell(s_1)$};
  % l(s_3): taglia il semipiano inferiore, si ferma su l(s_2)
  \draw[taglio] (2.3634,0.30) -- (3.5138,2.3506);
  \node[etg,anchor=north east] at (2.30,0.22) {$\ell(s_3)$};

  % ---- i tre segmenti (a due a due non intersecanti) ----------------------
  \draw[segmento] (2.70,0.90) -- (3.85,2.95);   % s_3, spezzato da l(s_2)
  \draw[segmento] (3.05,4.686) -- (4.42,2.735); % s_1
  \draw[segmento] (3.80,2.3459) -- (6.55,2.30);  % s_2, giace su l(s_2); l'estremo
                                                 % sinistro sta staccato dal punto
                                                 % in cui l(s_2) spezza s_3

  % punto in cui l(s_2) spezza s_3
  \filldraw[fill=white,draw=black,line width=0.5pt] (3.5138,2.3506) circle (1.7pt);

  % ---- etichette dei segmenti --------------------------------------------
  \node[et] at (2.80,1.90) {$s_3$};
  \node[et] at (4.34,3.30) {$s_1$};
  \node[et] at (5.00,1.98) {$s_2$};

  % ---- etichette delle regioni -------------------------------------------
  \node[et] at (2.00,3.90) {$r_4$};   % sopra l(s_2), a sinistra di l(s_1): porzione alta di s_3
  \node[et] at (5.60,4.10) {$r_2$};   % sopra l(s_2), a destra di l(s_1): contiene s_1
  \node[et] at (1.80,1.30) {$r_3$};   % sotto l(s_2), a sinistra di l(s_3): porzione bassa di s_3
  \node[et] at (5.10,0.85) {$r_1$};   % sotto l(s_2), a destra di l(s_3): contiene s_2

  % ---- annotazione --------------------------------------------------------
  \node[testo,anchor=west] at (6.55,1.55) {3 segmenti};
  \node[testo,anchor=west] at (6.55,1.05) {4 regioni};
\end{tikzpicture}

\caption{Tre segmenti bastano a impedire una partizione con tre sole regioni. Le tre rette
tratteggiate prolungano i tre segmenti; la prima, quella di $s_2$, spezza $s_3$, e le
regioni finali sono quattro: $r_1,\dots,r_4$.}
\end{figure}

Il numero di regioni, dunque, può eccedere il numero dei segmenti. La domanda interessante
diventa allora \emph{di quanto} lo può eccedere: bastano sempre $\Oh(n)$ regioni? È una
domanda tutt'altro che facile, e per porla in modo preciso conviene prima darsi una
rappresentazione formale delle partizioni.

\section{La BPP come albero binario}

\scan{61}

Rappresenteremo una BPP come un \emph{albero binario} i cui nodi interni sono coppie
\[
  (r,\ \ell)
\]
dove $r$ è la \emph{regione} associata al nodo e $\ell$ è il \emph{limite}, cioè la retta
che divide $r$: i due figli del nodo corrispondono alle due sottoregioni in cui $\ell$ taglia
$r$. Ogni foglia è etichettata dal segmento — o dalla porzione di segmento — contenuto nella
regione corrispondente. Chiamiamo \emph{dimensione} della partizione $P$ il numero delle sue
foglie, e la indichiamo con $\abs{P}$: è la quantità che vogliamo tenere piccola.

Le foglie sono almeno $n$, e possono essere di più.

\begin{esempio}
Sia $S=\set{s_1,s_2,s_3}$. La prima retta $\ell_1$ attraversa $s_3$ e lo spezza nelle due
porzioni $s_{3,1}$ e $s_{3,2}$: da una parte di $\ell_1$ restano $s_1$ e $s_{3,1}$, che
vengono separati da $\ell_2$; dall'altra restano $s_2$ e $s_{3,2}$, separati da $\ell_3$.
L'albero ha allora tre nodi interni e \emph{quattro} foglie, una in più dei tre segmenti di
partenza, proprio perché $s_3$ è stato spezzato.

\begin{figure}[H]
\centering
% Una BPP di S = {s_1,s_2,s_3} e l'albero binario che la rappresenta.
% A sinistra la scena: la prima retta l_1 taglia tutto il piano e attraversa
% s_3, spezzandolo in s_{3,1} (sopra l_1) e s_{3,2} (sotto l_1); l_2 e l_3
% sono i tagli delle due regioni figlie e sono quindi semirette che partono
% da l_1 e non la oltrepassano.
% A destra l'albero: nodi interni = coppie (r,l), foglie = segmenti o porzioni
% di segmento.  Tre nodi interni (tre tagli) e quattro foglie = 3 segmenti + 1
% spezzamento; le due foglie evidenziate sono le due porzioni di s_3.
\begin{tikzpicture}[
    segmento/.style={draw=black,line width=1.3pt,line cap=round},
    taglio/.style={draw=graydark,line width=0.7pt,line cap=round,
                   dash pattern=on 3.2pt off 2.4pt},
    et/.style={font=\footnotesize},
    etg/.style={font=\footnotesize,text=graydark},
    interno/.style={ellipse,draw=black,line width=0.6pt,inner sep=1.5pt,
                    minimum width=1.5cm,minimum height=0.74cm,font=\footnotesize},
    foglia/.style={circle,draw=black,line width=0.6pt,inner sep=0.5pt,
                   minimum size=0.76cm,font=\footnotesize},
    spezzata/.style={foglia,fill=accent!25},
    arco/.style={draw=black,line width=0.6pt},
  ]

  %% ======================= (a) la scena ==================================
  \begin{scope}[scale=0.78]

    % l_1: primo taglio, attraversa tutto il piano e spezza s_3
    \draw[taglio] (0.80,7.05) -- (6.50,2.10);
    \node[etg] at (6.90,1.80) {$\ell_1$};
    % l_2: taglia la regione sopra l_1 (semiretta uscente da l_1)
    \draw[taglio] (4.00,4.271) -- (6.10,6.85);
    \node[etg] at (6.42,7.10) {$\ell_2$};
    % l_3: taglia la regione sotto l_1 (semiretta uscente da l_1)
    \draw[taglio] (3.50,4.705) -- (3.85,1.15);
    \node[etg] at (3.92,0.78) {$\ell_3$};

    % i tre segmenti
    \draw[segmento] (2.30,6.00) -- (4.30,6.80);   % s_1
    \draw[segmento] (1.30,5.85) -- (2.40,3.80);   % s_2
    \draw[segmento] (4.40,3.20) -- (5.80,5.05);   % s_3, spezzato da l_1

    % punto in cui l_1 spezza s_3
    \filldraw[fill=white,draw=black,line width=0.5pt] (4.7305,3.6368) circle (2.2pt);

    \node[et] at (3.00,7.15) {$s_1$};
    \node[et] at (0.95,4.90) {$s_2$};
    \node[et] at (6.10,5.25) {$s_3$};
    \node[etg] at (5.72,4.00) {$s_{3,1}$};
    \node[etg] at (4.30,2.72) {$s_{3,2}$};
  \end{scope}

  %% ======================= (b) l'albero ==================================
  \begin{scope}[shift={(8.60,1.55)}]
    \node[interno] (n1) at ( 0.00,3.00) {$r_1\ \ \ell_1$};
    \node[interno] (n2) at (-1.50,1.50) {$r_2\ \ \ell_2$};
    \node[interno] (n3) at ( 1.50,1.50) {$r_3\ \ \ell_3$};

    \node[foglia]   (f1) at (-2.10,0.00) {$s_1$};
    \node[spezzata] (f2) at (-0.90,0.00) {$s_{3,1}$};
    \node[foglia]   (f3) at ( 0.90,0.00) {$s_2$};
    \node[spezzata] (f4) at ( 2.10,0.00) {$s_{3,2}$};

    \draw[arco] (n1) -- (n2);
    \draw[arco] (n1) -- (n3);
    \draw[arco] (n2) -- (f1);
    \draw[arco] (n2) -- (f2);
    \draw[arco] (n3) -- (f3);
    \draw[arco] (n3) -- (f4);
  \end{scope}
\end{tikzpicture}

\caption{Una BPP di $S=\set{s_1,s_2,s_3}$ (a sinistra la scena, a destra l'albero che la
rappresenta). La retta $\ell_1$ spezza $s_3$ in $s_{3,1}$ e $s_{3,2}$: le foglie sono
quattro, una più dei segmenti.}
\end{figure}
\end{esempio}

\begin{osservazione}
L'albero che rappresenta la $\mathrm{BPP}(S)$ è sempre un albero binario con un numero di
foglie pari a
\[
  \#\text{foglie} \;=\; n + \#\text{tagli},
\]
dove $\#\text{tagli}$ conta i tagli che spezzano un segmento: ogni taglio che spezza un
segmento produce una porzione in più, e quindi una foglia in più (un segmento attraversato da
$k$ tagli dà origine a $k+1$ porzioni). Questo numero non dipende solo da $S$, ma dalle scelte
fatte, cioè dalle rette con cui si decide di partizionare.
\end{osservazione}

\begin{center}
\fbox{\textbf{Goal:} tenere \emph{basso} il numero di tagli.}
\end{center}

\begin{domanda}[problema aperto]
Esiste sempre una $\mathrm{BPP}(S)$ con $\Oh(\abs{S})$ foglie, cioè con $\Oh(\abs{S})$
frammenti di segmento e quindi $\Oh(\abs{S})$ regioni?
\dubbioinl{negli appunti «Problema aperto» è seguito da una parola fra parentesi, di lettura
incerta: forse un nome proprio, «Clark»}
\end{domanda}

Negli appunti il problema è dato per aperto, come in Motwani--Raghavan da cui la domanda è
ripresa. Quello che sappiamo fare è $\Oh(\abs{S}\log\abs{S})$: diamo infatti un algoritmo
randomizzato, \textsc{RandAuto}, che dimostra che esiste sempre una $\mathrm{BPP}(S)$ con
quel numero di foglie.

\section{Applicazione: il \emph{painter's algorithm}}

\scan{62}

Una prima applicazione delle partizioni binarie del piano è il \emph{painter's algorithm}:
si usa l'albero della BPP come struttura dati per disegnare una scena.

Fissata la posizione dell'osservatore, l'albero viene visitato ricorsivamente: in ogni nodo
interno $(r,\ell)$ si disegna per prima la regione figlia che sta dalla parte di $\ell$
\emph{opposta} all'osservatore e poi l'altra. Si ottiene così un ordine di disegno dal più
lontano al più vicino, in cui gli oggetti vicini sovrascrivono quelli lontani e la rimozione
delle superfici nascoste è automatica. Il tempo di disegno è proporzionale a $\abs{P}$, cioè
al numero di foglie dell'albero: da qui l'interesse per partizioni di dimensione piccola.

\begin{figure}[H]
\centering
% p62-painter -- il painter's algorithm sull'albero di una BPP.
%
% La scena e' una regione irregolare del piano che contiene 7 segmenti disgiunti
% (tratto spesso).  Due tagli della BPP (tratteggiati, grigi) la suddividono in
% tre regioni:
%   ell_1  taglia la scena in una parte sinistra e una parte destra;
%   ell_2  parte da un punto interno di ell_1 (pallino grigio) e taglia la sola
%          parte destra in una porzione alta e una bassa.
% L'osservatore sta fuori dalla scena, in basso a sinistra: e' quindi a sinistra
% di ell_1 e sotto ell_2.  La visita ricorsiva dell'albero disegna per prima la
% regione dalla parte di ell opposta all'osservatore, quindi nell'ordine
%   (1) destra-alta, (2) destra-bassa, (3) sinistra,
% cioe' dal piu' lontano al piu' vicino.  I numeri cerchiati nelle regioni e la
% legenda in basso riportano quest'ordine; la gradazione del riempimento va dal
% piu' scuro (disegnato per primo) al piu' chiaro (disegnato per ultimo).
%
% Nota: nel manoscritto i due tagli sono curvi e partono entrambi dallo stesso
% punto del bordo; qui sono rettilinei (i tagli di una BPP sono rette) e ell_2
% nasce nell'interno di ell_1, cosi' da rendere visibile la gerarchia.
\begin{tikzpicture}[
    x=7.30cm, y=7.00cm,
    font=\footnotesize,
    >={Stealth[length=4.2pt,width=3.2pt]},
    bordo/.style={draw=black,line width=0.8pt,line join=round,line cap=round},
    taglio/.style={draw=graydark,line width=0.7pt,line cap=round,
                   dash pattern=on 3.2pt off 2.4pt},
    segmento/.style={draw=black,line width=1.4pt,line cap=round},
    etg/.style={text=graydark,inner sep=1.5pt},
    ordine/.style={circle,draw=graydark,fill=white,line width=0.5pt,
                   inner sep=0pt,minimum size=4.2mm,font=\scriptsize},
    occhio/.style={draw=black,line width=0.8pt},
    raggio/.style={->,draw=black,line width=0.6pt},
    frecciaord/.style={->,draw=graydark,line width=0.5pt},
  ]

  % ---- il bordo della scena (curva chiusa irregolare, schizzo a mano) ------
  \def\scenapath{plot[smooth cycle,tension=0.55] coordinates {
      (0.393,0.941) (0.457,0.927) (0.530,0.929) (0.590,0.948) (0.650,0.979)
      (0.713,0.996) (0.778,0.998) (0.838,0.984) (0.893,0.962) (0.918,0.949)
      (0.905,0.888) (0.893,0.815) (0.888,0.762) (0.910,0.710) (0.948,0.665)
      (0.996,0.626) (1.002,0.559) (0.996,0.488) (0.980,0.431) (0.957,0.357)
      (0.923,0.294) (0.877,0.247) (0.824,0.216) (0.767,0.201) (0.712,0.201)
      (0.666,0.182) (0.622,0.149) (0.575,0.100) (0.535,0.049) (0.471,0.011)
      (0.402,0.000) (0.347,0.008) (0.310,0.035) (0.295,0.078) (0.285,0.120)
      (0.243,0.116) (0.228,0.152) (0.194,0.172) (0.190,0.214) (0.157,0.240)
      (0.156,0.288) (0.124,0.318) (0.116,0.356) (0.085,0.388) (0.062,0.437)
      (0.043,0.457) (0.052,0.487) (0.030,0.532) (0.010,0.588) (0.002,0.645)
      (0.015,0.689) (0.032,0.723) (0.081,0.781) (0.118,0.793) (0.175,0.803)
      (0.246,0.853) (0.310,0.898)}}

  % estremi dei due tagli
  \coordinate (T)  at (0.3930,0.9410);   % ell_1 : bordo superiore
  \coordinate (Bs) at (0.5750,0.1000);   % ell_1 : bordo inferiore
  \coordinate (M)  at (0.4157,0.8360);   % ell_2 : punto interno di ell_1
  \coordinate (R)  at (0.9960,0.6260);   % ell_2 : bordo destro

  % ---- riempimento delle tre regioni (ritaglio scena + semipiani) ---------
  % (3) regione sinistra -- la piu' vicina all'osservatore, disegnata per ultima
  \begin{scope}
    \clip \scenapath;
    \clip (0.2838,1.4456) -- (0.6842,-0.4046) -- (-0.60,-0.4046) -- (-0.60,1.4456) -- cycle;
    \fill[accent!9] (-0.60,-0.50) rectangle (1.80,1.60);
  \end{scope}
  % (1) regione destra-alta -- la piu' lontana, disegnata per prima
  \begin{scope}
    \clip \scenapath;
    \clip (0.2838,1.4456) -- (0.6842,-0.4046) -- (1.80,-0.4046) -- (1.80,1.4456) -- cycle;
    \clip (-0.0485,1.0040) -- (1.4602,0.4580) -- (1.80,1.60) -- (-0.60,1.60) -- cycle;
    \fill[accent!38] (-0.60,-0.50) rectangle (1.80,1.60);
  \end{scope}
  % (2) regione destra-bassa
  \begin{scope}
    \clip \scenapath;
    \clip (0.2838,1.4456) -- (0.6842,-0.4046) -- (1.80,-0.4046) -- (1.80,1.4456) -- cycle;
    \clip (-0.0485,1.0040) -- (1.4602,0.4580) -- (1.80,-0.60) -- (-0.60,-0.60) -- cycle;
    \fill[accent!21] (-0.60,-0.50) rectangle (1.80,1.60);
  \end{scope}

  % ---- il bordo della scena ----------------------------------------------
  \draw[bordo] \scenapath;

  % ---- i due tagli della BPP ---------------------------------------------
  \draw[taglio] (T) -- (Bs);
  \draw[taglio] (M) -- (R);
  \fill[graydark] (M) circle (1.2pt);
  \node[etg,anchor=south east] at (0.352,0.960) {$\ell_1$};
  \node[etg,anchor=west]       at (1.012,0.646) {$\ell_2$};

  % ---- i sette segmenti della scena --------------------------------------
  \draw[segmento] (0.232,0.574) -- (0.356,0.662);
  \draw[segmento] (0.138,0.453) -- (0.219,0.395);
  \draw[segmento] (0.358,0.421) -- (0.416,0.540);
  \draw[segmento] (0.453,0.283) -- (0.501,0.212);
  \draw[segmento] (0.645,0.853) -- (0.828,0.779);
  \draw[segmento] (0.556,0.376) -- (0.727,0.580);
  \draw[segmento] (0.794,0.462) -- (0.915,0.431);

  % ---- ordine di disegno delle tre regioni -------------------------------
  \node[ordine] at (0.750,0.910) {1};
  \node[ordine] at (0.820,0.335) {2};
  \node[ordine] at (0.140,0.660) {3};

  % ---- l'osservatore ------------------------------------------------------
  \begin{scope}[shift={(-0.250,0.050)}]
    % occhio stilizzato: due archi che formano una lente, con la pupilla
    \draw[occhio] (-4.6mm,0mm) .. controls (-2.4mm, 3.5mm) and ( 2.4mm, 3.5mm)
               .. ( 4.6mm,0mm) .. controls ( 2.4mm,-3.5mm) and (-2.4mm,-3.5mm)
               .. cycle;
    \fill[black] (0mm,0mm) circle (1.15mm);
    % le tre direzioni di vista, a ventaglio verso la scena
    \draw[raggio] (55:6.4mm) -- (55:21mm);
    \draw[raggio] (33:6.4mm) -- (33:21mm);
    \draw[raggio] (14:6.4mm) -- (14:21mm);
  \end{scope}
  \node[anchor=north,yshift=-5.0mm] at (-0.250,0.050) {osservatore};

  % ---- legenda ------------------------------------------------------------
  \node[anchor=east] at (0.452,-0.092) {ordine di disegno:};
  \node[ordine] (o1) at (0.520,-0.092) {1};
  \node[ordine] (o2) at (0.630,-0.092) {2};
  \node[ordine] (o3) at (0.740,-0.092) {3};
  \draw[frecciaord] (o1) -- (o2);
  \draw[frecciaord] (o2) -- (o3);

\end{tikzpicture}

\caption{La scena, con i suoi segmenti, e l'osservatore. I tagli della BPP suddividono la
scena in regioni; l'albero della partizione fornisce, per ogni posizione dell'osservatore,
un ordine di disegno dal più lontano al più vicino.}
\end{figure}

\section{Autopartizioni e l'algoritmo \textsc{RandAuto}}

\begin{definizione}[autopartizione]\label{def:autopartizione}
Un'\emph{autopartizione} è una BPP tale che ogni retta $\ell_i$ che compare nei suoi nodi
interni $(r_i,\ell_i)$ è una delle rette su cui giace un segmento di $S$, o una parte di esso.
\end{definizione}

Per $s\in S$ scriviamo $\ell(s)$ per la retta su cui giace $s$, cioè il prolungamento di $s$
a tutta la retta.

\begin{figure}[H]
\centering
% p62-autopartizione -- il segmento s giace sulla retta ell = ell(s).
% La retta ell (tratteggiata) e' il prolungamento di s (tratto spesso, con le
% crocette agli estremi).  Tagliando la regione corrente con ell si ottengono
% le due regioni r_ell e r'_ell: s sta esattamente sulla loro frontiera comune
% e va assegnato all'una oppure all'altra.  Le due frecce, perpendicolari a
% ell e uscenti dall'estremo superiore di s, indicano proprio questa scelta.
\begin{tikzpicture}[
    font=\footnotesize,
    >={Stealth[length=4pt,width=3pt]},
    retta/.style={line width=0.5pt,dash pattern=on 2.2pt off 2.4pt},
    segmento/.style={line width=1.4pt},
    freccia/.style={->,line width=0.6pt},
  ]

  % ---- la retta ell : direzione (6.5,2.9), cioe' circa 24 gradi -----------
  \coordinate (P0) at (0,0.15);      % estremo sinistro del tratteggio
  \coordinate (P1) at (6.5,3.05);    % estremo destro del tratteggio
  \coordinate (A)  at (2.145,1.107); % estremo inferiore di s   (t = 0.33)
  \coordinate (B)  at (4.290,2.064); % estremo superiore di s   (t = 0.66)

  \draw[retta] (P0) -- (P1);
  \node[inner sep=1pt] at (6.78,3.16) {$\ell$};

  % ---- il segmento s -----------------------------------------------------
  \draw[segmento] (A) -- (B);
  \foreach \P in {A,B}{%
    \draw[line width=0.6pt] ($(\P)+(-1.1mm,-1.1mm)$) -- ($(\P)+(1.1mm,1.1mm)$);
    \draw[line width=0.6pt] ($(\P)+(-1.1mm, 1.1mm)$) -- ($(\P)+(1.1mm,-1.1mm)$);
  }
  \node[inner sep=1pt] at (3.05,1.99) {$s$};

  % ---- le due direzioni possibili, perpendicolari a ell ------------------
  % normale unitaria a ell:  n = (-0.4074, 0.9130)
  \draw[freccia] (4.209,2.247) -- (3.923,2.886);   % verso la regione r_ell
  \draw[freccia] (4.372,1.881) -- (4.657,1.242);   % verso la regione r'_ell

  % ---- le due regioni ----------------------------------------------------
  \node[inner sep=1pt] at (2.05,2.45) {$r_\ell$};
  \node[inner sep=1pt] at (4.95,1.08) {$r'_\ell$};

\end{tikzpicture}

\caption{Il segmento $s$ giace sulla retta $\ell=\ell(s)$: tagliando con $\ell$, $s$ finisce
sulla frontiera comune alle due regioni $r_\ell$ e $r'_\ell$ e va assegnato all'una oppure
all'altra.}
\end{figure}

\begin{osservazione}\label{oss:frontiera}
Definita la retta $\ell$ su cui giace $s$, il segmento $s$ si trova esattamente sulla
frontiera comune alle due regioni in cui $\ell$ divide la regione corrente. Sceglieremo
allora di collocare $s$ nella regione a destra oppure in quella a sinistra, a seconda di
dove conviene: per esempio nella regione con meno porzioni di segmenti. In particolare
$\ell(s)$ non taglia mai $s$, e quindi in un'autopartizione un segmento può essere spezzato
soltanto dalle rette \emph{degli altri} segmenti.
\end{osservazione}

L'algoritmo randomizzato che costruisce un'autopartizione è il seguente.

\begin{algorithm}[H]
\caption{\textsc{RandAuto}}
\begin{algorithmic}[1]
\Require $S=\set{s_1,\dots,s_n}$, segmenti a due a due non intersecanti.
\Ensure $P_\pi$, un'autopartizione binaria del piano (\emph{binary planar autopartition})
di $S$.
\State scegli uniformemente a caso una permutazione $\pi$ di $\set{1,\dots,n}$;
\While{esiste una regione $r$ che contiene più di una porzione di segmento}
    \State taglia $r$ usando $\ell(s_i)$, dove $s_i$ è il primo segmento nell'ordine dato
    da $\pi$ tale che $s_i$ (o una sua porzione) sia contenuto in $r$.
\EndWhile
\end{algorithmic}
\end{algorithm}

Tutte le scelte casuali sono concentrate nel passo~1: una volta estratta $\pi$, il resto
dell'esecuzione è completamente deterministico, e per questo la partizione prodotta si indica
con $P_\pi$.

\section{Analisi: la dimensione attesa di \texorpdfstring{$P_\pi$}{P pi}}

\scan{63}

\begin{teorema}[dimensione attesa di \textsc{RandAuto}]\label{thm:randauto-size}
Il valore atteso della dimensione dell'autopartizione prodotta da \textsc{RandAuto} su un
insieme $S$ di segmenti a due a due non intersecanti è
\[
  \E\bigl[\abs{P_\pi}\bigr]\;=\;\Oh\bigl(\abs{S}\log\abs{S}\bigr),
\]
dove la media è presa sulla scelta uniforme della permutazione $\pi$.
\end{teorema}

Il teorema riguarda la \emph{dimensione dell'output}, non il tempo di esecuzione.

\begin{esercizio}
Dimostrare che il \emph{tempo} di esecuzione di \textsc{RandAuto} su $S$ è, anch'esso in
valore atteso, $\Oh(n^2\log n)$, posto $n=\abs{S}$.
\end{esercizio}

\begin{corollario}[metodo probabilistico]\label{cor:esiste-autopartizione}
Dato $S$, esiste \emph{sempre} una BPP di $S$ di dimensione $\Oh(\abs{S}\log\abs{S})$.
\end{corollario}

\begin{dimostrazione}[del corollario]
Per il teorema esiste una costante $c$ tale che $\E[\abs{P_\pi}]\le c\,n\log n$, con
$n=\abs{S}$. Una variabile aleatoria non può assumere ovunque valori strettamente maggiori
della propria media: esiste quindi almeno una permutazione $\pi_0$ di $\set{1,\dots,n}$ con
$\abs{P_{\pi_0}}\le c\,n\log n$. La partizione $P_{\pi_0}$ è la BPP cercata (anzi, è
addirittura un'autopartizione). Si noti che l'argomento è puramente esistenziale: non
fornisce alcun modo di \emph{esibire} $\pi_0$.
\end{dimostrazione}

\bigskip
\textbf{Dimostrazione del Teorema~\ref{thm:randauto-size}.}
L'\emph{intuizione} è la seguente: dobbiamo stimare quanti tagli di segmenti vengono eseguiti
in un generico \emph{run} dell'algoritmo. Servono due nozioni preliminari.

\begin{definizione}[$\mathrm{index}$]\label{def:index}
Dati $u,v\in S$, definiamo
\begin{align*}
  \mathrm{index}(u,v)&=i
    &&\text{sse } \ell(u) \text{ interseca } i-1 \text{ segmenti prima di intersecare } v,\\
  \mathrm{index}(u,v)&=\infty
    &&\text{sse } \ell(u) \text{ non interseca } v.
\end{align*}
Il conteggio va inteso percorrendo $\ell(u)$ a partire da $u$ nella direzione di $v$; dunque
$\mathrm{index}(u,v)=1$ significa che $v$ è il primo segmento che $\ell(u)$ incontra uscendo
da $u$.
\end{definizione}

\begin{figure}[H]
\centering
% randauto-index -- la definizione di index(u,v) = i.
% Il tratto spesso a sinistra e' il segmento u (con i trattini agli estremi);
% il tratteggio che lo prolunga e' la retta ell(u).  Percorrendo ell(u) a
% partire da u si incontrano i segmenti u_1, u_2, ..., u_{i-1} -- contati
% dalla graffa, i-1 in tutto -- e solo dopo si incontra v.
\begin{tikzpicture}[
    scale=0.9,
    font=\small,
    retta/.style={line width=0.5pt,dash pattern=on 3.2pt off 3.2pt},
    segU/.style={line width=1.3pt},
    taglia/.style={line width=0.7pt},
  ]

  % ---- il segmento u e la retta ell(u) -----------------------------------
  \draw[segU] (0,0) -- (2.2,0);
  \draw[segU] (0,-0.15) -- (0,0.15);
  \draw[segU] (2.2,-0.15) -- (2.2,0.15);
  \draw[retta] (2.4,0) -- (12,0);

  \node[inner sep=1pt] at (1.1,0.42) {$u$};
  \node[inner sep=1pt] at (11.2,-0.45) {$\ell(u)$};

  % ---- i segmenti attraversati da ell(u) ---------------------------------
  \draw[taglia] (3.35,0.55) -- (3.80,-0.55);
  \node[inner sep=1pt] at (3.16,0.86) {$u_1$};

  \draw[taglia] (4.90,-0.55) -- (5.50,0.60);
  \node[inner sep=1pt] at (5.72,0.90) {$u_2$};

  \foreach \x in {6.4,6.7,7.0}{\fill (\x,-0.22) circle (1.3pt);}

  \draw[taglia] (8.30,-0.55) -- (8.50,0.60);
  \node[inner sep=1pt] at (8.66,0.90) {$u_{i-1}$};

  \draw[taglia] (9.85,-0.55) -- (10.00,0.78);
  \node[inner sep=1pt] at (10.22,1.05) {$v$};

  % ---- la graffa che conta gli i-1 segmenti intermedi (v escluso) --------
  \draw[line width=0.5pt,decorate,
        decoration={brace,amplitude=6pt,mirror,raise=1pt}]
       (3.35,-0.75) -- (8.55,-0.75);
  \node[inner sep=1pt] at (5.95,-1.42) {$i-1$};

\end{tikzpicture}

\caption{$\mathrm{index}(u,v)=i$: percorrendo $\ell(u)$ a partire da $u$ si incontrano i
segmenti $u_1,u_2,\dots,u_{i-1}$ ($i-1$ in tutto) prima di incontrare $v$.}
\label{fig:randauto-index}
\end{figure}

\begin{definizione}[notazione $u\dashv v$]\label{def:dashv}
Scriviamo
\[
  u\dashv v
\]
se $\ell(u)$ taglia $v$ durante l'esecuzione di \textsc{RandAuto}. È l'evento che corrisponde
all'aggiunta di una porzione di segmento in $P_\pi$.
\end{definizione}

\begin{osservazione}\label{oss:primo-in-pi}
Sia $\mathrm{index}(u,v)=i$ e siano $u_1,\dots,u_{i-1}$ i segmenti che $\ell(u)$ interseca
prima di $v$. Allora $u\dashv v$ si verifica \emph{solo se} $u$ occorre in $\pi$ prima di
$u_1,\dots,u_{i-1},v$. Se infatti uno degli $u_j$ venisse scelto prima di $u$, la retta
$\ell(u_j)$ attraverserebbe il tratto di $\ell(u)$ compreso fra $u$ e $v$, separandoli in due
regioni distinte, e $\ell(u)$ non potrebbe più tagliare $v$; se invece fosse $v$ a precedere
$u$, allora $\ell(v)$ verrebbe usata come taglio prima di $\ell(u)$ e $v$ si troverebbe sulla
frontiera della propria regione (Osservazione~\ref{oss:frontiera}), dove $\ell(u)$ non può
più spezzarlo.

Poiché $\pi$ è scelta uniformemente fra le $n!$ permutazioni, ciascuno degli $i+1$ segmenti
di $\set{u,u_1,\dots,u_{i-1},v}$ ha la stessa probabilità di precedere tutti gli altri. Che
$u$ sia il primo è dunque un evento di probabilità esattamente $\dfrac{1}{i+1}$.
\end{osservazione}

\scan{64}

Introduciamo, per ogni coppia di segmenti distinti $u,v\in S$, la \emph{variabile indicatrice}
dell'evento $u\dashv v$:
\[
  C_{u,v}\;\defeq\;\indic{u\dashv v}\;=\;
  \begin{cases}
    1 & \text{se }\ell(u)\text{ taglia }v\text{ durante l'esecuzione di \textsc{RandAuto}},\\[2pt]
    0 & \text{altrimenti.}
  \end{cases}
\]

Dall'Osservazione~\ref{oss:primo-in-pi} segue immediatamente
\[
  \E\bigl[C_{u,v}\bigr]\;=\;\Pr\bigl[u\dashv v\bigr]\;\le\;\frac{1}{\mathrm{index}(u,v)+1}.
\]
Si noti che la condizione dell'osservazione è soltanto necessaria: per questo si ottiene una
disuguaglianza e non un'uguaglianza.

\begin{osservazione}
Se $\mathrm{index}(u,v)=\infty$, cioè se $\ell(u)$ non interseca affatto $v$, il membro destro
va letto come $0$: in tal caso l'evento $u\dashv v$ è impossibile e $C_{u,v}=0$ con certezza.
\end{osservazione}

\paragraph{Dimensione attesa di $P_\pi$.}
La partizione parte dagli $n=\abs{S}$ segmenti e ogni taglio \emph{effettivo} spezza un
segmento in due, aggiungendo esattamente un frammento (e quindi una foglia occupata). Perciò
la dimensione attesa di $P_\pi$ è $n$ più il numero atteso di tagli:
\begin{align*}
  \E\bigl[\abs{P_\pi}\bigr]
    &\;=\; n+\E\Bigl[\;\sum_{u\in S}\;\sum_{v\in S,\;v\neq u} C_{u,v}\Bigr]\\[2pt]
    &\;=\; n+\sum_{u\in S}\;\sum_{v\neq u}\E\bigl[C_{u,v}\bigr]\\[2pt]
    &\;=\; n+\sum_{u\in S}\;\sum_{v\neq u}\Pr\bigl[u\dashv v\bigr]\\[2pt]
    &\;\le\; n+\sum_{u\in S}\;\sum_{v\neq u}\frac{1}{\mathrm{index}(u,v)+1}.
\end{align*}
Il secondo passaggio è la sola \emph{linearità del valore atteso}: le $C_{u,v}$ non sono
affatto indipendenti, ma qui l'indipendenza non serve.

\begin{osservazione}
Fissato $u$, ci sono \textbf{al più due} segmenti $v$ tali che $\mathrm{index}(u,v)=i$: uno a
destra e uno a sinistra di $u$ lungo $\ell(u)$. Infatti percorrendo $\ell(u)$ a partire da $u$
le direzioni possibili sono solo due e, in ciascuna di esse, il segmento preceduto da
esattamente $i-1$ segmenti è unico.
\end{osservazione}

Raggruppando allora gli addendi della somma interna secondo il valore $i=\mathrm{index}(u,v)$,
che varia fra $1$ e $n-1$ (gli altri segmenti sono $n-1$), si ottiene, per ogni $u\in S$
fissato,
\[
  \sum_{v\neq u}\frac{1}{\mathrm{index}(u,v)+1}\;\le\;\sum_{i=1}^{n-1}\frac{2}{i+1}.
\]

Sostituendo nella stima precedente:
\begin{align*}
  \E\bigl[\abs{P_\pi}\bigr]
    &\;\le\; n+\sum_{u\in S}\;\sum_{i=1}^{n-1}\frac{2}{i+1}
     \;=\; n+2n\sum_{i=1}^{n-1}\frac{1}{i+1}
     \;=\; n+2n\bigl(H_n-1\bigr)\\[2pt]
    &\;\le\; n+2n\,H_n\;\in\;\Oh(n\log n),
\end{align*}
dove $H_n=\sum_{k=1}^{n}\frac{1}{k}=\ln n+\Th(1)$ è l'$n$-esimo numero armonico.

Questo è esattamente l'enunciato del teorema: il valore atteso della dimensione
dell'autopartizione prodotta da \textsc{RandAuto} su un insieme $S$ di $n$ segmenti a due a
due non intersecanti è $\Oh(n\log n)$.
\hfill$\square$

\begin{osservazione}
Il risultato è di Paterson e Yao, e il fattore logaritmico non è un artefatto dell'analisi:
si sa che esistono insiemi di $n$ segmenti in cui \emph{ogni} autopartizione ha
$\Om(n\log n/\log\log n)$ foglie, e dunque per questa via non si arriva a $\Oh(n)$.
\end{osservazione}

\chapter{Random walk e algoritmo randomizzato per 2-SAT}

\section{La legge del disordine}

\scan{65}
Il riferimento è il libro divulgativo di George Gamow \emph{One, Two, Three\dots\ Infinity}
(1947), e in particolare il capitolo dedicato alla \emph{legge del disordine} (\emph{The Law
of Disorder}): anche da un comportamento completamente disordinato si riesce a estrarre una
legge quantitativa precisa.

L'esempio classico è il \emph{drunkard's walk}: un ubriaco $u$ parte da un lampione al centro
di una piazza e cammina senza meta, cambiando direzione a ogni passo in modo del tutto
casuale. È il prototipo di \textbf{random walk}; qui lo consideriamo in due dimensioni, su una
griglia.

\begin{figure}[H]
  \centering
  % ubriaco-griglia -- pag. 65
% Il "drunkard's walk" di Gamow: il piano xy visto in prospettiva (parallelogramma),
% i due lati obliqui fanno da assi cartesiani; al centro l'ubriaco u accanto al
% lampione da cui parte, e attorno il cammino irregolare (ellisse + curva tratteggiata).
% Il callout in alto a destra nomina la cosa: random walk, 2D su una griglia.
\begin{tikzpicture}[
    >=Latex,
    scale=0.95,
    font=\footnotesize,
    piano/.style={draw=black,line width=0.9pt},
    filo/.style={draw=black!75,line width=0.6pt},
    tratto/.style={draw=black!60,line width=0.7pt}]

  % ---- il piano (la piazza) in prospettiva ---------------------------------
  \coordinate (A) at (0.00,0.00);   % basso a sinistra
  \coordinate (B) at (4.00,0.00);   % basso a destra
  \coordinate (C) at (5.90,2.30);   % alto a destra
  \coordinate (D) at (1.90,2.30);   % alto a sinistra

  \draw[piano] (A) -- (B);
  \draw[piano] (D) -- (C);
  % lato AD prolungato oltre D: e' l'asse y
  \draw[piano,->] (A) -- (2.22,2.69);
  % lato BC con la punta di freccia a un terzo dell'altezza: e' l'asse x
  \draw[piano,postaction={decorate,decoration={markings,
        mark=at position 0.38 with {\arrow{Latex}}}}] (B) -- (C);

  \node[anchor=south east] at (2.16,2.66) {$y$};
  \node[anchor=north]      at (4.08,-0.10) {$x$};

  % ---- l'ubriaco ------------------------------------------------------------
  \draw[filo] (2.35,1.52) circle (0.13);            % testa
  \draw[filo] (2.35,1.39) -- (2.35,0.82);           % busto
  \draw[filo] (2.08,0.92) -- (2.35,1.25) -- (2.62,0.92);  % braccia
  \draw[filo] (2.35,0.82) -- (2.14,0.42);           % gambe
  \draw[filo] (2.35,0.82) -- (2.56,0.42);
  \node[anchor=east] at (2.00,1.20) {$u$};

  % ---- il lampione ----------------------------------------------------------
  \draw[filo,line width=1.0pt] (3.10,0.55) -- (3.10,1.76);
  \draw[filo] (3.10,2.32) -- (3.27,2.02) -- (3.10,1.72) -- (2.93,2.02) -- cycle;

  % ---- il cammino: prima l'orbita "regolare"... -----------------------------
  \draw[tratto,rotate around={10:(2.75,0.62)}] (2.75,0.62) ellipse (1.10 and 0.42);

  % ---- ...poi il vagabondaggio disordinato attorno ad essa -------------------
  \draw[tratto,densely dashed] plot[smooth cycle,tension=0.85] coordinates {
      (4.20,0.92) (3.96,1.32) (3.40,1.58) (2.70,1.66) (2.00,1.50) (1.38,1.18)
      (1.15,0.78) (1.35,0.42) (1.90,0.20) (2.60,0.14) (3.32,0.24) (3.96,0.50) };

  % ---- la domanda ------------------------------------------------------------
  \draw[filo,->] (2.45,0.02) to[out=-90,in=170] (3.52,-0.82);
  \node[anchor=west] at (3.60,-0.84) {\emph{arriver\`a a casa?}};

  % ---- il callout -------------------------------------------------------------
  \draw[filo] plot[smooth cycle,tension=0.85] coordinates {
      (8.45,3.32) (8.25,3.74) (7.58,3.94) (6.85,3.92) (6.15,3.70) (5.88,3.32)
      (6.18,2.94) (6.90,2.80) (7.72,2.82) (8.30,3.02) };
  \node at (7.15,3.36) {\textsf{random walk}};
  \draw[filo,->] (6.22,2.92) to[out=235,in=65] (5.96,2.40);
  \node[anchor=west] at (6.20,2.10) {($2$D su una griglia)};

\end{tikzpicture}

  \caption{Random walk in $2$D su una griglia: l'ubriaco $u$ parte dal lampione e vaga sul
  piano $xy$ seguendo un cammino del tutto irregolare. Arriverà a casa?}
\end{figure}

\begin{domanda}
L'ubriaco $u$ si allontanerà \emph{arbitrariamente} dal lampione?
\end{domanda}

La risposta è \textbf{sì}, e il conto che segue dice anche \emph{quanto}: dopo $N$ passi la
distanza dal lampione è dell'ordine di $\sqrt{N}$.

\begin{figure}[H]
  \centering
  % passo-griglia-2d -- pag. 65
% Un passo del random walk sulla griglia 2D: dal punto corrente si parte in una
% delle quattro direzioni cardinali, ciascuna con la stessa probabilita' p (=1/4).
% Le frecce non toccano il punto centrale, come nel manoscritto.
\begin{tikzpicture}[
    >=Latex,
    font=\footnotesize,
    passo/.style={line width=0.7pt,->}]

  % --- posizione corrente ---
  \fill (0,0) circle (1.5pt);

  % --- le quattro direzioni ---
  \draw[passo] (0, 0.20) -- (0, 1.00);
  \draw[passo] (0,-0.20) -- (0,-1.00);
  \draw[passo] ( 0.20,0) -- ( 1.00,0);
  \draw[passo] (-0.20,0) -- (-1.00,0);

  % --- probabilita', tutte uguali ---
  \node[anchor=west]  at ( 0.10, 1.00) {$p$};
  \node[anchor=north] at ( 1.00,-0.08) {$p$};
  \node[anchor=north] at (-1.00,-0.08) {$p$};
  \node[anchor=east]  at (-0.10,-1.00) {$p$};

\end{tikzpicture}

  \caption{Un passo del random walk sulla griglia bidimensionale: dal punto corrente ci si
  sposta in una delle quattro direzioni, ciascuna con la stessa probabilità $p$ (dunque
  $p = 1/4$).}
\end{figure}

Siano $x_i$ e $y_i$ le proiezioni sui due assi dell'$i$-esimo tratto, con segno positivo o
negativo a seconda del verso in cui l'ubriaco si muove lungo quell'asse. Se $R$ è la distanza
dal lampione dopo $N$ passi, per il teorema di Pitagora
\[
  R^{2} \;=\; (x_1 + x_2 + \dots + x_N)^{2} \;+\; (y_1 + y_2 + \dots + y_N)^{2}.
\]

Sviluppiamo il primo quadrato:
\begin{align*}
  (x_1 + x_2 + \dots + x_N)^{2}
    &\;=\; (x_1 + x_2 + \dots + x_N)\,(x_1 + x_2 + \dots + x_N)\\[2pt]
    &\;=\; x_1^{2} + x_2^{2} + x_3^{2} + \dots + x_N^{2}\\[4pt]
    &\qquad\;+\; 2\underbrace{\bigl(x_1x_2 + \dots + x_1x_N + x_2x_3 + \dots + x_{N-1}x_N\bigr)}_{\text{prodotti misti}} .
\end{align*}

\begin{osservazione}
Per ogni termine misto $x_i x_j$, al crescere di $N$ troverò lo stesso termine con segno
opposto con probabilità che tende a $1$: i passi sono indipendenti ed è ugualmente probabile
che la proiezione di un tratto sia positiva o negativa, quindi i prodotti misti si compensano
a due a due. Restano solo i quadrati, che sono tutti positivi.
\end{osservazione}

Dunque
\begin{align*}
  (x_1 + x_2 + \dots + x_N)^{2} &\;\longrightarrow\; x_1^{2} + x_2^{2} + \dots + x_N^{2} \;\approx\; N X^{2},\\
  (y_1 + y_2 + \dots + y_N)^{2} &\;\longrightarrow\; y_1^{2} + y_2^{2} + \dots + y_N^{2} \;\approx\; N Y^{2},
\end{align*}
dove $X^{2}$ e $Y^{2}$ denotano i valori medi di $x_i^{2}$ e di $y_i^{2}$, cioè le lunghezze
\emph{quadratiche} medie delle proiezioni di un passo, rispettivamente sull'asse $x$ e
sull'asse $y$. Sostituendo,
\[
  R \;=\; \sqrt{N\,(X^{2} + Y^{2})}
    \;=\; \sqrt{N}\cdot\underbrace{\sqrt{X^{2} + Y^{2}}}_{=\,1}
    \;=\; \sqrt{N},
\]
perché sulla griglia ogni passo ha lunghezza unitaria, cioè $x_i^{2}+y_i^{2}=1$ a ogni passo,
e dunque $X^{2}+Y^{2}=1$.

\begin{osservazione}
L'argomento della cancellazione è euristico: i prodotti misti non si annullano
algebricamente, quello che si annulla è il loro \emph{valore atteso}. Essendo i passi
indipendenti e a media nulla, $\E[x_i x_j] = \E[x_i]\,\E[x_j] = 0$ per $i \neq j$, e quindi
$\E[R^{2}] = N\bigl(\E[x^{2}] + \E[y^{2}]\bigr)$. Il $\sqrt{N}$ che si ottiene è perciò la
distanza \emph{quadratica media} $\sqrt{\E[R^{2}]}$, non la distanza «più probabile». Il
fatto qualitativo, che è quello che ci servirà, resta: la distanza cresce come $\sqrt{N}$ e
non come $N$.
\end{osservazione}

Il random walk, però, non ci interessa solo come curiosità fisica. Dal punto di vista
algoritmico la proprietà cruciale è un'altra: una passeggiata aleatoria è una \emph{visita} di
un grafo che non richiede di ricordare quali nodi siano già stati visitati, perché bastano la
posizione corrente e una sorgente di bit casuali. Nel resto del capitolo sfrutteremo
esattamente questo per ottenere un algoritmo per 2-SAT che lavora in spazio lineare nel numero
di variabili, laddove l'algoritmo deterministico costruisce e visita esplicitamente un grafo.

\section{SAT, 3-SAT, 2-SAT}

\scan{66}
Richiamiamo gli ingredienti del \emph{calcolo proposizionale}:
\[
\text{calcolo proposizionale}\quad
\begin{cases}
P,\; Q,\; R,\; X,\;\dots & \text{simboli proposizionali, con valori } 0/1,\\[6pt]
\wedge \qquad \vee \qquad \neg & \text{connettivi.}
\end{cases}
\]
Combinando simboli proposizionali e connettivi si ottengono le \emph{formule}, che possono
contenere variabili \emph{libere}; ammettendo anche i quantificatori $\forall$ ed $\exists$
sulle variabili proposizionali, una formula in cui nessuna variabile resta libera si dice
\emph{enunciato}.

\begin{esempio*}
\[
\forall Q\;\forall R\;\exists P\;\bigl[(P\wedge Q)\vee(\neg R\vee Q\vee\neg P)\bigr]
\]
è un enunciato: nessuna variabile resta libera. Formule di questa forma si chiamano
\emph{Quantified Boolean Formulae}; il problema di deciderne la verità si indica con QBF ed è
$\PSPACE$-completo.
\end{esempio*}

\begin{osservazione}
Passare da SAT (che è $\NP$-completo) alla sua versione \emph{quantificata} fa salire la
complessità da $\NP$ a $\PSPACE$: QBF è completo per $\PSPACE$, e sarebbe $\NP$-completo solo
se fosse $\NP = \PSPACE$.
\end{osservazione}

I tre problemi che ci interessano si scrivono allora nella forma seguente, in cui tutte le
variabili sono quantificate \emph{esistenzialmente}:
\[
\left\{
\begin{array}{l@{\qquad}l@{\quad}l}
\text{SAT}   & \exists x_1,\dots,x_n & \varphi\\[10pt]
\text{3-SAT} & \exists x_1,\dots,x_n &
   \displaystyle\bigwedge_{j=1}^{K}\bigl(\ell^{j}_{1}\vee \ell^{j}_{2}\vee \ell^{j}_{3}\bigr)\\[16pt]
\text{2-SAT} & \exists x_1,\dots,x_n &
   \displaystyle\bigwedge_{j=1}^{K}\bigl(\ell^{j}_{1}\vee \ell^{j}_{2}\bigr)
\end{array}
\right.
\]
dove ogni $\ell^{j}_{i}$ è un \emph{letterale}, cioè una variabile proposizionale oppure la
sua negazione, e $K$ è il numero di clausole dell'istanza; come d'uso, in $k$-SAT si richiede
che ogni clausola abbia \emph{al più} $k$ letterali, e i display qui sopra ne mostrano il caso
pieno. Le tre famiglie di istanze sono incluse l'una nell'altra,
\[
\text{2-SAT}\ \subseteq\ \text{3-SAT}\ \subseteq\ \text{SAT},
\]
perché una clausola di due letterali è in particolare una clausola di al più tre letterali, e
una congiunzione di clausole è in particolare una formula. In tutti e tre i casi si tratta del
\textbf{problema della soddisfacibilità} (per formule oppure per enunciati) di un dato tipo.

\paragraph{Versioni quantificate.}
Rimuovendo il vincolo che il prefisso sia puramente esistenziale, e ammettendo un'alternanza
arbitraria di $\forall$ ed $\exists$, si ottengono le corrispondenti versioni quantificate
\[
\text{QBF}\qquad \text{3-QBF}\qquad \text{2-QBF}.
\]

\begin{definizione}[Soddisfacibilità]
Risolvere un problema di soddisfacibilità significa cercare un \emph{assegnamento} di valori
di verità (cioè di valori $0/1$) alle variabili proposizionali che compaiono nella formula, in
modo da renderla vera.
\end{definizione}

\subsection{Forme normali}

Due forme normali ricorrenti per le formule proposizionali:
\begin{itemize}
  \item \textbf{CNF} (\emph{conjunctive normal form}): congiunzione di disgiunzioni di
        letterali;
  \item \textbf{DNF} (\emph{disjunctive normal form}): disgiunzione di congiunzioni di
        letterali.
\end{itemize}

\scan{67}
Il problema SAT si affronta dunque portando la formula in una delle due forme normali:
\[
  \text{SAT}\;
  \begin{cases}
    \;\longrightarrow\;\text{CNF}\\[4pt]
    \;\longrightarrow\;\text{DNF}
  \end{cases}
\]

\begin{osservazione}
Leggendo $\vee$ come somma e $\wedge$ come prodotto, una CNF è un \emph{prodotto di somme}
mentre una DNF è una \emph{somma di prodotti}:
\[
  \underbrace{(a+b)\cdot(c+d)\cdots}_{\text{CNF}}
  \qquad\qquad
  \underbrace{(a\cdot b)+(c\cdot d)+\cdots}_{\text{DNF}}
\]
\end{osservazione}

Una CNF ha la forma
\[
  (\ell_1\vee\ell_2)\wedge(\ell_3\vee\ell_4\vee\ell_5)\wedge\cdots
  \qquad\text{ossia}\qquad
  \bigwedge_{j=1}^{K}\Bigl(\bigvee_{i=1}^{k_j}\ell^{j}_{i}\Bigr),
\]
dove gli $\ell^{j}_{i}$ sono letterali e clausole diverse possono avere lunghezze diverse
($k_1$, $k_2$, \dots).

Per portare una formula in CNF si potrebbe procedere per pura distributività --- la
«moltiplicazione» dell'analogia aritmetica --- ma questa strada va scartata, perché fa
esplodere esponenzialmente il numero di clausole. Poiché però interessa soltanto
l'\emph{equisoddisfacibilità}, e non l'equivalenza logica, conviene invece
\textbf{introdurre nuove variabili}.

\begin{esempio}[da DNF a CNF introducendo una variabile]
Si parta dalla DNF
\[
  (\varphi_1\wedge\varphi_2)\vee(\varphi_3\wedge\varphi_4),
\]
in cui $\varphi_1,\dots,\varphi_4$ sono letterali. Introdotta una variabile fresca $u$ si
scrive
\[
  \bigl((\varphi_1\wedge\varphi_2)\vee u\bigr)\wedge\bigl((\varphi_3\wedge\varphi_4)\vee\neg u\bigr)
\]
e, distribuendo il $\vee$ dentro ciascuna congiunzione,
\[
  \bigl((\varphi_1\vee u)\wedge(\varphi_2\vee u)\bigr)\wedge\bigl((\varphi_3\vee\neg u)\wedge(\varphi_4\vee\neg u)\bigr)
\]
cioè la CNF
\[
  (\varphi_1\vee u)\wedge(\varphi_2\vee u)\wedge(\varphi_3\vee\neg u)\wedge(\varphi_4\vee\neg u).
\]
In notazione aritmetica lo stesso passaggio si legge
\[
  (a\cdot b)+(c\cdot d)\;\rightsquigarrow\;(a+u)\cdot(b+u)\cdot(c+\bar u)\cdot(d+\bar u).
\]
Con più di due termini si ripete il trucco introducendo altrettante variabili fresche.
\end{esempio}

\begin{osservazione}
Le due formule non sono logicamente equivalenti (la seconda parla anche di $u$), ma sono
equisoddisfacibili: se è vera $\varphi_1\wedge\varphi_2$ basta porre $u=0$, se è vera
$\varphi_3\wedge\varphi_4$ basta porre $u=1$; viceversa, in ogni assegnamento che soddisfa la
CNF il valore di $u$ seleziona quale dei due termini deve risultare vero.
\end{osservazione}

\subsection{Da CNF a 3-CNF}

Lo stesso trucco permette di accorciare le clausole:
\[
  \text{CNF}\;\longrightarrow\;\text{3-CNF}.
\]
Una clausola con $k>3$ letterali
\[
  \cdots\wedge(\ell_1\vee\ell_2\vee\cdots\vee\ell_k)\wedge\cdots
\]
\[\Downarrow\]
si sostituisce con la catena
\[
  \cdots\wedge(\ell_1\vee\ell_2\vee u_1)\wedge(\neg u_1\vee\ell_3\vee u_2)\wedge(\neg u_2\vee\ell_4\vee u_3)
  \wedge\cdots\wedge(\neg u_{k-3}\vee\ell_{k-1}\vee\ell_k)\wedge\cdots
\]
in cui $u_1,\dots,u_{k-3}$ sono variabili fresche. Ogni clausola prodotta ha esattamente $3$
letterali, e da una clausola di lunghezza $k$ se ne ottengono $k-2$: la trasformazione è
lineare nella dimensione della formula e preserva la soddisfacibilità.

\subsection{Aspvall--Plass--Tarjan: il grafo delle implicazioni}

Per il caso binario esiste un algoritmo deterministico lineare:
\[
  \Oh(n)\qquad\text{per 2-QBF.}
\]
Il risultato di Aspvall--Plass--Tarjan riguarda le formule booleane quantificate la cui
matrice è in 2-CNF (2-QBF); il caso puramente esistenziale è 2-SAT. L'idea è leggere ogni
clausola binaria come una coppia di implicazioni:
\[
  (\ell_1\vee\ell_2)
  \qquad\text{si legge}\qquad
  \overline{\ell_1}\to\ell_2,\qquad \overline{\ell_2}\to\ell_1 .
\]
Raccogliendo gli archi generati da tutte le clausole si ottiene il \emph{grafo delle
implicazioni} $G$, sul quale si calcolano le componenti fortemente connesse (CFC).

\begin{figure}[H]
\centering
% grafo-implicazioni-2sat -- pag. 67
% Il grafo delle implicazioni G di Aspvall--Plass--Tarjan.  I nodi sono letterali
% (senza cerchio, come nel manoscritto); la clausola (l_1 v l_2) genera i due archi
%       not l_1 --> l_2      e      not l_2 --> l_1 .
% La macchia irregolare rappresenta l'intero grafo G; fuori, a destra, la sigla CFC
% (componenti fortemente connesse), che e' cio' che su G si va a calcolare.
\begin{tikzpicture}[
    >=Latex,
    font=\small,
    bordo/.style={draw=black!70,line width=0.8pt},
    arco/.style={->,line width=0.7pt}]

  % ---- contorno del grafo (curva chiusa irregolare) ------------------------
  \draw[bordo] plot[smooth cycle,tension=0.85] coordinates {
      (-2.80, 0.10) (-2.40, 0.90) (-1.55, 1.30) (-0.40, 1.42)
      ( 0.80, 1.35) ( 1.85, 1.18) ( 2.58, 0.70) ( 2.80, 0.12)
      ( 2.58,-0.45) ( 1.75,-1.00) ( 0.60,-1.30) (-0.10,-1.45)
      (-1.10,-1.25) (-2.05,-0.95) (-2.65,-0.45) };

  % ---- nome del grafo, appoggiato al bordo in alto a destra ----------------
  \node[anchor=south west] at (1.72,0.98) {$G$};

  % ---- clausola (l_1 v l_2): primo arco,  not l_1 --> l_2 ------------------
  \node (nl1) at (-0.60, 0.62) {$\overline{\ell_1}$};
  \node (l2)  at (-1.85, 0.50) {$\ell_2$};
  \draw[arco] (nl1) to[bend left=15] (l2);

  % ---- clausola (l_1 v l_2): secondo arco,  not l_2 --> l_1 ----------------
  \node (nl2) at (0.15,-0.78) {$\overline{\ell_2}$};
  \node (l1)  at (1.45,-0.36) {$\ell_1$};
  \draw[arco] (nl2) to[bend right=15] (l1);

  % ---- cio' che si calcola su G -------------------------------------------
  \node[anchor=west] at (3.05,0.05) {CFC};

\end{tikzpicture}

\caption{Il grafo delle implicazioni $G$: la clausola $(\ell_1\vee\ell_2)$ produce i due archi
$\overline{\ell_1}\to\ell_2$ e $\overline{\ell_2}\to\ell_1$. Su $G$ si calcolano le componenti
fortemente connesse (CFC).}
\end{figure}

\begin{osservazione}
Nel caso puramente esistenziale il criterio è il seguente: la 2-CNF è insoddisfacibile se e
solo se esiste una variabile $x$ tale che $x$ e $\neg x$ cadono nella stessa componente
fortemente connessa di $G$. Poiché le CFC si calcolano con una visita in tempo lineare, si
ottiene l'algoritmo $\Oh(n)$ annunciato; per il caso quantificato (2-QBF) al criterio va
aggiunta una condizione sull'ordine dei quantificatori.
\end{osservazione}

\section{Un algoritmo randomizzato per 2-SAT}

\scan{69}
Riepilogando, per il caso binario abbiamo due algoritmi deterministici:
\[
  \text{CNF}\qquad \text{3-CNF}\qquad \text{2-CNF}
  \;\begin{cases}
    \;\longrightarrow\ \text{Aspvall--Plass--Tarjan (2-QBF)} & \Oh(n)\\[4pt]
    \;\longrightarrow\ \text{risoluzione} & \Oh(n^{2})
  \end{cases}
\]
dove la seconda riga è la chiusura per risoluzione dell'insieme delle clausole binarie. Il
punto debole di entrambi è lo stesso: costruiscono in modo esplicito una struttura di
dimensione almeno lineare nel numero dei letterali --- il grafo delle implicazioni ha un nodo
per letterale, mentre l'insieme dei risolventi può arrivare a $\Th(n^{2})$ clausole.
Possiamo semplificare ulteriormente, e usare la \emph{randomness}.

\scan{68}
Sia data una formula in 2-CNF e se ne consideri una clausola:
\[
  \cdots\;\wedge\;(\ell_1\vee\ell_2)\;\wedge\;\cdots
\]
Si parte da un \emph{assegnamento $f$ di valori di verità} e si itera il passo seguente: se
$(\ell_1\vee\ell_2)$ \emph{non} è soddisfatta da $f$, si fa il \textsc{flip} di \emph{uno dei
due valori di verità}, scelto \emph{a caso} fra i due.

\begin{osservazione}
Supponiamo che la formula di partenza sia soddisfacibile e fissiamo un assegnamento $f^{*}$
che la soddisfa. Se sotto $f$ i due letterali $\ell_1$ ed $\ell_2$ valgono entrambi $\bot$
--- cioè se la clausola $(\ell_1\vee\ell_2)$ non è soddisfatta --- allora $f^{*}$ ne rende
vero almeno uno, e quindi $f$ ed $f^{*}$ differiscono su almeno una delle due variabili
coinvolte. Invertendone una scelta a caso, con probabilità almeno $\tfrac12$ ci si
«avvicina» ad $f^{*}$: precisamente, il numero
\[
  k \;=\; \#\set{i \;:\; f(x_i)=f^{*}(x_i)}
\]
di variabili su cui $f$ coincide con $f^{*}$ aumenta di $1$ con probabilità almeno
$\tfrac12$; altrimenti diminuisce di $1$.
\end{osservazione}

Misurando così la distanza dall'obiettivo, ogni passo dell'algoritmo diventa uno spostamento
di una posizione lungo la linea degli stati $0,1,\dots,n$, dove $n$ è il numero delle
variabili proposizionali: da $k$ si passa a $k-1$ oppure a $k+1$, e a $k+1$ con probabilità
almeno $\tfrac12$.

\begin{figure}[H]
\centering
% random-walk-linea-2sat -- pagg. 68 e 69 (i due schizzi rappresentano la stessa cosa
% e qui sono fusi in un'unica figura).
% La linea degli stati 0,1,...,n: lo stato k conta le variabili su cui l'assegnamento
% corrente f coincide con l'assegnamento soddisfacente fissato f*.  Ogni FLIP sposta
% di +-1, e a k+1 si va con probabilita' almeno 1/2; n e' lo stato di arresto.
\begin{tikzpicture}[
    >=Latex,
    font=\footnotesize,
    asse/.style={line width=0.7pt},
    tacca/.style={line width=0.7pt},
    stato/.style={circle,fill,inner sep=0pt,minimum size=3.2pt},
    salto/.style={->,line width=0.6pt}]

  % ---- asse degli stati ----------------------------------------------------
  \draw[asse,->] (-0.55,0) -- (7.40,0);

  % ---- gli stati -----------------------------------------------------------
  \node[stato] at (0.00,0) {};
  \node[stato] at (0.75,0) {};
  \node[stato] at (1.45,0) {};
  \node[stato] at (2.95,0) {};
  \node[stato] at (3.95,0) {};
  \node[stato] at (4.95,0) {};
  \node[stato] at (6.55,0) {};

  % ---- tacche che attraversano l'asse: origine, stato corrente, arresto ----
  \draw[tacca] (0.00,-0.30) -- (0.00,0.30);
  \draw[tacca] (3.95,-0.26) -- (3.95,0.26);
  \draw[tacca] (6.55,-0.22) -- (6.55,0.22);

  % ---- etichette, tutte sulla stessa linea di base --------------------------
  \node[anchor=north] at (0.00,-0.34) {$0$};
  \node[anchor=north] at (0.75,-0.34) {$1$};
  \node[anchor=north] at (1.45,-0.34) {$2$};
  \node[anchor=north] at (2.15,-0.34) {$\cdots$};
  \node[anchor=north] at (2.95,-0.34) {$k-1$};
  \node[anchor=north] at (3.95,-0.34) {$k$};
  \node[anchor=north] at (4.95,-0.34) {$k+1$};
  \node[anchor=north] at (5.85,-0.34) {$\cdots$};
  \node[anchor=north] at (6.55,-0.34) {$n$};

  % ---- il passo: da k si va a k-1 oppure a k+1 ------------------------------
  \draw[salto] (3.88,0.32) to[out=145,in=55] (3.00,0.18);
  \draw[salto] (4.02,0.32) to[out=35,in=125] (4.90,0.18);
  \node[anchor=west,font=\scriptsize] at (5.05,0.46) {$\ge \tfrac12$};

  % ---- che cosa conta n -----------------------------------------------------
  \node[anchor=north west,align=left,font=\scriptsize,text=black!70]
    at (6.90,-0.26) {$n$ variabili\\proposizionali};

\end{tikzpicture}

\caption{Il passo dell'algoritmo visto come passeggiata aleatoria sulla linea degli stati
$0,1,\dots,n$: dallo stato $k$ si passa a $k-1$ oppure a $k+1$, e a $k+1$ con probabilità
almeno $\tfrac12$. Qui $n$ è il numero delle variabili proposizionali e lo stato $k$ conta le
variabili su cui l'assegnamento corrente $f$ coincide con l'assegnamento soddisfacente
fissato $f^{*}$.}
\end{figure}

Posso quindi vedere la mia ricerca come una \emph{visita randomizzata} di un \emph{grafo
lineare} (cammino): scriviamo $L_m$ per il cammino con $m$ nodi, cosicché la linea degli stati
$0,1,\dots,n$ è $L_{n+1}$.

\begin{figure}[H]
\centering
% grafo-lineare-L4 -- il grafo lineare (cammino) L_4 (quaderno, pag. 69).
% Cammino orizzontale con 4 nodi e 3 archi di lunghezza uguale, senza etichette
% sui nodi; a destra, staccata e leggermente piu' in basso, l'etichetta $L_4$.
\begin{tikzpicture}[
    font=\footnotesize,
    nodo/.style={circle,draw,line width=0.7pt,fill=white,
                 inner sep=0pt,minimum size=4.5pt},
    arco/.style={line width=0.7pt}]

  % --- i quattro nodi del cammino, equispaziati ---
  \foreach \i in {0,1,2,3}{
    \node[nodo] (v\i) at (\i*1.1,0) {};
  }

  % --- i tre archi ---
  \draw[arco] (v0) -- (v1);
  \draw[arco] (v1) -- (v2);
  \draw[arco] (v2) -- (v3);

  % --- etichetta del grafo ---
  \node[anchor=west] at (4.3,-0.22) {$L_4$};

\end{tikzpicture}

\caption{Il grafo lineare $L_4$: il cammino con $4$ nodi e $3$ archi.}
\end{figure}

Il guadagno rispetto agli algoritmi deterministici è tutto nello spazio: basta infatti tenere
in memoria l'assegnamento corrente, un bit per variabile, e non serve alcuna traccia dei nodi
già visitati.
\begin{center}
\fbox{\textsc{spazio}\quad $\Oh(n)$ bit}
\end{center}

Per il tempo, vedremo che il numero atteso di passi necessari a percorrere tutta la linea è
\[
  \Oh(n^{2}),
\]
e che il numero di passi che rende alta la probabilità di aver visitato tutto il grafo è di
poco superiore a questo. La giustificazione arriverà nella Sezione~\ref{sec:cover-lineare},
con il bound sul \emph{cover time} di un random walk su grafi.

\begin{osservazione}
Ogni volta che abbiamo un algoritmo che lavora su un grafo, nel random walk ritroviamo quello
stesso algoritmo incarnato in forma probabilistica.
\end{osservazione}

\section{Hitting time e cover time}

Per capire quanti passi servono, conviene svolgere due esercizi propedeutici sul grafo
completo $K_n$, dove i conti si fanno in modo esplicito.

\subsection{Hitting time su \texorpdfstring{$K_n$}{Kn}}

\scan{69}
\begin{esercizio}[Hitting time su $K_n$]\label{ex:hitting-kn}
Calcolare il valore atteso del numero di passi in $K_n$ (il \emph{complete graph}) per andare
da $u$ a $v$: è l'\emph{hitting time}.
\end{esercizio}

Il punto di partenza del conto è il numero $n-1$: tanti sono i vicini di ogni nodo di $K_n$,
e fra essi si sceglie uniformemente ad ogni passo. Le mosse possibili sono $n-1$ e non $n$
perché nei grafi non ci sono self-loop: un arco congiunge sempre due nodi distinti. Dunque
\begin{align*}
  \Pr[\,\text{trovo } v \text{ in } 1 \text{ passo}\,]  &= \frac{1}{n-1},\\[4pt]
  \Pr[\,\text{trovo } v \text{ in } 2 \text{ passi}\,]  &= \Bigl(1-\frac{1}{n-1}\Bigr)\frac{1}{n-1},\\[4pt]
  &\;\;\vdots\\[2pt]
  \Pr[\,\text{trovo } v \text{ in } t \text{ passi}\,]  &= \Bigl(1-\frac{1}{n-1}\Bigr)^{t-1}\frac{1}{n-1}.
\end{align*}

Si ricade quindi in uno schema di Bernoulli con probabilità di successo
\[
  p \;=\; \frac{1}{n-1}.
\]

\begin{osservazione}
Bernoulliane sono le singole prove (il singolo passo del cammino, che «riesce» se porta
proprio in $v$); la variabile aleatoria di cui si vuole il valore atteso --- il numero di
passi necessari per raggiungere $v$ --- è invece \emph{geometrica} di parametro $p$.
\end{osservazione}

\scan{70}
Detta $X$ la variabile aleatoria geometrica che conta il numero di tentativi fino al primo
successo --- qui: il numero di passi necessari a raggiungere $v$ --- quello che serve è dunque
il suo valore atteso:
\[
  \E[X] \;=\; \sum_{t=1}^{\infty} t\,(1-p)^{t-1}\,p ,
\]
dove l'indice $t$ conta il numero di passi. Il calcolo si ottiene derivando rispetto a $p$ la
serie geometrica; nel seguito l'apice denota $\frac{d}{dp}$:
\begin{align*}
  \sum_{t=1}^{\infty} t\,(1-p)^{t-1}\,p
    &= p\left(-\sum_{t=1}^{\infty}(1-p)^{t}\right)'
     \;=\; -p\left(\left(\sum_{t=0}^{\infty}(1-p)^{t}\right)-1\right)' \\[4pt]
    &= -p\left(\frac{1}{1-1+p}-1\right)'
     \;=\; -p\left(\frac{1-p}{p}\right)'
     \;=\; -p\left(-\frac{1}{p^{2}}\right)
     \;=\; \frac{1}{p} .
\end{align*}

Nel caso dell'Esercizio~\ref{ex:hitting-kn} si ha $p=\frac{1}{n-1}$: il valore atteso del numero di
passi per andare da $u$ a $v$ in $K_n$ (l'\emph{hitting time}) è quindi $n-1$.

\subsection{Cover time su \texorpdfstring{$K_n$}{Kn}}

\begin{esercizio}[Cover time su $K_n$]\label{ex:cover-kn}
Calcolare il valore atteso del numero di passi necessari per visitare tutto $K_n$, partendo da
un vertice $u$. Questo valore atteso è il \emph{cover time} di $K_n$.
\end{esercizio}

Sia $\tau_i$ la variabile aleatoria che conta il numero di passi effettuati per visitare $i$
nodi distinti, contando fra questi il vertice di partenza $u$. Per costruzione i suoi valori
crescono:
\[
  \tau_1 = 0 \;<\; \tau_2 = 1 \;<\; \tau_3 \;<\; \tau_4 \;<\; \cdots \;<\; \tau_n
\]
(il primo passo raggiunge sempre un nodo nuovo, quindi $\tau_2 = 1$ con certezza), e
$\tau_n$ è il numero, aleatorio, di passi necessari a coprire tutto il grafo: il cover time è
il suo valore atteso $\E[\tau_n]$.

La differenza $\tau_{i+1}-\tau_i$ misura il numero di passi necessari per visitare
l'$(i+1)$-esimo nodo distinto. Ragionando come nell'Esercizio~\ref{ex:hitting-kn}, la probabilità di
successo --- cioè di raggiungere in un passo un nodo non ancora visitato --- quando $i$ nodi
sono già stati visitati vale
\[
  p_i \;=\; \frac{n-i}{n-1}
\]
(dei $n-1$ vicini del nodo corrente, esattamente $n-i$ non sono ancora stati visitati: dei
nodi già visti, uno è quello su cui ci si trova e non è vicino di sé stesso). Conviene allora
calcolare $\E[\tau_{i+1}-\tau_i]$, in quanto, per telescopia e linearità del valore atteso,
\[
  \E[\tau_n] \;=\; \sum_{i=1}^{n-1} \E\bigl[\tau_{i+1}-\tau_i\bigr] .
\]

Ogni incremento è distribuito geometricamente con parametro $p_i$, quindi si ricade nel
calcolo appena svolto:
\[
  \E\bigl[\tau_{i+1}-\tau_i\bigr] \;=\; \sum_{t=1}^{\infty} t\,p_i\,(1-p_i)^{t-1}
  \;=\; \frac{1}{p_i} \;=\; \frac{n-1}{\,n-i\,} .
\]

\scan{71}
Sommando i contributi dei singoli incrementi si ottiene il valore atteso del numero di passi
necessari a visitare tutto $K_n$:
\[
  \E[\tau_n] \;=\; \sum_{i=1}^{n-1} \frac{1}{p_i}
  \;=\; \sum_{i=1}^{n-1} \frac{n-1}{n-i}
  \;=\; (n-1) \sum_{j=1}^{n-1} \frac{1}{j}
  \;=\; \underbrace{(n-1)\,H_{n-1}}_{\Oh(n\log n)},
\]
dove $H_n = \sum_{j=1}^{n} 1/j$ è l'$n$-esimo numero armonico. Il cover time di $K_n$ cresce
dunque come $\Th(n \log n)$: è il classico conto del \emph{coupon collector}.

\begin{osservazione}
Con la convenzione qui adottata il vertice di partenza conta come già visitato, per cui
$\tau_1 = 0$ e la somma parte da $i=1$. Se invece si fa partire la catena da $\tau_0 = 0$ e
si somma da $i=0$ --- come è naturale scrivere di getto --- il termine $i=0$ è spurio, perché
$p_0 = n/(n-1)$ non è una probabilità: il primo passo porta con certezza su un nodo non
ancora visitato. Quell'addendo vale $\frac{n-1}{n} < 1$, cioè meno di un passo, e la somma
diventa $(n-1)H_n \approx n H_n$, la forma in cui di solito si scrive il coupon collector:
l'ordine di grandezza $\Th(n\log n)$ non cambia.
\end{osservazione}

\section{Quando fermarsi: la disuguaglianza di Markov}

Il conto appena fatto è però un valore \emph{medio}, e questo è esattamente il punto delicato
della visita randomizzata:
\begin{center}
  Cover $\iff$ Visita $\iff$ \dots\ \emph{valore atteso}?
\end{center}
La copertura del grafo da parte del random walk fa le veci di una visita, ma di essa non
conosciamo il costo esatto: sappiamo soltanto quanti passi servono \emph{in media}. Nasce
allora la domanda:
\begin{center}
  \textbf{quando mi fermo?}
\end{center}
La risposta viene dalla disuguaglianza di Markov: è mediante Markov che siamo in grado di
stimare l'errore che si commette fermandosi dopo un numero prefissato di passi.

\begin{teorema}[Disuguaglianza di Markov]\label{thm:markov}
  Sia $X$ una variabile aleatoria a valori interi non negativi e sia $a > 0$. Allora
  \[
    \Pr[X \ge a] \;\le\; \frac{\E[X]}{a}.
  \]
\end{teorema}

\begin{dimostrazione}
  \begin{align*}
    \E[X] \;=\; \sum_{t=0}^{\infty} t \Pr[X = t]
      &\;\ge\; \sum_{t \ge a} t \Pr[X = t] \\
      &\;\ge\; \sum_{t \ge a} a \Pr[X = t]
      \;=\; a \sum_{t \ge a} \Pr[X = t]
      \;=\; a \Pr[X \ge a].
  \end{align*}
  La prima disuguaglianza vale perché si stanno buttando via addendi non negativi ($X$ è a
  valori $\ge 0$), la seconda perché nel dominio di somma si ha $t \ge a$. Dividendo per
  $a > 0$ si ottiene l'enunciato.
\end{dimostrazione}

Sulla base di quanto visto posso ridurre a piacere l'errore della mia visita, con un
trade-off fra tempo e correttezza. Basta leggere la disuguaglianza così:
\begin{itemize}
  \item $X$ è la complessità (numero di passi) della mia visita \emph{corretta}, cioè di
        quella che copre davvero tutto il grafo;
  \item fisso $a$ uguale al numero di passi eseguiti dalla mia visita, cioè al numero di
        passi che decido di spendere.
\end{itemize}
Allora l'evento $X \ge a$ è esattamente l'evento «mi fermo prima di aver coperto il grafo», e
Markov ne stima la probabilità:
\[
  \underbrace{\Pr[X \ge a]}_{\text{probabilità di errore}} \;\le\; \frac{\E[X]}{a}.
\]
Aumentando $a$ (cioè lavorando di più) la stima dell'errore scende quanto si vuole.

\section{Amplificazione dell'errore e cover time del grafo lineare}
\label{sec:cover-lineare}

\scan{72}
\begin{esempio}
Riprendiamo la visita randomizzata di cui sopra: $X$ è il numero di passi che essa impiega per
terminare \emph{correttamente} (cioè per visitare tutti i nodi) e supponiamo di sapere che
\[
  \E[X] \;=\; n^2 .
\]
Se decido di fermare la visita dopo $a$ passi, l'evento $X \ge a$ è esattamente l'evento «la
visita non era ancora terminata», cioè l'evento di errore; la disuguaglianza di Markov lo
limita immediatamente:
\[
  \Pr\Bigl[\;\underbrace{X \ge a}_{\text{errore}}\;\Bigr] \;\le\; \frac{\E[X]}{a}.
\]
Istanziando:
\begin{align*}
  a = 2n^2    &\;\;\Longrightarrow\;\; \Pr[X \ge a] \;\le\; \frac{n^2}{2n^2}   \;=\; \frac{1}{2},\\[2pt]
  a = 100\,n^2 &\;\;\Longrightarrow\;\; \Pr[X \ge a] \;\le\; \frac{n^2}{100\,n^2} \;=\; \frac{1}{100}.
\end{align*}
\end{esempio}

Allungando la singola visita fino ad $a = c\,n^2$ passi la probabilità di errore scende solo
come $1/c$. Conviene invece \emph{amplificare per ripetizione}: se eseguo $k$ visite
indipendenti da $2n^2$ passi ciascuna e tengo quella che ha visitato il massimo numero di
nodi, sbaglio solo se \emph{tutte} e $k$ le visite falliscono, e quindi
\[
  \Pr[\text{errore}] \;\le\; \Bigl(\frac{1}{2}\Bigr)^{k} \;=\; \frac{1}{2^{k}},
\]
cioè si guadagna esponenzialmente in $k$ a fronte di un costo solo lineare in $k$.

Resta da giustificare l'ordine di grandezza $\E[X] = \Oh(n^2)$ usato nell'esempio, che è poi
la stima di tempo annunciata per l'algoritmo randomizzato. Esso discende dal bound generale
sul \emph{cover time} di un random walk su grafi.

\begin{teorema}[Cover time di un random walk su grafi]\label{thm:cover}
  Sia $G$ un grafo connesso con $n$ nodi ed $m$ archi e sia $C(G)$ il suo \emph{cover time},
  cioè il numero atteso di passi che un random walk su $G$ impiega a visitare tutti i nodi.
  Allora
  \[
    C(G) \;\le\; 2m(n-1).
  \]
\end{teorema}

La dimostrazione, che richiede il linguaggio delle catene di Markov, è rimandata al capitolo
sui random walk su grafi (Teorema~\ref{thm:cover-time}).

È proprio questo il caso che interessa: la versione randomizzata di 2-SAT è un random walk
sulla retta $\set{0,1,\dots,n}$, dove la posizione conta quante variabili dell'assegnamento
corrente coincidono con quelle di un assegnamento soddisfacente fissato. Quella retta è il
cammino $L_{n+1}$, che ha $n+1$ nodi ed $m = n$ archi, e il teorema dà
\[
  C(L_{n+1}) \;\le\; 2\,n\,\bigl((n+1)-1\bigr) \;=\; 2n^2 .
\]

\begin{osservazione}
Le costanti vanno però maneggiate con cura. Il bound sul cover time dice $\E[X] \le 2n^2$, e
con $a = 2n^2$ la disuguaglianza di Markov restituirebbe la stima banale $1$. Per 2-SAT,
d'altra parte, non serve coprire tutta la linea: basta \emph{raggiungere} lo stato $n$. Detto
$h_k$ il numero atteso di passi per andare da $k$ ad $n$ nel caso peggiore --- probabilità
esattamente $\tfrac12$ in entrambe le direzioni, con il rimbalzo obbligato in $0$ --- si ha
\[
  h_n = 0,\qquad h_0 = h_1 + 1,\qquad
  h_k = 1 + \tfrac12 h_{k-1} + \tfrac12 h_{k+1}\quad (0<k<n),
\]
sistema risolto da $h_k = n^2 - k^2$. Partendo nel caso peggiore da $k=0$ si ottiene
$\E[X] \le h_0 = n^2$: è il valore usato nell'esempio, e con $a = 2n^2$ passi Markov dà
$\Pr[\text{errore}] \le \tfrac12$ per ogni singola esecuzione, che è esattamente l'ipotesi
dell'amplificazione per ripetizione.
\end{osservazione}

Mettendo insieme i pezzi si ottiene l'algoritmo nella sua forma finale.

\begin{algorithm}[H]
\caption{Algoritmo randomizzato per 2-SAT}
\begin{algorithmic}[1]
\Require una formula $\varphi$ in 2-CNF su $n$ variabili; un intero $k \ge 1$
\Ensure un assegnamento che soddisfa $\varphi$, oppure «insoddisfacibile»
\For{$r = 1,\dots,k$}
  \State $f \gets$ un assegnamento iniziale arbitrario (ad esempio $f(x_i)=1$ per ogni $i$)
  \For{$s = 1,\dots,2n^{2}$}
    \If{$f$ soddisfa $\varphi$}
      \State \Return $f$
    \EndIf
    \State scegli una clausola $(\ell_1\vee\ell_2)$ non soddisfatta da $f$
    \State scegli $\ell \in \set{\ell_1,\ell_2}$ uniformemente a caso
    \State inverti in $f$ il valore della variabile che occorre in $\ell$
  \EndFor
\EndFor
\State \Return «insoddisfacibile»
\end{algorithmic}
\end{algorithm}

\begin{osservazione}
L'errore è \emph{unilaterale}: se l'algoritmo restituisce un assegnamento, questo soddisfa
$\varphi$ con certezza, perché viene esibito; l'unica risposta che può essere sbagliata è
«insoddisfacibile». Poiché ogni \textsc{flip} avvicina all'assegnamento soddisfacente $f^{*}$
con probabilità almeno $\tfrac12$, ciascun round di $2n^{2}$ passi fallisce con probabilità al
più $\tfrac12$; essendo i $k$ round indipendenti, la risposta «insoddisfacibile» è sbagliata
con probabilità al più $1/2^{k}$. Si tratta dunque di un algoritmo \MonteCarlo, e il costo
complessivo è $\Oh(k\,n^{2})$ passi in spazio $\Oh(n)$ bit.
\end{osservazione}

\chapter{Catene di Markov}
\label{cap:catene-markov}

\scan{72}
Gli algoritmi randomizzati visti finora ci hanno portati più volte a studiare una
passeggiata aleatoria: la visita casuale di un grafo, di cui abbiamo stimato il
\emph{cover time}, e l'algoritmo per \textsc{2-cnf-sat}, che è un random walk sulla
retta $\{0,1,\dots,n\}$. La nozione generale di cui tutti questi processi sono casi
particolari è quella di \emph{catena di Markov}, alla quale è dedicato questo
capitolo. Il riferimento bibliografico classico è W.~Feller, \emph{An Introduction to
Probability Theory and Its Applications}, Vol.~I, cap.~XV; per il taglio algoritmico,
Motwani--Raghavan, \emph{Randomized Algorithms}, cap.~6.


\section{Definizione}

\scan{73}
Feller le introduce a partire da una successione di prove ripetute, e ne descrive la
probabilità di un'\emph{intera} sequenza di esiti.

\begin{definizione}[Catena di Markov, nella forma di Feller]\label{def:mc-feller}
  Si consideri una successione di prove (\emph{trials}, esperimenti) i cui esiti
  possibili sono $E_1, E_2, \dots$\,; le prove formano una \emph{catena di Markov} se la
  probabilità che gli esiti si verifichino, nell'ordine, $E_{j_0}, E_{j_1}, \dots,
  E_{j_t}$ si fattorizza come
  \[
    \Pr\bigl\{(E_{j_0}, E_{j_1}, \dots, E_{j_t})\bigr\}
      \;=\; a_{j_0} \cdot P_{j_0 j_1} \cdot P_{j_1 j_2} \cdot P_{j_2 j_3}
            \cdots P_{j_{t-1} j_t},
  \]
  dove
  \begin{itemize}
    \item $a_{j_0}$ è la probabilità che si verifichi $E_{j_0}$, cioè la distribuzione
          iniziale sugli esiti;
    \item $P_{j_{k-1} j_k}$ è la probabilità che si verifichi $E_{j_k}$ \emph{a
          condizione che} si sia verificato $E_{j_{k-1}}$.
  \end{itemize}
\end{definizione}

\paragraph{Memoryless (property).}
Tutto il contenuto della definizione sta nella forma dei fattori: la probabilità del
$k$-esimo esito è condizionata al solo esito immediatamente precedente, non all'intera
storia della successione. Oltre al presente, il passato non porta informazione.

D'ora in avanti chiameremo \emph{stati} gli esiti possibili e denoteremo con $S$
l'insieme che li raccoglie.

\subsection{Catene di Markov discrete}

\scan{75}
L'aggettivo \emph{discreta} si riferisce agli stati: sono «ben distinti», cioè $S$ è un
insieme (per noi \emph{finito}) di stati. In questa forma la definizione si dà
direttamente sulle variabili aleatorie.

\begin{definizione}[Catena di Markov discreta]\label{def:catena-markov}
  Una \emph{catena di Markov} su $S$ è un \emph{processo stocastico} su $S$, cioè una
  successione di variabili aleatorie
  \[
    X_0,\; X_1,\; \dots,\; X_t,\; X_{t+1},\; \dots
  \]
  a valori in $S$, che soddisfa la \emph{proprietà di Markov}: per ogni $t \ge 0$ e per
  ogni scelta di stati $i_0, \dots, i_t, j \in S$
  \[
    \Pr\bigl[X_{t+1} = j \;\big|\; X_0 = i_0,\, \dots,\, X_t = i_t\bigr]
    \;=\;
    \Pr\bigl[X_{t+1} = j \;\big|\; X_t = i_t\bigr].
  \]
\end{definizione}

\begin{osservazione}
  Le due definizioni dicono la stessa cosa. Sviluppando il condizionamento un passo
  alla volta, la proprietà di Markov dà infatti
  \[
    \Pr\bigl[X_0 = i_0,\, X_1 = i_1,\, \dots,\, X_t = i_t\bigr]
    \;=\; a_{i_0}\, P_{i_0 i_1} P_{i_1 i_2} \cdots P_{i_{t-1} i_t},
  \]
  che è esattamente la fattorizzazione della Definizione~\ref{def:mc-feller}, purché i fattori
  $P_{ij}$ non dipendano dall'istante $t$ in cui la transizione avviene: è l'ipotesi di
  \emph{omogeneità nel tempo}, implicita nella scrittura di Feller e assunta d'ora in
  poi. È proprio l'omogeneità che consente di raccogliere tutte le probabilità di
  transizione in un'unica matrice.
\end{osservazione}

\scan{73}
Tutta l'informazione sul passato che serve a predire il futuro è dunque contenuta nello
stato corrente:

\begin{center}
  \emph{Condizionatamente al presente, il futuro non dipende dal passato.}
\end{center}

\subsection{Due domande sulla proprietà di Markov}

\scan{74}
Lo slogan va però maneggiato con cura. Siano $X_0, X_1, X_2, \dots$ una catena di
Markov e siano $A$ e $B$ due sottoinsiemi di stati.

\begin{domanda}
\begin{enumerate}
  \item È vero che
  \[
    \Pr\bigl\{\, X_2 \in B \;\big|\; X_1 = x_1 \,\wedge\, X_0 \in A \,\bigr\}
    \;=\;
    \Pr\bigl\{\, X_2 \in B \;\big|\; X_1 = x_1 \,\bigr\}\;?
  \]
  \item È vero che
  \[
    \Pr\bigl\{\, X_2 \in B \;\big|\; X_1 \in A \,\wedge\, X_0 = x_0 \,\bigr\}
    \;=\;
    \Pr\bigl\{\, X_2 \in B \;\big|\; X_1 \in A \,\bigr\}\;?
  \]
\end{enumerate}
\end{domanda}

La prima uguaglianza è vera; la seconda, in generale, no.

\begin{osservazione}
Nella (1) il condizionamento fissa \emph{esattamente} lo stato presente
($X_1 = x_1$): per la proprietà di Markov l'informazione aggiuntiva sul passato
($X_0 \in A$) è irrilevante e l'uguaglianza vale.

Nella (2), invece, l'evento $X_1 \in A$ non individua un singolo stato presente.
Condizionando e poi decomponendo su quale stato di $A$ sia stato effettivamente
raggiunto,
\begin{align*}
  &\Pr\bigl\{X_2 \in B \mid X_1 \in A \wedge X_0 = x_0\bigr\} \\[2pt]
  &\qquad\qquad =\;
  \sum_{a \in A}
    \Pr\bigl\{X_1 = a \mid X_1 \in A \wedge X_0 = x_0\bigr\}\;
    \Pr\bigl\{X_2 \in B \mid X_1 = a\bigr\},
\end{align*}
si vede che sapere $X_0 = x_0$ cambia i \emph{pesi} con cui gli stati di $A$
entrano nella media (basta che da $x_0$ qualche stato di $A$ sia irraggiungibile
perché i pesi cambino): il passato torna a contare e l'uguaglianza cade. La
proprietà di Markov si applica al presente \emph{puntuale}, non al presente
«sfocato» in un insieme di stati.
\end{osservazione}

\begin{dubbio}
A margine della (2) il quaderno annota: «sì, nel caso in cui la probabilità di
passare da $x_0$ a qualche stato di $A$ è $0$» (la prima parola, sottolineata, si
legge «sì», ma potrebbe essere «se»). Le due letture di «qualche stato di $A$»
portano in direzioni opposte, e la nota non basta a decidere fra loro. Se la si
intende come $\Pr\{X_1 \in A \mid X_0 = x_0\} = 0$, allora l'evento condizionante
$\{X_1 \in A \wedge X_0 = x_0\}$ è impossibile e la (2) non è né vera né falsa:
il membro sinistro non è nemmeno definito. Se invece «qualche stato» va inteso
come «un particolare stato di $A$», la nota descrive proprio il meccanismo per cui
la (2) \emph{fallisce}, quello dell'osservazione qui sopra: se da $x_0$ un certo
$a \in A$ non è raggiungibile, i pesi con cui gli stati di $A$ entrano nella media
cambiano.
\end{dubbio}


\section{Matrice di transizione e vettore di stato}

\scan{73}
Supponiamo d'ora in avanti che gli stati siano in numero finito, $S = \{1, \dots, n\}$
(identifichiamo cioè l'esito $E_i$ con lo stato $i$), e raccogliamo le probabilità di
transizione nella \emph{matrice di transizione}
\[
  P \;=\;
  \begin{pmatrix}
    P_{11} & \cdots & P_{1n} \\
    \vdots &        & \vdots \\
    P_{n1} & \cdots & P_{nn}
  \end{pmatrix}.
\]
Ogni \emph{riga} di $P$ è una distribuzione di probabilità sugli $n$ stati — la riga
$i$-esima dice dove si può andare a partire da $i$, e con quale probabilità — sicché
\[
  \sum_{j=1}^{n} P_{ij} \;=\; 1 \qquad \text{per ogni } i = 1, \dots, n .
\]
Una matrice quadrata a entrate non negative con questa proprietà si dice
\emph{stocastica}.

\begin{osservazione}
  Il quaderno affianca a questa anche l'identità per colonne $\sum_{i=1}^{n} P_{ij} = 1$.
  Non è però una richiesta della definizione: vale soltanto per le matrici
  \emph{doppiamente stocastiche}, che sono un caso particolare. Lo è, per esempio, la
  passeggiata aleatoria simmetrica sul ciclo della Sezione~\ref{sec:mc-esempio-ciclo}, ed è
  esattamente il motivo per cui in quel caso la distribuzione limite risulta uniforme.
\end{osservazione}

\paragraph{Notazione.}
\scan{75}
\begin{itemize}
  \item $P_{ij}$ è la probabilità di transizione da $i$ a $j$ in \emph{un} passo,
        \[
          P_{ij} \;=\; \Pr[X_{t+1} = j \mid X_t = i],
        \]
        e per l'ipotesi di omogeneità non dipende da $t$;
  \item $P^{(t)}_{ij}$ è la probabilità di andare da $i$ a $j$ in \emph{esattamente}
        $t$ passi,
        \[
          P^{(t)}_{ij} \;=\; \Pr[X_t = j \mid X_0 = i].
        \]
\end{itemize}

\paragraph{Stato iniziale.} Per $X_0$ ci sono due possibilità: o è \emph{fisso},
oppure viene \emph{generato da una distribuzione} iniziale sugli stati.

\medskip

\scan{73}
Alla matrice $P$ si associa naturalmente un grafo orientato sull'insieme $S$ degli
stati: i nodi sono gli stati e c'è un arco da $i$ a $j$ quando $P_{ij} > 0$. È il punto
di vista che useremo sistematicamente più avanti.

\begin{figure}[H]
  \centering
  % p73-spazio-stati -- lo spazio degli stati S come "patata" con dentro un
% pezzo del grafo di transizione (quaderno, pag. 73).
% Tre stati s1, s2, s3; arco s1->s2, coppia s1<->s3 (lente), cappio su s3.
\begin{tikzpicture}[
    font=\footnotesize,
    stato/.style={circle,draw,line width=0.7pt,fill=white,
                  inner sep=0pt,minimum size=5.5pt},
    arco/.style={line width=0.7pt,-{Stealth[length=4.5pt,width=3.5pt]}}]

  %% ------------------- la "patata": l'insieme S degli stati -------------
  \draw[line width=0.8pt,gray]
    plot[smooth cycle,tension=0.75]
    coordinates {(0.30,1.80) (1.10,3.00) (2.40,3.35) (3.10,2.70)
                 (3.90,3.30) (5.00,3.10) (5.70,2.00) (5.20,0.90)
                 (4.20,0.25) (3.10,0.55) (2.40,0.15) (1.20,0.50)
                 (0.45,1.10)};

  %% ------------------- glossa esterna ----------------------------------
  \node[anchor=east,align=center] (etS) at (-0.70,2.15) {$S$\\[1pt]stati};
  \draw[line width=0.5pt,gray] (etS.east) -- (0.34,1.90);

  %% ------------------- gli stati ---------------------------------------
  \node[stato] (s1) at (1.80,1.30) {};
  \node[stato] (s2) at (2.65,2.25) {};
  \node[stato] (s3) at (3.95,1.55) {};

  \node[anchor=east]       at ([xshift=-2pt]s1.west)              {$s_1$};
  \node[anchor=south]      at ([yshift=2pt]s2.north)              {$s_2$};
  \node[anchor=north west] at ([shift={(1pt,-2pt)}]s3.south east) {$s_3$};

  %% ------------------- transizioni --------------------------------------
  \draw[arco] (s1) to[bend left=32] (s2);          % s1 -> s2
  \draw[arco] (s1) to[bend left=24] (s3);          % s1 -> s3
  \draw[arco] (s3) to[bend left=24] (s1);          % s3 -> s1
  \draw[arco] (s3) to[out=62,in=2,looseness=1,min distance=7mm] (s3); % cappio

\end{tikzpicture}

  \caption{L'insieme $S$ degli stati e alcune transizioni possibili fra di essi.}
  \label{fig:p73-spazio-stati}
\end{figure}

Letta così, una catena di Markov è \emph{una generalizzazione della nozione di
funzione}: a ogni stato non è associato un unico successore, ma una distribuzione di
probabilità sui successori possibili. Il caso deterministico si riottiene quando ogni
riga di $P$ contiene un solo $1$ e tutti gli altri elementi nulli.
L'osservazione è attribuita, nel quaderno, a \dubbioinl{J.~T.~Schwartz}.

\scan{75}
Quando gli stati sono pochi, il grafo di transizione si disegna per esteso, con le
probabilità come etichette degli archi.

\begin{figure}[H]
  \centering
  % p75-mc-due-stati -- diagramma di transizione di una catena a due stati
% (schizzo a matita nel margine, quaderno pag. 75).
% Le etichette sono quelle del quaderno: le uscite da 2 non sommano a 1,
% l'incongruenza e' discussa nell'ambiente "dubbio" subito dopo la figura.
\begin{tikzpicture}[
    font=\footnotesize,
    stato/.style={circle,draw,line width=0.7pt,fill=white,
                  inner sep=1pt,minimum size=7mm},
    arco/.style={line width=0.7pt,-{Stealth[length=5pt,width=4pt]}}]

  \node[stato] (u) at (0,0)   {$1$};
  \node[stato] (v) at (3.5,0) {$2$};

  % 1 -> 2 : arco che scavalca i due nodi
  \draw[arco] (u) to[bend left=45]  node[above,pos=0.5]{$0{,}8$} (v);
  % 2 -> 1 : arco di ritorno, sotto
  \draw[arco] (v) to[bend left=45]  node[below,pos=0.72]{$0{,}2$} (u);

  % cappi
  \draw[arco] (u) to[out=155,in=205,looseness=1,min distance=12mm]
              node[above left,pos=0.5,inner sep=1.5pt]{$0{,}2$} (u);
  \draw[arco] (v) to[out=-25,in=25,looseness=1,min distance=12mm]
              node[below right,pos=0.5,inner sep=1.5pt]{$0{,}2$} (v);

\end{tikzpicture}

  \caption{Catena di Markov a due stati, abbozzata a margine della pagina.}
  \label{fig:p75-mc-due-stati}
\end{figure}

\begin{dubbio}
  Nello schizzo le probabilità uscenti dallo stato $2$ sommano a $0{,}4$ e non a $1$:
  una delle due etichette $0{,}2$ che lo riguardano (il cappio su $2$ oppure l'arco di
  ritorno $2 \to 1$) dovrebbe valere $0{,}8$, ma dal quaderno non si capisce quale
  delle due.
\end{dubbio}

\medskip

Della catena interessa infine, a ogni istante, non tanto il singolo stato occupato
quanto la distribuzione di probabilità sugli stati.

\begin{definizione}[Vettore di stato, \emph{state probability vector}]\label{def:state-probability-vector}
  \scan{77}
  Il \emph{vettore di probabilità di stato} al tempo $t$ è il vettore riga
  \[
    \vec{q}^{\,(t)} \;=\; \bigl(q^{(t)}_{1},\, \dots,\, q^{(t)}_{n}\bigr),
  \]
  dove $q^{(t)}_{i}$ è la probabilità che la catena si trovi nello stato $i$ al tempo $t$.
\end{definizione}

Il vettore $\vec{q}^{\,(t)}$ è dunque la distribuzione dello stato al tempo $t$:
rappresenta la distribuzione della probabilità di visita di un random walk sul grafo di
transizione. Il passo della catena si ottiene moltiplicando il vettore (riga) di stato
per la matrice di transizione:
\[
  \vec{q}^{\,(t+1)} \;=\; \vec{q}^{\,(t)} \times P ,
  \qquad\text{e quindi}\qquad
  \vec{q}^{\,(t)} \;=\; \vec{q}^{\,(0)}\, P^{t} .
\]


\section{Un esempio: la passeggiata aleatoria sul ciclo}
\label{sec:mc-esempio-ciclo}

\scan{74}

\begin{esempio}\label{es:ciclo6}
Passeggiata aleatoria pigra sul ciclo a sei stati $\{0,1,2,3,4,5\}$: da ogni
stato ci si sposta al vicino precedente, si resta fermi, oppure ci si sposta al
vicino successivo, ciascuna cosa con probabilità $1/3$.

\begin{figure}[H]
\centering
% p74-ciclo6 -- passeggiata aleatoria "pigra" sul ciclo a sei stati
% (quaderno, pag. 74).  Sei stati 0..5 disposti in cerchio in senso orario
% a partire dall'alto; ogni stato ha un cappio e due archi (uno per verso)
% verso ciascun vicino.  Tutte le transizioni hanno probabilita' 1/3.
\begin{tikzpicture}[
    font=\footnotesize,
    stato/.style={circle,draw,line width=0.7pt,fill=white,
                  inner sep=1pt,minimum size=7mm},
    arco/.style={line width=0.7pt,-{Stealth[length=4.5pt,width=3.5pt]}}]

  %% ------------------- i sei stati --------------------------------------
  \foreach \n/\ang in {0/90, 1/30, 2/-30, 3/-90, 4/210, 5/150}
    \node[stato] (s\n) at (\ang:2.2) {$\n$};

  %% ------------------- archi fra vicini (uno per verso) -----------------
  \foreach \i/\j in {0/1, 1/2, 2/3, 3/4, 4/5, 5/0} {
    \draw[arco] (s\i) to[bend left=16] (s\j);
    \draw[arco] (s\j) to[bend left=16] (s\i);
  }

  %% ------------------- cappi, radiali verso l'esterno -------------------
  \foreach \n/\uscita/\entrata in {0/112/68, 1/52/8, 2/-8/-52,
                                   3/-68/-112, 4/232/188, 5/172/128}
    \draw[arco] (s\n) to[out=\uscita,in=\entrata,looseness=1,min distance=8mm]
                (s\n);

  %% ------------------- glossa al centro ---------------------------------
  \node[align=center] at (0,0) {ogni arco:\\[1pt]$\tfrac13$};

\end{tikzpicture}

\caption{La catena dell'esempio: sei stati disposti in cerchio, ciascuno con un
cappio e due archi (uno per verso) verso i due vicini; ogni arco ha probabilità
$1/3$.}
\label{fig:p74-ciclo6}
\end{figure}

Partendo da $i$, qual è la probabilità del passaggio $i \to j$?
\[
  P_{ij} \;=\;
  \begin{cases}
    \tfrac13 & j = i-1 \quad(\text{un passo indietro}),\\[2pt]
    \tfrac13 & j = i \phantom{{}-1} \quad(\text{fermo}),\\[2pt]
    \tfrac13 & j = i+1 \quad(\text{un passo avanti}),\\[2pt]
    0        & \text{altrimenti,}
  \end{cases}
\]
dove gli indici sono intesi modulo $6$.

Facendo partire la catena dallo stato $2$, il vettore di stato evolve così:
\begin{align*}
  \vec{q}^{\,(0)} &= \bigl(0,\;0,\;1,\;0,\;0,\;0\bigr), \\
  \vec{q}^{\,(1)} &= \bigl(0,\;\tfrac13,\;\tfrac13,\;\tfrac13,\;0,\;0\bigr)
               && \text{in $1$ passo,}\\
  \vec{q}^{\,(2)} &= \bigl(\tfrac19,\;\tfrac29,\;\tfrac13,\;\tfrac29,\;\tfrac19,\;0\bigr)
               && \text{in $2$ passi,}\\
  &\;\;\;\vdots \\
  \vec{q}^{\,(t)} &\;\xrightarrow[\;t\to\infty\;]{}\;
               \bigl(\tfrac16,\;\tfrac16,\;\tfrac16,\;\tfrac16,\;\tfrac16,\;\tfrac16\bigr).
\end{align*}
\end{esempio}

\begin{osservazione}
La matrice $P$ di questa catena è doppiamente stocastica (ogni riga e ogni colonna
somma a $1$); la catena è inoltre irriducibile, perché il ciclo è connesso, e
aperiodica, perché i cappi danno periodo $1$. Nel linguaggio della
Sezione~\ref{sec:mc-stazionaria} si tratta dunque di una catena ergodica: il
Teorema~\ref{thm:mc-fondamentale} ne garantisce la distribuzione stazionaria, e per le
catene ergodiche finite la distribuzione al tempo $t$ vi converge qualunque sia lo stato
di partenza --- è il limite scritto sopra. Che tale limite sia la distribuzione uniforme
lo dice la doppia stocasticità: se ogni colonna di $P$ somma a $1$, il vettore costante
$\vec{\pi} = \bigl(\tfrac16,\dots,\tfrac16\bigr)$ soddisfa $\vec{\pi} = \vec{\pi}P$.
\end{osservazione}

\begin{osservazione}
Il passo elementare $\vec{q}^{\,(t+1)} = \vec{q}^{\,(t)}P$ è, per questa catena, una
\emph{convoluzione}: per ogni stato $j$
\[
  q^{(t+1)}_{j} \;=\;
  \tfrac13\Bigl(q^{(t)}_{j-1} + q^{(t)}_{j} + q^{(t)}_{j+1}\Bigr),
\]
cioè si fa scorrere sulla distribuzione corrente la finestra di pesi
$\bigl(\tfrac13,\tfrac13,\tfrac13\bigr)$ e si sommano i tre contributi.

\begin{figure}[H]
\centering
% p74-convoluzione -- il passo q^(t+1) = q^(t) P della passeggiata sul ciclo
% letto come convoluzione (quaderno, pag. 74, schizzo a matita).
% Sopra: i tre valori adiacenti della distribuzione corrente, pesati 1/3 e
% sommati nel valore nuovo.  Sotto: la finestra (1/3,1/3,1/3) che scorre.
\begin{tikzpicture}[
    font=\footnotesize,
    linea/.style={line width=0.7pt},
    freccia/.style={line width=0.7pt,-{Stealth[length=4.5pt,width=3.5pt]}}]

  %% ---------- le tre posizioni adiacenti della distribuzione al tempo t --
  \foreach \i/\et in {0/{q^{(t)}_{j-1}}, 1/{q^{(t)}_{j}}, 2/{q^{(t)}_{j+1}}} {
    \node at ({1.7*\i},2.55) {$\et$};
    \draw[linea] ({1.7*\i},2.25) -- ({1.7*\i},1.95);   % trattino verticale
    \node at ({1.7*\i},1.72) {$\tfrac13$};             % peso
  }
  \node at (0.85,1.72) {$+$};
  \node at (2.55,1.72) {$+$};

  %% ---------- la graffa che raccoglie i tre contributi -------------------
  \draw[linea,decorate,decoration={brace,amplitude=5pt,mirror,raise=1pt}]
        (-0.55,1.40) -- (3.95,1.40);
  \draw[freccia] (1.7,1.02) -- (1.7,0.62);
  \node at (1.7,0.38) {$q^{(t+1)}_{j}$};

  %% ---------- la finestra di convoluzione --------------------------------
  \draw[linea] (-0.85,-1.20) rectangle (4.25,-0.50);
  \draw[linea] (0.85,-1.20) -- (0.85,-0.50);
  \draw[linea] (2.55,-1.20) -- (2.55,-0.50);
  \foreach \i in {0,1,2} \node at ({1.7*\i},-0.85) {$\tfrac13$};

  \node[anchor=east] at (-1.05,-0.85) {finestra};
  \draw[freccia] (4.45,-0.85) -- (5.15,-0.85);
  \node[anchor=west] at (5.25,-0.85) {scorre};

\end{tikzpicture}

\caption{Lo schema della convoluzione: i tre valori adiacenti, pesati $1/3$
ciascuno, si sommano nel valore nuovo; la finestra $(1/3,1/3,1/3)$ scorre lungo
la distribuzione.}
\label{fig:p74-convoluzione}
\end{figure}
\end{osservazione}

In termini matriciali il primo passo non è che l'estrazione di una riga: il vettore
iniziale $\vec{q}^{\,(0)}$ concentra tutta la massa nello stato $2$, dunque
\[
  \vec{q}^{\,(1)} \;=\; \vec{q}^{\,(0)} P \;=\; \bigl(0,\;\tfrac13,\;\tfrac13,\;\tfrac13,\;0,\;0\bigr)
\]
è la riga di indice $2$ di $P$ — e, essendo $P$ simmetrica, anche la sua colonna di
indice $2$.


\section{Raggiungibilità, tempi di attesa e classificazione degli stati}

\scan{75}
Per capire dove finisce, alla lunga, una catena servono due quantità: con che
probabilità uno stato viene raggiunto, e dopo quanto tempo. Entrambe si ricavano dalla
probabilità di raggiungerlo per la \emph{prima} volta a un dato istante.

\begin{definizione}[Prima transizione, \emph{first transition}]\label{def:first-transition}
  Per $i, j \in S$ e $t > 0$, sia $r^{(t)}_{ij}$ la probabilità che, partendo da $i$, la
  \emph{prima} transizione in $j$ avvenga al tempo $t$:
  \[
    r^{(t)}_{ij} \;=\;
    \Pr\Bigl[\; X_t = j \;\wedge\; \forall s \in \{1, \dots, t-1\}\;(X_s \ne j)
      \;\Big|\; X_0 = i \;\Bigr].
  \]
\end{definizione}

Da $r^{(t)}_{ij}$ si ricavano due quantità:
\[
  f_{ij} \;=\; \sum_{t > 0} r^{(t)}_{ij},
  \qquad\qquad
  h_{ij} \;=\; \sum_{t > 0} t\, r^{(t)}_{ij}.
\]
La prima è la probabilità che $j$ venga raggiunto, prima o poi, partendo da $i$; la
seconda è l'\emph{hitting time}, cioè il valore atteso del numero di passi necessari per
raggiungere $j$ a partire da $i$ (la prima volta). Detto $T_{ij}$ tale numero di passi,
con la convenzione $T_{ij} = +\infty$ sulle traiettorie che in $j$ non arrivano mai, si
pone $h_{ij} = \E[T_{ij}]$: la somma qui sopra lo calcola quando $f_{ij} = 1$, mentre
per $f_{ij} < 1$ si ha $h_{ij} = \infty$.

\begin{osservazione}
  La convenzione non è arbitraria. Se $f_{ij} < 1$, cioè se non è certo che $j$ venga
  raggiunto a partire da $i$, la serie $\sum_{t>0} t\, r^{(t)}_{ij}$ pesa soltanto le
  traiettorie che raggiungono $j$ e può benissimo convergere; ma la massa di probabilità
  residua $1 - f_{ij} > 0$ corrisponde a traiettorie che non lo raggiungono mai, e per le
  quali il numero di passi è infinito. Il valore atteso di $T_{ij}$ è dunque $\infty$, ed
  è questo il valore che si attribuisce a $h_{ij}$.
\end{osservazione}

Specializzando al caso $i = j$ si ottiene la classificazione degli stati.

\begin{definizione}[Stato transiente, stato persistente]\label{def:mc-transiente-persistente}
\scan{76}
Sia $i$ uno stato di una catena di Markov.
\begin{itemize}
  \item $i$ si dice \emph{transiente} se
    \[
      f_{ii} \;<\; 1 \qquad (\text{e dunque } h_{ii} = \infty);
    \]
  \item $i$ si dice \emph{persistente} se
    \[
      f_{ii} \;=\; 1 .
    \]
    Uno stato persistente si dice inoltre
    \begin{itemize}
      \item \emph{persistente non nullo} se $h_{ii} < \infty$;
      \item \emph{persistente nullo} se $h_{ii} = \infty$.
    \end{itemize}
\end{itemize}
\end{definizione}

\begin{osservazione}
Le tre situazioni si leggono così.
\begin{itemize}
  \item $i$ \emph{transiente}: non è certo che si ripassi da $i$. Con probabilità
    $1-f_{ii} > 0$ la catena, partita da $i$, non vi fa mai più ritorno, e per questo
    il tempo medio di ritorno vale $h_{ii} = \infty$.
  \item $i$ \emph{persistente non nullo}: il ritorno in $i$ è certo e avviene, in media,
    dopo un numero finito di passi.
  \item $i$ \emph{persistente nullo}: il ritorno in $i$ è certo ($f_{ii}=1$) eppure il
    numero atteso di passi diverge. Ciò accade quando esistono \emph{primi ritorni} in $i$
    arbitrariamente lunghi --- cammini da $i$ a $i$ che non passano per $i$ nei vertici
    intermedi --- con probabilità che non decresce abbastanza in fretta da rendere
    convergente la serie
    \[
      h_{ii} \;=\; \sum_{t>0} t\, r_{ii}^{(t)} .
    \]
\end{itemize}
\end{osservazione}

\begin{proposizione}\label{prop:mc-finita-no-nulli}
Nel caso che interessa a noi, cioè quando l'insieme degli stati è \emph{finito}, ogni stato è
transiente oppure persistente non nullo: non esistono stati persistenti nulli.
\end{proposizione}

\subsection{Il punto di vista dei grafi}

Riprendiamo il grafo di transizione della Figura~\ref{fig:p73-spazio-stati}: la struttura
combinatoria che governa la classificazione appena data è la sua decomposizione in
componenti fortemente connesse.

\begin{definizione}[Componenti fortemente connesse]\label{def:scc}
Le \emph{componenti fortemente connesse} (\emph{strongly connected components}, SCC) di un
grafo orientato sono le classi di equivalenza della relazione di \emph{mutua raggiungibilità}
\[
  i \sim j
  \quad\iff\quad
  \text{$j$ è raggiungibile da $i$\; e\; $i$ è raggiungibile da $j$.}
\]
\end{definizione}

\begin{figure}[H]
\centering
% p76-scc -- pag. 76
% Decomposizione di un grafo orientato nelle sue componenti fortemente connesse.
% Quattro bolle: dentro ciascuna un ciclo orientato che tocca tutti i suoi nodi
% (quindi ogni bolla e' davvero fortemente connessa); fra bolle diverse gli archi
% vanno tutti in un solo verso, cosi' che la condensazione
%   C1 -> C2 -> C4,  C1 -> C3
% sia un DAG (C3 e C4 sono le SCC finali).
\begin{tikzpicture}[x=1.1cm,y=1.05cm,line join=round,line cap=round,
    font=\footnotesize,
    nodo/.style={circle,fill=black,inner sep=0pt,minimum size=3.4pt},
    arco/.style={line width=0.5pt,-{Latex[length=1.5mm,width=1.2mm]}},
    arcone/.style={line width=0.7pt,-{Latex[length=2mm,width=1.5mm]}},
    bolla/.style={line width=0.6pt,black!45,fill=black!3},
    etic/.style={inner sep=1pt,text=black!60}]

  % ----------------------------------------------------------------- bolle --
  \draw[bolla] plot[smooth cycle,tension=0.75] coordinates {
    (1.05,3.02) (0.86,3.64) (0.28,3.86) (-0.34,3.52) (-0.56,2.86)
    (-0.40,2.24) (0.18,2.00) (0.70,2.36)};

  \draw[bolla] plot[smooth cycle,tension=0.75] coordinates {
    (6.60,1.55) (6.30,2.78) (5.40,3.48) (4.20,3.76) (3.00,3.62)
    (2.40,2.90) (2.28,1.90) (2.54,0.94) (3.30,0.34) (4.50,0.14)
    (5.70,0.20) (6.45,0.74)};

  \draw[bolla] plot[smooth cycle,tension=0.75] coordinates {
    (-0.66,0.86) (-0.80,1.56) (-1.36,1.82) (-1.86,1.46) (-2.02,0.76)
    (-1.86,0.10) (-1.36,-0.16) (-0.86,0.24)};

  \draw[bolla] plot[smooth cycle,tension=0.75] coordinates {
    (3.76,-1.44) (3.46,-0.74) (2.70,-0.48) (1.84,-0.54) (1.30,-1.04)
    (1.24,-1.80) (1.76,-2.36) (2.56,-2.56) (3.36,-2.20)};

  % ----------------------------------------------------------------- nodi ---
  \node[nodo] (n1) at ( 0.58, 3.38) {};
  \node[nodo] (n2) at (-0.04, 2.48) {};

  \node[nodo] (m1) at ( 3.30, 3.05) {};
  \node[nodo] (m2) at ( 5.60, 2.75) {};
  \node[nodo] (m3) at ( 5.95, 0.70) {};
  \node[nodo] (m4) at ( 3.15, 1.15) {};

  \node[nodo] (p1) at (-1.35, 0.30) {};
  \node[nodo] (p2) at (-1.05, 1.25) {};

  \node[nodo] (q1) at ( 1.95,-1.05) {};
  \node[nodo] (q2) at ( 3.15,-1.25) {};
  \node[nodo] (q3) at ( 2.45,-2.10) {};

  % --------------------------------------------------- cicli interni (SCC) --
  \draw[arco] (n2) to[bend left=36] (n1);
  \draw[arco] (n1) to[bend left=36] (n2);

  \draw[arco] (m1) to[bend left=12] (m2);
  \draw[arco] (m2) to[bend left=24] (m3);
  \draw[arco] (m3) to[bend left=20] (m4);
  \draw[arco] (m4) to[bend left=22] (m1);

  \draw[arco] (p1) to[bend left=36] (p2);
  \draw[arco] (p2) to[bend left=36] (p1);

  \draw[arco] (q1) to[bend left=20] (q2);
  \draw[arco] (q2) to[bend left=20] (q3);
  \draw[arco] (q3) to[bend left=20] (q1);

  % ------------------------------------------------ archi fra le componenti --
  \draw[arcone] (n1) to[bend right=10] (m1);
  \draw[arcone] (n2) to[bend right=14] (p2);
  \draw[arcone] (m4) to[bend left=12]  (q1);

  % ------------------------------------------------------------- etichette --
  \node[etic,anchor=south east] at (-0.34,3.56) {$C_1$};
  \node[etic,anchor=south west] at ( 6.24,2.90) {$C_2$};
  \node[etic,anchor=east]       at (-2.10,1.10) {$C_3$};
  \node[etic,anchor=west]       at ( 3.62,-2.10) {$C_4$};

\end{tikzpicture}

\caption{Le componenti fortemente connesse di un grafo orientato: ogni bolla è una classe di
equivalenza della relazione di mutua raggiungibilità. All'interno di una bolla ogni nodo
raggiunge ogni altro (e infatti un ciclo orientato li attraversa tutti); fra bolle distinte
gli archi vanno tutti in un solo verso, altrimenti le due bolle sarebbero un'unica SCC.}
\label{fig:p76-scc}
\end{figure}

\scan{77}
Se \emph{pensiamo} allo spazio degli stati come a un grafo, le componenti fortemente
connesse che ne otteniamo sono correlate con le nozioni di stato \emph{transiente} e di
stato \emph{persistente}.

\begin{definizione}[SCC finale]\label{def:scc-finale}
  Una componente fortemente connessa si dice \emph{finale} se è priva di archi uscenti.
\end{definizione}

\begin{definizione}[Catena irriducibile]\label{def:mc-irriducibile}
  Una catena di Markov è \emph{irriducibile} se ha una sola SCC --- che è allora,
  necessariamente, finale --- ossia se il suo grafo di transizione è fortemente connesso.
\end{definizione}

Le catene irriducibili sono le catene di Markov che sappiamo trattare meglio: è per esse
che vale il teorema fondamentale enunciato nel seguito.

\begin{osservazione}
  In una catena finita il legame fra le SCC e la classificazione degli stati è questo: gli
  stati che stanno in una SCC finale sono persistenti (non nulli, per la
  Proposizione~\ref{prop:mc-finita-no-nulli}), tutti gli altri sono transienti. Il quaderno
  definisce l'irriducibilità chiedendo che vi sia \emph{una sola SCC finale}; la
  condizione corretta è più forte, cioè che vi sia una sola SCC in tutto. La differenza
  non è innocua: il teorema fondamentale afferma che $\pi_i > 0$ per \emph{ogni} stato
  $i$, e questo è falso non appena esistono SCC non finali, i cui stati sono transienti e
  hanno probabilità stazionaria nulla.
\end{osservazione}


\section{Distribuzione stazionaria e teorema fondamentale}
\label{sec:mc-stazionaria}

\scan{77}

\begin{definizione}[Distribuzione stazionaria]\label{def:distribuzione-stazionaria}
  Un vettore di probabilità $\vec{\pi}$ è una \emph{distribuzione stazionaria} della catena
  di matrice di transizione $P$ se
  \[
    \vec{\pi} \;=\; \vec{\pi} P .
  \]
\end{definizione}

\scan{78}
Quando raggiungiamo la distribuzione stazionaria non ci sono più «sorprese»: la catena è
entrata nel suo \emph{steady state}.

\begin{definizione}[Periodicità di uno stato]\label{def:mc-periodicita}
  La \emph{periodicità} di uno stato $i$ è
  \[
    \max
    \Bigl\{\;
      T\in\NN \;:\;
      \exists\,\vec{q}^{\,(0)}\;\;\exists\,a\in\NN\;\;\forall t\;\;
      \bigl(q_i^{(t)}>0 \;\to\; t\equiv a \mod T\bigr)
    \;\Bigr\}.
  \]
\end{definizione}

In parole: il valore $T$ è ammissibile quando esistono una distribuzione iniziale
$\vec{q}^{\,(0)}$ e una classe di resto $a$ modulo $T$ tali che lo stato $i$ possa essere
occupato con probabilità non nulla soltanto negli istanti $t\equiv a \pmod{T}$; la
periodicità di $i$ è il più grande fra questi $T$.

\begin{figure}[H]
  \centering
  % p78-periodicita -- pag. 78
% La periodicita' letta sull'asse dei tempi: i soli istanti in cui lo stato i puo'
% essere occupato sono a, a+T, a+2T, ... (pallini pieni); fra un pallino e il
% successivo passano T passi e in tutti quegli istanti q_i^{(t)} = 0.
\begin{tikzpicture}[x=1cm,y=1cm,font=\footnotesize,line cap=round,
    asse/.style={line width=0.6pt,dash pattern=on 2.4pt off 2.4pt},
    pallino/.style={circle,fill=black,inner sep=0pt,minimum size=3.6pt}]

  % ------------------------------------------------------- asse dei tempi ---
  \draw[asse] (0.15,0) -- (10.15,0);
  \node[anchor=west,inner sep=2.5pt] at (10.15,0) {$\cdots$};

  % --------------------------------------- gli istanti a, a+T, a+2T, ... ----
  \node[pallino] at (2.80,0) {};
  \node[pallino] at (5.80,0) {};
  \node[pallino] at (8.80,0) {};

  % ------------------------------------ graffa dello sfasamento iniziale ----
  \draw[line width=0.5pt,decorate,
        decoration={brace,mirror,amplitude=4pt,raise=2.5pt}]
       (0.15,0) -- (2.80,0);
  \node[anchor=north,inner sep=1.5pt] at (1.475,-0.28) {$a$};

  % ----------------------------- etichette degli istanti «buoni» (pallini) --
  \node[anchor=north,inner sep=1.5pt] at (2.80,-0.32) {$q_i^{(t)}>0$};
  \node[anchor=north,inner sep=1.5pt] at (5.80,-0.32) {$q_i^{(t)}>0$};
  \node[anchor=north,inner sep=1.5pt] at (8.80,-0.32) {$q_i^{(t)}>0$};

  % -------------------------------- etichette degli intervalli fra pallini --
  \node[anchor=south,inner sep=1.5pt] at (4.30,0.14) {$T$ passi};
  \node[anchor=south,inner sep=1.5pt] at (7.30,0.14) {$T$ passi};

  \node[anchor=north,inner sep=1.5pt] at (4.30,-1.02) {$q_i^{(t)}=0$};
  \node[anchor=north,inner sep=1.5pt] at (7.30,-1.02) {$q_i^{(t)}=0$};

\end{tikzpicture}

  \caption{La periodicità letta sull'asse dei tempi: lo stato $i$ può essere occupato
  ($q_i^{(t)}>0$) soltanto negli istanti $a,\,a+T,\,a+2T,\dots$, mentre in tutti gli altri
  istanti si ha $q_i^{(t)}=0$.}
  \label{fig:p78-periodicita}
\end{figure}

\begin{osservazione}
  Prendendo $\vec{q}^{\,(0)}$ concentrata in $i$ si ha $q_i^{(0)} = 1 > 0$, il che forza
  $a \equiv 0$, e la condizione diventa: $T$ divide ogni $t$ per cui $P_{ii}^{(t)}>0$. Con
  questa scelta la periodicità di $i$ è il massimo comune divisore dell'insieme
  $\set{t\ge 1 \;:\; P_{ii}^{(t)}>0}$ dei possibili tempi di ritorno in $i$, e il massimo
  esiste davvero non appena in $i$ si può tornare: $T=1$ è sempre ammissibile, e ogni $T$
  ammissibile divide --- dunque non supera --- un qualunque tempo di ritorno. La
  quantificazione esistenziale su $\vec{q}^{\,(0)}$ non cambia il valore per le catene
  irriducibili, che sono poi le sole per cui useremo la nozione; in generale, invece, la
  renderebbe inservibile, perché una distribuzione iniziale da cui $i$ non è raggiungibile
  soddisfa la condizione a vuoto, per ogni $T$.
\end{osservazione}

\begin{definizione}[Stati e catene aperiodiche]\label{def:mc-aperiodica}
  Lo stato $i$ è \emph{periodico} se ha periodicità $>1$; altrimenti è \emph{aperiodico}.
  Una catena è \emph{aperiodica} se ogni suo stato lo è.
\end{definizione}

\begin{definizione}[Stati e catene ergodiche]\label{def:mc-ergodico}
  Uno stato è \emph{ergodico} se è
  \begin{itemize}
    \item aperiodico;
    \item persistente non nullo.
  \end{itemize}
  Una catena di Markov è \emph{ergodica} se ogni suo stato lo è.
\end{definizione}

\begin{osservazione}
  I due requisiti, messi insieme, dicono che un random walk sulla catena ritorna in quello
  stato con probabilità $1$ (persistenza), dopo un tempo di attesa di valore atteso finito
  (non nullità) e senza che gli istanti di visita siano confinati in una progressione
  aritmetica (aperiodicità): il random walk continua dunque a visitare, indefinitamente nel
  tempo, ogni suo stato.
\end{osservazione}

\begin{teorema}[fondamentale sulle catene di Markov]\label{thm:mc-fondamentale}
  Ogni catena di Markov irriducibile, finita e aperiodica gode delle seguenti proprietà.
  \begin{enumerate}
    \item Tutti i suoi stati sono ergodici.
    \item Esiste una distribuzione stazionaria $\vec{\pi}$ tale che $\pi_i>0$ per ogni $i$.
    \item Per ogni $i$ si ha
      \[
        f_{ii}=1 \qquad\text{e}\qquad h_{ii}=\frac{1}{\pi_i}.
      \]
    \item \scan{79}%
      Detto $N(i,t)$ il numero medio di volte in cui la catena visita lo stato $i$ nei
      primi $t$ passi, vale
      \[
        \lim_{t\to\infty}\frac{N(i,t)}{t}=\pi_i .
      \]
  \end{enumerate}
\end{teorema}

\begin{osservazione}
  La distribuzione stazionaria del punto~2 è anche \emph{unica}, e i punti~3 e~4 ne danno
  due letture concrete. Il punto~3 dice che $\pi_i$ è il reciproco del tempo medio di
  ritorno in $i$; il punto~4 è la lettura «frequentista»: $\pi_i$ è la frazione di tempo
  che la catena trascorre, alla lunga, nello stato $i$. Le due letture sono coerenti fra
  loro: se il tempo medio di ritorno in $i$ è $h_{ii}=1/\pi_i$, in $t$ passi lo stato $i$
  viene visitato circa $t\,\pi_i$ volte.

  Sotto le stesse ipotesi vale inoltre il \emph{basic limit theorem}: qualunque sia la
  distribuzione iniziale $\vec{q}^{\,(0)}$ si ha $\vec{q}^{\,(t)} \to \vec{\pi}$ per
  $t\to\infty$. È il fatto usato nell'Esempio~\ref{es:ciclo6}, e va tenuto distinto dal
  punto~4: quello è un limite alla Cesàro e continua a valere per le catene irriducibili
  \emph{periodiche}, per le quali invece $\vec{q}^{\,(t)}$ in generale non converge.
\end{osservazione}

\chapter{Random walk su grafi: hitting, commute e cover time}

Il \emph{random walk} su un grafo è l'esempio più naturale di catena di Markov: gli stati
sono i vertici e a ogni passo si sceglie a caso uno degli archi incidenti al vertice
corrente. In questo capitolo vediamo che cosa il teorema fondamentale delle catene di
Markov dice di questa catena --- la distribuzione stazionaria è proporzionale al grado ---
e introduciamo le tre grandezze con cui se ne misura il costo: \emph{hitting time},
\emph{commute time} e \emph{cover time}.

\section{La catena di Markov associata a un grafo}

\scan{83}

\begin{osservazione}[Richiamo: il teorema fondamentale delle catene di Markov]
  Ogni catena di Markov irriducibile, finita e aperiodica gode delle seguenti proprietà:
  tutti i suoi stati sono ergodici (e lo è dunque la catena stessa); esiste un'unica
  distribuzione stazionaria $\vec{\pi}$, e per essa $\pi_i > 0$ per ogni
  $i \in \{1,\dots,n\}$; per ogni stato $i$ valgono $f_{i,i} = 1$ e $h_{i,i} = 1/\pi_i$;
  infine, detto $N(i,t)$ il numero medio di volte in cui la catena visita lo stato $i$ nei
  primi $t$ passi, $\lim_{t\to\infty} N(i,t)/t = \pi_i$.

  Il Teorema~\ref{thm:mc-fondamentale} è stato enunciato e discusso nel capitolo precedente;
  qui ne useremo due conseguenze: l'\emph{unicità} della distribuzione stazionaria e l'identità
  $h_{i,i} = 1/\pi_i$ fra tempo medio di ritorno in uno stato e probabilità stazionaria
  dello stato.
\end{osservazione}

Sia dunque $G = (V,E)$ un grafo \emph{connesso}, \emph{non orientato} e \emph{non
bipartito}, con
\[
  \abs{V} = n , \qquad \abs{E} = m .
\]
Queste tre ipotesi restano in vigore per tutto il capitolo. Per $u \in V$ indichiamo con
$\Gamma(u)$ l'insieme dei vicini di $u$ e con $d(u) = \abs{\Gamma(u)}$ il suo grado.
L'ultima ipotesi --- la non bipartitezza --- serve a garantire la \emph{non periodicità}
(cioè l'aperiodicità) della catena che stiamo per costruire, ed è l'oggetto
dell'osservazione che chiude la sezione; la connessione serve all'irriducibilità e la
finitezza è ovvia.

\begin{definizione}[Random walk su un grafo]\label{def:random-walk-grafo}
  La catena di Markov $M_G$ associata a $G$ è definita così:
  \begin{itemize}
    \item gli \textbf{stati} sono i \textbf{vertici} del grafo;
    \item per ogni $u, v \in V$ la probabilità di transizione è
      \[
        P_{u,v} \;=\;
        \begin{cases}
          \dfrac{1}{d(u)} & \text{se } (u,v) \in E,\\[8pt]
          0 & \text{altrimenti.}
        \end{cases}
      \]
  \end{itemize}
  Il passo elementare è dunque: dal vertice corrente $u$ si sceglie uniformemente a caso
  uno dei $d(u)$ archi incidenti e lo si attraversa.
\end{definizione}

\begin{figure}[H]
  \centering
  % p83-transizioni-mg -- pag. 83
% Un passo del random walk M_G a partire dal vertice u: dei d(u) archi incidenti
% se ne sceglie uno uniformemente a caso, dunque P_{u,v} = 1/d(u).
% Il quaderno disegna tre archi a ventaglio; qui i tre archi restano, ma fra l'uno
% e l'altro si aggiungono i puntini che segnalano i vicini omessi, cosi' che l'arco
% di cerchio possa davvero contare d(u) archi.  Grafo non orientato: niente frecce.
\begin{tikzpicture}[
    x=1cm, y=1cm, line join=round, line cap=round,
    font=\footnotesize,
    nodo/.style   = {circle, draw, line width=0.7pt, fill=white,
                     inner sep=0pt, minimum size=5pt},
    nodov/.style  = {circle, draw, line width=0.8pt, fill=white,
                     inner sep=0pt, minimum size=7pt},
    arco/.style   = {line width=0.7pt},
    conta/.style  = {gray!70!black, line width=0.6pt},
  ]

  % ---- il vertice corrente ------------------------------------------------
  \node[nodov] (u) at (0,0) {};
  \node[anchor=north east, inner sep=2pt] at (u.south west) {$u$};

  % ---- il ventaglio dei vicini (tre archi, campione dei d(u)) -------------
  \def\Ra{2.55}
  \node[nodo]  (w1) at (48:\Ra) {};
  \node[nodov] (v)  at (-2:\Ra) {};
  \node[nodo]  (w2) at (-52:\Ra) {};

  \node[anchor=west, inner sep=3pt] at (v.east) {$v$};

  \draw[arco] (u) -- (w1);
  \draw[arco] (u) -- (v);
  \draw[arco] (u) -- (w2);

  % ---- i vicini omessi: puntini fra un arco e l'altro ---------------------
  \foreach \ang in {32,23,14}{\fill[gray!70!black] (\ang:2.45) circle (0.65pt);}
  \foreach \ang in {-18,-27,-36}{\fill[gray!70!black] (\ang:2.45) circle (0.65pt);}

  % ---- la probabilita' di scegliere l'arco u--v ---------------------------
  \node[anchor=south, inner sep=2.5pt] at (-2:1.5) {$\dfrac{1}{d(u)}$};

  % ---- arco di cerchio che "conta" gli archi incidenti a u ----------------
  \draw[conta] (58:1.05) arc[start angle=58, end angle=-62, radius=1.05];
  \node[anchor=north east, inner sep=1pt] at (-64:1.05) {$d(u)$};

\end{tikzpicture}

  \caption{Un passo di $M_G$ a partire da $u$: dei $d(u)$ archi incidenti se ne sceglie
  uno uniformemente a caso, dunque $P_{u,v} = 1/d(u)$ per ogni vicino $v$ di $u$.}
  \label{fig:transizioni-mg}
\end{figure}

\begin{osservazione}[$M_G$ è aperiodica]
  La \emph{periodicità} di un vertice (cioè di uno stato) $v$ è
  \[
    T_v \;=\; \gcd\bigl\{\, \text{lunghezze di tutti i cammini chiusi (cicli) da } v \,\bigr\}.
  \]
  Nel caso di $M_G$:
  \begin{itemize}
    \item ogni arco di $G$, essendo non orientato, dà luogo alle due transizioni opposte
      $u \to v$ e $v \to u$: abbiamo dunque cicli di lunghezza $2$;
    \item $G$ non è bipartito, quindi ci sono anche cicli di lunghezza dispari, diciamo
      $2k+1$.
  \end{itemize}
  Ne segue
  \[
    T_v \;=\; \gcd(2,\, 2k+1) \;=\; 1
  \]
  per ogni $v$, cioè $M_G$ è aperiodica.

  Perché l'argomento sia completo occorre però che entrambi i cammini chiusi passino per lo
  \emph{stesso} vertice $v$. Per il primo non c'è nulla da dire: $G$ è connesso, quindi ogni
  vertice ha almeno un arco incidente. Per il secondo: se il ciclo dispari non passa per
  $v$, si raggiunge un suo vertice $w$ con un cammino di lunghezza $\ell$, si percorre il
  ciclo e si torna indietro lungo lo stesso cammino, ottenendo un cammino chiuso da $v$ di
  lunghezza $2\ell + (2k+1)$, ancora dispari. La connessione di $G$ garantisce dunque che
  ogni vertice stia su cammini chiusi di lunghezza $2$ e di lunghezza dispari.
\end{osservazione}

\begin{figure}[H]
  \centering
  % p83-arco-ciclo2 -- pag. 83
% Ogni arco non orientato di G diventa, in M_G, la coppia di transizioni opposte
% u -> v e v -> u: un cammino chiuso di lunghezza 2.  Due pannelli separati da =>.
% Nel quaderno i due punti non sono etichettati: qui si aggiungono u, v e i nomi
% delle due probabilita' di transizione.
\begin{tikzpicture}[
    x=1cm, y=1cm, line join=round, line cap=round,
    font=\footnotesize,
    punto/.style = {circle, fill, inner sep=0pt, minimum size=3pt},
    arco/.style  = {line width=0.7pt},
    tran/.style  = {line width=0.7pt, -{Latex[length=1.9mm,width=1.5mm]}},
    didasc/.style= {gray!70!black, anchor=north, inner sep=1pt},
  ]

  % ================= pannello di sinistra: l'arco di G =====================
  \node[punto] (a) at (0,0)   {};
  \node[punto] (b) at (1.7,0) {};
  \node[anchor=east, inner sep=3pt] at (a.west) {$u$};
  \node[anchor=west, inner sep=3pt] at (b.east) {$v$};

  \draw[arco] (a) to[bend left=45] (b);

  \node[didasc] at (0.85,-0.95) {in $G$: un arco};

  % ================= freccia di implicazione ==============================
  \node at (2.95,0.14) {$\Longrightarrow$};

  % ================= pannello di destra: il 2-ciclo di M_G ================
  \node[punto] (c) at (4.2,0) {};
  \node[punto] (d) at (5.9,0) {};
  \node[anchor=east, inner sep=3pt] at (c.west) {$u$};
  \node[anchor=west, inner sep=3pt] at (d.east) {$v$};

  \draw[tran] (c) to[bend left=45] node[above,inner sep=1.5pt] {$P_{u,v}$} (d);
  \draw[tran] (d) to[bend left=45] node[below,inner sep=1.5pt] {$P_{v,u}$} (c);

  \node[didasc] at (5.05,-0.95) {in $M_G$: un ciclo di lunghezza $2$};

\end{tikzpicture}

  \caption{Ogni arco non orientato di $G$ diventa, in $M_G$, una coppia di transizioni
  opposte, cioè un cammino chiuso di lunghezza $2$.}
  \label{fig:arco-ciclo2}
\end{figure}

\scan{84}

Le tre ipotesi del teorema fondamentale sono dunque tutte soddisfatte da $M_G$:
\[
  \left.
  \begin{array}{l}
    M_G \text{ è aperiodica} \\[2pt]
    M_G \text{ è finita} \\[2pt]
    M_G \text{ è irriducibile}
  \end{array}
  \right\}
  \;\Longrightarrow\;
  \text{posso applicare il teorema.}
\]
L'aperiodicità è quella appena verificata, la finitezza viene dal fatto che gli stati sono
gli $n$ vertici del grafo e l'irriducibilità dalla connessione di $G$.

\section{Intermezzo: mescolare un mazzo di carte}

\begin{esercizio}[Mescolare un mazzo di carte]\label{ex:mazzo-carte}
  Come si mescola un mazzo di carte? Ovvero: quanto vale il \emph{mixing time} della catena
  di Markov che descrive il mescolamento?
\end{esercizio}

Gli stati della catena sono le permutazioni del mazzo: per un mazzo di $K = 52$ carte sono
\[
  K! \;=\; 52! \;\approx\; 10^{68}
\]
possibili stati.

Il mescolamento che consideriamo è il \emph{top-in-at-random}: a ogni passo si toglie la
carta in cima al mazzo e la si reinserisce in una delle $K$ posizioni possibili, scelta
uniformemente a caso.

\begin{figure}[H]
  \centering
  % top-in-at-random -- pag. 84
% Il mescolamento top-in-at-random: il mazzo e' visto di taglio (una carta = un
% segmento orizzontale).  Pannello superiore: la carta in cima viene tolta e
% reinserita in una posizione scelta a caso; B e' la carta in fondo al mazzo.
% Pannello inferiore: dopo un certo numero di passi B e' risalita e le carte che
% le sono passate sotto sono in ordine uniformemente casuale (l'invariante).
% Nel quaderno i due mazzi hanno 6 e 7 carte: qui entrambi ne hanno 7, perche' il
% numero di carte non cambia lungo il mescolamento.
\begin{tikzpicture}[
    x=1cm, y=1cm, line join=round, line cap=round,
    font=\footnotesize,
    carta/.style  = {line width=0.8pt},
    cartaB/.style = {line width=1.7pt},
    frec/.style   = {line width=0.7pt, -{Latex[length=2mm,width=1.6mm]}},
    graffa/.style = {decorate, decoration={brace, amplitude=4pt},
                     gray!70!black, line width=0.6pt},
  ]

  \def\W{1.9}     % larghezza del mazzo
  \def\P{0.3}     % passo verticale fra due carte

  % =================== intestazione =======================================
  \node[anchor=south] at (1.55, 0.45) {\emph{top-in-at-random}};

  % =================== pannello superiore: il mazzo iniziale ==============
  % sette carte, y = 0, -0.3, ..., -1.8;  B e' l'ultima
  \foreach \i in {0,...,5}{\draw[carta] (0,-\i*\P) -- (\W,-\i*\P);}
  \draw[cartaB] (0,-6*\P) -- (\W,-6*\P);
  \node[anchor=east, inner sep=4pt] at (0,-6*\P) {$B$};

  % la carta in cima esce e rientra in una posizione a caso
  \draw[frec] (\W,0.03) .. controls (2.85,0.14) and (2.95,-1.02) .. (1.62,-1.06);
  \node[anchor=west, inner sep=2pt] at (2.62,-0.42) {a caso};

  % =================== transizione fra i due pannelli =====================
  \draw[line width=1.5pt, -{Latex[length=3mm,width=2.6mm]}]
        (\W/2,-2.24) -- (\W/2,-2.86);
  \node[anchor=west, inner sep=5pt, align=left, gray!70!black]
        at (\W/2,-2.55) {dopo un certo\\numero di passi};

  % =================== pannello inferiore: dopo alcuni passi ==============
  % sette carte, y = -3.3, ..., -5.1;  B e' la quarta (tre sopra, tre sotto)
  \def\Y{-3.3}
  \foreach \i in {0,1,2}{\draw[carta] (0,\Y-\i*\P) -- (\W,\Y-\i*\P);}
  \draw[cartaB] (0,\Y-3*\P) -- (\W+0.22,\Y-3*\P);
  \node[anchor=east, inner sep=4pt] at (0,\Y-3*\P) {$B$};
  \foreach \i in {4,5,6}{\draw[carta] (0,\Y-\i*\P) -- (\W,\Y-\i*\P);}

  % le carte finite sotto B sono in ordine uniformemente casuale
  \draw[graffa] (\W+0.42,\Y-3.6*\P) -- (\W+0.42,\Y-6.4*\P);
  \node[anchor=west, inner sep=8pt] at (\W+0.42,\Y-5*\P) {a caso};

\end{tikzpicture}

  \caption{Il mescolamento \emph{top-in-at-random}: la carta in cima al mazzo viene
  reinserita in una posizione scelta uniformemente a caso. In alto la situazione iniziale,
  con la carta $B$ in fondo al mazzo; in basso, dopo un certo numero di passi, $B$ è
  risalita e le carte che le sono passate sotto si trovano in ordine uniformemente
  casuale.}
  \label{fig:top-in-at-random}
\end{figure}

Sia $B$ la carta che all'inizio si trova in fondo al mazzo. Vogliamo valutare il numero $T$
di operazioni che servono per portare $B$ in prima posizione.

Supponiamo che sotto $B$ ci siano già $i-1$ carte. La carta reinserita finisce sotto $B$
esattamente quando la posizione estratta è una delle $i$ che stanno al di sotto di $B$, il
che accade con probabilità
\[
  p_i \;=\; \frac{i}{K}
  \qquad\text{cioè, via via,}\quad \frac{1}{K},\ \frac{2}{K},\ \dots,\ \frac{K-1}{K}.
\]
Detto $T_i$ il numero di passi necessari per passare da $i-1$ a $i$ carte sotto $B$, la
variabile $T_i$ è geometrica di parametro $p_i$, e quindi, per il conto già fatto sul valore
atteso di una variabile geometrica,
\[
  \E[T_i] \;=\; \frac{1}{p_i} \;=\; \frac{K}{i}.
\]

La carta $B$ arriva in prima posizione quando tutte le altre $K-1$ carte le sono passate
sotto, dunque
\[
  T \;=\; T_1 + T_2 + \dots + T_{K-1}
\]
e, per linearità del valore atteso,
\[
  \E[T] \;=\; \sum_{i=1}^{K-1} \E[T_i]
        \;=\; K \sum_{i=1}^{K-1} \frac{1}{i}
        \;=\; K\,H_{K-1}
        \;=\; \Oh(K \log K),
\]
dove $H_K = \sum_{i=1}^{K} 1/i$ è il $K$-esimo numero armonico: è di nuovo il conto del
\emph{coupon collector}.

\begin{osservazione}
  Un solo passo in più basta a concludere il mescolamento: nell'istante in cui $B$, arrivata
  in cima, viene a sua volta reinserita in una posizione casuale, tutte le $K$ carte sono
  state collocate a caso e il mazzo è uniformemente mescolato. Contando anche quel passo si
  ottiene
  \[
    \sum_{i=1}^{K} \frac{K}{i} \;=\; K\,H_K \;\approx\; K \ln K
  \]
  (l'addendo con $i=K$ vale esattamente $1$ ed è proprio il passo finale). Per $K = 52$
  servono dunque in media $52\,H_{52} \approx 236$ mescolate: l'ordine di grandezza è
  $K \ln K \approx 205$, e i circa $30$ passi di differenza sono il termine $\gamma K$ che
  l'approssimazione $H_K \approx \ln K$ trascura ($\gamma \approx 0{,}577$ è la costante di
  Eulero--Mascheroni).
\end{osservazione}

\section{La distribuzione stazionaria del random walk}

\begin{lemma}[Distribuzione stazionaria di un random walk]\label{lem:stazionaria-grado}
  Sia $M_G$ la catena associata al grafo $G = (V,E)$, con $\abs{V} = n$ e $\abs{E} = m$.
  Allora
  \[
    \forall v \in V \qquad \pi_v \;=\; \frac{d(v)}{2m}.
  \]
\end{lemma}

\begin{dimostrazione}
  Dobbiamo trovare $\vec{\pi}$ soluzione di
  \[
    \vec{\pi} \;=\; \vec{\pi}\,P ,
  \]
  e per farlo basta verificare che il vettore di componenti $d(v)/2m$, al variare di
  $v \in V$, soddisfa questa equazione, cioè che per ogni $u \in V$ vale
  $\pi_u = \sum_{v \in V} \pi_v P_{v,u}$.

  \scan{85}
  Poiché $P_{v,u} \neq 0$ soltanto se $v \in \Gamma(u)$, e in tal caso
  $P_{v,u} = 1/d(v)$,
  \[
    \frac{d(u)}{2m}
    \;\overset{?}{=}\;
    \sum_{v \in \Gamma(u)} \frac{d(v)}{2m} \cdot P_{v,u}
    \;=\;
    \sum_{v \in \Gamma(u)} \frac{d(v)}{2m} \cdot \frac{1}{d(v)}
    \;=\;
    \frac{1}{2m} \sum_{v \in \Gamma(u)} 1
    \;=\;
    \frac{d(u)}{2m} \,,
  \]
  perché la somma ha esattamente $\abs{\Gamma(u)} = d(u)$ addendi, uno per ogni vicino di $u$.

  Il vettore così definito è inoltre una distribuzione di probabilità: ogni arco contribuisce
  al grado di entrambi i suoi estremi, quindi $\sum_{v \in V} d(v) = 2m$ e di conseguenza
  $\sum_{v \in V} \pi_v = 1$.
\end{dimostrazione}

Ergo, per ogni $v \in V$, il tempo medio di ritorno in $v$ --- che per il teorema
fondamentale è il reciproco della probabilità stazionaria --- vale
\[
  h_{v,v} \;=\; \frac{1}{\pi_v} \;=\; \frac{2m}{d(v)} \,.
\]

\section{Hitting time, commute time e cover time}

\begin{definizione}[Hitting time, commute time, cover time]\label{def:tempi}
  Si consideri il random walk su $G = (V,E)$.
  \begin{itemize}
    \item $h_{u,v}$ --- \emph{hitting time}: numero medio di passi per andare da $u$ a $v$.
    \item $C_{u,v}$ --- \emph{commute time}:
      \[
        h_{u,v} + h_{v,u} \;=\; C_{u,v} \;=\; C_{v,u} \,.
      \]
    \item $C_u(G)$ --- tempo medio per tornare a $u$ dopo aver visitato l'intero $G$
      partendo da $u$.
    \item $C(G)$ --- \emph{cover time} di $G$:
      \[
        C(G) \;=\; \max_{u} \, C_u(G) \,.
      \]
  \end{itemize}
\end{definizione}

\begin{domanda}\label{q:simmetria-hitting}
  Vale $h_{u,v} = h_{v,u}$?
\end{domanda}

Si consideri il \emph{lollipop graph} $L_n$: una clique $K_{n/2}$ alla quale è incollato,
nel vertice $u$, un cammino di $n/2$ vertici che termina in $v$.

\begin{figure}[H]
  \centering
  % p85-lollipop -- pag. 85 (la stessa figura torna a p. 88, deduplicata)
% Il lollipop graph L_n: la clique K_{n/2} e' rappresentata astrattamente da un
% cerchio (vertici e archi interni non si disegnano); nel vertice u -- l'unico
% pieno, sul bordo destro del cerchio -- e' incollato un cammino di n/2 vertici
% che termina in v.  In tutto n vertici.
% Il quaderno e' tutto a penna blu: qui si usa il nero come nelle altre figure.
\begin{tikzpicture}[
    x=1cm, y=1cm, line join=round, line cap=round,
    font=\footnotesize,
    nodo/.style   = {circle, draw, line width=0.8pt, fill=white,
                     inner sep=0pt, minimum size=5pt},
    nodou/.style  = {circle, draw, fill=black, line width=0.8pt,
                     inner sep=0pt, minimum size=5pt},
    arco/.style   = {line width=0.8pt},
    graffa/.style = {decorate, decoration={brace, amplitude=5pt},
                     gray!70!black, line width=0.7pt},
  ]

  % =================== la clique K_{n/2} ==================================
  \draw[arco, line width=1pt] (0,0) circle (1.15);
  \node at (0,0) {$K_{n/2}$};

  % =================== il bastoncino: cammino di n/2 vertici ==============
  \node[nodou] (u)  at (1.15,0) {};
  \node[nodo]  (p1) at (2.05,0) {};
  \node[nodo]  (p2) at (2.95,0) {};
  \node[nodo]  (p3) at (3.85,0) {};
  \node[nodo]  (v)  at (5.35,0) {};

  \draw[arco] (u)  -- (p1);
  \draw[arco] (p1) -- (p2);
  \draw[arco] (p2) -- (p3);

  % i vertici omessi
  \draw[arco] (p3) -- (4.15,0);
  \foreach \x in {4.38,4.60,4.82}{\fill (\x,0) circle (0.7pt);}
  \draw[arco] (5.05,0) -- (v);

  \node[anchor=north, inner sep=4pt] at (1.30,-0.09) {$u$};
  \node[anchor=north, inner sep=4pt] at (5.48,-0.09) {$v$};

  % =================== la graffa che misura il bastoncino =================
  \draw[graffa] (1.18,0.42) -- (5.35,0.42)
        node[midway, above=6pt, black] {$n/2$};

\end{tikzpicture}

  \caption{Il \emph{lollipop graph} $L_n$: una clique $K_{n/2}$ alla quale è incollato,
  nel vertice $u$, un cammino di $n/2$ vertici che termina in $v$; in tutto $n$ vertici.}
  \label{fig:lollipop}
\end{figure}

\scan{86}

Dimostreremo che su $L_n$ l'hitting time non è affatto simmetrico:
\[
  h_{u,v} \;=\; \Th(n^{3}),
  \qquad\text{mentre}\qquad
  h_{v,u} \;=\; \Th(n^{2}).
\]
I due versi differiscono dunque di un fattore $n$: la risposta alla
Domanda~\ref{q:simmetria-hitting} è negativa.

Dimostreremo anche che sul grafo completo $K_n$ il cover time vale
\[
  C_u(K_n) \;=\; \Th(n \log n) \qquad\text{per ogni } u \in V .
\]

\begin{osservazione}
  Su $K_n$ ogni vertice ha $n-1$ vicini, tutti equiprobabili, dunque il numero di passi
  necessari per raggiungere un fissato $v \neq u$ è una variabile geometrica di parametro
  $1/(n-1)$ e
  \[
    h_{u,v} \;=\; n-1 \;=\; \Th(n) \qquad \text{per ogni } u \neq v :
  \]
  sul grafo completo l'hitting time è lineare e uguale per tutte le coppie. È invece il
  numero di passi necessari a \emph{visitare tutti} i vertici a valere
  $(n-1)H_{n-1} = \Th(n \log n)$: è il conto del collezionista di figurine.
\end{osservazione}

\paragraph{Morale.} Non è ovvio come ragionare per stimare $C_u(G)$ e $C(G)$. Lo strumento
classico è l'analogia fra passeggiate aleatorie su grafi e reti di resistenze, esposta in
«\emph{Random Walks and Electric Networks}» di P.~G.~Doyle e J.~L.~Snell (1984); qui
seguiamo invece una strada elementare, che parte dal commute time di un singolo arco.

\section{Il commute time lungo un arco}

\begin{lemma}[Commute time lungo un arco]\label{lem:commute-arco}
  Sia $G = (V,E)$ come sopra, con $\abs{E} = m$. Se $u$ e $v$ sono \emph{adiacenti}, cioè
  $(u,v) \in E$, allora
  \[
    C_{u,v} \;=\; h_{u,v} + h_{v,u} \;\le\; 2m .
  \]
\end{lemma}

L'ipotesi di adiacenza è essenziale: se cade, la disuguaglianza può essere falsa. Sul
lollipop graph, per esempio, $h_{u,v} + h_{v,u} = \Th(n^{3})$ mentre $2m = \Th(n^{2})$,
perché la clique $K_{n/2}$ ha già $\Th(n^{2})$ archi.

\begin{dimostrazione}
  Introduciamo una catena di Markov ausiliaria i cui stati sono gli \emph{archi orientati}
  di $G$: ogni arco non orientato di $G$ dà luogo ai due stati $(u,v)$ e $(v,u)$, per un
  totale di $2m$ stati. In questa dimostrazione, dunque, la scrittura $(u,v)$ denota un arco
  \emph{orientato}; lo stato $(u,v)$ va letto come «la passeggiata ha appena attraversato
  l'arco da $u$ a $v$», e si trova quindi nel vertice $v$.

  La matrice di transizione $Q$ della catena così introdotta è
  \[
    Q_{(u,v),(v,w)} \;=\; P_{v,w} \;=\; \frac{1}{d(v)}
    \qquad (w \in \Gamma(v)),
  \]
  e tutte le altre entrate sono nulle, perché dallo stato $(u,v)$ si può passare soltanto a
  uno stato della forma $(v,w)$, e lo si fa con la stessa probabilità con cui il random walk
  su $G$, trovandosi in $v$, sceglie il vicino $w$.

  \medskip
  \textbf{Osservazione: $Q$ è \emph{doubly stochastic}}, cioè righe e colonne sommano a $1$.
  \begin{itemize}
    \item \emph{Righe.} Fissato lo stato $(u,v)$,
      \[
        \sum_{w \in \Gamma(v)} Q_{(u,v),(v,w)}
        \;=\; \sum_{w \in \Gamma(v)} \frac{1}{d(v)}
        \;=\; \frac{d(v)}{d(v)} \;=\; 1 .
      \]
    \item \emph{Colonne.} Fissato lo stato $(v,w)$, nella somma su tutti gli stati $(x,y)$
      sopravvivono soltanto i termini con $y = v$ e $x \in \Gamma(v)$, e quindi
      \[
        \sum_{x \in V} \; \sum_{y \in \Gamma(x)} Q_{(x,y),(v,w)}
        \;=\; \sum_{x \in \Gamma(v)} Q_{(x,v),(v,w)}
        \;=\; \sum_{x \in \Gamma(v)} \frac{1}{d(v)}
        \;=\; \frac{d(v)}{d(v)} \;=\; 1 .
      \]
  \end{itemize}

  \medskip
  \scan{87}
  Anche questa catena soddisfa le ipotesi del teorema fondamentale: è \emph{finita}, perché
  ha $2m$ stati; è \emph{irriducibile}, perché nello stato $(u,v)$ la passeggiata si trova in
  $v$ e, essendo $G$ connesso, può raggiungere qualunque vertice $x$ e attraversare poi un
  qualunque arco $x \to y$, arrivando nello stato $(x,y)$; ed è \emph{aperiodica}, perché
  $(u,v) \to (v,u) \to (u,v)$ è un cammino chiuso di lunghezza $2$, mentre un cammino chiuso
  di lunghezza dispari da $v$ a $v$ --- che esiste, per quanto visto sopra, perché $G$ non è
  bipartito --- seguito da $v \to u \to v$ riporta in $(u,v)$ dopo un numero dispari di passi.

  La distribuzione stazionaria di una matrice \emph{doubly stochastic} è quella uniforme.
  Basta dunque verificare che $\vec{\pi} = \vec{\pi}\,Q$ per $\vec{\pi}$ uniforme e poi
  invocare il teorema fondamentale delle catene di Markov, che garantisce l'unicità della
  distribuzione stazionaria.

  Poiché le colonne di $Q$ sommano a $1$, «mi piacerebbe» che fosse $\pi_x = 1$ per ogni
  stato $x$: il vettore di tutti $1$ è in effetti fissato da $Q$. Ma così
  $\sum_x \pi_x \neq 1$, e $\vec{\pi}$ non sarebbe una distribuzione di probabilità. Il rimedio
  è normalizzare: poniamo
  \[
    \pi_x \;=\; \frac{1}{\text{numero degli stati}} \;=\; \frac{1}{2m},
  \]
  visto che gli stati della catena sono i $2m$ archi orientati di $G$.

  Le due verifiche sono ora immediate. La normalizzazione:
  \[
    \sum_x \pi_x \;=\; \sum_x \frac{1}{2m} \;=\; \frac{2m}{2m} \;=\; 1 ;
  \]
  la stazionarietà: per ogni stato $x$,
  \[
    \sum_y \pi_y\, Q_{y,x}
    \;=\; \frac{1}{2m} \underbrace{\Bigl(\sum_y Q_{y,x}\Bigr)}_{=\;1\ (Q \text{ doubly stochastic})}
    \;=\; \frac{1}{2m}
    \;=\; \pi_x .
  \]

  Per l'unicità della distribuzione stazionaria si conclude allora che, per ogni arco
  orientato $(u,v)$, si ha $\pi_{(u,v)} = 1/(2m)$ e quindi, per l'identità fra tempo medio di
  ritorno e probabilità stazionaria,
  \[
    h_{(u,v),(u,v)} \;=\; \frac{1}{\pi_{(u,v)}} \;=\; 2m .
  \]

  Resta da collegare questo tempo di ritorno ai due hitting time $h_{u,v}$ e $h_{v,u}$ sui
  \emph{vertici}. La quantità $h_{(u,v),(u,v)}$ è il tempo medio che intercorre fra una visita
  dello stato $(u,v)$ e la visita successiva dello stesso stato: appena attraversato l'arco
  $u \to v$ la passeggiata si trova in $v$, vaga per il grafo, prima o poi ritorna in $u$ e
  solo allora può attraversare una seconda volta l'arco $u \to v$ (Figura~\ref{fig:ritorno-arco}).

  Spezzando il tragitto nel primo arrivo in $u$ e in ciò che segue si ottiene
  \[
    2m \;=\; h_{(u,v),(u,v)} \;\ge\; h_{v,u} + h_{u,v} :
  \]
  tornare in $u$ partendo da $v$ costa in media $h_{v,u}$ e, una volta in $u$, per rivedere
  lo stato $(u,v)$ la passeggiata deve come minimo ritrovarsi in $v$, il che costa in media
  almeno $h_{u,v}$. Dunque
  \[
    h_{u,v} + h_{v,u} \;\le\; 2m
    \qquad\text{per } (u,v) \in E ,
  \]
  che è la tesi.
\end{dimostrazione}

\begin{figure}[H]
  \centering
  % p87-ritorno-arco -- pag. 87
% Primo ritorno nello stato (u,v) della catena sugli archi orientati:
%   u -> v  (prima traversata)  ~>  vagabondaggio nel grafo fino a rientrare in u
%   ~>  u -> v  (seconda traversata).
% Visualizza la decomposizione  h_{(u,v),(u,v)} >= h_{v,u} + h_{u,v}.
% Il grafo resta implicito: la curva chiusa e' la traiettoria della passeggiata,
% non un arco di G.
\begin{tikzpicture}[
    x=1cm, y=1cm,
    vert/.style     = {circle, fill=graydark, inner sep=0pt, minimum size=4pt},
    traversa/.style = {draw=accent, line width=1.1pt, line cap=round,
                       -{Stealth[length=5pt,width=4pt]}},
    vaga/.style     = {draw=graydark, line width=0.6pt, line cap=round,
                       line join=round},
    et/.style       = {font=\footnotesize},
  ]

  \coordinate (u) at (0.00, 0.00);
  \coordinate (v) at (2.40, 1.05);

  % ---- il vagabondaggio: un'unica traiettoria che esce da v, gira a lungo
  %      per il grafo e rientra in u ------------------------------------------
  \draw[vaga, -{Stealth[length=5pt,width=4pt]}]
       (2.50, 1.09)
       .. controls ( 4.70, 1.30) and ( 5.60,-0.30) .. ( 4.90,-1.60)
       .. controls ( 4.20,-2.90) and ( 2.20,-3.40) .. ( 0.40,-3.00)
       .. controls (-1.30,-2.60) and (-2.10,-1.30) .. (-1.50,-0.55)
       .. controls (-1.05,-0.40) and ( 0.90,-0.60) .. ( 1.90,-0.85)
       .. controls ( 3.10,-1.55) and ( 3.00,-2.45) .. ( 1.70,-2.55)
       .. controls ( 0.55,-2.65) and (-0.62,-2.05) .. (-0.78,-1.35)
       .. controls (-0.88,-0.80) and (-0.50,-0.35) .. (-0.10,-0.10);

  % ---- le due traversate dell'arco u -> v ------------------------------------
  \draw[traversa] (u) to[bend left=16]  (v);
  \draw[traversa] (u) to[bend right=16] (v);

  \node[vert] at (u) {};
  \node[vert] at (v) {};

  \node[et, anchor=east,  xshift=-4pt] at (u) {$u$};
  \node[et, anchor=south, yshift=4pt]  at (v) {$v$};

  \node[et, rotate=23.6, anchor=south] at (1.13, 0.76) {1\textsuperscript{a} traversata};
  \node[et, rotate=23.6, anchor=north] at (1.30, 0.29) {2\textsuperscript{a} traversata};

  \node[et, anchor=north, align=center] at (1.70,-3.30)
       {vagabondaggio nel grafo\\ prima del rientro in $u$ \ ($h_{v,u}$)};

\end{tikzpicture}

  \caption{Il ritorno nello stato $(u,v)$ della catena sugli archi orientati: dopo la prima
  traversata di $u \to v$ la passeggiata vaga per il grafo, rientra in $u$ e attraversa una
  seconda volta $u \to v$.}
  \label{fig:ritorno-arco}
\end{figure}

\section{Una stima del cover time}

\scan{88}

Il lemma appena dimostrato --- $C_{u,v} = h_{u,v} + h_{v,u} \le 2m$ per ogni arco
$(u,v) \in E$ --- si propaga da un singolo arco a tutto il grafo e fornisce la stima del
cover time.

\begin{teorema}[Stima del cover time]\label{thm:cover-time}
  Sia $G = (V,E)$ come sopra, con $n = \abs{V}$ e $m = \abs{E}$. Allora
  \[
    C(G) \;\le\; 2m\,(n-1) \,.
  \]
\end{teorema}

\begin{dimostrazione}
  Si fissi uno spanning tree $T$ di $G$ e lo si percorra con un tour euleriano dell'albero
  raddoppiato: partendo da un vertice qualsiasi $u$ si toccano tutti i vertici e si torna in
  $u$ attraversando ciascuno degli $n-1$ archi di $T$ una volta per verso. Il costo atteso di
  questo tour è la somma dei commute time dei suoi archi, e ogni arco di $T$ è un arco di
  $G$, al quale si applica il Lemma~\ref{lem:commute-arco}: per ogni $u \in V$
  \[
    C_u(G) \;\le\; \sum_{(x,y) \in T} C_{x,y} \;\le\; (n-1) \cdot 2m \,,
  \]
  e passando al massimo su $u$ si ottiene la tesi.
\end{dimostrazione}

\begin{esempio}[I due casi estremi]
  \textbf{Cammino.} Se $G$ è un cammino su $n$ vertici (Figura~\ref{fig:cammino}), allora
  $m = n-1$ e la stima dà
  \[
    2(n-1)(n-1) \;\in\; \Th(n^2) \,.
  \]
  Il cammino è bipartito, e cade quindi fuori dalle ipotesi permanenti fissate all'inizio del
  capitolo; la stima vale lo stesso, perché di quelle ipotesi la catena di deduzioni usa
  soltanto l'unicità della distribuzione stazionaria e l'identità $h_{i,i} = 1/\pi_i$, che
  valgono per ogni catena finita e irriducibile, anche periodica: la non bipartitezza serve a
  garantire la convergenza alla distribuzione stazionaria, che qui non viene mai invocata.

  \medskip
  \textbf{Lollipop.} Se $G$ è il lollipop graph $L_n$ (Figura~\ref{fig:lollipop}), la sola
  clique $K_{n/2}$ porta $\binom{n/2}{2} \in \Th(n^2)$ archi, dunque $m \in \Th(n^2)$ e la
  stima dà
  \[
    2m\,(n-1) \;\in\; \Th(n^3) \,.
  \]
\end{esempio}

\begin{figure}[H]
  \centering
  % p88-cammino -- pag. 88
% Il cammino su n vertici: l'albero con il minimo numero di archi, m = n-1.
% E' il caso in cui la stima C(G) <= 2m(n-1) vale Theta(n^2).
\begin{tikzpicture}[
    x=1cm, y=1cm,
    filo/.style = {draw=graydark, line width=0.7pt, line cap=round},
    vert/.style = {circle, fill=graydark, inner sep=0pt, minimum size=4pt},
    punto/.style= {circle, fill=graydark, inner sep=0pt, minimum size=2pt},
    et/.style   = {font=\footnotesize},
  ]

  % ---- i primi quattro vertici del cammino --------------------------------
  \foreach \i in {0,1,2,3} { \coordinate (n\i) at (1.55*\i, 0); }
  \draw[filo] (n0) -- (n1) -- (n2) -- (n3);

  % ---- lo stacco: il cammino prosegue -------------------------------------
  \foreach \x in {5.60, 6.05, 6.50} { \node[punto] at (\x, 0) {}; }

  % ---- gli ultimi due vertici ---------------------------------------------
  \coordinate (m0) at (7.45, 0);
  \coordinate (m1) at (9.00, 0);
  \draw[filo] (m0) -- (m1);

  % ---- i pallini ----------------------------------------------------------
  \foreach \n in {n0,n1,n2,n3,m0,m1} { \node[vert] at (\n) {}; }

  % ---- le etichette -------------------------------------------------------
  \node[et, below=4pt] at (n0) {$v_1$};
  \node[et, below=4pt] at (n1) {$v_2$};
  \node[et, below=4pt] at (n2) {$v_3$};
  \node[et, below=4pt] at (n3) {$v_4$};
  \node[et, below=4pt] at (m0) {$v_{n-1}$};
  \node[et, below=4pt] at (m1) {$v_n$};

\end{tikzpicture}

  \caption{Il cammino su $n$ vertici: $m = n-1$ archi.}
  \label{fig:cammino}
\end{figure}

\begin{osservazione}
  In entrambi i casi la stima è dell'ordine giusto --- il cover time del cammino è
  $\Th(n^2)$, quello del lollipop è $\Th(n^3)$ --- e quindi $2m(n-1)$ non è migliorabile in
  generale. Non è però sempre stretta: su $K_n$ dà $\Th(n^3)$, mentre il cover time vero è
  $\Th(n \log n)$ (collezionista di figurine).

  Resta invece impregiudicato il conto annunciato sopra e mai svolto: la stima puntuale
  $h_{u,v} = \Th(n^3)$ e $h_{v,u} = \Th(n^2)$ sul lollipop, che risponde negativamente alla
  Domanda~\ref{q:simmetria-hitting}.
\end{osservazione}


\appendix
\part{Appendice}
\chapter{Graph Isomorphism e la gerarchia di Weisfeiler--Leman}

% ---------------------------------------------------------------------------
% Nel quaderno il capitolo si apre a metà di p.79, con il titolo sottolineato
% «Graph Isomorphism»; accanto, fra parentesi, un nome proprio di lettura
% incerta («Francesco Nassimbeni»), verosimilmente chi ha tenuto questa parte
% del corso. Il nome non è riportato nel testo perché la lettura resta dubbia.
% ---------------------------------------------------------------------------

Le tre pagine di quaderno da cui nasce questo capitolo stanno fuori dal programma
ordinario del corso --- sono gli appunti di un seminario --- ed è per questo che le si
raccoglie in appendice. Il quaderno le lascia allo stato di traccia: titoli, formule
chiave, qualche disegno, e l'ultima pagina si interrompe a metà. Qui se ne dà una
versione leggibile, colmando le lacune con i risultati standard della letteratura;
tutte le integrazioni sono registrate nelle note editoriali.

Il problema di partenza è il \emph{graph isomorphism}: dati due grafi $G$ e $H$,
stabilire se esiste una biiezione fra i loro vertici che preservi gli archi, nel qual
caso si scrive $G \cong H$. Non se ne conosce un algoritmo polinomiale, né lo si sa
$\NP$-completo. La famiglia di test di Weisfeiler--Leman è la risposta pratica al
problema: una gerarchia di procedure polinomiali, indicizzate da un parametro $k$, che
non lo risolvono ma lo attaccano dal lato facile --- sanno certificare il \emph{non}
isomorfismo, mai l'isomorfismo --- e che diventano tanto più fini quanto più $k$ cresce.

\section{WL-1: il raffinamento dei colori}

\scan{79}

Il test di Weisfeiler--Leman di dimensione~$1$ --- il classico \emph{raffinamento dei
colori} (\emph{color refinement}) --- assegna un colore a ogni vertice e lo raffina
iterativamente: a ogni passo il nuovo colore di un vertice è determinato dal suo colore
corrente e dal multi-insieme dei colori dei suoi vicini. La colorazione partiziona
l'insieme dei vertici in classi di colore, e il procedimento si arresta quando la
partizione non si raffina più.

\begin{figure}[H]
  \centering
  % wl1-partizione-colori -- pag. 79
% Il disco e' l'insieme dei vertici V; i settori sono le classi di colore
% prodotte dal raffinamento e le crocette i vertici che vi cadono dentro.
% NB: nel quaderno la crocetta manca nel settore piu' piccolo (in alto a
% sinistra) -- quattro settori, tre crocette -- e qui la si riproduce cosi'
% com'e', come da specifica registrata nelle note editoriali.
\begin{tikzpicture}[x=1cm,y=1cm,line join=round,line cap=round,font=\footnotesize]

  % ---- il disco V ----------------------------------------------------------
  \draw[line width=0.9pt] (0,0) circle (1.5);

  % ---- il diametro, leggermente inclinato ----------------------------------
  \draw[line width=0.8pt] (185:1.5) -- (5:1.5);

  % ---- i due raggi che suddividono la meta' superiore ----------------------
  \draw[line width=0.8pt] (0,0) -- (90:1.5);
  \draw[line width=0.8pt] (0,0) -- (145:1.5);

  % ---- i vertici, come nel manoscritto (tre crocette in quattro settori) ---
  \foreach \ang/\rad in {266/0.80, 47/0.80, 117/0.78}{
    \draw[line width=0.7pt]
      ([shift={(\ang:\rad)}]-0.09,-0.09) -- ([shift={(\ang:\rad)}]0.09,0.09)
      ([shift={(\ang:\rad)}]-0.09, 0.09) -- ([shift={(\ang:\rad)}]0.09,-0.09);
  }

  % ---- nome dell'insieme ---------------------------------------------------
  \node[anchor=south west,inner sep=1pt] at (38:1.60) {$V$};

\end{tikzpicture}

  \caption{Il raffinamento dei colori partiziona l'insieme dei vertici in classi di
  colore.}
  \label{fig:wl1-partizione}
\end{figure}

Il colore che WL-1 assegna a un vertice dopo $r$ raffinamenti codifica il suo
\emph{albero di srotolamento} (\emph{unrolling}) di profondità $r$: l'albero che ha per
radice il vertice, per figli della radice i suoi vicini, per figli di questi i vicini
dei vicini, e così via.

\begin{figure}[H]
  \centering
  % wl1-grafo-e-unrolling -- pag. 79
% Due pannelli: a sinistra il grafo G (un albero con due centri, sei vertici
% colorati R, B, G, N); a destra il suo srotolamento di profondita' 2 radicato
% nel vertice N.
% Controllo semantico: i figli della radice sono i vicini di N, cioe' R, B
% (foglia) e B (il secondo centro); sotto ciascuno compaiono i suoi vicini,
% compreso il ritorno al padre. Le due foglie R e B hanno il solo vicino N; il
% secondo centro B ha per vicini G, R e N.
% NB: nel quaderno i pallini non sono colorati, solo etichettati con l'iniziale
% del colore (Rosso, Blu, Giallo, Nero): si mantiene questa scelta.
\begin{tikzpicture}[x=1cm,y=1cm,line join=round,line cap=round,font=\footnotesize,
    nodo/.style={circle,fill,inner sep=0pt,minimum size=3.6pt},
    arco/.style={line width=0.7pt},
    col/.style={inner sep=2.5pt}]

  % ======================= pannello sinistro: il grafo G ===================
  \begin{scope}[shift={(0,0)}]
    \node[nodo] (c1) at ( 0.00, 0.00) {};   % centro superiore, colore N
    \node[nodo] (r1) at (-1.10, 1.05) {};   % foglia R
    \node[nodo] (b1) at ( 1.10, 1.05) {};   % foglia B
    \node[nodo] (c2) at ( 0.00,-1.60) {};   % centro inferiore, colore B
    \node[nodo] (g)  at (-0.90,-2.60) {};   % foglia G
    \node[nodo] (r2) at ( 0.90,-2.60) {};   % foglia R

    \draw[arco] (c1) -- (r1)  (c1) -- (b1)  (c1) -- (c2)
                (c2) -- (g)   (c2) -- (r2);

    \node[col,anchor=west,inner sep=4pt] at (c1.east) {$N$};
    \node[col,anchor=south east] at (r1.north west) {$R$};
    \node[col,anchor=south west] at (b1.north east) {$B$};
    \node[col,anchor=west]       at (c2.east)       {$B$};
    \node[col,anchor=north east] at (g.south west)  {$G$};
    \node[col,anchor=north west] at (r2.south east) {$R$};
  \end{scope}

  % ============ pannello destro: lo srotolamento radicato in N =============
  \begin{scope}[shift={(5.30,0.45)}]
    \node[nodo] (t)   at ( 0.00, 0.00) {};   % radice: il vertice N
    \node[nodo] (tR)  at (-1.60,-1.20) {};   % vicino R (foglia)
    \node[nodo] (tB)  at ( 0.00,-1.20) {};   % vicino B (secondo centro)
    \node[nodo] (tB2) at ( 1.60,-1.20) {};   % vicino B (foglia)
    \node[nodo] (l1)  at (-2.60,-2.50) {};   % vicino di R: N
    \node[nodo] (l2)  at (-0.80,-2.50) {};   % vicini del secondo centro: G
    \node[nodo] (l3)  at ( 0.00,-2.50) {};   %                            R
    \node[nodo] (l4)  at ( 0.80,-2.50) {};   %                            N
    \node[nodo] (l5)  at ( 2.60,-2.50) {};   % vicino di B (foglia): N

    \draw[arco] (t) -- (tR)  (t) -- (tB)  (t) -- (tB2)
                (tR) -- (l1)
                (tB) -- (l2)  (tB) -- (l3)  (tB) -- (l4)
                (tB2) -- (l5);

    \node[col,anchor=south]      at (t.north)      {$N$};
    \node[col,anchor=south east] at (tR.north west) {$R$};
    \node[col,anchor=south east] at (tB.north west) {$B$};
    \node[col,anchor=south west] at (tB2.north east){$B$};
    \node[col,anchor=north]      at (l1.south)     {$N$};
    \node[col,anchor=north]      at (l2.south)     {$G$};
    \node[col,anchor=north]      at (l3.south)     {$R$};
    \node[col,anchor=north]      at (l4.south)     {$N$};
    \node[col,anchor=north]      at (l5.south)     {$N$};
  \end{scope}

\end{tikzpicture}

  \caption{A sinistra un grafo i cui vertici portano quattro colori distinti $R$, $B$,
  $G$, $N$; a destra l'albero di srotolamento di profondità~$2$ radicato nel vertice
  $N$.}
  \label{fig:wl1-unrolling}
\end{figure}

Per decidere l'isomorfismo di due grafi $G$ e $H$ si esegue il raffinamento su entrambi
e si confrontano i multi-insiemi dei colori finali: il test ha esito \emph{vero} quando
i due multi-insiemi coincidono. Valgono allora
\begin{align*}
  G \cong H &\;\Longrightarrow\; \mathrm{WL}\text{-}1(G,H) = \text{vero},\\
  \mathrm{WL}\text{-}1(G,H) = \text{vero} &\;\nRightarrow\; G \cong H,\\
  \mathrm{WL}\text{-}1(G,H) = \text{falso} &\;\Longrightarrow\; G \ncong H.
\end{align*}

\begin{osservazione}
  WL-1 è dunque un test con errore \emph{unilatero}: la risposta negativa è una prova di
  non isomorfismo, mentre la risposta positiva non permette di concludere nulla. La
  terza implicazione è la contronominale della prima.
\end{osservazione}

L'esempio minimale di fallimento è la coppia formata dal ciclo $C_6$ e dall'unione
disgiunta di due triangoli: entrambi i grafi sono $2$-regolari, tutti i vertici restano
dello stesso colore a ogni iterazione, e WL-1 risponde \emph{vero} pur trattandosi di
grafi non isomorfi. È questa cecità che si cerca di curare salendo di dimensione.

\section{WL-\texorpdfstring{$k$}{k}: colorare le \texorpdfstring{$k$}{k}-uple}

Il potere discriminante del test si accresce colorando non più i singoli vertici, ma le
$k$-uple di vertici: ogni $\bar v = (v_1,\dots,v_k) \in V^{k}$ viene trattata come un
unico oggetto da colorare. Servono due ingredienti --- i \emph{tipi atomici}, che danno
la colorazione iniziale, e i \emph{vicinati}, che danno il passo di raffinamento.

\subsection{Tipi atomici e colorazione iniziale}

Il colore iniziale di una $k$-upla è il suo \emph{tipo atomico}
$\operatorname{atp}(\bar v)$, cioè l'informazione puramente locale sulla $k$-upla:
quali componenti coincidono ($v_i = v_j$) e quali coppie ordinate $(v_i,v_j)$ sono archi
del grafo. In altri termini $\operatorname{atp}(\bar v)$ è il tipo di isomorfismo del
sottografo \emph{ordinato} indotto da $\bar v$: due $k$-uple $\bar v$ e $\bar v'$
partono con lo stesso colore se e solo se la corrispondenza $v_i \mapsto v_i'$,
componente per componente, è un isomorfismo fra i due sottografi indotti.

\begin{figure}[H]
  \centering
  % wlk-tipo-atomico -- pag. 79
% Il tipo atomico di una k-upla e' la sua "fotografia locale": quali componenti
% coincidono e quali coppie ordinate sono archi. Qui la tripla (u,v,w) con
% u -> v e v -> w, e nessun arco fra u e w.
% NB: nel quaderno le parentesi sono angolari e la freccia v -> w si ferma poco
% prima del pallino; qui parentesi tonde (come nel testo) e freccia che tocca
% il nodo.
\begin{tikzpicture}[x=1cm,y=1cm,line join=round,line cap=round,font=\footnotesize,
    nodo/.style={circle,fill,inner sep=0pt,minimum size=3.6pt},
    arco/.style={line width=0.8pt,-{Stealth[length=4pt,width=3pt]},shorten >=0.6pt}]

  % ---- la scritta a sinistra, come nel quaderno ----------------------------
  \node[anchor=east,inner sep=0pt] at (-0.75,0.28) {$\operatorname{atp}(u,v,w,\dots)$};

  % ---- i tre vertici della tripla -----------------------------------------
  \node[nodo] (u) at (0.00, 0.00) {};
  \node[nodo] (v) at (1.10, 0.90) {};
  \node[nodo] (w) at (1.15,-0.35) {};

  \node[anchor=south east,inner sep=2pt] at (u.north west) {$u$};
  \node[anchor=south,inner sep=2.5pt]    at (v.north)      {$v$};
  \node[anchor=west,inner sep=3pt]       at (w.east)       {$w$};

  % ---- gli archi orientati -------------------------------------------------
  \draw[arco] (u) -- (v);
  \draw[arco] (v) -- (w);

\end{tikzpicture}

  \caption{Il tipo atomico della $k$-upla $(u,v,w,\dots)$ registra le coincidenze fra le
  componenti e le adiacenze: qui $u \to v$ e $v \to w$.}
  \label{fig:wlk-tipo-atomico}
\end{figure}

\scan{80}

La colorazione iniziale è dunque
\[
  C_0^k(\bar v) \;=\; \operatorname{atp}(\bar v)
  \qquad \forall\, \bar v \in V^{k} .
\]

\begin{figure}[H]
  \centering
  % wlk-tupla -- pag. 80
% Le k componenti v_1,...,v_k sono raccolte in un unico oggetto vbar: e' a
% questo, e non ai singoli vertici, che WL-k assegna un colore.
% NB: nel quaderno i nodi in alto sono quattro (v_1..v_4) e il nodo in basso e'
% spostato verso sinistra; qui l'ultimo nodo e' etichettato v_k, con i puntini
% di sospensione, per accordarsi alla notazione del testo, e il nodo in basso
% e' centrato sotto la fila.
\begin{tikzpicture}[x=1cm,y=1cm,line join=round,line cap=round,font=\footnotesize,
    nodo/.style={circle,fill,inner sep=0pt,minimum size=3.4pt},
    tupla/.style={circle,fill,inner sep=0pt,minimum size=5pt},
    arco/.style={line width=0.7pt}]

  % ---- le componenti della k-upla -----------------------------------------
  \node[nodo] (v1) at (0.0,0) {};
  \node[nodo] (v2) at (1.2,0) {};
  \node[nodo] (v3) at (2.4,0) {};
  \node[nodo] (vk) at (3.6,0) {};

  \node[anchor=south,inner sep=2.5pt] at (v1.north) {$v_1$};
  \node[anchor=south,inner sep=2.5pt] at (v2.north) {$v_2$};
  \node[anchor=south,inner sep=2.5pt] at (v3.north) {$v_3$};
  \node[anchor=south,inner sep=2.5pt] at (vk.north) {$v_k$};

  \node[inner sep=1pt] at (3.0,0) {$\cdots$};

  % ---- la k-upla, vista come un solo oggetto ------------------------------
  \node[tupla] (vb) at (1.8,-1.8) {};
  \node[anchor=north,inner sep=4pt] at (vb.south) {$\bar v$};

  \draw[arco] (v1) -- (vb);
  \draw[arco] (v2) -- (vb);
  \draw[arco] (v3) -- (vb);
  \draw[arco] (vk) -- (vb);

\end{tikzpicture}

  \caption{La $k$-upla $\bar v = (v_1,\dots,v_k)$ è vista come un solo oggetto: è ad
  essa, e non ai vertici che la compongono, che WL-$k$ assegna un colore.}
  \label{fig:wlk-tupla}
\end{figure}

\subsection{Vicinati e passo di raffinamento}

I \emph{vicinati} di una $k$-upla sono le $k$-uple che si ottengono da essa sostituendo
una componente alla volta con un qualunque vertice del grafo. Scriviamo dunque
$\bar v[w/j]$ per la $k$-upla ottenuta da $\bar v$ rimpiazzando la $j$-esima componente
con $w$:
\[
  \bar v[w/j] \;=\; (v_1,\dots,v_{j-1},\,w,\,v_{j+1},\dots,v_k).
\]
Il vicinato di $\bar v$ al passo $i$ è il multi-insieme
\[
  M_i^k(\bar v) \;=\;
  \Bigl\{\!\Bigl\{\;
     \bigl( C_{i-1}^k(\bar v[w/1]),\; \dots,\; C_{i-1}^k(\bar v[w/k]) \bigr)
  \;:\; w \in V \;\Bigr\}\!\Bigr\}
\]
e il colore aggiornato è
\[
  C_i^k(\bar v) \;=\; \bigl( C_{i-1}^k(\bar v),\; M_i^k(\bar v) \bigr).
\]
Poiché il nuovo colore contiene il vecchio, ogni passo può solo raffinare la partizione
di $V^{k}$ indotta dal colore: dopo al più $n^{k}$ iterazioni --- tante quante sono le
$k$-uple --- la colorazione si stabilizza, e il test confronta i multi-insiemi dei
colori stabili dei due grafi esattamente come faceva WL-1.

\begin{esempio}[WL-2]
  Per $k = 2$ la tupla è una coppia $(u,v)$ e il vertice $w$, che scorre su tutto $V$,
  la spezza nei due «mezzi cammini» $(u,w)$ e $(w,v)$:
  \[
    M_i^2(u,v) \;=\;
    \Bigl\{\!\Bigl\{\; \bigl( C_{i-1}^2(u,w),\; C_{i-1}^2(w,v) \bigr) \;:\; w \in V
    \;\Bigr\}\!\Bigr\}.
  \]
  L'ordine delle due componenti della coppia è convenzionale --- qui è quello del
  cammino $u \to w \to v$, mentre la formula generale elenca le posizioni nell'ordine
  $1,\dots,k$ --- e basta fissarlo una volta per tutte.
\end{esempio}

\begin{figure}[H]
  \centering
  % wl2-esempio -- pag. 80
% WL-2: la coppia (u,v) viene guardata «attraverso» un terzo vertice w.
% Frecce continue = archi del grafo diretto (w->v, x->v, u->v, u->w, u->x);
% triangolo tratteggiato = le tre COPPIE in gioco nel passo di raffinamento,
% cioe' (u,v), (u,w) e (w,v).  Il tratteggio non e' un arco del grafo.
\begin{tikzpicture}[
    >=Latex,
    font=\footnotesize,
    nodo/.style={circle,fill=black,inner sep=0pt,minimum size=3.4pt},
    arco/.style={->,line width=0.7pt},
    coppia/.style={dashed,dash pattern=on 2.4pt off 1.8pt,
                   line width=0.6pt,draw=black!55},
    et/.style={font=\footnotesize,text=black!60,inner sep=1.5pt}]

  %% ---- vertici ------------------------------------------------------------
  \node[nodo,label={[inner sep=2pt]above left:$v$}]  (v) at ( 0.00, 2.00) {};
  \node[nodo,label={[inner sep=2pt]right:$w$}]       (w) at ( 2.80, 1.60) {};
  \node[nodo,label={[inner sep=3pt]below:$u$}]       (u) at ( 1.25, 0.00) {};
  \node[nodo,label={[inner sep=2pt]left:$x$}]        (x) at (-1.40,-0.10) {};

  %% ---- archi del grafo (continui, orientati) ------------------------------
  \draw[arco] (w) -- (v);
  \draw[arco] (x) -- (v);
  \draw[arco] (u) -- (v);
  \draw[arco] (u) -- (w);
  \draw[arco] (u) -- (x);

  %% ---- le tre coppie coinvolte nel passo (tratteggio, senza punte) --------
  \draw[coppia] (v) to[bend left=22]  node[et,above=1pt]           {$(w,v)$} (w);
  \draw[coppia] (w) to[bend left=22]  node[et,pos=0.45,right=1pt]  {$(u,w)$} (u);
  \draw[coppia] (v) to[bend right=18] node[et,pos=0.55,left=1pt]   {$(u,v)$} (u);

\end{tikzpicture}

  \caption{WL-2: alla coppia $(u,v)$ si guarda attraverso un terzo vertice $w$. Le
  frecce continue sono archi del grafo, il triangolo tratteggiato evidenzia le tre
  \emph{coppie} coinvolte nel passo di raffinamento: $(u,v)$, $(u,w)$ e $(w,v)$.}
  \label{fig:wl2-esempio}
\end{figure}

\subsection{FWL contro OWL}

La definizione appena data è quella \emph{folklore} (FWL): per ogni $w$ si raccoglie
un'unica $k$-upla di colori, quindi il multi-insieme ricorda \emph{quali} colori
provengono dallo stesso $w$. Nella variante \emph{oblivious} (OWL) si tengono invece
$k$ multi-insiemi separati, uno per posizione,
\[
  M_i^{k,j}(\bar v) \;=\;
  \bigl\{\!\bigl\{\, C_{i-1}^k(\bar v[w/j]) \;:\; w \in V \,\bigr\}\!\bigr\},
  \qquad j = 1,\dots,k,
\]
\[
  C_i^k(\bar v) \;=\;
  \bigl( C_{i-1}^k(\bar v),\; M_i^{k,1}(\bar v),\; \dots,\; M_i^{k,k}(\bar v) \bigr),
\]
e la correlazione fra le posizioni va perduta: FWL legge i due mezzi cammini
\emph{insieme}, per lo stesso $w$; OWL li legge \emph{separatamente}.

\begin{figure}[H]
  \centering
  % fwl-vs-owl -- pag. 80
% Trittico.  (a) configurazione di partenza: la coppia (u,v) da colorare
% (tratteggiata, come in wl2-esempio) e i due mezzi cammini u->w, w->v.
% (b) FWL: i due mezzi cammini stanno nello STESSO disegno, con un solo w:
%     il multi-insieme ricorda che i due colori vengono dallo stesso w.
% (c) OWL: i due mezzi cammini vivono in due schemi staccati, ciascuno con la
%     propria copia di w: la correlazione fra le due posizioni e' persa.
\begin{tikzpicture}[
    >=Latex,
    font=\footnotesize,
    nodo/.style={circle,fill=black,inner sep=0pt,minimum size=3.2pt},
    arco/.style={->,line width=0.7pt},
    coppia/.style={dashed,dash pattern=on 2.4pt off 1.8pt,
                   line width=0.6pt,draw=black!55},
    titolo/.style={font=\footnotesize\scshape},
    glossa/.style={font=\footnotesize,align=center,text width=3.5cm}]

  %% =================== (a) configurazione di partenza ======================
  \begin{scope}[shift={(0,0)}]
    \node[nodo,label={[inner sep=2pt]above left:$v$}] (av) at (0.00,1.45) {};
    \node[nodo,label={[inner sep=2pt]right:$w$}]      (aw) at (1.70,1.30) {};
    \node[nodo,label={[inner sep=3pt]below:$u$}]      (au) at (0.75,0.00) {};
    \draw[arco]   (aw) -- (av);
    \draw[arco]   (au) -- (aw);
    \draw[coppia] (au) -- (av);
    \node[glossa] at (0.85,-0.85) {configurazione\\di partenza};
  \end{scope}

  %% ------------------------------ separatore -------------------------------
  \draw[draw=black!25,line width=0.5pt] (2.95,-1.15) -- (2.95,2.15);

  %% ============================== (b) FWL ==================================
  \begin{scope}[shift={(3.75,0)}]
    \node[titolo] at (0.85,2.00) {FWL};
    \node[nodo,label={[inner sep=2pt]above left:$v$}] (bv) at (0.00,1.45) {};
    \node[nodo,label={[inner sep=2pt]right:$w$}]      (bw) at (1.70,1.30) {};
    \node[nodo,label={[inner sep=3pt]below:$u$}]      (bu) at (0.75,0.00) {};
    \draw[arco] (bw) -- (bv);
    \draw[arco] (bu) -- (bw);
    \node[glossa] at (0.85,-0.85) {considera la coppia:\\lo stesso $w$};
  \end{scope}

  %% ------------------------------ separatore -------------------------------
  \draw[draw=black!25,line width=0.5pt] (6.70,-1.15) -- (6.70,2.15);

  %% ============================== (c) OWL ==================================
  \begin{scope}[shift={(7.50,0)}]
    \node[titolo] at (1.05,2.00) {OWL};
    % primo mezzo cammino: w -> v
    \node[nodo,label={[inner sep=2pt]above left:$v$}] (cv)  at (0.00,1.65) {};
    \node[nodo,label={[inner sep=2pt]right:$w$}]      (cw1) at (1.00,1.20) {};
    \draw[arco] (cw1) -- (cv);
    % secondo mezzo cammino: u -> w  (copia distinta di w, schema staccato)
    \node[nodo,label={[inner sep=3pt]below:$u$}]      (cu)  at (0.85,0.00) {};
    \node[nodo,label={[inner sep=2pt]right:$w$}]      (cw2) at (1.85,0.65) {};
    \draw[arco] (cu) -- (cw2);
    \node[glossa] at (1.05,-0.85) {considera i cammini\\separatamente};
  \end{scope}

\end{tikzpicture}

  \caption{FWL contro OWL sulla coppia $(u,v)$. A sinistra la configurazione di
  partenza; al centro FWL, che considera la coppia di colori nello stesso disegno
  (stesso $w$); a destra OWL, che considera i due cammini separatamente, in due
  multi-insiemi distinti.}
  \label{fig:fwl-owl}
\end{figure}

\begin{osservazione}
  Le due varianti non hanno lo stesso potere distintivo a parità di $k$: per $k \ge 2$,
  FWL-$k$ distingue esattamente ciò che distingue OWL-$(k+1)$. La gerarchia è quindi la
  stessa, ma sfasata di uno: quando in questo capitolo si scrive WL-$k$ senza altra
  specificazione si intende la versione folklore, che è quella definita qui sopra.
\end{osservazione}

\begin{osservazione}[Il caso $k = 1$ va letto a parte]
  Presa alla lettera, la formula del passo di raffinamento è vuota per $k = 1$: si ha
  $\bar v = (v)$ e $\bar v[w/1] = (w)$, quindi
  \[
    M_i^{1}(v) \;=\; \bigl\{\!\bigl\{\, C_{i-1}^{1}(w) \;:\; w \in V \,\bigr\}\!\bigr\}
  \]
  non dipende affatto da $v$, e la colorazione non si raffina mai. È il motivo per cui
  WL-1 --- il raffinamento dei colori della sezione precedente --- si definisce a parte,
  facendo scorrere $w$ sui soli \emph{vicini} di $v$ anziché su tutto $V$. Per $k \ge 2$
  la restrizione non serve: il tipo atomico di $\bar v[w/j]$ registra già se $w$ sia o
  no adiacente alle altre componenti, e l'informazione di adiacenza non va persa.
  Vale inoltre che OWL-2 ha esattamente il potere distintivo di WL-1: convenendo che
  FWL-1 sia per definizione WL-1, lo sfasamento di uno dell'osservazione precedente si
  estende anche al primo livello, e la gerarchia comincia a essere davvero più fine da
  OWL-3, ossia da FWL-2, in poi.
\end{osservazione}

\section{Che cosa distingue WL-\texorpdfstring{$k$}{k}}

Fissato $k$, quali coppie di grafi sfuggono al test? La domanda ha tre risposte, di
natura molto diversa, ed è a queste che il quaderno dedica le ultime righe:
una \emph{combinatoria} --- su quali classi di grafi un $k$ fissato basta ---, una
\emph{algebrica} --- il rilassamento continuo dell'isomorfismo ---, e una \emph{logica}
--- il frammento del prim'ordine che ha lo stesso potere separante del test.

\subsection{Minori esclusi}

Escludere un minore fisso comporta dimensione di WL limitata: per ogni grafo $H$ esiste
un $k = k(H)$ tale che ogni grafo che non contiene $H$ come minore viene identificato da
WL-$k$, cioè non è confuso con nessun grafo non isomorfo. Il caso emblematico è quello
dei grafi planari, che per il teorema di Wagner sono esattamente i grafi privi di $K_5$ e
di $K_{3,3}$ come minori: sui planari, dunque, una dimensione costante basta.

\begin{figure}[H]
  \centering
  % minori-esclusi-wagner -- pag. 80
% I due minori proibiti del teorema di Wagner.  Nel quaderno sono due schizzi a
% matita appena accennati (K_{3,3} ridotto a due terne di vertici, K_5 cerchiato
% a mano): qui sono ridisegnati per intero e corretti.
%   K_{3,3}: bipartito completo, due classi da tre vertici, tutti i 9 archi.
%   K_5:     completo su 5 vertici disposti a pentagono, tutti i 10 archi.
\begin{tikzpicture}[
    font=\footnotesize,
    nodo/.style={circle,fill=black,inner sep=0pt,minimum size=3.2pt},
    arco/.style={line width=0.55pt}]

  %% ============================== K_{3,3} ==================================
  \begin{scope}[shift={(0,0)}]
    \foreach \i in {1,2,3} {
      \node[nodo] (a\i) at ({(\i-1)*0.95}, 1.50) {};
      \node[nodo] (b\i) at ({(\i-1)*0.95}, 0.00) {};
    }
    \foreach \i in {1,2,3}
      \foreach \j in {1,2,3}
        \draw[arco] (a\i) -- (b\j);
    \node at (0.95,-0.68) {$K_{3,3}$};
  \end{scope}

  %% ================================ K_5 ====================================
  \begin{scope}[shift={(4.30,0.75)}]
    \foreach \i/\ang in {1/90, 2/162, 3/234, 4/306, 5/18}
      \node[nodo] (c\i) at (\ang:1.10) {};
    \foreach \i/\j in {1/2,1/3,1/4,1/5,2/3,2/4,2/5,3/4,3/5,4/5}
      \draw[arco] (c\i) -- (c\j);
    \node at (0,-1.43) {$K_5$};
  \end{scope}

\end{tikzpicture}

  \caption{I due minori proibiti del teorema di Wagner: $K_{3,3}$ e $K_5$.}
  \label{fig:minori-esclusi}
\end{figure}

\subsection{Isomorfismo frazionario}

\scan{81}

Se fra due grafi sussiste un isomorfismo frazionario, il raffinamento dei colori
fallisce: i due grafi ricevono la stessa colorazione e WL-1 non riesce a separarli.

\begin{definizione}[Isomorfismo frazionario]
  Siano $G$ e $H$ grafi su $n$ nodi, con matrici di adiacenza $A_G$ e $A_H$. Diciamo che
  $G$ e $H$ sono \emph{frazionariamente isomorfi}, e scriviamo $G \cong_f H$, se esiste
  una matrice doppiamente stocastica $S \in \RR^{n \times n}$ --- cioè con entrate non
  negative e con tutte le righe e tutte le colonne di somma $1$ --- tale che
  \[
    A_G\,S \;=\; S\,A_H .
  \]
\end{definizione}

\begin{osservazione}
  Un isomorfismo vero e proprio è una matrice di \emph{permutazione} $P$ con
  $A_G P = P A_H$; l'isomorfismo frazionario è il rilassamento continuo di quella
  condizione, cioè l'ammissibilità del programma lineare che si ottiene sostituendo le
  matrici di permutazione con quelle doppiamente stocastiche. Il legame con
  Weisfeiler--Leman è esatto, ma vale al primo livello: $G \cong_f H$ se e solo se il
  raffinamento dei colori, cioè WL-1, produce sui due grafi lo stesso multi-insieme di
  colori. È questo il senso della frase qui sopra --- quando l'isomorfismo frazionario
  c'è, la colorazione non ha più nulla su cui fare presa. Ai livelli superiori il
  fallimento non è più garantito: il ciclo $C_6$ e l'unione di due triangoli sono
  frazionariamente isomorfi, essendo entrambi $2$-regolari, eppure WL-2 li distingue,
  perché il suo passo di raffinamento conta i vicini comuni di ogni coppia di vertici
  ($0$ nel ciclo, $1$ per gli archi dei triangoli). Salendo nella gerarchia, la
  controparte di WL-$k$ non è più il programma lineare «nudo», ma il suo livello $k$
  nella gerarchia di Sherali--Adams.
\end{osservazione}

\subsection{La logica con conteggio}

\begin{definizione}[Il frammento $C^{k}$]
  $C^{k}$ è il frammento della logica del prim'ordine in cui si usano
  \begin{itemize}
    \item al più $k$ variabili distinte;
    \item quantificatori esistenziali \emph{con conteggio}
      \[
        \exists^{\ge m} x \; \varphi(x), \qquad m \in \NN,
      \]
      da leggersi «esistono almeno $m$ elementi $x$ per cui vale $\varphi(x)$».
  \end{itemize}
\end{definizione}

\begin{lacuna}
  Nel quaderno, sotto la definizione di $C^{k}$, compare un secondo punto elenco lasciato
  in bianco.
\end{lacuna}

\begin{osservazione}[Cai--Fürer--Immerman, 1992]
  Il punto rimasto in sospeso è quasi certamente l'enunciato che giustifica
  l'introduzione di $C^{k}$: WL-$k$ (nella versione \emph{folklore} adottata qui)
  distingue due grafi $G$ e $H$ se e solo se esiste un enunciato di $C^{k+1}$ vero in uno
  e falso nell'altro. Lo sfasamento di uno è lo stesso già incontrato fra FWL e OWL: è
  OWL-$k$, non FWL-$k$, ad avere il potere distintivo di $C^{k}$. Per $k = 1$ si ritrova
  il caso noto: il raffinamento dei colori ha esattamente il potere distintivo di $C^{2}$,
  la logica a due variabili con conteggio.
\end{osservazione}

\section{Il kernel di Weisfeiler--Leman}

L'algoritmo si può usare anche, al di fuori del problema dell'isomorfismo, come
estrattore di caratteristiche per il machine learning su grafi: si eseguono $r$
iterazioni di Weisfeiler--Leman sul grafo e si conta quante volte compare ciascun colore.
Il vettore dei conteggi $\varphi_r(G)$ è la rappresentazione del grafo, e il kernel è il
prodotto scalare fra le rappresentazioni,
\[
  K(G,H) \;=\; \langle\, \varphi_r(G),\; \varphi_r(H) \,\rangle .
\]

\begin{figure}[H]
  \centering
  % wl-kernel-diamante -- pag. 81
% Il grafo a matita sotto il titolo «WL-kernel»: ciclo di lunghezza 4 con una
% corda, cioe' K_4 privato di un lato (il «diamante»).  I quattro lati del
% quadrato piu' la sola diagonale alto-destra / basso-sinistra: i due vertici
% che essa collega hanno grado 3, gli altri due grado 2.
% Il quadrato e' leggermente irregolare, come nel manoscritto; la corda e'
% tracciata un filo piu' sottile perche' nella scansione ha tratto piu' leggero.
\begin{tikzpicture}[
    font=\footnotesize,
    nodo/.style={circle,fill=black,inner sep=0pt,minimum size=3.2pt},
    arco/.style={line width=0.65pt},
    corda/.style={line width=0.5pt}]

  \node[nodo] (a) at (0.00, 1.85) {};   % alto a sinistra   (grado 2)
  \node[nodo] (b) at (1.90, 1.95) {};   % alto a destra     (grado 3)
  \node[nodo] (c) at (0.10, 0.00) {};   % basso a sinistra  (grado 3)
  \node[nodo] (d) at (1.85,-0.05) {};   % basso a destra    (grado 2)

  \draw[arco]  (a) -- (b);
  \draw[arco]  (a) -- (c);
  \draw[arco]  (b) -- (d);
  \draw[arco]  (c) -- (d);
  \draw[corda] (b) -- (c);

\end{tikzpicture}

  \caption{Il grafo disegnato a matita sotto il titolo «WL-kernel»: un ciclo di lunghezza
  $4$ con una corda, cioè $K_4$ privato di un lato.}
  \label{fig:wl-kernel-diamante}
\end{figure}

\begin{osservazione}
  Annotazione a matita, rimasta a metà: «grafi indistinguibili da WL-$k$
  $\Longrightarrow$ \dots». La conclusione non è stata scritta, ma è obbligata: se
  WL-$k$ non distingue $G$ da $H$, i due grafi hanno lo stesso multi-insieme di colori,
  dunque $\varphi_r(G) = \varphi_r(H)$ e il kernel non li separa. I limiti di
  espressività del test si trasferiscono di peso al kernel, e per lo stesso motivo ---
  l'aggiornamento è quello di Weisfeiler--Leman --- anche alle reti neurali su grafi a
  passaggio di messaggi, che non superano il potere distintivo di WL-1.
\end{osservazione}

% --- Gli appunti si interrompono qui: le pagine successive del quaderno sono bianche.


\end{document}
